% ================================================================
%  main.tex  —  Compilador maestro
%  Almería en Cifras y Formas · Presentaciones Beamer
%  ---------------------------------------------------------------
%  Este archivo reúne todas las sesiones en un único PDF.
%  Para compilar una sesión individual, abre directamente
%  sesionXX.tex y compila con pdflatex (dos pasadas).
%
%  Para Overleaf: sube toda la carpeta como proyecto.
%  El archivo principal de Overleaf debe ser main.tex.
%
%  Compilación local (requiere TeX Live o MiKTeX):
%    pdflatex main.tex
%    pdflatex main.tex   (segunda pasada para referencias)
% ================================================================
\documentclass[aspectratio=169, 12pt, spanish]{beamer}
\usepackage{almeria-beamer}

% ---------------------------------------------------------------
%  Datos del proyecto (comunes a todas las sesiones)
% ---------------------------------------------------------------
\renewcommand{\ProjectName}{Almería en Cifras y Formas}
\renewcommand{\AuthorName}{Juan María Gámez Ortiz}
\renewcommand{\InstituteName}{IES Al-Ándalus}
\renewcommand{\SchoolYear}{2025--2026}
\renewcommand{\SchoolSubject}{Matemáticas 1.º ESO}

% ---------------------------------------------------------------
%  Portada del proyecto completo
% ---------------------------------------------------------------
\renewcommand{\SessionNumber}{\phantom{0}}
\renewcommand{\SessionTitle}{Presentaciones de aula}
\renewcommand{\SessionName}{Colección completa}
\renewcommand{\SessionObjective}{%
  Guía teórica completa con ejemplos, soluciones paso a paso
  y representaciones gráficas para las 20 sesiones de la
  situación de aprendizaje «Almería en Cifras y Formas».}
\renewcommand{\BlockName}{Bloques 1--4 · 20 sesiones}

\begin{document}

\begin{frame}[plain]
  % Portada especial de proyecto completo
  \begin{tikzpicture}[remember picture, overlay]
    \fill[AlmAzul] (current page.south west) rectangle (current page.north east);
    \fill[AlmAzulM, opacity=0.4]
      (current page.south west) --
      ([yshift=3.5cm]current page.south west) ..controls
      ([xshift=5cm,yshift=4cm]current page.south west) and
      ([xshift=-5cm,yshift=3cm]current page.south east)..
      ([yshift=3.5cm]current page.south east) --
      (current page.south east) -- cycle;
    \fill[AlmAmbar]
      ([yshift=-3pt]current page.north west)
      rectangle
      ([yshift=-10pt]current page.north east);
  \end{tikzpicture}
  \vspace{0.6cm}
  \begin{center}
    {\color{AlmAmbar}\normalsize\bfseries
      \faMapMarkerAlt\quad\ProjectName}\\[0.5cm]
    {\color{AlmBlanco}\Huge\bfseries Presentaciones de aula}\\[0.3cm]
    {\color{AlmBlanco!70}\large 20 sesiones · Guía teórica completa}\\[0.8cm]
    \begin{tcolorbox}[
      enhanced, width=0.72\textwidth,
      colback=AlmAzulM!50!AlmAzul,
      colframe=AlmAmbar, boxrule=1pt, arc=4pt,
      halign=center]
      {\color{AlmBlanco}\small
        Bloques 1--4 · \SchoolSubject · \SchoolYear\\
        \AuthorName · \InstituteName}
    \end{tcolorbox}
  \end{center}
\end{frame}

% ---------------------------------------------------------------
%  Índice general
% ---------------------------------------------------------------
\begin{frame}{Índice general de sesiones}
  \begin{columns}[T]
    \begin{column}{0.48\textwidth}
      {\color{AlmAzulM}\bfseries Bloque 1: Mapa y coordenadas}
      \begin{enumerate}
        \item ¿Qué es una guía turística matemática?
        \item ¿Qué se puede medir en una ciudad?
        \item Plano por cuadrantes
        \item Coordenadas y localización
      \end{enumerate}
      \vspace{6pt}
      {\color{AlmTerra}\bfseries Bloque 2: Geometría urbana}
      \begin{enumerate}\setcounter{enumi}{4}
        \item Rectas, segmentos y ángulos
        \item Figuras planas en la ciudad
        \item Perímetro de zonas urbanas
        \item Área de zonas urbanas
        \item Perímetro y área: no son lo mismo
        \item Rutas y distancias por Almería
      \end{enumerate}
    \end{column}
    \begin{column}{0.48\textwidth}
      {\color{AlmVerde}\bfseries Bloque 3: Datos y gráficas}
      \begin{enumerate}\setcounter{enumi}{10}
        \item Variables urbanas
        \item De enunciado a tabla
        \item De tabla a gráfica
        \item Leer e interpretar gráficas
        \item Comparar gráficas
        \item Relaciones funcionales y proporcionalidad
      \end{enumerate}
      \vspace{6pt}
      {\color{BCrear}\bfseries Bloque 4: Producto final}
      \begin{enumerate}\setcounter{enumi}{16}
        \item Seleccionar la información de la guía
        \item Construimos la guía
        \item Ensayo y mejora
        \item Exposición y cierre del proyecto
      \end{enumerate}
    \end{column}
  \end{columns}
\end{frame}

% ================================================================
%  INCLUSIÓN DE TODAS LAS SESIONES
%  Cada \input carga el cuerpo de la sesión (sin \documentclass).
%  NOTA: Para incluir sesiones individuales descomenta las líneas.
%  En Overleaf, puedes compilar el main.tex completo o
%  cada sesionXX.tex de forma independiente.
% ================================================================

% Descomenta las líneas siguientes para generar el PDF completo.
% Con todas activadas, el PDF resultante tiene ~200+ diapositivas.

% % ================================================================
%  sesion01.tex  —  Sesión 1 de 20
%  ¿Qué es una guía turística matemática?
%  Almería en Cifras y Formas · Matemáticas 1.º ESO
%  Compilar con:  pdflatex sesion01.tex  (dos pasadas)
% ================================================================
\documentclass[aspectratio=169, 10pt, spanish]{beamer}
\usepackage{almeria-beamer}

% ---------------------------------------------------------------
%  DATOS DE LA SESIÓN
% ---------------------------------------------------------------
\renewcommand{\SessionNumber}{01}
\renewcommand{\SessionTitle}{¿Qué es una guía turística matemática?}
\renewcommand{\SessionName}{Sesión 01}
\renewcommand{\SessionObjective}{%
  Identificar qué tipo de información puede tener carácter
  matemático en una ciudad y distinguir datos descriptivos
  de datos cuantitativos usando Almería como contexto.}
\renewcommand{\BlockName}{Bloque 1: Mapa y coordenadas}

% ================================================================
\begin{document}

% ---------------------------------------------------------------
%  PORTADA
% ---------------------------------------------------------------
\begin{frame}[plain]
  \titlepage
\end{frame}

% ---------------------------------------------------------------
%  ÍNDICE DE LA SESIÓN
% ---------------------------------------------------------------
\begin{frame}{Índice de la sesión}
  \begin{columns}[t]
    \begin{column}{0.5\textwidth}
      \begin{enumerate}
        \item ¿Qué es una guía turística?
        \item ¿Qué la hace \textit{matemática}?
        \item Tipos de información: descriptiva vs.\ cuantitativa
        \item Datos matemáticos ocultos en Almería
        \item Ejemplo paso a paso: el Cable Inglés
      \end{enumerate}
    \end{column}
    \begin{column}{0.5\textwidth}
      \begin{enumerate}
        \setcounter{enumi}{5}
        \item Ejercicios multinivel (Bloom)
        \item Evidencia del día
        \item Error frecuente
        \item Resumen
      \end{enumerate}
    \end{column}
  \end{columns}
\end{frame}

% ---------------------------------------------------------------
\seccionframe{1. ¿Qué es una guía turística?}
% ---------------------------------------------------------------

\begin{frame}{¿Qué es una guía turística?}
  \begin{columns}[c]
    \begin{column}{0.52\textwidth}
      \begin{definicionbox}{Guía turística}
        Publicación (papel o digital) que describe los
        lugares de interés de una ciudad o región:
        monumentos, museos, rutas, horarios, precios\ldots
      \end{definicionbox}
      
      \begin{itemize}
        \item Cuenta \textbf{qué} hay y \textbf{cómo llegar}.
        \item Usa palabras, fotos y mapas.
        \item Va dirigida al \textbf{turista}.
      \end{itemize}
      
      \begin{notabox}[Ejemplos conocidos]
        
        Lonely Planet · Guía Repsol · Tripadvisor
      \end{notabox}
    \end{column}
    \begin{column}{0.44\textwidth}
      \begin{center}
        \begin{tikzpicture}[scale=0.9]
          % Portada de libro estilizada
          \fill[AlmAzul, rounded corners=4pt]
            (0,0) rectangle (3.2,4.4);
          \fill[AlmAmbar]
            (0.15,3.6) rectangle (3.05,4.2);
          \node[text=AlmBlanco, font=\bfseries] at (1.6,3.9)
            {ALMERÍA};
          \fill[AlmAzulC!40]
            (0.15,2.1) rectangle (3.05,3.5);
          \node[text=AlmAzul, font=] at (1.6,2.8)
            {mapa de la ciudad};
          \foreach \y in {1.2,0.85,0.5}{
            \fill[AlmBlanco!30, rounded corners=1pt]
              (0.3,\y) rectangle (2.8,\y+0.22);
          }
          \node[text=AlmAmbar, font=\bfseries] at (1.6,0.05)
            {Guía Turística};
        \end{tikzpicture}
      \end{center}
    \end{column}
  \end{columns}
\end{frame}

% ---------------------------------------------------------------
\seccionframe{2. ¿Qué la hace \textit{matemática}?}
% ---------------------------------------------------------------

\begin{frame}{¿Qué hace matemática a una guía?}
  \begin{columns}[c]
    \begin{column}{0.55\textwidth}
      \begin{teoriabox}
        Una guía turística \textbf{matemática} incorpora,
        además de texto descriptivo, datos numéricos,
        medidas y representaciones cuantitativas que
        permiten \textit{analizar} y \textit{comparar}
        la realidad de la ciudad.
      \end{teoriabox}
      {8pt}
      ¿Qué puede aparecer?
      \begin{itemize}
        \item \textbf{Coordenadas} de monumentos
        \item \textbf{Distancias} entre puntos turísticos
        \item \textbf{Alturas}, superficies, aforos
        \item \textbf{Gráficas} de visitantes o temperatura
        \item \textbf{Estadísticas} históricas
      \end{itemize}
    \end{column}
    \begin{column}{0.42\textwidth}
      \begin{tikzpicture}[every node/.style={font=}]
        % Diagrama Venn
        \fill[AlmAzulC!25, opacity=0.7] (-0.6,0) ellipse (2cm and 1.4cm);
        \fill[AlmAmbar!35, opacity=0.7] (0.6,0) ellipse (2cm and 1.4cm);
        \node[text=AlmAzul, align=center] at (-1.4,0) {Guía\\clásica};
        \node[text=AlmAzul] at (2.0,0) {Matemáticas};
        \node[text=AlmNegro, align=center] at (0,0)
          {Guía\\matemática};
        \node[text=AlmGris, font=, align=center] at (-1.5,-0.55)
          {fotos\\textos};
        \node[text=AlmGris, font=, align=center] at (1.6,-0.55)
          {cálculos\\gráficas};
        \node[text=AlmAzul, font=, align=center] at (0,0.55)
          {datos\\medidas};
      \end{tikzpicture}
    \end{column}
  \end{columns}
\end{frame}

% ---------------------------------------------------------------
\seccionframe{3. Información descriptiva vs.\ cuantitativa}
% ---------------------------------------------------------------

\begin{frame}{Dos tipos de información}
  \begin{columns}[T]
    \begin{column}{0.48\textwidth}
      \begin{definicionbox}{Información descriptiva}
        \begin{itemize}
          \item Usa palabras, no números.
          \item Describe \textbf{cómo es} algo.
          \item Ejemplo: «La Alcazaba es una fortaleza árabe
                de impresionante belleza.»
        \end{itemize}
      \end{definicionbox}
    \end{column}
    \begin{column}{0.48\textwidth}
      \begin{definicionbox}{Información cuantitativa}
        \begin{itemize}
          \item Usa \textbf{números y unidades}.
          \item Dice \textbf{cuánto} o \textbf{cuántos}.
          \item Ejemplo: «La Alcazaba ocupa
                $43\,000\ \text{m}^2$ y
                recibe $300\,000$ visitantes/año.»
        \end{itemize}
      \end{definicionbox}
    \end{column}
  \end{columns}
  
  \begin{center}
    \begin{tabular}{lll}
      \toprule
      \textbf{Monumento} & \textbf{Descriptivo} & \textbf{Cuantitativo}\\
      \midrule
      Alcazaba   & Fortaleza islámica medieval & $43\,000\ \text{m}^2$, s.\ X\\
      Cable Inglés & Viaducto de hierro único   & $1{,}1\ \text{km}$, 1904\\
      Catedral   & Estilo gótico-renacentista   & 55 m altura, 1524\\
      \bottomrule
    \end{tabular}
  \end{center}
\end{frame}

% ---------------------------------------------------------------
\seccionframe{4. Datos matemáticos ocultos en Almería}
% ---------------------------------------------------------------

\begin{frame}{Datos matemáticos ocultos}
  \begin{columns}[T]
    \begin{column}{0.6\textwidth}
      Cualquier elemento urbano esconde información
      matemática. Hay que saber \textbf{hacer preguntas}:
      
      \begin{tabular}{@{}p{2.8cm}p{5cm}@{}}
        \toprule
        \textbf{Lugar} & \textbf{Preguntas matemáticas}\\
        \midrule
        Alcazaba
          & ¿Cuántos metros de muralla tiene?\\[2pt]
        Puerto de Almería
          & ¿Qué área ocupa? ¿Cuántos barcos atraca?\\[2pt]
        Paseo Marítimo
          & ¿Cuántos km mide de extremo a extremo?\\[2pt]
        Catedral
          & ¿Cuánto tarda en dar vuelta la sombra de la torre?\\
        \bottomrule
      \end{tabular}
    \end{column}
    \begin{column}{0.37\textwidth}
      \begin{notabox}[¿Sabías que\ldots?]
        
        El \textbf{Puerto de Almería} ocupa
        \textbf{35 ha}, lo que equivale a
        $35\times 10\,000 = 350\,000\ \text{m}^2$.
        ¡Más de 49 campos de fútbol!
      \end{notabox}
      
      \begin{tikzpicture}[scale=0.75]
        \fill[AlmAzulC!30, rounded corners=3pt]
          (0,0) rectangle (3.8,2.2);
        \fill[AlmVerde!40]
          (0.2,0.2) rectangle (1.6,0.6);
        \node[font=, color=AlmAzul] at (0.9,0.4)
          {campo fútbol};
        \node[font=\bfseries, color=AlmAzul] at (1.9,1.2)
          {$35\ \text{ha}$};
        \node[font=, color=AlmGris] at (1.9,0.7)
          {Puerto};
      \end{tikzpicture}
    \end{column}
  \end{columns}
\end{frame}

% ---------------------------------------------------------------
\seccionframe{5. Ejemplo: el Cable Inglés}
% ---------------------------------------------------------------

\begin{frame}{El Cable Inglés — Ejemplo paso a paso}

  \framesubtitle{Longitud 1,1 km · Velocidad media de paseo 5 km/h}
  
        \begin{ejemplobox}
          
          El Cable Inglés es un viaducto de hierro de
          $1{,}1\ \text{km}$ de longitud que cruza la
          desembocadura de la Rambla.\\[4pt]
          \textbf{Pregunta:} ¿Cuánto tarda en cruzarlo
          un peatón a $5\ \text{km/h}$?
        \end{ejemplobox}
  \begin{columns}[c]
    \begin{column}{0.5\textwidth}
      
      \begin{pasobox}[1]
        
        \textbf{Datos:}
        $d = 1{,}1\ \text{km}$;\quad
        $v = 5\ \text{km/h}$
      \end{pasobox}
      
            \begin{pasobox}[2]
              
              \textbf{Fórmula tiempo--distancia--velocidad:}
              \[
                t = \frac{d}{v}
              \]
            \end{pasobox}
    \end{column}
    \begin{column}{0.46\textwidth}
      
      \begin{pasobox}[3]
        
        \textbf{Sustitución:}
        \[
          t = \frac{1{,}1\ \text{km}}{5\ \text{km/h}}
            = 0{,}22\ \text{h}
        \]
      \end{pasobox}
      
      \begin{pasobox}[4]
        
        \textbf{Conversión a minutos:}
        \[
          t = 0{,}22 \times 60 \approx \mathbf{13\ \text{min}}
        \]
      \end{pasobox}
    \end{column}
  \end{columns}
\end{frame}

% ---------------------------------------------------------------
\seccionframe{6. Ejercicios multinivel — Bloom}
% ---------------------------------------------------------------

\begin{frame}{Ejercicios multinivel}
  \begin{columns}[T]
    \begin{column}{0.31\textwidth}
      \begin{recordarbox}
        
        \textbf{Define} con tus palabras qué es
        una guía turística.\par
        Escribe \textbf{3 monumentos} de Almería
        y anota 2 datos matemáticos que podrías
        buscar de cada uno.
      \end{recordarbox}
    \end{column}
    \begin{column}{0.31\textwidth}
      \begin{comprenderbox}
        
        El Paseo de Almería mide $370\ \text{m}$
        y se recorre a $4\ \text{km/h}$.\par
        \textbf{Explica} qué tipo de información es
        cada dato y calcula cuántos minutos dura
        el paseo.
      \end{comprenderbox}
    \end{column}
    \begin{column}{0.31\textwidth}
      \begin{analizarbox}
        
        Observa tres monumentos de tu elección.\par
        \textbf{Extrae} al menos 2 preguntas
        matemáticas que podrían incluirse en
        una guía turística para cada uno.\par
        ¿Qué tipo de matemáticas necesitarías
        para responderlas?
      \end{analizarbox}
    \end{column}
  \end{columns}
\end{frame}

% ---------------------------------------------------------------
%  EVIDENCIA
% ---------------------------------------------------------------
\begin{frame}{Evidencia del día}
  \begin{evidenciabox}
    \begin{itemize}
      \item Entrega tu \textbf{tarjeta de trabajo} con
            al menos \textbf{5 elementos} que incluirás
            en vuestra guía turística matemática.
      \item Para cada elemento indica:
            \begin{enumerate}
              \item Nombre del monumento o lugar
              \item Un dato \textit{descriptivo} y uno \textit{cuantitativo}
              \item Una pregunta matemática que se podría resolver
            \end{enumerate}
      \item Incluye las operaciones intermedias, no solo el resultado.
    \end{itemize}
  \end{evidenciabox}
\end{frame}

% ---------------------------------------------------------------
%  ERROR FRECUENTE
% ---------------------------------------------------------------
\begin{frame}{¡Cuidado con este error!}
  \begin{errorbox}
    % --- CORRECCIÓN: Uso de minipage en lugar de columns ---
    \begin{minipage}[t]{0.48\textwidth}
      \correcto\ \textbf{Correcto} \\[4pt]
      Las matemáticas sirven para \textbf{describir,
      medir y comparar} la realidad:\par
      «La Alcazaba recibe
      $\mathbf{300\,000}$ visitas/año
      y tiene $\mathbf{43\,000\ \text{m}^2}$.»
    \end{minipage}%
    \hfill%
    \begin{minipage}[t]{0.48\textwidth}
      \incorrecto\ \textbf{Error frecuente}\\[4pt]
      Pensar que las matemáticas son
      \textbf{solo cálculo}:\par
      «Las matemáticas no tienen nada que ver
      con describir ciudades.»
    \end{minipage}
    % -------------------------------------------------------
  \end{errorbox}
  \begin{notabox}[Recuerda]
    
    Todo lo que se puede \textbf{medir, contar, localizar
    o comparar} en una ciudad es material matemático.
  \end{notabox}
\end{frame}

% ---------------------------------------------------------------
%  RESUMEN
% ---------------------------------------------------------------
\begin{frame}{Resumen de la sesión 01}
  \begin{resumenbox}
    \begin{itemize}
      \item Una \textbf{guía turística matemática} combina
            descripciones con datos numéricos, medidas y
            representaciones cuantitativas.
      \item Los datos pueden ser \textbf{descriptivos}
            (palabras) o \textbf{cuantitativos} (números).
      \item Cualquier monumento o espacio urbano esconde
            información matemática si sabemos \textbf{hacer
            las preguntas correctas}.
      \item Hemos aplicado la fórmula
            $t = \dfrac{d}{v}$ al Cable Inglés.
    \end{itemize}
  \end{resumenbox}
  
  \begin{center}
    {\color{AlmAzulM}
    \faArrowCircleRight\enskip
    \textbf{Próxima sesión:}
    ¿Qué se puede medir y representar en una ciudad?}
  \end{center}
\end{frame}

\end{document}



% % ================================================================
%  sesion02.tex  —  Sesión 2 de 20
%  ¿Qué se puede medir y representar en una ciudad?
%  Almería en Cifras y Formas · Matemáticas 1.º ESO
%  Compilar con:  pdflatex sesion02.tex  (dos pasadas)
% ================================================================
\documentclass[aspectratio=169, 10pt, spanish]{beamer}
\usepackage{almeria-beamer}

% ---------------------------------------------------------------
%  DATOS DE LA SESIÓN
% ---------------------------------------------------------------
\renewcommand{\SessionNumber}{02}
\renewcommand{\SessionTitle}{¿Qué se puede medir y representar en una ciudad?}
\renewcommand{\SessionName}{Sesión 02}
\renewcommand{\SessionObjective}{%
  Clasificar los tipos de información matemática presentes en una ciudad
  (ubicación, medida, dato) e identificar ejemplos reales en Almería.}
\renewcommand{\BlockName}{Bloque 1: Mapa y coordenadas}

% ================================================================
\begin{document}

% ---------------------------------------------------------------
%  PORTADA
% ---------------------------------------------------------------
\begin{frame}[plain]
  \titlepage
\end{frame}

% ---------------------------------------------------------------
%  ÍNDICE DE LA SESIÓN
% ---------------------------------------------------------------
\begin{frame}{Índice de la sesión}
  \begin{columns}[t]
    \begin{column}{0.5\textwidth}
      \begin{enumerate}
        \item Tres tipos de información matemática
        \item Magnitudes y unidades
        \item Tabla de conversión de unidades
        \item Ejemplos reales de Almería
        \item Conversiones con la Rambla y el Puerto
      \end{enumerate}
    \end{column}
    \begin{column}{0.5\textwidth}
      \begin{enumerate}
        \setcounter{enumi}{5}
        \item Ejercicios multinivel (Bloom)
        \item Evidencia del día
        \item Error frecuente
        \item Resumen
      \end{enumerate}
    \end{column}
  \end{columns}
\end{frame}

% ---------------------------------------------------------------
\seccionframe{1. Tres tipos de información matemática}
% ---------------------------------------------------------------

\begin{frame}{¿Qué información matemática hay en una ciudad?}
  \begin{columns}[c]
    \begin{column}{0.52\textwidth}
      \begin{teoriabox}
        Toda la información matemática de una ciudad
        puede clasificarse en \textbf{tres grandes tipos}:
        \begin{itemize}
          \item \textbf{Ubicación:} ¿Dónde está?
          \item \textbf{Medida:} ¿Cuánto mide/pesa/dura?
          \item \textbf{Dato:} ¿Cuántos/cuántas en un momento?
        \end{itemize}
      \end{teoriabox}
      \vspace{6pt}
      Reconocer el tipo correcto nos ayuda a elegir
      la \textbf{herramienta matemática} adecuada.
    \end{column}
    \begin{column}{0.44\textwidth}
      \begin{center}
        \begin{tikzpicture}[scale=0.88,
            nodo/.style={rounded corners=5pt, minimum width=2.5cm,
                         minimum height=0.7cm, align=center,
                         font=\AlmFontSmall\bfseries}]
          % Nodo raíz
          \node[nodo, fill=AlmAzul, text=AlmBlanco]
            (inf) at (0,2.8) {Información\\matemática};
          % Tres hijos
          \node[nodo, fill=AlmAzulM, text=AlmBlanco]
            (ubi) at (-2.6,1.2) {Ubicación};
          \node[nodo, fill=AlmAmbar!80!black, text=AlmBlanco]
            (med) at (0,1.2)    {Medida};
          \node[nodo, fill=AlmVerde, text=AlmBlanco]
            (dat) at (2.6,1.2)  {Dato};
          % Flechas
          \draw[AlmAmbar, line width=1.4pt, ->]
            (inf) -- (ubi);
          \draw[AlmAmbar, line width=1.4pt, ->]
            (inf) -- (med);
          \draw[AlmAmbar, line width=1.4pt, ->]
            (inf) -- (dat);
          % Ejemplos pequeños
          \node[font=\AlmFontTiny, color=AlmGris, align=center]
            at (-2.6,0.45) {coordenadas\\referencias};
          \node[font=\AlmFontTiny, color=AlmGris, align=center]
            at (0,0.45)    {longitud\\superficie\\masa};
          \node[font=\AlmFontTiny, color=AlmGris, align=center]
            at (2.6,0.45)  {población\\temperatura\\turistas};
        \end{tikzpicture}
      \end{center}
    \end{column}
  \end{columns}
\end{frame}

% -------
\begin{frame}{Los tres tipos en detalle}
  \begin{columns}[T]
    \begin{column}{0.32\textwidth}
      \begin{definicionbox}{Ubicación}
        \AlmFontFootnote
        Indica \textbf{dónde} se encuentra algo
        mediante coordenadas o referencias espaciales.\\[4pt]
        Ejemplo en Almería:\\
        \datoalm{Alcazaba}{36,8409°\,N\;/\;2,4699°\,O}
      \end{definicionbox}
    \end{column}
    \begin{column}{0.32\textwidth}
      \begin{definicionbox}{Medida}
        \AlmFontFootnote
        Expresa una \textbf{magnitud física} con un valor
        \emph{fijo} (no cambia con el tiempo).\\[4pt]
        Ejemplo en Almería:\\
        \datoalm{Rambla}{370 m de longitud}\\
        \datoalm{Puerto}{35 ha de superficie}
      \end{definicionbox}
    \end{column}
    \begin{column}{0.32\textwidth}
      \begin{definicionbox}{Dato}
        \AlmFontFootnote
        Información numérica \textbf{variable}: puede
        cambiar según el momento o el año.\\[4pt]
        Ejemplo en Almería:\\
        \datoalm{Temperatura}{máx. 38 °C en agosto}\\
        \datoalm{Turistas}{5,8 millones en 2023}
      \end{definicionbox}
    \end{column}
  \end{columns}
  \vspace{8pt}
  \begin{notabox}[Clave para distinguirlos]
    \AlmFontSmall
    \correcto\ La \textbf{medida} no cambia: la longitud de la Rambla
    siempre es 370 m.\\
    \correcto\ El \textbf{dato} cambia: la temperatura de mañana puede ser distinta.
  \end{notabox}
\end{frame}

% ---------------------------------------------------------------
\seccionframe{2. Magnitudes y unidades}
% ---------------------------------------------------------------

\begin{frame}{Magnitudes físicas y sus unidades}
  \begin{columns}[c]
    \begin{column}{0.52\textwidth}
      \begin{teoriabox}
        Una \textbf{magnitud} es todo aquello que se puede
        medir. Toda medida consta de:
        \[
          \text{medida} = \underbrace{\text{número}}_{\text{valor}}
            \;\times\;
          \underbrace{\text{unidad}}_{\text{patrón}}
        \]
        Ejemplo: $370\ \unidad{m}$ — el número es 370
        y la unidad es el metro (\unidad{m}).
      \end{teoriabox}
      \vspace{6pt}
      \begin{notabox}[Regla de oro]
        \AlmFontSmall
        Nunca escribas un resultado \textbf{sin unidades}.\\
        «370» no significa nada; «$370\ \text{m}$» sí.
      \end{notabox}
    \end{column}
    \begin{column}{0.44\textwidth}
      \vspace{4pt}
      {\AlmFontSmall
      \begin{tabular}{@{}lll@{}}
        \toprule
        \textbf{Magnitud} & \textbf{Unidad base} & \textbf{Símbolo}\\
        \midrule
        Longitud     & metro        & \unidad{m}\\
        Superficie   & metro cuadrado & \unidad{m\textsuperscript{2}}\\
        Masa         & kilogramo    & \unidad{kg}\\
        Tiempo       & segundo      & \unidad{s}\\
        Temperatura  & grado Celsius & \unidad{°C}\\
        \bottomrule
      \end{tabular}}
      \vspace{8pt}
      \begin{ejemplobox}
        \AlmFontFootnote
        La \textbf{Catedral de Almería} tiene
        55 \unidad{m} de altura. Su torre mide
        un \textbf{tiempo} de sombra diferente
        según la estación del año.
      \end{ejemplobox}
    \end{column}
  \end{columns}
\end{frame}

% ---------------------------------------------------------------
\seccionframe{3. Tabla de conversión de unidades}
% ---------------------------------------------------------------

\begin{frame}{Conversiones de unidades más usadas}
  \begin{columns}[c]
    \begin{column}{0.56\textwidth}
      \begin{center}
        {\AlmFontSmall
        \begin{tabular}{@{}llll@{}}
          \toprule
          \textbf{Magnitud} & \textbf{Unidad} & \textbf{Equivalencia}\\
          \midrule
          \multirow{3}{*}{Longitud}
            & \unidad{km}  & $1\ \text{km} = 1\,000\ \text{m}$\\
            & \unidad{m}   & $1\ \text{m} = 100\ \text{cm}$\\
            & \unidad{cm}  & $1\ \text{cm} = 10\ \text{mm}$\\[4pt]
          \multirow{3}{*}{Superficie}
            & \unidad{km\textsuperscript{2}}
                           & $1\ \text{km}^2 = 1\,000\,000\ \text{m}^2$\\
            & \unidad{ha}  & $1\ \text{ha} = 10\,000\ \text{m}^2$\\
            & \unidad{m\textsuperscript{2}}
                           & unidad base\\[4pt]
          \multirow{2}{*}{Masa}
            & \unidad{t}   & $1\ \text{t} = 1\,000\ \text{kg}$\\
            & \unidad{g}   & $1\ \text{kg} = 1\,000\ \text{g}$\\[4pt]
          \multirow{2}{*}{Tiempo}
            & \unidad{h}   & $1\ \text{h} = 3\,600\ \text{s}$\\
            & \unidad{min} & $1\ \text{min} = 60\ \text{s}$\\
          \bottomrule
        \end{tabular}}
      \end{center}
    \end{column}
    \begin{column}{0.40\textwidth}
      \begin{formulabox}[AlmAzulM]
        \AlmFontSmall
        \[
          1\ \text{ha} = 10\,000\ \text{m}^2
        \]
        \[
          1\ \text{km} = 1\,000\ \text{m}
        \]
        \[
          1\ \text{km}^2 = 100\ \text{ha}
        \]
      \end{formulabox}
      \vspace{6pt}
      \begin{tikzpicture}[scale=0.80]
        % Escalera de conversión de longitud
        \foreach \i/\lbl/\col in {%
          0/{mm}/AlmGris,
          1/{cm}/AlmAzulC,
          2/{dm}/AlmAzulC,
          3/{m}/AlmAzulM,
          4/{dam}/AlmAzulM,
          5/{hm}/AlmAzulM,
          6/{km}/AlmAzul}{
          \fill[\col, rounded corners=2pt]
            (\i*0.82, \i*0.38)
            rectangle (\i*0.82+2.6, \i*0.38+0.34);
          \node[font=\AlmFontTiny\bfseries, text=white]
            at (\i*0.82+1.3, \i*0.38+0.17) {\lbl};
        }
        \node[font=\AlmFontTiny, color=AlmGris] at (1.3,-0.28)
          {$\times 10$\quad $\times 100$\quad $\times 1000$};
      \end{tikzpicture}
    \end{column}
  \end{columns}
\end{frame}

% ---------------------------------------------------------------
\seccionframe{4 y 5. Ejemplos reales de Almería}
% ---------------------------------------------------------------

\begin{frame}{La Rambla de Almería — Longitud y anchura}
  \framesubtitle{Magnitud: Longitud · Unidades: m, km}
  \begin{columns}[c]
    \begin{column}{0.52\textwidth}
      \begin{ejemplobox}
        \AlmFontSmall
        La \textbf{Rambla de Almería} es el paseo central
        de la ciudad.\\[6pt]
        \datoalm{Longitud}{370 m}\\
        \datoalm{Anchura máxima}{50 m}\\[6pt]
        Ambas medidas expresan la magnitud
        \textbf{longitud} en metros (\unidad{m}).
      \end{ejemplobox}
      \vspace{6pt}
      \begin{pasobox}[1]
        \AlmFontFootnote
        Convertir longitud a kilómetros:
        \[
          370\ \text{m} \div 1\,000 = \mathbf{0{,}37\ \text{km}}
        \]
      \end{pasobox}
    \end{column}
    \begin{column}{0.44\textwidth}
    
          \begin{pasobox}[2]
            \AlmFontFootnote
            Estimar área (rectángulo aproximado):
            \[
              A \approx 370 \times 50 = \mathbf{18\,500\ \text{m}^2}
            \]
          \end{pasobox}
      \begin{center}
        \begin{tikzpicture}[scale=0.82]
          % Representación esquemática de la Rambla
          \fill[AlmAzulC!20, rounded corners=3pt]
            (0,0) rectangle (5.2,1.4);
          \fill[AlmVerde!35, rounded corners=2pt]
            (0.12,0.12) rectangle (5.08,1.28);
          % Acotación longitud
          \draw[acota] (0,-0.28) -- (5.2,-0.28);
          \node[font=\AlmFontTiny\bfseries, color=AlmAzulM]
            at (2.6,-0.52) {370 m};
          % Acotación anchura
          \draw[acota] (5.38,0) -- (5.38,1.4);
          \node[font=\AlmFontTiny\bfseries, color=AlmAzulM,
                rotate=90, anchor=south]
            at (5.62,0.7) {50 m};
          % Etiqueta
          \node[font=\AlmFontScript\bfseries, color=AlmAzul]
            at (2.6,0.7) {Rambla de Almería};
          % Flecha norte
          \node[font=\AlmFontTiny, color=AlmGris] at (0.35,1.15)
            {N$\uparrow$};
        \end{tikzpicture}
      \end{center}
      \vspace{4pt}
      \begin{notabox}[Tipo de información]
        \AlmFontFootnote
        \textbf{Longitud} (370 m) $\rightarrow$ \emph{Medida}\\
        \textbf{Anchura} (50 m) $\rightarrow$ \emph{Medida}\\
        \textbf{Visitantes diarios} $\rightarrow$ \emph{Dato}
      \end{notabox}
    \end{column}
  \end{columns}
\end{frame}

\begin{frame}{El Puerto de Almería — Conversión de superficie}
  \framesubtitle{Magnitud: Superficie · 1 ha = 10\,000 m²}
  \begin{columns}[c]
    \begin{column}{0.50\textwidth}
      \begin{ejemplobox}
        \AlmFontSmall
        El \textbf{Puerto de Almería} ocupa
        \textbf{35 ha} de superficie.\\[6pt]
        ¿Cuántos metros cuadrados son?\\
        ¿Cuántos km\textsuperscript{2} son?
      \end{ejemplobox}
      \vspace{6pt}
      \begin{pasobox}[1]
        \AlmFontFootnote
        Convertir a \textbf{m\textsuperscript{2}}:
        \[
          35\ \text{ha} \times 10\,000 = \mathbf{350\,000\ \text{m}^2}
        \]
      \end{pasobox}
      \vspace{4pt}
      \begin{pasobox}[2]
        \AlmFontFootnote
        Convertir a \textbf{km\textsuperscript{2}}:
        \[
          350\,000\ \text{m}^2 \div 1\,000\,000 = \mathbf{0{,}35\ \text{km}^2}
        \]
      \end{pasobox}
    \end{column}
    \begin{column}{0.46\textwidth}
      \begin{center}
        \begin{tikzpicture}[scale=0.78]
          % Cuadrado grande = Puerto (35 ha)
          \fill[AlmAzulC!30, rounded corners=4pt]
            (0,0) rectangle (4.5,4.5);
          \node[font=\AlmFontScript\bfseries, color=AlmAzul]
            at (2.25,4.1) {Puerto de Almería};
          \node[font=\AlmFontSmall\bfseries, color=AlmAzulM]
            at (2.25,2.4) {35 ha};
          \node[font=\AlmFontFootnote, color=AlmGris]
            at (2.25,1.85) {$= 350\,000\ \text{m}^2$};
          \node[font=\AlmFontFootnote, color=AlmGris]
            at (2.25,1.3)  {$= 0{,}35\ \text{km}^2$};
          % Cuadradito = campo de fútbol ~0,7 ha
          \fill[AlmVerde!50, rounded corners=2pt]
            (0.15,0.15) rectangle (1.35,1.0);
          \node[font=\AlmFontTiny, color=AlmVerde!70!black]
            at (0.75,0.57) {\(\approx\)campo};
          % Acotación
          \draw[acota] (0,-0.28) -- (4.5,-0.28);
          \node[font=\AlmFontTiny, color=AlmGris] at (2.25,-0.52)
            {escala aprox.};
        \end{tikzpicture}
      \end{center}
      \vspace{2pt}
      \begin{notabox}[Curiosidad]
        \AlmFontFootnote
        35 ha equivalen a unos
        \textbf{49 campos de fútbol}.
      \end{notabox}
    \end{column}
  \end{columns}
\end{frame}

% ---------------------------------------------------------------
\seccionframe{6. Ejercicios multinivel — Bloom}
% ---------------------------------------------------------------

\begin{frame}{Ejercicio RECORDAR}
  \bloomtag{RECORDAR}\quad\dificultad{1}
  \vspace{4pt}
  \begin{recordarbox}
    \AlmFontSmall
    \textbf{Tarea:} Completa la siguiente tabla de memoria
    (o con tus apuntes).\\[6pt]
    \begin{tabular}{@{}lll@{}}
      \toprule
      \textbf{Magnitud} & \textbf{Unidad base} & \textbf{Ejemplo en Almería}\\
      \midrule
      Longitud     & \ldots\ldots & \ldots\ldots\ldots\ldots\\[4pt]
      Superficie   & \ldots\ldots & \ldots\ldots\ldots\ldots\\[4pt]
      Masa         & \ldots\ldots & \ldots\ldots\ldots\ldots\\[4pt]
      Tiempo       & \ldots\ldots & \ldots\ldots\ldots\ldots\\[4pt]
      Temperatura  & \ldots\ldots & \ldots\ldots\ldots\ldots\\
      \bottomrule
    \end{tabular}
    \\[6pt]
    Relaciona cada magnitud con un lugar o elemento real de Almería.
  \end{recordarbox}
\end{frame}

\begin{frame}{Ejercicio COMPRENDER}
  \bloomtag{COMPRENDER}\quad\dificultad{2}
  \vspace{4pt}
  \begin{comprenderbox}
    \AlmFontSmall
    La \textbf{playa de El Zapillo} mide 2,5 km de longitud.\\[6pt]
    Responde razonando cada respuesta:
    \begin{enumerate}
      \item ¿Qué \textbf{magnitud} representa esa medida?
            ¿A qué tipo de información pertenece
            (ubicación / medida / dato)?
      \item Convierte 2,5 km a \textbf{metros}.
            Muestra la operación completa.
      \item Para calcular la \textbf{superficie} de la playa,
            ¿qué otra medida necesitarías?
            ¿Qué operación harías?
    \end{enumerate}
  \end{comprenderbox}
  \vspace{4pt}
  \begin{notabox}[Pista]
    \AlmFontFootnote
    Longitud $\times$ anchura $=$ superficie (si es rectangular).
  \end{notabox}
\end{frame}

\begin{frame}{Ejercicio APLICAR}
  \bloomtag{APLICAR}\quad\dificultad{3}
  \vspace{4pt}
  \begin{aplicarbox}
    \AlmFontFootnote
    Clasifica cada dato de la tabla como \textbf{Ubicación},
    \textbf{Medida} o \textbf{Dato}. Convierte las unidades
    cuando se indique.\\[4pt]
    \begin{tabular}{@{}p{4.2cm}lp{2.4cm}@{}}
      \toprule
      \textbf{Información} & \textbf{Tipo} & \textbf{Conversión}\\
      \midrule
      Alcazaba: 36,84°\,N / 2,47°\,O
        & \ldots & —\\[3pt]
      Catedral: 55 m de altura
        & \ldots & \ldots en km\\[3pt]
      Temperatura máx. agosto: 38 °C
        & \ldots & —\\[3pt]
      Puerto: 35 ha
        & \ldots & \ldots en m\textsuperscript{2}\\[3pt]
      Turistas 2023: 5\,800\,000
        & \ldots & —\\[3pt]
      Paseo Marítimo: 3 km
        & \ldots & \ldots en m\\
      \bottomrule
    \end{tabular}
  \end{aplicarbox}
\end{frame}

% ---------------------------------------------------------------
%  EVIDENCIA
% ---------------------------------------------------------------
\begin{frame}{Evidencia del día}
  \begin{evidenciabox}
    Entrega tu \textbf{Tabla Clasificada} con al menos
    \textbf{5 elementos reales de Almería},
    una fila por elemento:
    \vspace{6pt}
    \begin{center}
      \AlmFontSmall
      \begin{tabular}{@{}lll@{}}
        \toprule
        \textbf{Tipo (Ubicación / Medida / Dato)}
          & \textbf{Ejemplo de Almería}
          & \textbf{Unidad}\\
        \midrule
        Medida & Longitud de la Rambla: 370 m & \unidad{m}\\
        \ldots & \ldots & \ldots\\
        \bottomrule
      \end{tabular}
    \end{center}
    \vspace{6pt}
    \begin{itemize}
      \item Incluye al menos \textbf{una conversión de unidades}
            con los pasos completos.
      \item Justifica en una frase por qué cada elemento
            es Ubicación, Medida o Dato.
    \end{itemize}
  \end{evidenciabox}
\end{frame}

% ---------------------------------------------------------------
%  ERROR FRECUENTE
% ---------------------------------------------------------------
\begin{frame}{¡Cuidado con este error!}
  \begin{errorbox}
    \begin{columns}[c]
      \begin{column}{0.50\textwidth}
        \incorrecto\ \textbf{Error frecuente}\\[4pt]
        Confundir un \textbf{dato} (varía con el tiempo)
        con una \textbf{medida} (valor fijo).\\[6pt]
        «La temperatura de Almería es
        \cancel{una medida fija}» — \textbf{incorrecto}.\\
        La temperatura cambia: es un \textbf{dato}.
      \end{column}
      \begin{column}{0.46\textwidth}
        \correcto\ \textbf{Correcto}\\[4pt]
        La \textbf{altura} de la Catedral (55 m)
        \emph{no cambia con el tiempo}:
        es una \textbf{medida}.\\[6pt]
        La \textbf{temperatura} máxima de un día
        \emph{sí cambia} cada día:
        es un \textbf{dato}.
      \end{column}
    \end{columns}
  \end{errorbox}
  \vspace{8pt}
  \begin{notabox}[Regla rápida]
    \AlmFontSmall
    ¿Cambiaría si volvemos a medirlo mañana?
    \correcto\;Sí $\Rightarrow$ \textbf{dato}.\quad
    \correcto\;No $\Rightarrow$ \textbf{medida}.
  \end{notabox}
\end{frame}

% ---------------------------------------------------------------
%  RESUMEN
% ---------------------------------------------------------------
\begin{frame}{Resumen de la sesión 02}
  \begin{resumenbox}
    \begin{itemize}
      \item La información matemática de una ciudad se clasifica
            en \textbf{Ubicación}, \textbf{Medida} y \textbf{Dato}.
      \item Toda medida tiene \textbf{valor numérico} y \textbf{unidad};
            escribir la unidad es obligatorio.
      \item Las magnitudes más frecuentes son longitud (\unidad{m}, \unidad{km}),
            superficie (\unidad{m\textsuperscript{2}}, \unidad{ha}),
            masa (\unidad{kg}), tiempo (\unidad{s}, \unidad{h}) y
            temperatura (\unidad{°C}).
      \item Conversiones clave:
            $1\ \text{km}=1\,000\ \text{m}$;
            $1\ \text{ha}=10\,000\ \text{m}^2$.
      \item Un dato \textbf{cambia} con el tiempo;
            una medida \textbf{no cambia}.
    \end{itemize}
  \end{resumenbox}
  \vspace{6pt}
  \begin{center}
    {\AlmFontSmall\color{AlmAzulM}
    \faArrowCircleRight\enskip
    \textbf{Próxima sesión:}
    Plano por cuadrantes — localizar lugares con letra y número.}
  \end{center}
\end{frame}

\end{document}




% % ================================================================
%  sesion03.tex  —  Sesión 3 de 20
%  Plano por cuadrantes
%  Almería en Cifras y Formas · Matemáticas 1.º ESO
%  Compilar con:  pdflatex sesion03.tex  (dos pasadas)
% ================================================================
\documentclass[aspectratio=169, 12pt, spanish]{beamer}
\usepackage{almeria-beamer}

% ---------------------------------------------------------------
%  DATOS DE LA SESIÓN
% ---------------------------------------------------------------
\renewcommand{\SessionNumber}{03}
\renewcommand{\SessionTitle}{Plano por cuadrantes}
\renewcommand{\SessionName}{Sesión 03}
\renewcommand{\SessionObjective}{%
  Localizar puntos de interés en un mapa cuadriculado de Almería usando
  un sistema de cuadrantes (letra + número) y compararlo con el sistema
  de coordenadas cartesianas.}
\renewcommand{\BlockName}{Bloque 1: Mapa y coordenadas}

% ================================================================
\begin{document}

% ---------------------------------------------------------------
%  PORTADA
% ---------------------------------------------------------------
\begin{frame}[plain]
  \titlepage
\end{frame}

% ---------------------------------------------------------------
%  ÍNDICE DE LA SESIÓN
% ---------------------------------------------------------------
\begin{frame}{Índice de la sesión}
  \begin{columns}[t]
    \begin{column}{0.5\textwidth}
      \begin{enumerate}
        \item ¿Qué es un plano cuadriculado?
        \item Cómo leer un código de cuadrante
        \item Mapa de Almería: localizar puntos
        \item Pasos para encontrar un lugar
      \end{enumerate}
    \end{column}
    \begin{column}{0.5\textwidth}
      \begin{enumerate}
        \setcounter{enumi}{4}
        \item Comparación: cuadrantes vs.\ coordenadas vs.\ GPS
        \item Ejercicios multinivel (Bloom)
        \item Evidencia del día
        \item Error frecuente
        \item Resumen
      \end{enumerate}
    \end{column}
  \end{columns}
\end{frame}

% ---------------------------------------------------------------
\seccionframe{1. ¿Qué es un plano cuadriculado?}
% ---------------------------------------------------------------

\begin{frame}{El plano cuadriculado}
  \begin{columns}[c]
    \begin{column}{0.52\textwidth}
      \begin{definicionbox}{Plano cuadriculado}
        Un mapa dividido en \textbf{filas} (letras A, B, C\ldots)
        y \textbf{columnas} (números 1, 2, 3\ldots) que forma una
        rejilla. Cada celda recibe un código único
        \textbf{letra + número}.
      \end{definicionbox}
      \vspace{6pt}
      \begin{itemize}
        \item Fácil de usar: no necesita tecnología.
        \item Estándar en guías de calles y mapas turísticos.
        \item La letra identifica la \textbf{fila horizontal}.
        \item El número identifica la \textbf{columna vertical}.
      \end{itemize}
      \vspace{4pt}
      \begin{notabox}[Regla de lectura]
        \small
        \textbf{Primero la letra} (fila),
        \textbf{después el número} (columna):
        \quad B3 = fila B, columna 3.
      \end{notabox}
    \end{column}
    \begin{column}{0.44\textwidth}
      \begin{center}
        \begin{tikzpicture}[scale=0.85]
          % Cuadrícula 4x4 básica explicativa
          \foreach \c in {1,2,3,4}{
            \foreach \f in {1,2,3,4}{
              \fill[AlmFondo, draw=AlmAzulC!50, line width=0.5pt]
                (\c-1, \f-1) rectangle (\c, \f);
            }
          }
          % Etiquetas de filas (letras)
          \foreach \f/\ltr in {4/A, 3/B, 2/C, 1/D}{
            \node[font=\small\bfseries, color=AlmAzul]
              at (-0.35, \f-0.5) {\ltr};
          }
          % Etiquetas de columnas (números)
          \foreach \c/\num in {1/1, 2/2, 3/3, 4/4}{
            \node[font=\small\bfseries, color=AlmAzul]
              at (\c-0.5, 4.35) {\num};
          }
          % Celda destacada B3
          \fill[AlmAmbar!45, draw=AlmAmbar, line width=1.5pt]
            (2, 2) rectangle (3, 3);
          \node[font=\small\bfseries, color=AlmAzul]
            at (2.5, 2.5) {B3};
          % Flechas indicadoras
          \draw[AlmTerra, ->, line width=1pt]
            (-0.35, 2.5) -- (0.05, 2.5);
          \draw[AlmAzulM, ->, line width=1pt]
            (2.5, 4.35) -- (2.5, 3.05);
        \end{tikzpicture}
      \end{center}
    \end{column}
  \end{columns}
\end{frame}

% ---------------------------------------------------------------
\seccionframe{2. Cómo leer un código de cuadrante}
% ---------------------------------------------------------------

\begin{frame}{Lectura de un código: paso a paso}
  \begin{columns}[c]
    \begin{column}{0.50\textwidth}
      \begin{pasobox}[1]
        \small
        \textbf{Lee la letra.}\\
        Identifica la \textbf{fila horizontal}
        en el margen izquierdo del mapa.
      \end{pasobox}
      \vspace{6pt}
      \begin{pasobox}[2]
        \small
        \textbf{Lee el número.}\\
        Identifica la \textbf{columna vertical}
        en el margen superior del mapa.
      \end{pasobox}
      \vspace{6pt}
      \begin{pasobox}[3]
        \small
        \textbf{Marca la intersección.}\\
        El lugar está \emph{dentro de la celda}
        donde se cruzan esa fila y esa columna.
      \end{pasobox}
    \end{column}
    \begin{column}{0.46\textwidth}
      \begin{ejemplobox}
        \small
        \textbf{Código: C4}\\[4pt]
        \begin{enumerate}
          \item Busco la fila \textbf{C}
                (tercera fila desde arriba).
          \item Busco la columna \textbf{4}
                (cuarta columna desde la izquierda).
          \item Marco donde se cruzan: ahí está el lugar.
        \end{enumerate}
      \end{ejemplobox}
      \vspace{6pt}
      \begin{formulabox}[AlmAzulM]
        \small
        Código \;=\; \textbf{letra} + \textbf{número}\\[4pt]
        \large B3 \;=\; fila\,B\;,\;columna\,3
      \end{formulabox}
    \end{column}
  \end{columns}
\end{frame}

% ---------------------------------------------------------------
\seccionframe{3. Mapa de Almería: localizar puntos}
% ---------------------------------------------------------------

\begin{frame}{Mapa cuadriculado de Almería}
  \framesubtitle{Cuadrícula 4 × 5 — puntos de interés reales}
  \begin{center}
    \begin{tikzpicture}[scale=1.05]
      % ------- Fondo del mapa -------
      \fill[AlmAzulC!12, rounded corners=4pt]
        (-0.5,-0.5) rectangle (5.5,4.5);
      % ------- Cuadrícula 4 filas × 5 columnas -------
      \foreach \c in {0,...,4}{
        \foreach \f in {0,...,3}{
          \fill[AlmBlanco!70, draw=AlmAzulC!60, line width=0.5pt]
            (\c,\f) rectangle (\c+1,\f+1);
        }
      }
      % ------- Etiquetas filas (A=top=fila 3, D=bottom=fila 0) -------
      \foreach \f/\ltr in {3/A, 2/B, 1/C, 0/D}{
        \node[font=\small\bfseries, color=AlmAzul]
          at (-0.32, \f+0.5) {\ltr};
      }
      % ------- Etiquetas columnas -------
      \foreach \c in {1,...,5}{
        \node[font=\small\bfseries, color=AlmAzul]
          at (\c-0.5, 4.28) {\c};
      }
      % ------- Puntos de interés -------
      % Plaza Catedral: B2 -> col 2, fila B(=2) -> (1.5, 2.5)
      \node[punto] (cat) at (1.5, 2.5) {};
      \node[font=\tiny\bfseries, color=AlmAzulM, above right=-1pt and 1pt of cat]
        {Catedral};
      \node[font=\tiny, color=AlmGris, below right=-1pt and 1pt of cat]
        {B2};
      % Alcazaba: A2 -> col 2, fila A(=3) -> (1.5, 3.5)
      \node[puntot] (alc) at (1.5, 3.5) {};
      \node[font=\tiny\bfseries, color=AlmTerra, above right=-1pt and 1pt of alc]
        {Alcazaba};
      \node[font=\tiny, color=AlmGris, below right=-1pt and 1pt of alc]
        {A2};
      % Mercado Central: B3 -> col 3, fila B(=2) -> (2.5, 2.5)
      \node[puntov] (mer) at (2.5, 2.5) {};
      \node[font=\tiny\bfseries, color=AlmVerde, above right=-1pt and 1pt of mer]
        {Mercado};
      \node[font=\tiny, color=AlmGris, below right=-1pt and 1pt of mer]
        {B3};
      % Puerto: C4 -> col 4, fila C(=1) -> (3.5, 1.5)
      \node[punto, fill=BAplicar] (puer) at (3.5, 1.5) {};
      \node[font=\tiny\bfseries, color=BAplicar!70!black,
            above right=-1pt and 1pt of puer]
        {Puerto};
      \node[font=\tiny, color=AlmGris, below right=-1pt and 1pt of puer]
        {C4};
      % Parque Nicolás Salmerón: C3 -> col 3, fila C(=1) -> (2.5, 1.5)
      \node[puntov] (par) at (2.5, 1.5) {};
      \node[font=\tiny\bfseries, color=AlmVerde,
            above left=-1pt and 1pt of par]
        {P. Salmerón};
      \node[font=\tiny, color=AlmGris, below left=-1pt and 1pt of par]
        {C3};
      % ------- Leyenda -------
      \node[font=\tiny, color=AlmAzul, align=left,
            draw=AlmAzulC!50, fill=AlmBlanco!80,
            rounded corners=2pt, inner sep=3pt]
        at (4.35, 3.5)
        {\textbf{Leyenda}\\
         \textcolor{AlmAzulM}{$\bullet$} Catedral\\
         \textcolor{AlmTerra}{$\bullet$} Alcazaba\\
         \textcolor{AlmVerde}{$\bullet$} Mercado\\
         \textcolor{BAplicar}{$\bullet$} Puerto\\
         \textcolor{AlmVerde}{$\bullet$} P.\,Salmerón};
    \end{tikzpicture}
  \end{center}
\end{frame}

% ---------------------------------------------------------------
\seccionframe{4. Pasos para encontrar un lugar}
% ---------------------------------------------------------------

\begin{frame}{Localización paso a paso en el mapa}
  \begin{columns}[c]
    \begin{column}{0.50\textwidth}
      \begin{teoriabox}
        \small
        Para localizar un lugar en un plano cuadriculado:
        \begin{enumerate}
          \item Busca la \textbf{letra} en el margen izquierdo
                y traza mentalmente una línea horizontal.
          \item Busca el \textbf{número} en el margen superior
                y traza una línea vertical.
          \item La celda de la \textbf{intersección}
                contiene el lugar buscado.
          \item Coloca un \textbf{punto} en el centro
                de esa celda.
        \end{enumerate}
      \end{teoriabox}
    \end{column}
    \begin{column}{0.46\textwidth}
      \begin{ejemplobox}
        \small
        \textbf{Localizar: Puerto (C4)}\\[4pt]
        \begin{enumerate}
          \item Fila \textbf{C} $\rightarrow$ tercera fila.
          \item Columna \textbf{4} $\rightarrow$ cuarta columna.
          \item Intersección $\rightarrow$ celda C4.
          \item Marca el punto \textcolor{BAplicar}{\faMapMarkerAlt}.
        \end{enumerate}
      \end{ejemplobox}
      \vspace{6pt}
      \begin{notabox}[El código siempre es único]
        \footnotesize
        Dos lugares distintos \textbf{no pueden}
        tener el mismo código.\\
        Dos códigos distintos \textbf{nunca} señalan
        al mismo lugar.
      \end{notabox}
    \end{column}
  \end{columns}
\end{frame}

% ---------------------------------------------------------------
\seccionframe{5. Comparación de sistemas de localización}
% ---------------------------------------------------------------

\begin{frame}{Cuadrantes · Coordenadas cartesianas · GPS}
  \begin{center}
    \small
    \begin{tabular}{@{}p{3.0cm}p{3.4cm}p{3.4cm}p{3.4cm}@{}}
      \toprule
      \textbf{Aspecto}
        & \textbf{Cuadrantes}\newline(letra + número)
        & \textbf{Coordenadas}\newline cartesianas $(x, y)$
        & \textbf{GPS}\newline (latitud, longitud)\\
      \midrule
      \textbf{Precisión}
        & Baja: sitúa en una celda ($\sim$km\textsuperscript{2})
        & Media: punto exacto en el mapa
        & Muy alta: $< 5$ m\\[4pt]
      \textbf{Facilidad}
        & Muy fácil;\newline sin cálculos
        & Media;\newline requiere entender ejes
        & Muy fácil con dispositivo\\[4pt]
      \textbf{Tecnología}
        & Ninguna (papel y lápiz)
        & Ninguna (cuadrícula)
        & Satélite / smartphone\\[4pt]
      \textbf{Uso típico}
        & Guías turísticas,\newline callejeros
        & Geometría, planos de arquitectura
        & Navegación, \newline cartografía digital\\[4pt]
      \textbf{Ejemplo}\newline \textbf{Almería}
        & Catedral: B2
        & Catedral: $(1{,}5;\; 2{,}5)$
        & 36,8368°\,N\newline 2,4695°\,O\\
      \bottomrule
    \end{tabular}
  \end{center}
  \vspace{4pt}
  \begin{notabox}[Conclusión]
    \footnotesize
    Cada sistema tiene su contexto ideal. Para esta SA usaremos
    \textbf{cuadrantes} y \textbf{coordenadas cartesianas}.
  \end{notabox}
\end{frame}

% ---------------------------------------------------------------
\seccionframe{6. Ejercicios multinivel — Bloom}
% ---------------------------------------------------------------

\begin{frame}{Ejercicio RECORDAR}
  \bloomtag{RECORDAR}\quad\dificultad{1}
  \vspace{4pt}
  \begin{recordarbox}
    \small
    En tu cuaderno, \textbf{dibuja} una cuadrícula 4 × 4:
    \begin{itemize}
      \item Rotula las filas con las letras A, B, C, D
            (de arriba abajo) en el margen izquierdo.
      \item Rotula las columnas con los números 1, 2, 3, 4
            (de izquierda a derecha) en el margen superior.
    \end{itemize}
    \vspace{4pt}
    Responde:
    \begin{enumerate}
      \item ¿Cuántas celdas tiene la cuadrícula en total?
      \item Colorea la celda \textbf{A3} de azul
            y la celda \textbf{D2} de rojo.
      \item ¿Qué código tiene la celda de la esquina
            superior derecha?
    \end{enumerate}
  \end{recordarbox}
\end{frame}

\begin{frame}{Ejercicio APLICAR}
  \bloomtag{APLICAR}\quad\dificultad{3}
  \vspace{4pt}
  \begin{aplicarbox}
    \footnotesize
    Usa el mapa de Almería de la diapositiva anterior.\\[4pt]
    \textbf{Parte A — Dibujar desde códigos:}
    \begin{enumerate}
      \item Localiza y marca en el mapa los siguientes lugares:\\
        Rambla de Almería: B2 \quad
        Parque Nicolás Salmerón: C3\\
        Teatro Cervantes: B3 \quad
        Playa de Almería: D4
    \end{enumerate}
    \vspace{4pt}
    \textbf{Parte B — Leer códigos del mapa:}
    \begin{enumerate}
      \setcounter{enumi}{1}
      \item El mapa muestra tres puntos sin nombre.\\
        Escribe el código de cuadrante de cada uno.
    \end{enumerate}
    \vspace{4pt}
    \textbf{Parte C — Distancia aproximada:}
    \begin{enumerate}
      \setcounter{enumi}{2}
      \item Si cada celda representa 1 km, ¿cuántos km
        hay entre la Alcazaba (A2) y el Puerto (C4)?\\
        Explica cómo has calculado esa distancia.
    \end{enumerate}
  \end{aplicarbox}
\end{frame}

\begin{frame}{Ejercicio ANALIZAR}
  \bloomtag{ANALIZAR}\quad\dificultad{4}
  \vspace{4pt}
  \begin{analizarbox}
    \small
    Compara el sistema de \textbf{cuadrantes} con el de
    \textbf{coordenadas cartesianas}.\\[6pt]
    Escribe un argumento de \textbf{5--6 líneas} que responda:
    \begin{itemize}
      \item ¿En qué situaciones es más útil el sistema de cuadrantes?
      \item ¿Cuándo preferirías coordenadas cartesianas?
      \item ¿Cuál de los dos usarías para diseñar un mapa
            turístico impreso de Almería? Justifica tu elección.
    \end{itemize}
    \vspace{4pt}
    Usa \textbf{ejemplos concretos} de Almería en tu argumento.
  \end{analizarbox}
\end{frame}

% ---------------------------------------------------------------
%  EVIDENCIA
% ---------------------------------------------------------------
\begin{frame}{Evidencia del día}
  \begin{evidenciabox}
    Entrega tu \textbf{Mapa base de Almería completado}:
    \begin{itemize}
      \item Cuadrícula de al menos 4 × 4 con líneas rectas
            y perpendiculares.
      \item Ejes correctamente nombrados
            (letras a la izquierda, números arriba).
      \item Al menos \textbf{5 puntos de interés} de Almería
            correctamente situados con su código.
      \item \textbf{Leyenda} con el nombre de cada punto.
      \item Una \textbf{frase de reflexión}:
            «El sistema de cuadrantes me sirve para\ldots»
    \end{itemize}
  \end{evidenciabox}
\end{frame}

% ---------------------------------------------------------------
%  ERROR FRECUENTE
% ---------------------------------------------------------------
\begin{frame}{¡Cuidado con este error!}
  \begin{errorbox}
    \begin{columns}[c]
      \begin{column}{0.50\textwidth}
        \incorrecto\ \textbf{Error frecuente}\\[6pt]
        Leer el código en el \textbf{orden incorrecto}:
        número antes que letra.\\[8pt]
        \begin{center}
          \large
          \textbf{\sout{3B}}~$\neq$~\textbf{B3}
        \end{center}
        \vspace{4pt}
        \footnotesize
        Si dices «3B» estás buscando la columna 3
        y la fila B, lo cual puede interpretarse
        de forma errónea y llevar a una celda
        \textbf{completamente distinta}.
      \end{column}
      \begin{column}{0.46\textwidth}
        \correcto\ \textbf{Correcto}\\[6pt]
        \textbf{Siempre primero la letra, luego el número.}\\[8pt]
        \begin{itemize}
          \item B3 $\rightarrow$ fila B, columna 3 \correcto
          \item 3B $\rightarrow$ orden incorrecto \incorrecto
        \end{itemize}
        \vspace{6pt}
        \footnotesize
        \textbf{Truco:} piensa en el alfabeto antes que en los números.
      \end{column}
    \end{columns}
  \end{errorbox}
  \vspace{6pt}
  \begin{notabox}[Analogía]
    \small
    En una hoja de cálculo, la celda B3 está
    en la columna B y la fila 3 — el mismo orden.
  \end{notabox}
\end{frame}

% ---------------------------------------------------------------
%  RESUMEN
% ---------------------------------------------------------------
\begin{frame}{Resumen de la sesión 03}
  \begin{resumenbox}
    \begin{itemize}
      \item Un \textbf{plano cuadriculado} divide el mapa en filas
            (letras) y columnas (números).
      \item El código de un lugar siempre es
            \textbf{letra primero}, número segundo: B3, C4\ldots
      \item Para localizar: \textbf{1) fila, 2) columna,
            3) intersección}.
      \item Este sistema es sencillo pero tiene
            \textbf{poca precisión} comparado con GPS.
      \item Las \textbf{coordenadas cartesianas} dan más precisión
            y las estudiaremos en la próxima sesión.
    \end{itemize}
  \end{resumenbox}
  \vspace{6pt}
  \begin{center}
    {\small\color{AlmAzulM}
    \faArrowCircleRight\enskip
    \textbf{Próxima sesión:}
    Coordenadas cartesianas — leer y escribir pares $(x, y)$.}
  \end{center}
\end{frame}

\end{document}

% % ================================================================
%  sesion04.tex  —  Sesión 4 de 20
%  Coordenadas y localización
%  Almería en Cifras y Formas · Matemáticas 1.º ESO
%  Compilar con:  pdflatex sesion04.tex  (dos pasadas)
% ================================================================
\documentclass[aspectratio=169, 10pt, spanish]{beamer}
\usepackage{almeria-beamer}

% ---------------------------------------------------------------
%  DATOS DE LA SESIÓN
% ---------------------------------------------------------------
\renewcommand{\SessionNumber}{04}
\renewcommand{\SessionTitle}{Coordenadas y localización}
\renewcommand{\SessionName}{Sesión 04}
\renewcommand{\SessionObjective}{%
  Leer, escribir y localizar lugares exactos en un mapa de Almería
  usando pares de coordenadas cartesianas $(x, y)$.}
\renewcommand{\BlockName}{Bloque 1: Mapa y coordenadas}

% ================================================================
\begin{document}

% ---------------------------------------------------------------
%  PORTADA
% ---------------------------------------------------------------
\begin{frame}[plain]
  \titlepage
\end{frame}

% ---------------------------------------------------------------
%  ÍNDICE DE LA SESIÓN
% ---------------------------------------------------------------
\begin{frame}{Índice de la sesión}
  \begin{columns}[t]
    \begin{column}{0.5\textwidth}
      \begin{enumerate}
        \item El sistema de coordenadas cartesianas
        \item Leer un par $(x, y)$
        \item Los cuatro cuadrantes y sus signos
        \item Representar puntos en el plano
      \end{enumerate}
    \end{column}
    \begin{column}{0.5\textwidth}
      \begin{enumerate}
        \setcounter{enumi}{4}
        \item Escala: de la cuadrícula al mapa real
        \item Ejercicios multinivel (Bloom)
        \item Evidencia del día
        \item Error frecuente
        \item Resumen
      \end{enumerate}
    \end{column}
  \end{columns}
\end{frame}

% ---------------------------------------------------------------
\seccionframe{1. El sistema de coordenadas cartesianas}
% ---------------------------------------------------------------

\begin{frame}{Ejes cartesianos: eje X y eje Y}
  \begin{columns}[c]
    \begin{column}{0.50\textwidth}
      \begin{definicionbox}{Sistema de coordenadas cartesianas}
        Dos rectas perpendiculares que se cortan en el
        \textbf{origen} $(0,0)$:
        \begin{itemize}
          \item \textbf{Eje X} (horizontal): valores positivos
                a la derecha, negativos a la izquierda.
          \item \textbf{Eje Y} (vertical): valores positivos
                hacia arriba, negativos hacia abajo.
        \end{itemize}
      \end{definicionbox}
      \vspace{6pt}
      \begin{teoriabox}
        \AlmFontSmall
        \textbf{Regla mnemotécnica:}\\
        «Primero camina por la \textbf{calle} (eje X),
        luego sube o baja en el \textbf{ascensor} (eje Y).»
      \end{teoriabox}
    \end{column}
    \begin{column}{0.46\textwidth}
      \begin{center}
        \begin{tikzpicture}[scale=0.72]
          % Ejes
          \draw[recto, ->] (-3.4,0) -- (3.6,0)
            node[right, font=\AlmFontSmall\bfseries, color=AlmAzulM] {$x$};
          \draw[recto, ->] (0,-3.4) -- (0,3.6)
            node[above, font=\AlmFontSmall\bfseries, color=AlmAzulM] {$y$};
          % Cuadrícula ligera
          \foreach \i in {-3,-2,-1,1,2,3}{
            \draw[AlmAzulC!25, line width=0.3pt]
              (\i,-3.3) -- (\i,3.3);
            \draw[AlmAzulC!25, line width=0.3pt]
              (-3.3,\i) -- (3.3,\i);
          }
          % Marcas en los ejes
          \foreach \i in {-3,-2,-1,1,2,3}{
            \draw[AlmAzulM, line width=0.6pt]
              (\i,-0.08) -- (\i,0.08);
            \draw[AlmAzulM, line width=0.6pt]
              (-0.08,\i) -- (0.08,\i);
            \node[font=\AlmFontTiny, color=AlmGris] at (\i,-0.28) {\i};
            \node[font=\AlmFontTiny, color=AlmGris] at (-0.28,\i) {\i};
          }
          % Origen
          \node[font=\AlmFontTiny, color=AlmGris] at (-0.28,-0.28) {0};
          % Indicadores de dirección
          \node[font=\AlmFontTiny\bfseries, color=AlmTerra] at (2.2,-0.45)
            {$+$};
          \node[font=\AlmFontTiny\bfseries, color=AlmTerra] at (-2.2,-0.45)
            {$-$};
          \node[font=\AlmFontTiny\bfseries, color=AlmTerra] at (0.35,2.3)
            {$+$};
          \node[font=\AlmFontTiny\bfseries, color=AlmTerra] at (0.35,-2.3)
            {$-$};
        \end{tikzpicture}
      \end{center}
    \end{column}
  \end{columns}
\end{frame}

% ---------------------------------------------------------------
\seccionframe{2. Leer un par $(x, y)$}
% ---------------------------------------------------------------

\begin{frame}{¿Qué significa el par $(x, y)$?}
  \begin{columns}[c]
    \begin{column}{0.50\textwidth}
      \begin{definicionbox}{Par de coordenadas $(x,\,y)$}
        \begin{itemize}
          \item \textbf{Primera coordenada} $x$:\\
                desplazamiento \emph{horizontal}
                desde el origen.\\
                $(+)$ = derecha, $(-)$ = izquierda.
          \item \textbf{Segunda coordenada} $y$:\\
                desplazamiento \emph{vertical}
                desde el origen.\\
                $(+)$ = arriba, $(-)$ = abajo.
        \end{itemize}
      \end{definicionbox}
      \vspace{6pt}
      \begin{pasobox}[1]
        \AlmFontFootnote
        Para el punto $A(3,\,2)$:\\
        \quad$x = 3$: muévete 3 unidades a la \textbf{derecha}.\\
        \quad$y = 2$: muévete 2 unidades hacia \textbf{arriba}.
      \end{pasobox}
    \end{column}
    \begin{column}{0.46\textwidth}
      \begin{center}
        \begin{tikzpicture}[scale=0.75]
          % Ejes
          \draw[recto, ->] (-0.5,0) -- (4.5,0)
            node[right, font=\AlmFontSmall\bfseries, color=AlmAzulM] {$x$};
          \draw[recto, ->] (0,-0.5) -- (0,3.5)
            node[above, font=\AlmFontSmall\bfseries, color=AlmAzulM] {$y$};
          % Cuadrícula
          \foreach \i in {1,2,3,4}{
            \draw[AlmAzulC!25, line width=0.3pt] (\i,-0.4)--(\i,3.3);
          }
          \foreach \j in {1,2,3}{
            \draw[AlmAzulC!25, line width=0.3pt] (-0.4,\j)--(4.3,\j);
          }
          % Marcas
          \foreach \i in {1,2,3,4}{
            \node[font=\AlmFontTiny, color=AlmGris] at (\i,-0.25) {\i};
          }
          \foreach \j in {1,2,3}{
            \node[font=\AlmFontTiny, color=AlmGris] at (-0.22,\j) {\j};
          }
          \node[font=\AlmFontTiny, color=AlmGris] at (-0.22,-0.22) {0};
          % Punto A(3,2)
          \node[punto] (A) at (3,2) {};
          \node[font=\AlmFontSmall\bfseries, color=AlmAzulM, above right=1pt]
            at (A) {$A(3,\,2)$};
          % Flechas de desplazamiento
          \draw[AlmTerra, ->, line width=1.2pt, dashed]
            (0,0) -- (3,0)
            node[midway, below, font=\AlmFontTiny\bfseries, color=AlmTerra]
            {$x=3$};
          \draw[AlmVerde, ->, line width=1.2pt, dashed]
            (3,0) -- (3,2)
            node[midway, right, font=\AlmFontTiny\bfseries, color=AlmVerde]
            {$y=2$};
        \end{tikzpicture}
      \end{center}
    \end{column}
  \end{columns}
\end{frame}

% ---------------------------------------------------------------
\seccionframe{3. Los cuatro cuadrantes}
% ---------------------------------------------------------------

\begin{frame}{Los cuatro cuadrantes y sus signos}
  \begin{columns}[c]
    \begin{column}{0.44\textwidth}
      \begin{center}
        \begin{tikzpicture}[scale=0.80]
          % Zonas sombreadas
          \fill[AlmAzulC!22] (0,0) rectangle (2.8,2.8);   % I
          \fill[AlmAmbar!22] (-2.8,0) rectangle (0,2.8);  % II
          \fill[AlmTerra!18] (-2.8,-2.8) rectangle (0,0); % III
          \fill[AlmVerde!18] (0,-2.8) rectangle (2.8,0);  % IV
          % Ejes
          \draw[recto, ->] (-3.1,0) -- (3.2,0)
            node[right, font=\AlmFontSmall\bfseries, color=AlmAzulM] {$x$};
          \draw[recto, ->] (0,-3.1) -- (0,3.2)
            node[above, font=\AlmFontSmall\bfseries, color=AlmAzulM] {$y$};
          % Etiquetas de cuadrante
          \node[font=\AlmFontNormal\bfseries, color=AlmAzul]
            at (1.5, 1.5)  {I};
          \node[font=\AlmFontNormal\bfseries, color=AlmAmbar!80!black]
            at (-1.5, 1.5) {II};
          \node[font=\AlmFontNormal\bfseries, color=AlmTerra]
            at (-1.5,-1.5) {III};
          \node[font=\AlmFontNormal\bfseries, color=AlmVerde]
            at (1.5,-1.5)  {IV};
          % Signos en cuadrantes
          \node[font=\AlmFontTiny, color=AlmAzul]
            at (1.5, 0.9) {$(+,+)$};
          \node[font=\AlmFontTiny, color=AlmAmbar!80!black]
            at (-1.5, 0.9) {$(-,+)$};
          \node[font=\AlmFontTiny, color=AlmTerra]
            at (-1.5,-0.9) {$(-,-)$};
          \node[font=\AlmFontTiny, color=AlmVerde]
            at (1.5,-0.9)  {$(+,-)$};
          % Origen
          \node[font=\AlmFontTiny, color=AlmGris] at (-0.22,-0.22) {O};
        \end{tikzpicture}
      \end{center}
    \end{column}
    \begin{column}{0.52\textwidth}
      \vspace{4pt}
      {\AlmFontSmall
      \begin{tabular}{@{}cllc@{}}
        \toprule
        \textbf{Cuadrante} & \textbf{Signo $x$}
          & \textbf{Signo $y$} & \textbf{Color}\\
        \midrule
        \textbf{I}   & positivo $(+)$ & positivo $(+)$
          & \textcolor{AlmAzulM}{azul}\\
        \textbf{II}  & negativo $(-)$ & positivo $(+)$
          & \textcolor{AlmAmbar!70!black}{ámbar}\\
        \textbf{III} & negativo $(-)$ & negativo $(-)$
          & \textcolor{AlmTerra}{terracota}\\
        \textbf{IV}  & positivo $(+)$ & negativo $(-)$
          & \textcolor{AlmVerde}{verde}\\
        \bottomrule
      \end{tabular}}
      \vspace{8pt}
      \begin{notabox}[Truco para memorizar]
        \AlmFontSmall
        Empieza en el \textbf{cuadrante I} (arriba-derecha)
        y gira en sentido \textbf{contrario a las agujas del reloj}:
        I → II → III → IV.
      \end{notabox}
    \end{column}
  \end{columns}
\end{frame}

% ---------------------------------------------------------------
\seccionframe{4. Representar puntos en el plano}
% ---------------------------------------------------------------

\begin{frame}{Cinco puntos en el plano de Almería}
  \framesubtitle{Cada punto representa un lugar de la ciudad}
  \begin{columns}[c]
    \begin{column}{0.50\textwidth}
      \begin{center}
        \begin{tikzpicture}[scale=0.72]
          % Cuadrícula
          \foreach \i in {-4,-3,-2,-1,0,1,2,3,4,5}{
            \draw[AlmAzulC!20, line width=0.3pt]
              (\i,-4.3)--(\i,4.5);
          }
          \foreach \j in {-3,-2,-1,0,1,2,3,4}{
            \draw[AlmAzulC!20, line width=0.3pt]
              (-4.3,\j)--(5.3,\j);
          }
          % Ejes
          \draw[recto, ->] (-4.5,0) -- (5.5,0)
            node[right, font=\AlmFontSmall\bfseries, color=AlmAzulM] {$x$};
          \draw[recto, ->] (0,-3.8) -- (0,4.8)
            node[above, font=\AlmFontSmall\bfseries, color=AlmAzulM] {$y$};
          % Marcas numéricas
          \foreach \i in {-4,-3,-2,-1,1,2,3,4,5}{
            \node[font=\AlmFontTiny, color=AlmGris] at (\i,-0.3) {\i};
          }
          \foreach \j in {-3,-2,-1,1,2,3,4}{
            \node[font=\AlmFontTiny, color=AlmGris] at (-0.3,\j) {\j};
          }
          \node[font=\AlmFontTiny, color=AlmGris] at (-0.28,-0.28) {0};
          % Punto A(3,2)
          \node[punto] (A) at (3,2) {};
          \node[font=\AlmFontScript\bfseries, color=AlmAzulM,
                above right=0pt] at (A) {$A(3,2)$};
          % Punto B(-4,1)
          \node[puntot] (B) at (-4,1) {};
          \node[font=\AlmFontScript\bfseries, color=AlmTerra,
                above right=0pt] at (B) {$B(-4,1)$};
          % Punto C(-2,-3)
          \node[puntov] (C) at (-2,-3) {};
          \node[font=\AlmFontScript\bfseries, color=AlmVerde,
                above right=0pt] at (C) {$C(-2,-3)$};
          % Punto D(5,-1)
          \node[circle, fill=BCrear, inner sep=2.2pt,
                draw=AlmBlanco, line width=0.8pt] (D) at (5,-1) {};
          \node[font=\AlmFontScript\bfseries, color=BCrear,
                above left=0pt] at (D) {$D(5,-1)$};
          % Punto E(0,4)
          \node[circle, fill=BAplicar, inner sep=2.2pt,
                draw=AlmBlanco, line width=0.8pt] (E) at (0,4) {};
          \node[font=\AlmFontScript\bfseries, color=BAplicar!80!black,
                right=2pt] at (E) {$E(0,4)$};
        \end{tikzpicture}
      \end{center}
    \end{column}
    \begin{column}{0.46\textwidth}
      \vspace{4pt}
      {\AlmFontSmall
      \begin{tabular}{@{}clcc@{}}
        \toprule
        \textbf{Punto} & \textbf{Coordenadas}
          & \textbf{Cuadrante} & \textbf{Color}\\
        \midrule
        $A$ & $(3,\;2)$    & I   & \textcolor{AlmAzulM}{azul}\\
        $B$ & $(-4,\;1)$   & II  & \textcolor{AlmTerra}{rojo}\\
        $C$ & $(-2,\;-3)$  & III & \textcolor{AlmVerde}{verde}\\
        $D$ & $(5,\;-1)$   & IV  & \textcolor{BCrear}{púrpura}\\
        $E$ & $(0,\;4)$    & eje $y$ & \textcolor{BAplicar!80!black}{ámbar}\\
        \bottomrule
      \end{tabular}}
      \vspace{8pt}
      \begin{notabox}[Punto en el eje]
        \AlmFontFootnote
        Si $x = 0$, el punto está \textbf{sobre el eje Y}.\\
        Si $y = 0$, el punto está \textbf{sobre el eje X}.\\
        No pertenece a ningún cuadrante.
      \end{notabox}
    \end{column}
  \end{columns}
\end{frame}

% ---------------------------------------------------------------
\seccionframe{5. Escala: de la cuadrícula al mapa real}
% ---------------------------------------------------------------

\begin{frame}{Coordenadas y escala en el mapa de Almería}
  \begin{columns}[c]
    \begin{column}{0.52\textwidth}
      \begin{teoriabox}
        \AlmFontSmall
        En nuestro \textbf{mapa de Almería},
        asignamos la escala:
        \[
          1\ \text{unidad} = 100\ \text{m}
        \]
        Así, el punto $P(3,\,5)$ significa:
        \begin{itemize}
          \item $x=3$: 300 m hacia el \textbf{este}.
          \item $y=5$: 500 m hacia el \textbf{norte}.
        \end{itemize}
      \end{teoriabox}
      \vspace{6pt}
      \begin{ejemplobox}
        \AlmFontFootnote
        La \textbf{Alcazaba} de Almería tiene
        coordenadas GPS reales:\\
        \[
          36{,}8409°\,\text{N}, \quad 2{,}4699°\,\text{O}
        \]
        Estas coordenadas son \textbf{latitud/longitud},
        un sistema diferente al $(x,y)$ cartesiano.
      \end{ejemplobox}
    \end{column}
    \begin{column}{0.44\textwidth}
      \begin{center}
        \begin{tikzpicture}[scale=0.82]
          % Mini mapa con escala
          \fill[AlmAzulC!15, rounded corners=3pt]
            (0,0) rectangle (4.2,3.5);
          % Cuadrícula
          \foreach \i in {1,2,3,4}{
            \draw[AlmAzulC!40, line width=0.4pt] (\i,0)--(\i,3.5);
          }
          \foreach \j in {1,2,3}{
            \draw[AlmAzulC!40, line width=0.4pt] (0,\j)--(4.2,\j);
          }
          % Ejes
          \draw[recto, ->] (0,0)--(4.5,0)
            node[right,font=\AlmFontTiny\bfseries, color=AlmAzulM]{$x$};
          \draw[recto, ->] (0,0)--(0,3.8)
            node[above,font=\AlmFontTiny\bfseries, color=AlmAzulM]{$y$};
          % Marcas
          \foreach \i in {1,2,3,4}{
            \node[font=\AlmFontTiny, color=AlmGris] at (\i,-0.22) {\i};
          }
          \foreach \j in {1,2,3}{
            \node[font=\AlmFontTiny, color=AlmGris] at (-0.22,\j) {\j};
          }
          \node[font=\AlmFontTiny, color=AlmGris] at (-0.18,-0.18){0};
          % Punto P(3,2) — Alcazaba aprox.
          \node[puntot] (P) at (3,2) {};
          \node[font=\AlmFontTiny\bfseries, color=AlmTerra, above right=0pt]
            at (P) {Alcazaba};
          \node[font=\AlmFontTiny, color=AlmTerra, below right=0pt]
            at (P) {$(3,2)$};
          % Acotación x
          \draw[acota] (0,-0.48) -- (3,-0.48);
          \node[font=\AlmFontTiny, color=AlmAzulM] at (1.5,-0.72) {$300\ \text{m}$};
          % Acotación y
          \draw[acota] (3.5,0) -- (3.5,2);
          \node[font=\AlmFontTiny, color=AlmAzulM, rotate=90, anchor=south]
            at (3.75,1) {$200\ \text{m}$};
          % Escala
          \node[font=\AlmFontTiny, color=AlmGris, align=center]
            at (2.1,3.3) {escala: 1 ud = 100 m};
        \end{tikzpicture}
      \end{center}
    \end{column}
  \end{columns}
\end{frame}

% ---------------------------------------------------------------
\seccionframe{6. Ejercicios multinivel — Bloom}
% ---------------------------------------------------------------

\begin{frame}{Ejercicio COMPRENDER}
  \bloomtag{COMPRENDER}\quad\dificultad{2}
  \vspace{4pt}
  \begin{comprenderbox}
    \AlmFontSmall
    Responde razonando cada respuesta:
    \begin{enumerate}
      \item ¿Qué significa el par $(-3,\,4)$?
            Explica el movimiento en el plano:
            ¿hacia dónde y cuánto nos desplazamos en cada eje?
      \item ¿Por qué el par $(3,\,5)$ es diferente al par $(5,\,3)$?
            Representa \textbf{ambos} en un mismo sistema de ejes
            y comprueba que son puntos distintos.
      \item ¿Qué pares de coordenadas $(x, y)$ con $x = y$
            están a la misma distancia del origen
            sobre la recta diagonal?
            Escribe tres ejemplos.
    \end{enumerate}
  \end{comprenderbox}
\end{frame}

\begin{frame}{Ejercicio APLICAR}
  \bloomtag{APLICAR}\quad\dificultad{3}
  \vspace{4pt}
  \begin{aplicarbox}
    \AlmFontFootnote
    \textbf{Parte A — Representa los puntos:}\\
    Dibuja un sistema de ejes de $-2$ a $6$ en $x$
    y de $0$ a $6$ en $y$. Sitúa los puntos:
    \[
      A(1,\,4), \quad B(3,\,2), \quad C(5,\,5), \quad D(0,\,3)
    \]
    Une los puntos en orden $A \to B \to C \to D \to A$.
    ¿Qué forma geométrica se obtiene? Nómbrala.\\[6pt]
    \textbf{Parte B — Lee coordenadas del mapa:}\\
    Tu profesor proyectará un mapa con 4 puntos marcados.
    Escribe las coordenadas $(x, y)$ de cada uno.
    Compara con tus compañeros. ¿Todos obtuvisteis
    los mismos resultados?\\[6pt]
    \textbf{Parte C — Aplicación a Almería:}\\
    Si el origen $(0,0)$ es la Plaza de la Catedral y
    la escala es 1 ud = 100 m, escribe las coordenadas
    aproximadas de un monumento que conozcas.
  \end{aplicarbox}
\end{frame}

\begin{frame}{Ejercicio EVALUAR}
  \bloomtag{EVALUAR}\quad\dificultad{4}
  \vspace{4pt}
  \begin{evaluarbox}
    \AlmFontSmall
    \textbf{Debate y argumenta:}\\[6pt]
    El Ayuntamiento de Almería quiere crear una nueva
    guía turística digital.\\[4pt]
    \begin{itemize}
      \item ¿Deberían usar \textbf{coordenadas cartesianas}
            $(x, y)$ o \textbf{coordenadas GPS} (latitud/longitud)?
    \end{itemize}
    \vspace{6pt}
    Escribe un argumento de \textbf{5--6 líneas} que incluya:
    \begin{enumerate}
      \item Ventajas e inconvenientes de cada sistema.
      \item Tu recomendación justificada.
      \item Al menos un ejemplo concreto de Almería.
    \end{enumerate}
  \end{evaluarbox}
  \vspace{4pt}
  \begin{notabox}[Pista]
    \AlmFontFootnote
    Piensa en el público de la guía: ¿turistas con smartphone,
    escolares con papel, o profesionales de cartografía?
  \end{notabox}
\end{frame}

% ---------------------------------------------------------------
%  EVIDENCIA
% ---------------------------------------------------------------
\begin{frame}{Evidencia del día}
  \begin{evidenciabox}
    Entrega tu \textbf{Ficha de Coordenadas} con:
    \begin{itemize}
      \item Sistema de ejes claramente dibujado con regla,
            etiquetado $x$ e $y$.
      \item Al menos \textbf{6 puntos} correctamente colocados
            en las intersecciones de la cuadrícula
            (no entre líneas).
      \item Cada punto identificado con su \textbf{nombre}
            y su par $(x, y)$.
      \item El \textbf{cuadrante} al que pertenece indicado
            para cada punto.
      \item Una nota explicando con tus palabras
            la diferencia entre $(x, y)$ y GPS.
    \end{itemize}
  \end{evidenciabox}
\end{frame}

% ---------------------------------------------------------------
%  ERROR FRECUENTE
% ---------------------------------------------------------------
\begin{frame}{¡Cuidado con este error!}
  \begin{errorbox}
    \begin{columns}[c]
      \begin{column}{0.50\textwidth}
        \incorrecto\ \textbf{Error frecuente}\\[6pt]
        Confundir la coordenada $x$ con la $y$:
        \begin{center}
          \AlmFontLarge
          $B(1,\,-4) \;\neq\; B(-4,\,1)$
        \end{center}
        \AlmFontFootnote
        Si marcas el punto $B(-4,\,1)$ en la posición
        de $B(1,\,-4)$, estarás a \textbf{cientos de metros}
        del lugar real en el mapa de Almería.
      \end{column}
      \begin{column}{0.46\textwidth}
        \correcto\ \textbf{Correcto}\\[6pt]
        \AlmFontSmall
        \textbf{Recuerda el orden estricto:}
        \begin{enumerate}
          \item \textbf{Primera} $\to$ eje $x$ (horizontal).
          \item \textbf{Segunda} $\to$ eje $y$ (vertical).
        \end{enumerate}
        \vspace{4pt}
        \AlmFontFootnote
        $B(-4,\,1)$:\\
        \quad$x = -4$: 4 unidades a la \textbf{izquierda}.\\
        \quad$y = 1$: 1 unidad hacia \textbf{arriba}.
      \end{column}
    \end{columns}
  \end{errorbox}
  \vspace{6pt}
  \begin{notabox}[Regla de oro]
    \AlmFontSmall
    «Primero la \textbf{calle} ($x$ horizontal),
    luego el \textbf{piso} ($y$ vertical).»
  \end{notabox}
\end{frame}

% ---------------------------------------------------------------
%  RESUMEN
% ---------------------------------------------------------------
\begin{frame}{Resumen de la sesión 04}
  \begin{resumenbox}
    \begin{itemize}
      \item El sistema cartesiano tiene un \textbf{eje X}
            (horizontal) y un \textbf{eje Y} (vertical)
            que se cruzan en el \textbf{origen} $(0,0)$.
      \item Un par $(x,\,y)$: primero la coordenada
            \textbf{horizontal}, luego la \textbf{vertical}.
      \item Los cuatro cuadrantes tienen signos:
            I $(+,+)$, II $(-,+)$, III $(-,-)$, IV $(+,-)$.
      \item $(3,\,5) \neq (5,\,3)$:
            el \textbf{orden importa}.
      \item Con escala $1\ \text{ud} = 100\ \text{m}$
            podemos usar $(x,y)$ para el mapa de Almería.
    \end{itemize}
  \end{resumenbox}
  \vspace{6pt}
  \begin{center}
    {\AlmFontSmall\color{AlmAzulM}
    \faArrowCircleRight\enskip
    \textbf{Próxima sesión:}
    Rectas, segmentos y ángulos en la ciudad.}
  \end{center}
\end{frame}

\end{document}




% % ================================================================
%  sesion05.tex  —  Sesión 5 de 20
%  Rectas, segmentos y ángulos en la ciudad
%  Almería en Cifras y Formas · Matemáticas 1.º ESO
%  Compilar con:  pdflatex sesion05.tex  (dos pasadas)
% ================================================================
\documentclass[aspectratio=169, 10pt, spanish]{beamer}
\usepackage{almeria-beamer}

% ---------------------------------------------------------------
%  DATOS DE LA SESIÓN
% ---------------------------------------------------------------
\renewcommand{\SessionNumber}{05}
\renewcommand{\SessionTitle}{Rectas, segmentos y ángulos en la ciudad}
\renewcommand{\SessionName}{Sesión 05}
\renewcommand{\SessionObjective}{%
  Identificar y trazar puntos, rectas, segmentos y ángulos en planos
  urbanos, reconociendo rectas paralelas y perpendiculares en las
  calles de Almería.}
\renewcommand{\BlockName}{Bloque 2: Geometría urbana}

% ================================================================
\begin{document}

% ---------------------------------------------------------------
%  PORTADA
% ---------------------------------------------------------------
\begin{frame}[plain]
  \titlepage
\end{frame}

% ---------------------------------------------------------------
%  ÍNDICE DE LA SESIÓN
% ---------------------------------------------------------------
\begin{frame}{Índice de la sesión}
  \begin{columns}[t]
    \begin{column}{0.5\textwidth}
      \begin{enumerate}
        \item Punto, recta, semirrecta y segmento
        \item Ángulos y sus tipos
        \item Rectas paralelas y perpendiculares
        \item Almería: geometría en las calles
      \end{enumerate}
    \end{column}
    \begin{column}{0.5\textwidth}
      \begin{enumerate}
        \setcounter{enumi}{4}
        \item Diagrama completo de elementos
        \item Ejercicios multinivel (Bloom)
        \item Evidencia del día
        \item Error frecuente
        \item Resumen
      \end{enumerate}
    \end{column}
  \end{columns}
\end{frame}

% ---------------------------------------------------------------
\seccionframe{1. Punto, recta, semirrecta y segmento}
% ---------------------------------------------------------------

\begin{frame}{Elementos geométricos básicos (1/2)}
  \begin{columns}[T]
    \begin{column}{0.50\textwidth}
      \begin{definicionbox}{Punto}
        \AlmFontFootnote
        Posición exacta en el plano, \textbf{sin dimensiones}.\\
        Se representa con una letra mayúscula: $A$, $B$, $P$\ldots
      \end{definicionbox}
      \vspace{5pt}
      \begin{definicionbox}{Recta}
        \AlmFontFootnote
        Se extiende \textbf{infinitamente} en ambas direcciones;
        no tiene extremos. Se representa con una letra minúscula: $r$, $s$\ldots
      \end{definicionbox}
    \end{column}
    \begin{column}{0.46\textwidth}
      \begin{center}
        \begin{tikzpicture}[scale=0.88]
          % ---------- Punto ----------
          \node[punto] (P) at (0, 5.2) {};
          \node[font=\AlmFontSmall\bfseries, color=AlmAzulM, right=3pt of P] {$P$};
          \node[font=\AlmFontTiny, color=AlmGris, right=22pt of P] {punto};

          % ---------- Recta ----------
          \draw[recto, <->] (-1.8, 4.2) -- (2.2, 4.2);
          \node[font=\AlmFontTiny, color=AlmAzulM] at (-1.5, 4.45) {$r$};
          \node[font=\AlmFontTiny, color=AlmGris] at (2.8, 4.2) {recta};
        \end{tikzpicture}
      \end{center}
    \end{column}
  \end{columns}
\end{frame}

\begin{frame}{Elementos geométricos básicos (2/2)}
  \begin{columns}[T]
    \begin{column}{0.50\textwidth}
      \begin{definicionbox}{Semirrecta}
        \AlmFontFootnote
        Mitad de una recta: tiene \textbf{un extremo} (origen)
        y se extiende hacia un solo lado.
      \end{definicionbox}
      \vspace{5pt}
      \begin{definicionbox}{Segmento}
        \AlmFontFootnote
        Porción \textbf{finita} de recta entre dos puntos $A$ y $B$.
        Se escribe $\overline{AB}$. Tiene longitud medible.
      \end{definicionbox}
    \end{column}
    \begin{column}{0.46\textwidth}
      \begin{center}
        \begin{tikzpicture}[scale=0.88]
          % ---------- Semirrecta ----------
          \node[punto] at (-1.8, 3.2) {};
          \draw[recto, ->] (-1.8, 3.2) -- (2.2, 3.2);
          \node[font=\AlmFontTiny, color=AlmGris] at (2.8, 3.2) {semirrecta};

          % ---------- Segmento ----------
          \node[punto] (Sa) at (-1.8, 2.2) {};
          \node[punto] (Sb) at (2.2, 2.2) {};
          \draw[recto] (Sa) -- (Sb);
          \node[font=\AlmFontSmall\bfseries, color=AlmAzulM, below=2pt] at (Sa) {$A$};
          \node[font=\AlmFontSmall\bfseries, color=AlmAzulM, below=2pt] at (Sb) {$B$};
          \node[font=\AlmFontTiny, color=AlmGris] at (2.8, 2.2) {$\overline{AB}$};
          \draw[acota] (-1.8,1.78) -- (2.2,1.78);
          \node[font=\AlmFontTiny, color=AlmGris] at (0.2, 1.52) {longitud medible};

          \draw[AlmAmbar!60, line width=0.6pt, dashed]
            (-2.1,0) -- (3.2,0);
          \node[font=\AlmFontTiny\bfseries, color=AlmAmbar!80!black]
            at (0.5, -0.22) {elementos geométricos básicos};
        \end{tikzpicture}
      \end{center}
    \end{column}
  \end{columns}
\end{frame}

% ---------------------------------------------------------------
\seccionframe{2. Ángulos y sus tipos}
% ---------------------------------------------------------------

\begin{frame}{Qué es un ángulo}
  \begin{columns}[c]
    \begin{column}{0.48\textwidth}
      \begin{definicionbox}{Ángulo}
        \AlmFontSmall
        \textbf{Abertura} entre dos semirrectas
        que parten del mismo punto (vértice).\\
        Se mide en \textbf{grados} (°).
      \end{definicionbox}
      \vspace{8pt}
      \begin{notabox}[Cómo clasificarlo]
        \AlmFontFootnote
        Basta comparar su medida con 90° y 180°:
        antes de 90° es agudo, en 90° es recto,
        entre 90° y 180° es obtuso y en 180° es llano.
      \end{notabox}
    \end{column}
    \begin{column}{0.46\textwidth}
      {\AlmFontSmall
      \begin{tabular}{@{}lll@{}}
        \toprule
        \textbf{Tipo} & \textbf{Condición} & \textbf{Color}\\
        \midrule
        Agudo   & $0° < \alpha < 90°$
          & \textcolor{AlmAzulM}{azul}\\
        Recto   & $\alpha = 90°$
          & \textcolor{AlmAmbar!80!black}{ámbar}\\
        Obtuso  & $90° < \alpha < 180°$
          & \textcolor{AlmTerra}{rojo}\\
        Llano   & $\alpha = 180°$
          & \textcolor{AlmVerde}{verde}\\
        \bottomrule
      \end{tabular}}
      \vspace{8pt}
      \begin{ejemplobox}
        \AlmFontFootnote
        Una puerta poco abierta suele formar un
        \textbf{ángulo agudo}; la esquina de un folio
        es un \textbf{ángulo recto}.
      \end{ejemplobox}
    \end{column}
  \end{columns}
\end{frame}

\begin{frame}{Galería de tipos de ángulos}
  \begin{center}
    \begin{tikzpicture}[scale=0.86]
      % ---- Agudo ~55° ----
      \draw[recto, color=AlmAzulM]
        (0,0) -- (1.3,0);
      \draw[recto, color=AlmAzulM]
        (0,0) -- (0.74,1.06);
      \draw[AlmAzulM, line width=0.7pt]
        (0.45,0) arc (0:55:0.45);
      \node[font=\AlmFontTiny\bfseries, color=AlmAzulM] at (0.7,0.25)
        {$55°$};
      \node[font=\AlmFontTiny, color=AlmGris] at (0.65,-0.28)
        {agudo};

      % ---- Recto 90° ----
      \draw[recto, color=AlmAmbar!80!black]
        (2.5,0) -- (3.8,0);
      \draw[recto, color=AlmAmbar!80!black]
        (2.5,0) -- (2.5,1.3);
      \draw[rectoang]
        (2.5,0) ++(0.28,0) -- ++(0,0.28) -- ++(-0.28,0);
      \node[font=\AlmFontTiny\bfseries, color=AlmAmbar!80!black]
        at (3.2,0.45) {$90°$};
      \node[font=\AlmFontTiny, color=AlmGris] at (3.15,-0.28)
        {recto};

      % ---- Obtuso ~120° ----
      \draw[recto, color=AlmTerra]
        (4.8,0) -- (6.1,0);
      \draw[recto, color=AlmTerra]
        (4.8,0) -- (4.17,1.06);
      \draw[AlmTerra, line width=0.7pt]
        (5.22,0) arc (0:120:0.42);
      \node[font=\AlmFontTiny\bfseries, color=AlmTerra] at (5.05,0.35)
        {$120°$};
      \node[font=\AlmFontTiny, color=AlmGris] at (5.45,-0.28)
        {obtuso};

      % ---- Llano 180° ----
      \draw[recto, color=AlmVerde]
        (0,-1.5) -- (3,-1.5);
      \node[punto, fill=AlmVerde] at (1.5,-1.5) {};
      \draw[AlmVerde, line width=0.7pt]
        (1.92,-1.5) arc (0:180:0.42);
      \node[font=\AlmFontTiny\bfseries, color=AlmVerde] at (1.5,-1.1)
        {$180°$};
      \node[font=\AlmFontTiny, color=AlmGris] at (1.5,-1.82)
        {llano};
    \end{tikzpicture}
  \end{center}
  \vspace{6pt}
  \begin{center}
    \begin{notabox}[Idea clave]
      \AlmFontFootnote
      Cuanto mayor es la abertura, mayor es la medida del ángulo.
    \end{notabox}
  \end{center}
\end{frame}

% ---------------------------------------------------------------
\seccionframe{3. Rectas paralelas y perpendiculares}
% ---------------------------------------------------------------

\begin{frame}{Rectas paralelas y perpendiculares}
  \begin{columns}[T]
    \begin{column}{0.48\textwidth}
      \begin{definicionbox}{Rectas paralelas}
        \AlmFontSmall
        Dos rectas son \textbf{paralelas} si
        \textbf{nunca se cortan}, por mucho que
        las prolonguemos.\\[4pt]
        Siempre mantienen la \textbf{misma distancia}
        entre sí.\\[4pt]
        Notación: $r \parallel s$
      \end{definicionbox}
      \vspace{6pt}
      \begin{center}
        \begin{tikzpicture}[scale=0.75]
          \draw[recto, color=AlmAzulM]
            (0,0.5) -- (3.5,0.5)
            node[right, font=\AlmFontTiny\bfseries, color=AlmAzulM]{$r$};
          \draw[recto, color=AlmAzulM]
            (0,1.6) -- (3.5,1.6)
            node[right, font=\AlmFontTiny\bfseries, color=AlmAzulM]{$s$};
          % Marcas de paralelismo
          \draw[AlmAmbar, line width=1pt]
            (1.5,0.38) -- (1.5,0.62);
          \draw[AlmAmbar, line width=1pt]
            (1.5,1.48) -- (1.5,1.72);
          % Distancia
          \draw[acota] (2.8,0.5) -- (2.8,1.6);
          \node[font=\AlmFontTiny, color=AlmGris, right=2pt]
            at (2.8,1.05) {dist. cte.};
        \end{tikzpicture}
      \end{center}
    \end{column}
    \begin{column}{0.48\textwidth}
      \begin{definicionbox}{Rectas perpendiculares}
        \AlmFontSmall
        Dos rectas son \textbf{perpendiculares} si
        se cortan formando un \textbf{ángulo de 90°}
        (ángulo recto).\\[4pt]
        Notación: $r \perp s$
      \end{definicionbox}
      \vspace{6pt}
      \begin{center}
        \begin{tikzpicture}[scale=0.75]
          \draw[recto, color=AlmTerra]
            (0,1.05) -- (3.5,1.05)
            node[right, font=\AlmFontTiny\bfseries, color=AlmTerra]{$r$};
          \draw[recto, color=AlmTerra]
            (1.75,0) -- (1.75,2.3)
            node[above, font=\AlmFontTiny\bfseries, color=AlmTerra]{$s$};
          % Cuadradito ángulo recto
          \draw[rectoang]
            (1.75,1.05) ++(0.22,0) -- ++(0,0.22) -- ++(-0.22,0);
          \node[font=\AlmFontSmall\bfseries, color=AlmAmbar!80!black]
            at (2.35,1.5) {$90°$};
        \end{tikzpicture}
      \end{center}
    \end{column}
  \end{columns}
\end{frame}

% ---------------------------------------------------------------
\seccionframe{4. Almería: geometría en las calles}
% ---------------------------------------------------------------

\begin{frame}{Geometría en las calles de Almería}
  \begin{columns}[c]
    \begin{column}{0.50\textwidth}
      \begin{ejemplobox}
        \AlmFontSmall
        \textbf{Paralelas en Almería:}\\
        La \textbf{Avenida del Mediterráneo} y la
        \textbf{Calle del Obispo} discurren en la misma
        dirección; nunca se encuentran. Son \textbf{paralelas}.
        \vspace{4pt}\\
        \textbf{Perpendiculares en Almería:}\\
        La \textbf{Calle Reyes Católicos} corta a la
        \textbf{Rambla de Almería} formando un ángulo
        de aproximadamente \textbf{90°}. Son \textbf{perpendiculares}.
      \end{ejemplobox}
    \end{column}
    \begin{column}{0.46\textwidth}
      \begin{center}
        \begin{tikzpicture}[scale=0.85]
          % Fondo de manzana
          \fill[AlmFondo, rounded corners=3pt]
            (0,0) rectangle (4.4,3.8);
          % Calle horizontal 1 — Av. Mediterráneo (paralela)
          \fill[AlmAzulC!30]
            (0,2.6) rectangle (4.4,3.1);
          \node[font=\AlmFontTiny\bfseries, color=AlmAzul, align=center]
            at (2.2,2.85) {Av. Mediterráneo};
          % Calle horizontal 2 — C. del Obispo (paralela)
          \fill[AlmAzulC!30]
            (0,0.6) rectangle (4.4,1.1);
          \node[font=\AlmFontTiny\bfseries, color=AlmAzul, align=center]
            at (2.2,0.85) {C. del Obispo};
          % Marcas de paralelismo
          \draw[AlmAmbar, line width=1.2pt]
            (2.0,2.62) -- (2.0,3.08);
          \draw[AlmAmbar, line width=1.2pt]
            (2.0,0.62) -- (2.0,1.08);
          % Calle vertical — C. Reyes Católicos (perpendicular a Rambla)
          \fill[AlmTerra!25]
            (1.8,0) rectangle (2.3,3.8);
          \node[font=\AlmFontTiny\bfseries, color=AlmTerra,
                rotate=90, align=center]
            at (2.05,1.85) {C. Reyes Católicos};
          % Cuadradito ángulo recto en intersección Rambla
          \draw[rectoang, draw=AlmTerra, line width=0.9pt]
            (2.3,1.1) ++(0,0) -- ++(0.2,0) -- ++(0,0.2);
          % Leyenda
          \node[font=\AlmFontTiny, color=AlmAzulM] at (3.6,3.45)
            {\textcolor{AlmAzulM}{$\parallel$} paralelas};
          \node[font=\AlmFontTiny, color=AlmTerra] at (3.65,3.15)
            {\textcolor{AlmTerra}{$\perp$} perpendicular};
        \end{tikzpicture}
      \end{center}
    \end{column}
  \end{columns}
\end{frame}

\begin{frame}{Las calles son segmentos}
  \begin{columns}[c]
    \begin{column}{0.56\textwidth}
      \begin{notabox}[Idea clave]
        \AlmFontSmall
        Cualquier calle tiene un inicio y un final:
        en el plano urbano la representamos como un
        \textbf{segmento}, no como una recta infinita.
      \end{notabox}
      \vspace{8pt}
      \begin{ejemplobox}
        \AlmFontFootnote
        Si una calle va desde la plaza A hasta el cruce B,
        podemos escribirla como $\overline{AB}$ y medir
        su longitud aproximada.
      \end{ejemplobox}
    \end{column}
    \begin{column}{0.38\textwidth}
      \begin{center}
        \begin{tikzpicture}[scale=0.9]
          \node[punto] (A) at (0,0) {};
          \node[punto] (B) at (3.2,0) {};
          \draw[recto] (A) -- (B);
          \node[font=\AlmFontSmall\bfseries, color=AlmAzulM, below=2pt] at (A) {$A$};
          \node[font=\AlmFontSmall\bfseries, color=AlmAzulM, below=2pt] at (B) {$B$};
          \draw[acota] (0,-0.45) -- (3.2,-0.45);
          \node[font=\AlmFontTiny, color=AlmGris] at (1.6,-0.72) {tramo de calle};
        \end{tikzpicture}
      \end{center}
    \end{column}
  \end{columns}
\end{frame}

% ---------------------------------------------------------------
\seccionframe{5. Diagrama completo de elementos}
% ---------------------------------------------------------------

\begin{frame}{Todos los elementos en un mapa urbano}
  \framesubtitle{Plano esquemático de un barrio de Almería}
  \begin{center}
    \begin{tikzpicture}[scale=1.0]
      % Fondo
      \fill[AlmFondo, rounded corners=4pt]
        (-0.6,-0.6) rectangle (9.6,5.0);

      % ---- Calles paralelas horizontales ----
      \draw[recto, color=AlmAzulM, line width=1.8pt]
        (0,4.0) -- (9,4.0);
      \draw[recto, color=AlmAzulM, line width=1.8pt]
        (0,2.2) -- (9,2.2);
      \draw[recto, color=AlmAzulM, line width=1.8pt]
        (0,0.4) -- (9,0.4);

      % Marcas de paralelismo en las tres rectas
      \foreach \y in {4.0,2.2,0.4}{
        \draw[AlmAmbar, line width=1.0pt]
          (4.3,\y-0.12) -- (4.3,\y+0.12);
        \draw[AlmAmbar, line width=1.0pt]
          (4.6,\y-0.12) -- (4.6,\y+0.12);
      }
      \node[font=\AlmFontTiny\bfseries, color=AlmAzulM] at (-0.3,4.0)
        {$r_1$};
      \node[font=\AlmFontTiny\bfseries, color=AlmAzulM] at (-0.3,2.2)
        {$r_2$};
      \node[font=\AlmFontTiny\bfseries, color=AlmAzulM] at (-0.3,0.4)
        {$r_3$};
      \node[font=\AlmFontTiny, color=AlmAzulM] at (9.3,4.0) {Av.\ Mar};
      \node[font=\AlmFontTiny, color=AlmAzulM] at (9.3,2.2) {Rambla};
      \node[font=\AlmFontTiny, color=AlmAzulM] at (9.3,0.4) {C.\ Obispo};

      % ---- Calle perpendicular ----
      \draw[rectot, line width=1.8pt]
        (2.8,-0.4) -- (2.8,4.8);
      \node[font=\AlmFontTiny\bfseries, color=AlmTerra, rotate=90]
        at (2.55,2.2) {C.\ Reyes Católicos};
      % Cuadradito ángulo recto
      \draw[rectoang, draw=AlmTerra, line width=0.9pt]
        (2.8,2.2) ++(0.18,0) -- ++(0,0.18) -- ++(-0.18,0);

      % ---- Segmentos verdes ----
      \draw[rectov, line width=2pt]
        (5.5,0.4) -- (5.5,2.2);
      \node[font=\AlmFontTiny\bfseries, color=AlmVerde, rotate=90]
        at (5.3,1.3) {$\overline{CD}$};
      \draw[rectov, line width=2pt]
        (7.0,2.2) -- (8.5,4.0);
      \node[font=\AlmFontTiny\bfseries, color=AlmVerde]
        at (8.0,3.35) {$\overline{EF}$};
      \draw[rectov, line width=2pt]
        (1.0,0.4) -- (2.8,0.4);
      \node[font=\AlmFontTiny\bfseries, color=AlmVerde]
        at (1.9,0.62) {$\overline{AB}$};

      % ---- Puntos notables ----
      \node[punto] (Pt) at (2.8,4.0) {};
      \node[font=\AlmFontTiny\bfseries, color=AlmAzulM, above right=0pt]
        at (Pt) {$P$};
      \node[punto] (Qt) at (5.5,2.2) {};
      \node[font=\AlmFontTiny\bfseries, color=AlmAzulM, above right=0pt]
        at (Qt) {$Q$};

      % ---- Ángulo agudo entre segmento oblicuo y horizontal ----
      \draw[AlmAmbar!80!black, line width=0.7pt]
        (7.28,2.2) arc (0:50:0.28);
      \node[font=\AlmFontTiny\bfseries, color=AlmAmbar!80!black]
        at (7.7,2.45) {agudo};

      % ---- Leyenda ----
      \node[draw=AlmAzulC!50, fill=AlmBlanco, rounded corners=2pt,
            inner sep=4pt, font=\AlmFontTiny, align=left]
        at (1.5,3.35) {
          \textcolor{AlmAzulM}{—} paralelas ($r_1 \parallel r_2 \parallel r_3$)\\
          \textcolor{AlmTerra}{—} perpendicular ($\perp$)\\
          \textcolor{AlmVerde}{—} segmentos ($\overline{AB}$, $\overline{CD}$, $\overline{EF}$)
        };
    \end{tikzpicture}
  \end{center}
\end{frame}

% ---------------------------------------------------------------
\seccionframe{6. Ejercicios multinivel — Bloom}
% ---------------------------------------------------------------

\begin{frame}{Ejercicio RECORDAR}
  \bloomtag{RECORDAR}\quad\dificultad{1}
  \vspace{4pt}
  \begin{recordarbox}
    \AlmFontFootnote
    Con \textbf{regla y compás} (o solo regla), dibuja en
    tu cuaderno cada uno de los siguientes elementos:
    \begin{enumerate}
      \item Una \textbf{recta} $r$ (indica la doble flecha
            para mostrar que no tiene extremos).
      \item Un \textbf{segmento} $\overline{AB}$ de 5 cm.
            Mide y anota la longitud.
      \item Un \textbf{ángulo recto} de 90°.
            Marca el cuadradito del ángulo recto.
      \item Un \textbf{ángulo agudo} de aproximadamente 60°.
      \item Un \textbf{ángulo obtuso} de aproximadamente 110°.
    \end{enumerate}
    \vspace{4pt}
    Escribe junto a cada dibujo el nombre del elemento
    y su definición en una frase breve.
  \end{recordarbox}
\end{frame}

\begin{frame}{Ejercicio COMPRENDER}
  \bloomtag{COMPRENDER}\quad\dificultad{2}
  \vspace{4pt}
  \begin{comprenderbox}
    \AlmFontFootnote
    Observa una fotografía del centro histórico de Almería
    (o el plano que facilita el profesor).\\[6pt]
    Identifica y \textbf{justifica} con una frase cada uno:
    \begin{enumerate}
      \item \textbf{Dos calles paralelas.}\\
            Justificación: «Son paralelas porque\ldots»
      \item \textbf{Una intersección perpendicular.}\\
            Justificación: «Son perpendiculares porque\ldots»
      \item \textbf{Un ángulo obtuso} visible en un cruce o edificio.\\
            Justificación: «Es obtuso porque mide más de 90°\ldots»
    \end{enumerate}
    \vspace{4pt}
    \begin{notabox}[Pista]
      \AlmFontFootnote
      Para justificar el paralelismo: comprueba que
      las calles \emph{nunca} se cruzan en el mapa.
    \end{notabox}
  \end{comprenderbox}
\end{frame}

\begin{frame}{Ejercicio APLICAR}
  \bloomtag{APLICAR}\quad\dificultad{3}
  \vspace{4pt}
  \begin{aplicarbox}
    \AlmFontFootnote
    En el \textbf{plano simplificado de Almería} que te
    entrega el profesor:
    \begin{itemize}
      \item Traza con \textcolor{AlmAzulM}{\textbf{azul}}
            las \textbf{2 calles paralelas} que identifiques.
            Añade las marcas de paralelismo
            (\,$\shortmid$\;$\shortmid$\,) sobre ellas.
      \item Traza con \textcolor{AlmTerra}{\textbf{rojo}}
            la \textbf{intersección perpendicular} que
            identifiques. Dibuja el cuadradito de 90°.
      \item Marca con \textcolor{AlmVerde}{\textbf{verde}}
            \textbf{3 segmentos} que representen tramos de calles.
            Escribe su nombre ($\overline{AB}$, $\overline{CD}$,
            $\overline{EF}$) y mide su longitud aproximada en cm.
      \item Señala \textbf{2 ángulos agudos} y \textbf{1 ángulo obtuso}
            que aparezcan en algún cruce o edificio del plano.
    \end{itemize}
  \end{aplicarbox}
\end{frame}

% ---------------------------------------------------------------
%  EVIDENCIA
% ---------------------------------------------------------------
\begin{frame}{Evidencia del día}
  \begin{evidenciabox}
    \AlmFontFootnote
    Entrega tu \textbf{Plano anotado de Almería} con:
    \begin{itemize}
      \item \textcolor{AlmAzulM}{\textbf{Paralelas en azul}}
            con marcas de paralelismo y nombre de las calles.
      \item \textcolor{AlmTerra}{\textbf{Perpendiculares en rojo}}
            con el cuadradito de ángulo recto dibujado.
      \item \textcolor{AlmVerde}{\textbf{Segmentos en verde}}
            con sus extremos nombrados y longitud aproximada.
      \item Al menos \textbf{2 ángulos} etiquetados con su tipo
            (agudo / recto / obtuso / llano).
    \end{itemize}
  \end{evidenciabox}
\end{frame}

\begin{frame}{Evidencia del día: tabla resumen}
  \begin{evidenciabox}
    \AlmFontFootnote
    Añade al final una \textbf{tabla resumen} como esta:
    \vspace{6pt}
    \begin{center}
      \begin{tabular}{@{}lll@{}}
        \toprule
        \textbf{Elemento} & \textbf{Ejemplo en Almería} & \textbf{Justificación}\\
        \midrule
        Paralelas & Av.\,Mediterráneo y C.\,Obispo & Nunca se cortan\\
        Perpendiculares & C.\,Reyes Católicos $\perp$ Rambla & Forman 90°\\
        Segmento & Tramo $\overline{AB}$ & Tiene inicio y final\\
        \bottomrule
      \end{tabular}
    \end{center}
  \end{evidenciabox}
\end{frame}

% ---------------------------------------------------------------
%  ERROR FRECUENTE
% ---------------------------------------------------------------
\begin{frame}{¡Cuidado con este error!}
  \begin{errorbox}
    \begin{columns}[c]
      \begin{column}{0.50\textwidth}
        \incorrecto\ \textbf{Error frecuente}\\[6pt]
        Llamar «recta» a una calle porque parece
        infinitamente larga.\\[6pt]
        \begin{center}
          \cancel{«La Calle Mayor es una \textbf{recta}»}
        \end{center}
        \vspace{4pt}
        \AlmFontFootnote
        Toda calle tiene un inicio y un final:
        \textbf{siempre es un segmento}.
        Una recta matemática no tiene extremos
        y se extiende al infinito en ambos sentidos,
        lo que es imposible en el mundo real.
      \end{column}
      \begin{column}{0.46\textwidth}
        \correcto\ \textbf{Correcto}\\[6pt]
        \AlmFontSmall
        «La Calle Mayor es el
        \textbf{segmento} $\overline{AB}$,
        donde $A$ es su inicio y $B$ su final.»\\[8pt]
        \begin{center}
          \begin{tikzpicture}[scale=0.8]
            \node[punto] (A) at (0,0) {};
            \node[font=\AlmFontSmall\bfseries, color=AlmAzulM,
                  below=2pt] at (A) {$A$};
            \node[punto] (B) at (3.5,0) {};
            \node[font=\AlmFontSmall\bfseries, color=AlmAzulM,
                  below=2pt] at (B) {$B$};
            \draw[recto] (A) -- (B);
            \draw[acota] (0,-0.42) -- (3.5,-0.42);
            \node[font=\AlmFontTiny, color=AlmGris] at (1.75,-0.65)
              {longitud finita};
          \end{tikzpicture}
        \end{center}
      \end{column}
    \end{columns}
  \end{errorbox}
\end{frame}

\begin{frame}{Regla de oro}
  \begin{notabox}[Para no confundirlo]
    \AlmFontSmall
    \textbf{Recta} $\to$ sin extremos, infinita.\\[4pt]
    \textbf{Segmento} $\to$ dos extremos, longitud medible.
  \end{notabox}
  \vspace{10pt}
  \begin{center}
    \begin{tikzpicture}[scale=0.9]
      \draw[recto,<->] (-1.8,0.8) -- (1.8,0.8);
      \node[font=\AlmFontTiny, color=AlmGris] at (0,1.1) {recta};
      \node[punto] (A) at (-1.4,-0.6) {};
      \node[punto] (B) at (1.4,-0.6) {};
      \draw[recto] (A) -- (B);
      \node[font=\AlmFontTiny, color=AlmGris] at (0,-1.0) {segmento};
    \end{tikzpicture}
  \end{center}
\end{frame}

% ---------------------------------------------------------------
%  RESUMEN
% ---------------------------------------------------------------
\begin{frame}{Resumen de la sesión 05}
  \begin{resumenbox}
    \begin{itemize}
      \item \textbf{Punto}: posición sin dimensiones. \textbf{Recta}: sin extremos, infinita. \textbf{Semirrecta}: un extremo. \textbf{Segmento}: dos extremos, longitud medible.
      \item Un \textbf{ángulo} es la abertura entre dos semirrectas
            con el mismo origen; se mide en grados (°).
      \item Tipos: \textbf{agudo} ($< 90°$),
            \textbf{recto} ($= 90°$),
            \textbf{obtuso} ($90°$--$180°$),
            \textbf{llano} ($= 180°$).
      \item \textbf{Paralelas}: nunca se cortan ($r \parallel s$).
      \item \textbf{Perpendiculares}: se cortan en 90° ($r \perp s$).
      \item Las calles de Almería son \textbf{segmentos},
            no rectas, porque tienen inicio y final.
    \end{itemize}
  \end{resumenbox}
  \vspace{6pt}
  \begin{center}
    {\AlmFontSmall\color{AlmAzulM}
    \faArrowCircleRight\enskip
    \textbf{Próxima sesión:}
    Polígonos en la ciudad — triángulos y cuadriláteros en Almería.}
  \end{center}
\end{frame}

\end{document}




% % ================================================================
%  SESIÓN 06 — Figuras planas en la ciudad
%  Almería en Cifras y Formas · Matemáticas 1.º ESO
%  IES Al-Ándalus · 2025-2026
% ================================================================
\documentclass[aspectratio=169, 12pt, spanish]{beamer}
\usepackage{almeria-beamer}

\renewcommand{\SessionNumber}{06}
\renewcommand{\SessionTitle}{Figuras planas en la ciudad}
\renewcommand{\SessionName}{Sesión 06}
\renewcommand{\SessionObjective}{%
  Identificar y clasificar figuras planas (polígonos) en elementos
  arquitectónicos de Almería, reconociendo sus propiedades básicas
  (lados, vértices, ángulos).%
}
\renewcommand{\BlockName}{Bloque 2: Geometría urbana}

\begin{document}

% ---------------------------------------------------------------
%  PORTADA
% ---------------------------------------------------------------
\begin{frame}[plain]
  \titlepage
\end{frame}

% ---------------------------------------------------------------
%  ÍNDICE
% ---------------------------------------------------------------
\begin{frame}{Contenidos de hoy}
  \begin{columns}[t]
    \begin{column}{0.48\textwidth}
      \begin{itemize}
        \item Figura plana y polígono
        \item Clasificación de polígonos
        \item Propiedades: lados, vértices, ángulos
        \item Polígonos regulares e irregulares
        \item Cuadriláteros especiales
      \end{itemize}
    \end{column}
    \begin{column}{0.48\textwidth}
      \begin{itemize}
        \item Tipos de triángulos
        \item Suma de ángulos interiores
        \item Polígonos en Almería
        \item Ejercicios Bloom
        \item Evidencia del día
      \end{itemize}
    \end{column}
  \end{columns}
  \vspace{0.4cm}
  \begin{notabox}[Niveles de Bloom]
    \bloomtag{Recordar}\quad
    \bloomtag{Comprender}\quad
    \bloomtag{Aplicar}\quad
    \bloomtag{Analizar}
  \end{notabox}
\end{frame}

% ---------------------------------------------------------------
\seccionframe{Figuras planas y polígonos}
% ---------------------------------------------------------------

% ---------------------------------------------------------------
%  FRAME 3 — Figura plana y polígono
% ---------------------------------------------------------------
\begin{frame}{¿Qué es una figura plana?}
  \begin{columns}[c]
    \begin{column}{0.54\textwidth}
      \begin{definicionbox}{Figura plana}
        Figura geométrica que se puede dibujar completamente en
        un plano (en dos dimensiones: largo y ancho).
        No tiene profundidad.
      \end{definicionbox}
      \vspace{0.3cm}
      \begin{definicionbox}{Polígono}
        Figura plana \textbf{cerrada} formada por segmentos de recta
        (llamados \textbf{lados}). Los segmentos se unen en puntos
        llamados \textbf{vértices}.
      \end{definicionbox}
    \end{column}
    \begin{column}{0.42\textwidth}
      \begin{center}
      \begin{tikzpicture}[scale=0.85]
        % polígono ejemplo (pentágono irregular)
        \draw[zona] (0,0) -- (2,0) -- (2.6,1.4) -- (1.3,2.2) -- (-0.4,1.3) -- cycle;
        \foreach \p/\lbl/\pos in {
          (0,0)/A/below left,
          (2,0)/B/below right,
          (2.6,1.4)/C/right,
          (1.3,2.2)/D/above,
          (-0.4,1.3)/E/left}{
          \node[punto] at \p {};
          \node[\pos, font=\bfseries\scriptsize, color=AlmAzulM] at \p {\lbl};
        }
        % etiqueta lados
        \node[color=AlmTerra, font=\tiny\bfseries] at (1,-0.3) {lado};
        % etiqueta vértice
        \draw[->, AlmTerra, thin] (0.6,-0.5) -- (-0.05,-0.05);
        \node[color=AlmTerra, font=\tiny\bfseries] at (0.9,-0.55) {vértice};
      \end{tikzpicture}
      \end{center}
      \begin{notabox}[Clave]
        \scriptsize Un círculo \textbf{NO} es un polígono:
        no tiene lados rectos.
      \end{notabox}
    \end{column}
  \end{columns}
\end{frame}

% ---------------------------------------------------------------
%  FRAME 4 — Galería de polígonos (TikZ)
% ---------------------------------------------------------------
\begin{frame}{Galería de polígonos: de 3 a 8 lados}
  \begin{center}
  \begin{tikzpicture}[scale=0.78]
    % ---- Triángulo ----
    \draw[zona] (0,0) -- (1.2,0) -- (0.6,1.0) -- cycle;
    \node[below, font=\tiny\bfseries, color=AlmAzulM] at (0.6,-0.12) {Triángulo};
    \node[below, font=\tiny, color=AlmGris] at (0.6,-0.38) {3 lados / 180°};
    % ---- Cuadrilátero ----
    \draw[zona] (1.7,0) -- (3.1,0) -- (3.1,1.0) -- (1.7,1.0) -- cycle;
    \node[below, font=\tiny\bfseries, color=AlmAzulM] at (2.4,-0.12) {Cuadrilátero};
    \node[below, font=\tiny, color=AlmGris] at (2.4,-0.38) {4 lados / 360°};
    % ---- Pentágono ----
    \begin{scope}[shift={(4.0,0.5)}]
      \draw[zona] \foreach \i in {0,1,2,3,4}{
        -- ({cos(90+\i*72)*0.6},{sin(90+\i*72)*0.6})
      } -- cycle;
    \end{scope}
    \node[below, font=\tiny\bfseries, color=AlmAzulM] at (4.0,-0.12) {Pentágono};
    \node[below, font=\tiny, color=AlmGris] at (4.0,-0.38) {5 lados / 540°};
    % ---- Hexágono ----
    \begin{scope}[shift={(5.7,0.5)}]
      \draw[zona] \foreach \i in {0,1,2,3,4,5}{
        -- ({cos(90+\i*60)*0.6},{sin(90+\i*60)*0.6})
      } -- cycle;
    \end{scope}
    \node[below, font=\tiny\bfseries, color=AlmAzulM] at (5.7,-0.12) {Hexágono};
    \node[below, font=\tiny, color=AlmGris] at (5.7,-0.38) {6 lados / 720°};
    % ---- Heptágono ----
    \begin{scope}[shift={(7.4,0.5)}]
      \draw[zona] \foreach \i in {0,1,2,3,4,5,6}{
        -- ({cos(90+\i*51.43)*0.6},{sin(90+\i*51.43)*0.6})
      } -- cycle;
    \end{scope}
    \node[below, font=\tiny\bfseries, color=AlmAzulM] at (7.4,-0.12) {Heptágono};
    \node[below, font=\tiny, color=AlmGris] at (7.4,-0.38) {7 lados / 900°};
    % ---- Octógono ----
    \begin{scope}[shift={(9.1,0.5)}]
      \draw[zona] \foreach \i in {0,1,2,3,4,5,6,7}{
        -- ({cos(90+\i*45)*0.6},{sin(90+\i*45)*0.6})
      } -- cycle;
    \end{scope}
    \node[below, font=\tiny\bfseries, color=AlmAzulM] at (9.1,-0.12) {Octógono};
    \node[below, font=\tiny, color=AlmGris] at (9.1,-0.38) {8 lados / 1080°};
  \end{tikzpicture}
  \end{center}
  \vspace{0.15cm}
  \begin{teoriabox}
    \small\textbf{Regla clave:} en todo polígono, el número de lados
    $=$ número de vértices $=$ número de ángulos interiores.
    La suma de ángulos interiores de un polígono de $n$ lados es
    $(n-2)\times 180°$.
  \end{teoriabox}
\end{frame}

% ---------------------------------------------------------------
%  FRAME 5 — Regular vs. Irregular
% ---------------------------------------------------------------
\begin{frame}{Polígono regular vs. irregular}
  \begin{columns}[c]
    \begin{column}{0.48\textwidth}
      \begin{definicionbox}{Polígono regular}
        Todos los lados son \textbf{iguales}
        \textbf{y} todos los ángulos son \textbf{iguales}.
        \\[4pt]
        Ejemplos: triángulo equilátero, cuadrado, hexágono regular.
      \end{definicionbox}
      \vspace{0.3cm}
      \begin{center}
      \begin{tikzpicture}[scale=0.8]
        \begin{scope}[shift={(0,0)}]
          \draw[rectov, fill=AlmVerde!15] \foreach \i in {0,1,2}{
            -- ({cos(90+\i*120)*0.7},{sin(90+\i*120)*0.7})
          } -- cycle;
          \node[below, font=\scriptsize\bfseries, color=AlmVerde] at (0,-0.8) {Regular};
        \end{scope}
        \begin{scope}[shift={(1.9,0)}]
          \draw[rectov, fill=AlmVerde!15] (0,0) -- (1.2,0) -- (1.2,1.2) -- (0,1.2) -- cycle;
          \node[below, font=\scriptsize\bfseries, color=AlmVerde] at (0.6,-0.12) {Regular};
        \end{scope}
      \end{tikzpicture}
      \end{center}
    \end{column}
    \begin{column}{0.48\textwidth}
      \begin{definicionbox}{Polígono irregular}
        Los lados \textbf{o} los ángulos son diferentes entre sí.
        \\[4pt]
        Muchas figuras de la arquitectura real son irregulares.
      \end{definicionbox}
      \vspace{0.3cm}
      \begin{center}
      \begin{tikzpicture}[scale=0.8]
        \begin{scope}[shift={(0,0)}]
          \draw[rectot, fill=AlmTerra!12]
            (0,0) -- (1.4,0) -- (1.0,0.8) -- (0.3,1.1) -- cycle;
          \node[below, font=\scriptsize\bfseries, color=AlmTerra] at (0.7,-0.12) {Irregular};
        \end{scope}
        \begin{scope}[shift={(2.2,0)}]
          \draw[rectot, fill=AlmTerra!12]
            (0,0.3) -- (1.1,0) -- (1.3,0.9) -- (0.8,1.3) -- (0.1,1.0) -- cycle;
          \node[below, font=\scriptsize\bfseries, color=AlmTerra] at (0.65,-0.12) {Irregular};
        \end{scope}
      \end{tikzpicture}
      \end{center}
    \end{column}
  \end{columns}
\end{frame}

% ---------------------------------------------------------------
%  FRAME 6 — Cuadriláteros especiales
% ---------------------------------------------------------------
\begin{frame}{Cuadriláteros especiales}
  \begin{center}
  \begin{tikzpicture}[scale=0.78, every node/.style={font=\tiny}]
    % Rectángulo
    \draw[zona] (0,0) rectangle (1.8,1.0);
    \node[below=2pt] at (0.9,-0.0) {\textbf{Rectángulo}};
    \node[below=14pt, color=AlmGris] at (0.9,0) {4 ángulos rectos};
    \node[below=24pt, color=AlmGris] at (0.9,0) {lados opuestos =};
    % Cuadrado
    \draw[zona] (2.4,0) rectangle (3.6,1.2);
    \node[below=2pt] at (3.0,-0.0) {\textbf{Cuadrado}};
    \node[below=14pt, color=AlmGris] at (3.0,0) {4 ángulos rectos};
    \node[below=24pt, color=AlmGris] at (3.0,0) {4 lados iguales};
    % Rombo
    \draw[zona] (4.6,0.5) -- (5.4,0) -- (6.2,0.5) -- (5.4,1.0) -- cycle;
    \node[below=2pt] at (5.4,-0.0) {\textbf{Rombo}};
    \node[below=14pt, color=AlmGris] at (5.4,0) {4 lados iguales};
    \node[below=24pt, color=AlmGris] at (5.4,0) {ángulos ≠ 90°};
    % Trapecio
    \draw[zona] (7.0,0) -- (8.6,0) -- (8.2,1.0) -- (7.4,1.0) -- cycle;
    \node[below=2pt] at (7.8,-0.0) {\textbf{Trapecio}};
    \node[below=14pt, color=AlmGris] at (7.8,0) {1 par de lados};
    \node[below=24pt, color=AlmGris] at (7.8,0) {paralelos};
    % Paralelogramo
    \draw[zona] (9.4,0) -- (11.0,0) -- (10.6,1.0) -- (9.0,1.0) -- cycle;
    \node[below=2pt] at (10.0,-0.0) {\textbf{Paralelogramo}};
    \node[below=14pt, color=AlmGris] at (10.0,0) {2 pares paralelos};
    \node[below=24pt, color=AlmGris] at (10.0,0) {lados opuestos =};
  \end{tikzpicture}
  \end{center}
  \vspace{0.1cm}
  \begin{notabox}[Jerarquía]
    \small Cuadrado $\subset$ Rectángulo $\subset$ Paralelogramo $\subset$ Cuadrilátero
  \end{notabox}
\end{frame}

% ---------------------------------------------------------------
%  FRAME 7 — Tipos de triángulos
% ---------------------------------------------------------------
\begin{frame}{Tipos de triángulos}
  \begin{columns}[t]
    \begin{column}{0.50\textwidth}
      \textbf{\color{AlmAzulM}Por sus lados:}
      \vspace{0.2cm}
      \begin{center}
      \begin{tikzpicture}[scale=0.75]
        % Equilátero
        \draw[rectov, fill=AlmVerde!12]
          (0,0) -- (1.0,0) -- (0.5,0.87) -- cycle;
        \node[below, font=\tiny\bfseries, color=AlmVerde] at (0.5,-0.05) {Equilátero};
        \node[below, font=\tiny, color=AlmGris] at (0.5,-0.25) {3 lados iguales};
        % Isósceles
        \draw[zona]
          (1.8,0) -- (3.2,0) -- (2.5,1.1) -- cycle;
        \node[below, font=\tiny\bfseries, color=AlmAzulM] at (2.5,-0.05) {Isósceles};
        \node[below, font=\tiny, color=AlmGris] at (2.5,-0.25) {2 lados iguales};
        % Escaleno
        \draw[rectot, fill=AlmTerra!10]
          (4.0,0) -- (5.6,0) -- (5.2,1.0) -- cycle;
        \node[below, font=\tiny\bfseries, color=AlmTerra] at (4.9,-0.05) {Escaleno};
        \node[below, font=\tiny, color=AlmGris] at (4.9,-0.25) {todos distintos};
      \end{tikzpicture}
      \end{center}
    \end{column}
    \begin{column}{0.46\textwidth}
      \textbf{\color{AlmAzulM}Por sus ángulos:}
      \vspace{0.2cm}
      \begin{center}
      \begin{tikzpicture}[scale=0.75]
        % Acutángulo
        \draw[rectov, fill=AlmVerde!12]
          (0,0) -- (1.4,0) -- (0.8,1.0) -- cycle;
        \node[below, font=\tiny\bfseries, color=AlmVerde] at (0.7,-0.05) {Acutángulo};
        \node[below, font=\tiny, color=AlmGris] at (0.7,-0.25) {todos $<$ 90°};
        % Rectángulo
        \draw[zona]
          (1.9,0) -- (3.3,0) -- (1.9,1.0) -- cycle;
        \draw[rectoang] (1.9,0) -- (2.2,0) -- (2.2,0.3) -- (1.9,0.3);
        \node[below, font=\tiny\bfseries, color=AlmAzulM] at (2.6,-0.05) {Rectángulo};
        \node[below, font=\tiny, color=AlmGris] at (2.6,-0.25) {un ángulo $=90°$};
        % Obtusángulo
        \draw[rectot, fill=AlmTerra!10]
          (4.0,0.5) -- (5.6,0) -- (4.3,1.1) -- cycle;
        \node[below, font=\tiny\bfseries, color=AlmTerra] at (4.9,-0.05) {Obtusángulo};
        \node[below, font=\tiny, color=AlmGris] at (4.9,-0.25) {un ángulo $>90°$};
      \end{tikzpicture}
      \end{center}
    \end{column}
  \end{columns}
  \vspace{0.2cm}
  \begin{formulabox}
    Suma de ángulos de un triángulo $= \alpha + \beta + \gamma = 180°$
    \quad\quad
    Suma de ángulos de un cuadrilátero $= 360°$
  \end{formulabox}
\end{frame}

% ---------------------------------------------------------------
\seccionframe[BAnalizar]{Polígonos en Almería}
% ---------------------------------------------------------------

% ---------------------------------------------------------------
%  FRAME 8 — Almería y polígonos
% ---------------------------------------------------------------
\begin{frame}{Polígonos en la arquitectura almeriense}
  \begin{columns}[t]
    \begin{column}{0.50\textwidth}
      \begin{ejemplobox}
        \small
        \begin{itemize}
          \item \textbf{Catedral de Almería:} fachada con rectángulos,
                arcos y triángulos en las cornisas.
          \item \textbf{Plazas del centro:} cuadriláteros (cuadrados
                y trapecios) en el trazado urbano.
          \item \textbf{Torres de la Alcazaba:} polígonos irregulares
                en la planta de cada torre.
          \item \textbf{Puerto de Almería:} edificios con ventanas
                hexagonales regulares.
        \end{itemize}
      \end{ejemplobox}
    \end{column}
    \begin{column}{0.46\textwidth}
      \begin{center}
      \begin{tikzpicture}[scale=0.72]
        % Silueta estilizada de edificio con formas
        % Base rectangular
        \draw[zona] (0,0) rectangle (3.6,1.8);
        % Ventana rectangular
        \draw[rectov, fill=AlmAzulC!25] (0.3,0.4) rectangle (1.0,1.2);
        \draw[rectov, fill=AlmAzulC!25] (1.4,0.4) rectangle (2.1,1.2);
        \draw[rectov, fill=AlmAzulC!25] (2.5,0.4) rectangle (3.2,1.2);
        % Tejado triangular
        \draw[rectot, fill=AlmTerra!20] (0,1.8) -- (1.8,2.8) -- (3.6,1.8) -- cycle;
        % Hexágono (ventana especial)
        \node[font=\tiny\bfseries, color=AlmAzul] at (1.8,0.2) {Rectángulos};
        \node[font=\tiny\bfseries, color=AlmTerra] at (1.8,2.55) {Triángulo};
      \end{tikzpicture}
      \end{center}
      \begin{notabox}[Dato]
        \scriptsize La Catedral de Almería (s. XVI) es una de las
        pocas catedrales-fortaleza de España: sus muros combinan
        rectángulos y arcos de medio punto.
      \end{notabox}
    \end{column}
  \end{columns}
\end{frame}

% ---------------------------------------------------------------
\seccionframe[BRecordar]{Ejercicios — RECORDAR}
% ---------------------------------------------------------------

% ---------------------------------------------------------------
%  FRAME 9 — Ejercicio RECORDAR
% ---------------------------------------------------------------
\begin{frame}{Ejercicio 1 — RECORDAR \dificultad{1}}
  \begin{recordarbox}
    \begin{enumerate}
      \item ¿Cómo se llama el polígono de 5 lados?
      \item ¿Cuántos vértices tiene un hexágono?
      \item Completa la tabla:

      \vspace{0.2cm}
      \begin{center}
      \small
      \begin{tabular}{|>{\bfseries\color{AlmAzulM}}l|c|c|c|}
        \hline
        \rowcolor{AlmAzul!15}
        \textbf{Figura} & \textbf{Lados} & \textbf{Vértices} &
        \textbf{Suma ángulos} \\
        \hline
        Triángulo  & 3 & \phantom{000} & \phantom{000°} \\
        \hline
        Cuadrilátero & \phantom{0} & 4 & \phantom{000°} \\
        \hline
        Pentágono  & \phantom{0} & \phantom{0} & 540° \\
        \hline
        Hexágono   & \phantom{0} & \phantom{0} & \phantom{000°} \\
        \hline
        Octógono   & \phantom{0} & \phantom{0} & \phantom{000°} \\
        \hline
      \end{tabular}
      \end{center}
    \end{enumerate}
  \end{recordarbox}
  \vspace{0.15cm}
  \begin{notabox}[Pista]
    \small Suma de ángulos $= (n-2)\times 180°$, donde $n$ es el número de lados.
  \end{notabox}
\end{frame}

% ---------------------------------------------------------------
\seccionframe[BComprender]{Ejercicios — COMPRENDER}
% ---------------------------------------------------------------

% ---------------------------------------------------------------
%  FRAME 10 — Ejercicio COMPRENDER
% ---------------------------------------------------------------
\begin{frame}{Ejercicio 2 — COMPRENDER \dificultad{2}}
  \begin{comprenderbox}
    Una plaza de Almería tiene cuatro lados de longitudes:
    \textbf{25 m, 30 m, 25 m y 30 m}. Los cuatro ángulos son de 90°.
    \begin{enumerate}
      \item ¿Qué tipo de cuadrilátero es?
      \item ¿Es regular o irregular? Razona tu respuesta.
      \item ¿Qué condición adicional necesitaría para ser un cuadrado?
    \end{enumerate}
  \end{comprenderbox}
  \vspace{0.2cm}
  \begin{columns}[t]
    \begin{column}{0.48\textwidth}
      \begin{pasobox}[1]
        Identificar el tipo por sus propiedades:
        lados opuestos iguales + 4 ángulos rectos.
      \end{pasobox}
    \end{column}
    \begin{column}{0.48\textwidth}
      \begin{pasobox}[2]
        Comparar con la definición de polígono regular:
        ¿todos los lados son iguales?
      \end{pasobox}
    \end{column}
  \end{columns}
\end{frame}

% ---------------------------------------------------------------
\seccionframe[BAplicar]{Ejercicios — APLICAR}
% ---------------------------------------------------------------

% ---------------------------------------------------------------
%  FRAME 11 — Ejercicio APLICAR
% ---------------------------------------------------------------
\begin{frame}{Ejercicio 3 — APLICAR \dificultad{3}}
  \begin{aplicarbox}
    Identifica los polígonos que aparecen en estos 5 elementos
    de Almería y clasifícalos (regular/irregular):

    \vspace{0.15cm}
    \small
    \begin{tabular}{@{}p{3.2cm}p{3.5cm}p{2.8cm}@{}}
      \toprule
      \textbf{Elemento} & \textbf{Polígono(s)} & \textbf{Reg. / Irreg.} \\
      \midrule
      Fachada Catedral & \phantom{xxxxxxxxxx} & \phantom{xxxxx} \\
      Ventana hexagonal & \phantom{xxxxxxxxxx} & \phantom{xxxxx} \\
      Torre Alcazaba & \phantom{xxxxxxxxxx} & \phantom{xxxxx} \\
      Plaza del Ayuntamiento & \phantom{xxxxxxxxxx} & \phantom{xxxxx} \\
      Panel solar (tejado) & \phantom{xxxxxxxxxx} & \phantom{xxxxx} \\
      \bottomrule
    \end{tabular}
  \end{aplicarbox}
  \vspace{0.15cm}
  \begin{notabox}[Recuerda]
    \small Cuenta los lados visibles. Comprueba si todos los lados
    y ángulos son iguales para determinar si es regular.
  \end{notabox}
\end{frame}

% ---------------------------------------------------------------
\seccionframe[BAnalizar]{Ejercicios — ANALIZAR}
% ---------------------------------------------------------------

% ---------------------------------------------------------------
%  FRAME 12 — Ejercicio ANALIZAR
% ---------------------------------------------------------------
\begin{frame}{Ejercicio 4 — ANALIZAR \dificultad{4}}
  \begin{analizarbox}
    Compara un \textbf{cuadrado} y un \textbf{rectángulo}:
    \begin{enumerate}
      \item ¿Qué tienen en común (número de lados, vértices,
            suma de ángulos)?
      \item ¿En qué se diferencian?
      \item Si ves una figura con 4 ángulos rectos pero sin
            ninguna medida, ¿puedes saber si es cuadrado
            o rectángulo? ¿Por qué?
    \end{enumerate}
  \end{analizarbox}
  \vspace{0.2cm}
  \begin{columns}[c]
    \begin{column}{0.30\textwidth}
      \begin{center}
      \begin{tikzpicture}[scale=0.65]
        \draw[zona] (0,0) rectangle (1.4,1.4);
        \draw[rectoang] (0,0) -- (0.25,0) -- (0.25,0.25) -- (0,0.25);
        \node[below, font=\tiny\bfseries] at (0.7,-0.1) {Cuadrado};
        \node[below, font=\tiny, color=AlmGris] at (0.7,-0.35) {$l = l$};
      \end{tikzpicture}
      \end{center}
    \end{column}
    \begin{column}{0.30\textwidth}
      \begin{center}
      \begin{tikzpicture}[scale=0.65]
        \draw[zona] (0,0) rectangle (2.0,1.1);
        \draw[rectoang] (0,0) -- (0.25,0) -- (0.25,0.25) -- (0,0.25);
        \node[below, font=\tiny\bfseries] at (1.0,-0.1) {Rectángulo};
        \node[below, font=\tiny, color=AlmGris] at (1.0,-0.35) {$l \neq w$};
      \end{tikzpicture}
      \end{center}
    \end{column}
    \begin{column}{0.34\textwidth}
      \small\color{AlmAzulM}
      Ambos: 4 lados, 4 vértices, 360°. \\[4pt]
      Diferencia: ¿todos los lados iguales?
    \end{column}
  \end{columns}
\end{frame}

% ---------------------------------------------------------------
%  FRAME 13 — Error frecuente
% ---------------------------------------------------------------
\begin{frame}{¡Cuidado con estos errores!}
  \begin{errorbox}
    \begin{itemize}
      \item \incorrecto\ \textbf{«El círculo es un polígono»}
            \\
            \correcto\ Un círculo tiene un contorno \emph{curvo}, no lados rectos.
            \textbf{No es un polígono.}
      \vspace{0.3cm}
      \item \incorrecto\ \textbf{«El hexágono tiene 5 lados»}
            \\
            \correcto\ HEXA = 6. \quad
            \textbf{HEXÁgono $\rightarrow$ 6 lados}.
            \\
            Pentágono (PENTA $= 5$) es el de 5 lados.
      \vspace{0.3cm}
      \item \incorrecto\ \textbf{«Todo cuadrilátero es un cuadrado»}
            \\
            \correcto\ Solo si los 4 lados son iguales
            \textbf{y} los 4 ángulos son rectos.
    \end{itemize}
  \end{errorbox}
\end{frame}

% ---------------------------------------------------------------
%  FRAME 14 — Resumen
% ---------------------------------------------------------------
\begin{frame}{Resumen de la sesión}
  \begin{resumenbox}
    \begin{itemize}
      \item Un \textbf{polígono} es una figura plana cerrada con lados rectos.
      \item Se clasifica por número de lados: triángulo (3) hasta octógono (8).
      \item En todo polígono: \#lados $=$ \#vértices $=$ \#ángulos.
      \item \textbf{Regular}: todos los lados y ángulos iguales.
            \textbf{Irregular}: no se cumple.
      \item Triángulos por lados: equilátero, isósceles, escaleno.
      \item Triángulos por ángulos: acutángulo, rectángulo, obtusángulo.
      \item Suma de ángulos: triángulo $= 180°$, cuadrilátero $= 360°$,
            en general $(n-2)\times 180°$.
    \end{itemize}
  \end{resumenbox}
\end{frame}

% ---------------------------------------------------------------
%  FRAME 15 — Evidencia del día
% ---------------------------------------------------------------
\begin{frame}{Evidencia del día}
  \begin{evidenciabox}
    \textbf{Catálogo de figuras planas de Almería}

    \vspace{0.3cm}
    Elabora un catálogo con al menos \textbf{5 edificios o elementos}
    de Almería. Para cada uno indica:

    \vspace{0.15cm}
    \begin{enumerate}
      \item Nombre del edificio o elemento.
      \item Polígono(s) identificado(s).
      \item Número de lados y vértices.
      \item Regular o irregular — justifica con las propiedades.
      \item Dibuja el polígono (a mano o con regla y compás).
    \end{enumerate}

    \vspace{0.2cm}
    \textbf{Criterio de éxito:} los 5 polígonos están correctamente
    identificados, clasificados y dibujados.
  \end{evidenciabox}
\end{frame}

\end{document}

% % ================================================================
%  SESIÓN 07 — Perímetro de zonas urbanas
%  Almería en Cifras y Formas · Matemáticas 1.º ESO
%  IES Al-Ándalus · 2025-2026
% ================================================================
\documentclass[aspectratio=169, 12pt, spanish]{beamer}
\usepackage{almeria-beamer}

\renewcommand{\SessionNumber}{07}
\renewcommand{\SessionTitle}{Perímetro de zonas urbanas}
\renewcommand{\SessionName}{Sesión 07}
\renewcommand{\SessionObjective}{%
  Calcular el perímetro de diferentes figuras geométricas y aplicarlo
  a zonas urbanas de Almería modelizadas como polígonos.%
}
\renewcommand{\BlockName}{Bloque 2: Geometría urbana}

\begin{document}

% ---------------------------------------------------------------
%  PORTADA
% ---------------------------------------------------------------
\begin{frame}[plain]
  \titlepage
\end{frame}

% ---------------------------------------------------------------
%  ÍNDICE
% ---------------------------------------------------------------
\begin{frame}{Contenidos de hoy}
  \begin{columns}[t]
    \begin{column}{0.48\textwidth}
      \begin{itemize}
        \item Definición de perímetro
        \item Analogía: la valla del parque
        \item Fórmulas de perímetro
        \item Pasos para calcular
        \item Unidades de longitud
      \end{itemize}
    \end{column}
    \begin{column}{0.48\textwidth}
      \begin{itemize}
        \item Ejemplos en Almería
        \item Problema paso a paso
        \item Ejercicios Bloom
        \item Error frecuente
        \item Evidencia del día
      \end{itemize}
    \end{column}
  \end{columns}
  \vspace{0.4cm}
  \begin{notabox}[Niveles de Bloom]
    \bloomtag{Recordar}\quad
    \bloomtag{Comprender}\quad
    \bloomtag{Aplicar}
  \end{notabox}
\end{frame}

% ---------------------------------------------------------------
\seccionframe{¿Qué es el perímetro?}
% ---------------------------------------------------------------

% ---------------------------------------------------------------
%  FRAME 3 — Definición
% ---------------------------------------------------------------
\begin{frame}{Perímetro: el borde de la figura}
  \begin{columns}[c]
    \begin{column}{0.54\textwidth}
      \begin{definicionbox}{Perímetro}
        El \textbf{perímetro} de una figura plana es la
        \textbf{suma de las longitudes de todos sus lados}.
        Mide el \emph{contorno}, no el interior.
      \end{definicionbox}

      \vspace{0.4cm}
      \begin{notabox}[Analogía]
        \small El perímetro es como la \textbf{valla} que rodea un parque:
        mides cuánta valla necesitas, no cuánto césped hay dentro.
      \end{notabox}

      \vspace{0.3cm}
      \textbf{Unidades:} m, km, cm, mm \ldots
      \\
      \textbf{Nunca:} m$^2$, km$^2$ (eso sería área).
    \end{column}
    \begin{column}{0.42\textwidth}
      \begin{center}
      \begin{tikzpicture}[scale=0.80]
        % Parque (rectángulo)
        \draw[zona] (0,0) rectangle (3.2,2.0);
        % Valla (borde grueso destacado)
        \draw[AlmTerra, line width=3pt] (0,0) rectangle (3.2,2.0);
        % Césped interior
        \fill[AlmVerde!20] (0.1,0.1) rectangle (3.1,1.9);
        \node[font=\small\bfseries, color=AlmVerde] at (1.6,1.0) {ÁREA};
        \node[font=\small\bfseries, color=AlmTerra] at (1.6,-0.35) {PERÍMETRO (valla)};
        % Flechas de las dimensiones
        \draw[acota] (0,-0.6) -- (3.2,-0.6)
          node[midway, below, font=\scriptsize] {$l$};
        \draw[acota] (3.6,0) -- (3.6,2.0)
          node[midway, right, font=\scriptsize] {$w$};
      \end{tikzpicture}
      \end{center}
    \end{column}
  \end{columns}
\end{frame}

% ---------------------------------------------------------------
%  FRAME 4 — Fórmulas
% ---------------------------------------------------------------
\begin{frame}{Fórmulas de perímetro}
  \begin{columns}[t]
    \begin{column}{0.50\textwidth}
      \formuladest[AlmAzulM]{P_{\text{rect}} = 2l + 2w = 2(l+w)}
      \vspace{0.1cm}
      \formuladest[AlmVerde]{P_{\text{cuad}} = 4l}
      \vspace{0.1cm}
      \formuladest[AlmTerra]{P_{\text{triáng}} = a + b + c}
    \end{column}
    \begin{column}{0.46\textwidth}
      \formuladest[BAnalizar]{P_{\text{reg.}} = n \cdot s}
      \vspace{0.1cm}
      \begin{teoriabox}
        \small
        Donde $n$ = número de lados y $s$ = longitud de cada lado.\\[6pt]
        Para cualquier polígono:
        \[P = \sum_i l_i\]
        (suma de todos los lados).
      \end{teoriabox}
    \end{column}
  \end{columns}
\end{frame}

% ---------------------------------------------------------------
%  FRAME 5 — Pasos para calcular
% ---------------------------------------------------------------
\begin{frame}{Cómo calcular el perímetro: paso a paso}
  \begin{columns}[t]
    \begin{column}{0.50\textwidth}
      \begin{pasobox}[1]
        Identifica \textbf{todos los lados exteriores}
        de la figura. No olvides ninguno.
      \end{pasobox}
      \vspace{0.3cm}
      \begin{pasobox}[2]
        Asegúrate de que \textbf{todas las medidas}
        están en la \textbf{misma unidad}.
        Si hay metros y centímetros, convierte primero.
      \end{pasobox}
    \end{column}
    \begin{column}{0.46\textwidth}
      \begin{pasobox}[3]
        \textbf{Suma} todas las longitudes.
      \end{pasobox}
      \vspace{0.3cm}
      \begin{pasobox}[4]
        Escribe el resultado con su \textbf{unidad}
        (m, km, cm…).
        El perímetro es siempre una \emph{longitud}.
      \end{pasobox}
      \vspace{0.3cm}
      \begin{notabox}[Nunca olvides]
        \small La unidad del perímetro es de
        \textbf{una sola dimensión}: m, no m².
      \end{notabox}
    \end{column}
  \end{columns}
\end{frame}

% ---------------------------------------------------------------
\seccionframe[BAplicar]{Perímetros en Almería}
% ---------------------------------------------------------------

% ---------------------------------------------------------------
%  FRAME 6 — Ejemplos en Almería
% ---------------------------------------------------------------
\begin{frame}{Ejemplos en Almería}
  \begin{columns}[t]
    \begin{column}{0.50\textwidth}
      \begin{ejemplobox}
        \small
        \textbf{Parque Nicolás Salmerón}
        (rectángulo aprox. 400 m $\times$ 80 m):
        \[P = 2(400 + 80) = 2 \times 480 = \mathbf{960\ m}\]

        \vspace{0.15cm}
        \textbf{Plaza triangular hipotética}
        (lados 120 m, 85 m, 100 m):
        \[P = 120 + 85 + 100 = \mathbf{305\ m}\]

        \vspace{0.15cm}
        \textbf{Plaza cuadrada} (lado 60 m):
        \[P = 4 \times 60 = \mathbf{240\ m}\]
      \end{ejemplobox}
    \end{column}
    \begin{column}{0.46\textwidth}
      \begin{center}
      \begin{tikzpicture}[scale=0.70]
        % Rectángulo (parque)
        \draw[zona] (0,0) rectangle (4.0,0.8);
        \node[above, font=\tiny\bfseries, color=AlmAzulM] at (2.0,0.8) {Parque};
        \draw[acota] (0,-0.35) -- (4.0,-0.35)
          node[midway,below,font=\tiny]{400 m};
        \draw[acota] (4.3,0) -- (4.3,0.8)
          node[midway,right,font=\tiny]{80 m};
        % Triángulo
        \draw[zonat] (0,1.5) -- (2.2,1.5) -- (1.5,2.5) -- cycle;
        \node[above, font=\tiny\bfseries, color=AlmTerra] at (1.1,2.5) {Plaza $\blacktriangle$};
        \node[below, font=\tiny, color=AlmGris] at (1.1,1.45) {120 m};
        % Cuadrado
        \draw[rectov, fill=AlmVerde!15] (2.8,1.5) rectangle (4.0,2.7);
        \node[above, font=\tiny\bfseries, color=AlmVerde] at (3.4,2.7) {Plaza $\square$};
        \node[below, font=\tiny, color=AlmGris] at (3.4,1.45) {60 m};
      \end{tikzpicture}
      \end{center}
    \end{column}
  \end{columns}
\end{frame}

% ---------------------------------------------------------------
%  FRAME 7 — Problema completo paso a paso
% ---------------------------------------------------------------
\begin{frame}{Problema resuelto: manzana de Almería}
  \begin{definicionbox}{Enunciado}
    \small Una manzana rectangular del centro de Almería mide
    \textbf{180 m} de largo y \textbf{120 m} de ancho.
    Calcula su perímetro y el tiempo necesario para darle
    la vuelta caminando a 5 km/h.
  \end{definicionbox}

  \vspace{0.2cm}
  \begin{columns}[t]
    \begin{column}{0.48\textwidth}
      \begin{pasobox}[1]
        \small Identificar la fórmula:
        \[P = 2(l + w)\]
      \end{pasobox}
      \vspace{0.2cm}
      \begin{pasobox}[2]
        \small Sustituir valores:
        \[P = 2(180 + 120) = 2 \times 300 = \mathbf{600\ m}\]
      \end{pasobox}
    \end{column}
    \begin{column}{0.48\textwidth}
      \begin{pasobox}[3]
        \small Convertir a km para calcular el tiempo:
        \[600\ \text{m} = 0{,}6\ \text{km}\]
      \end{pasobox}
      \vspace{0.2cm}
      \begin{pasobox}[4]
        \small Calcular el tiempo:
        \[t = \frac{d}{v} = \frac{0{,}6\ \text{km}}{5\ \text{km/h}}
          = 0{,}12\ \text{h} = \mathbf{7{,}2\ \text{min}}\]
      \end{pasobox}
    \end{column}
  \end{columns}
\end{frame}

% ---------------------------------------------------------------
%  FRAME 8 — Comparativa de perímetros (PGFPlots)
% ---------------------------------------------------------------
\begin{frame}{Comparativa: perímetros de zonas de Almería}
  \begin{center}
  \begin{tikzpicture}
    \begin{axis}[
      almeria bar,
      height=5.5cm, width=10cm,
      ymin=0, ymax=1100,
      ylabel={Perímetro (m)},
      symbolic x coords={Parque Salmerón, Plaza cuadrada, Plaza triangular},
      xtick=data,
      xticklabel style={font=\scriptsize, align=center, text width=2.5cm},
      ytick={0,200,400,600,800,1000},
      bar width=1.2cm,
    ]
      \addplot[fill=AlmAzulM, draw=AlmAzul] coordinates {
        (Parque Salmerón, 960)
        (Plaza cuadrada, 240)
        (Plaza triangular, 305)
      };
    \end{axis}
  \end{tikzpicture}
  \end{center}
  \vspace{-0.2cm}
  \begin{notabox}[Observa]
    \small El Parque Nicolás Salmerón tiene el mayor perímetro (960 m)
    porque sus dimensiones totales son mucho mayores, aunque su forma
    sea rectangular.
  \end{notabox}
\end{frame}

% ---------------------------------------------------------------
\seccionframe[BRecordar]{Ejercicios — RECORDAR}
% ---------------------------------------------------------------

% ---------------------------------------------------------------
%  FRAME 9 — Ejercicio RECORDAR
% ---------------------------------------------------------------
\begin{frame}{Ejercicio 1 — RECORDAR \dificultad{1}}
  \begin{recordarbox}
    Calcula el perímetro de cada figura:
    \begin{enumerate}
      \item Rectángulo de \textbf{8 cm} de largo y \textbf{5 cm} de ancho.
      \item Cuadrado de lado \textbf{7 m}.
      \item Triángulo con lados \textbf{6 m}, \textbf{8 m} y \textbf{10 m}.
    \end{enumerate}
  \end{recordarbox}

  \vspace{0.25cm}
  \begin{columns}[t]
    \begin{column}{0.32\textwidth}
      \begin{center}
      \begin{tikzpicture}[scale=0.7]
        \draw[zona] (0,0) rectangle (2.4,1.5);
        \draw[acota] (0,-0.45) -- (2.4,-0.45)
          node[midway,below,font=\tiny]{8 cm};
        \draw[acota] (2.75,0) -- (2.75,1.5)
          node[midway,right,font=\tiny]{5 cm};
      \end{tikzpicture}
      \end{center}
    \end{column}
    \begin{column}{0.32\textwidth}
      \begin{center}
      \begin{tikzpicture}[scale=0.7]
        \draw[rectov, fill=AlmVerde!15] (0,0) rectangle (1.8,1.8);
        \draw[acota] (0,-0.45) -- (1.8,-0.45)
          node[midway,below,font=\tiny]{7 m};
      \end{tikzpicture}
      \end{center}
    \end{column}
    \begin{column}{0.32\textwidth}
      \begin{center}
      \begin{tikzpicture}[scale=0.7]
        \draw[zonat] (0,0) -- (2.4,0) -- (1.8,1.5) -- cycle;
        \node[below,font=\tiny,color=AlmGris] at (1.2,-0.1) {6 m};
        \node[right,font=\tiny,color=AlmGris] at (2.1,0.75) {8 m};
        \node[left,font=\tiny,color=AlmGris] at (0.9,0.75) {10 m};
      \end{tikzpicture}
      \end{center}
    \end{column}
  \end{columns}
\end{frame}

% ---------------------------------------------------------------
\seccionframe[BComprender]{Ejercicios — COMPRENDER}
% ---------------------------------------------------------------

% ---------------------------------------------------------------
%  FRAME 10 — Ejercicio COMPRENDER
% ---------------------------------------------------------------
\begin{frame}{Ejercicio 2 — COMPRENDER \dificultad{2}}
  \begin{comprenderbox}
    Dos parques de Almería:
    \begin{itemize}
      \item \textbf{Parque A:} rectángulo de \textbf{200 m $\times$ 100 m}.
      \item \textbf{Parque B:} cuadrado de lado \textbf{150 m}.
    \end{itemize}
    \begin{enumerate}
      \item Calcula el perímetro de cada parque.
      \item ¿Cuál tiene mayor perímetro?
      \item ¿Por qué el área \textbf{no} es relevante para saber
            cuánta valla se necesita?
    \end{enumerate}
  \end{comprenderbox}
  \vspace{0.2cm}
  \begin{columns}[c]
    \begin{column}{0.46\textwidth}
      \begin{tikzpicture}[scale=0.7]
        \draw[zona] (0,0) rectangle (3.2,1.6);
        \node[font=\scriptsize\bfseries, color=AlmAzulM] at (1.6,0.8) {Parque A};
        \draw[acota] (0,-0.4) -- (3.2,-0.4)
          node[midway,below,font=\tiny]{200 m};
        \draw[acota] (3.55,0) -- (3.55,1.6)
          node[midway,right,font=\tiny]{100 m};
      \end{tikzpicture}
    \end{column}
    \begin{column}{0.46\textwidth}
      \begin{tikzpicture}[scale=0.7]
        \draw[rectov, fill=AlmVerde!15] (0,0) rectangle (2.4,2.4);
        \node[font=\scriptsize\bfseries, color=AlmVerde] at (1.2,1.2) {Parque B};
        \draw[acota] (0,-0.4) -- (2.4,-0.4)
          node[midway,below,font=\tiny]{150 m};
      \end{tikzpicture}
    \end{column}
  \end{columns}
\end{frame}

% ---------------------------------------------------------------
\seccionframe[BAplicar]{Ejercicios — APLICAR}
% ---------------------------------------------------------------

% ---------------------------------------------------------------
%  FRAME 11 — Ejercicio APLICAR
% ---------------------------------------------------------------
\begin{frame}{Ejercicio 3 — APLICAR \dificultad{3}}
  \begin{aplicarbox}
    \small El \textbf{Parque de las Familias} de Almería tiene forma
    de polígono irregular de 6 lados con las medidas:
    \textbf{250 m, 180 m, 300 m, 200 m, 150 m, 220 m}.

    \begin{enumerate}
      \item Calcula su perímetro total.
      \item Rodear el parque contabiliza como \textbf{1 vuelta}.
            ¿Cuántas vueltas completas caben en \textbf{2 km}?
            ¿Sobra algún metro?
    \end{enumerate}
  \end{aplicarbox}

  \vspace{0.2cm}
  \begin{columns}[c]
    \begin{column}{0.52\textwidth}
      \begin{center}
      \begin{tikzpicture}[scale=0.55]
        \draw[zona]
          (0,1.5) -- (2.5,0) -- (5.5,0.3) --
          (6.5,2.0) -- (5.0,3.5) -- (1.5,3.2) -- cycle;
        \node[font=\tiny, color=AlmGris] at (1.25,0.55) {250 m};
        \node[font=\tiny, color=AlmGris] at (4.0,-0.25) {300 m};
        \node[font=\tiny, color=AlmGris] at (6.25,1.15) {180 m};
        \node[font=\tiny, color=AlmGris] at (5.95,2.85) {200 m};
        \node[font=\tiny, color=AlmGris] at (3.15,3.55) {150 m};
        \node[font=\tiny, color=AlmGris] at (0.45,2.45) {220 m};
        \node[font=\scriptsize\bfseries, color=AlmAzulM] at (3.25,1.75) {Parque};
      \end{tikzpicture}
      \end{center}
    \end{column}
    \begin{column}{0.44\textwidth}
      \begin{pasobox}[1]
        \scriptsize $P = 250 + 180 + 300 + 200 + 150 + 220 = ?$
      \end{pasobox}
      \vspace{0.2cm}
      \begin{pasobox}[2]
        \scriptsize Convertir 2 km a metros: $2\ \text{km} = 2000\ \text{m}$
        \\Vueltas $= 2000 \div P$
      \end{pasobox}
    \end{column}
  \end{columns}
\end{frame}

% ---------------------------------------------------------------
%  FRAME 12 — Error frecuente
% ---------------------------------------------------------------
\begin{frame}{¡Cuidado con este error!}
  \begin{errorbox}
    \begin{itemize}
      \item \incorrecto\
            \textbf{«El perímetro del parque es 960 metros cuadrados»}
            \\[6pt]
            \correcto\
            El perímetro es una \textbf{longitud} (1 dimensión),
            su unidad es \textbf{m, km, cm…}
            \\
            Nunca se mide en m$^2$ ni cm$^2$
            (eso sería \emph{área}).
      \vspace{0.4cm}
      \item \incorrecto\
            \textbf{Sumar lados con distintas unidades sin convertir:}
            \\
            $P = 80\ \text{cm} + 1{,}2\ \text{m} + 0{,}5\ \text{m}$
            sin pasar todo a la misma unidad.
            \\[6pt]
            \correcto\
            Convierte primero: $80\ \text{cm} = 0{,}8\ \text{m}$,
            \\
            $P = 0{,}8 + 1{,}2 + 0{,}5 = \mathbf{2{,}5\ \text{m}}$
    \end{itemize}
  \end{errorbox}
\end{frame}

% ---------------------------------------------------------------
%  FRAME 13 — Resumen
% ---------------------------------------------------------------
\begin{frame}{Resumen de la sesión}
  \begin{resumenbox}
    \begin{itemize}
      \item El \textbf{perímetro} es la suma de todos los lados:
            mide el \emph{contorno}, como una valla.
      \item Fórmulas clave:
        $P_{\text{rect}} = 2(l+w)$, \quad
        $P_{\text{cuad}} = 4l$, \quad
        $P_{\text{triáng}} = a+b+c$, \quad
        $P_{\text{reg}} = n\cdot s$.
      \item 4 pasos: identificar lados → unidad única →
            sumar → escribir con unidad.
      \item Unidad del perímetro: \textbf{longitud} (m, km, cm…),
            \textbf{nunca} m$^2$.
      \item Parque Salmerón: $P \approx 960\ \text{m}$.
    \end{itemize}
  \end{resumenbox}
\end{frame}

% ---------------------------------------------------------------
%  FRAME 14 — Evidencia del día
% ---------------------------------------------------------------
\begin{frame}{Evidencia del día}
  \begin{evidenciabox}
    \textbf{Ficha de Perímetros de Almería}

    \vspace{0.2cm}
    Elige \textbf{3 zonas} de Almería (parques, plazas o manzanas
    del callejero). Para cada una:
    \begin{enumerate}
      \item Dibuja la forma aproximada con las medidas.
      \item Muestra \textbf{todos los pasos} del cálculo
            (no solo el resultado final).
      \item Escribe el resultado con su unidad correcta.
      \item Calcula cuánto tiempo tardarías en rodearla
            caminando a 4 km/h.
    \end{enumerate}
    \vspace{0.2cm}
    \textbf{Criterio de éxito:} los 3 perímetros están bien calculados,
    los pasos son visibles y las unidades son correctas.
  \end{evidenciabox}
\end{frame}

\end{document}

% % ================================================================
%  SESIÓN 08 — Área de zonas urbanas
%  Almería en Cifras y Formas · Matemáticas 1.º ESO
%  IES Al-Ándalus · 2025-2026
% ================================================================
\documentclass[aspectratio=169, 10pt, spanish]{beamer}
\usepackage{almeria-beamer}

\renewcommand{\SessionNumber}{08}
\renewcommand{\SessionTitle}{Área de zonas urbanas}
\renewcommand{\SessionName}{Sesión 08}
\renewcommand{\SessionObjective}{%
  Calcular el área de figuras planas básicas y aplicarla a zonas verdes
  y plazas de Almería, interpretando correctamente las unidades
  de superficie.%
}
\renewcommand{\BlockName}{Bloque 2: Geometría urbana}

\begin{document}

% ---------------------------------------------------------------
%  PORTADA
% ---------------------------------------------------------------
\begin{frame}[plain]
  \titlepage
\end{frame}

% ---------------------------------------------------------------
%  ÍNDICE
% ---------------------------------------------------------------
\begin{frame}{Contenidos de hoy}
  \begin{columns}[t]
    \begin{column}{0.48\textwidth}
      \begin{itemize}
        \item Definición de área
        \item Analogía: el césped del parque
        \item Unidades de superficie
        \item Fórmulas de las 6 figuras básicas
        \item Método de cuadrícula
      \end{itemize}
    \end{column}
    \begin{column}{0.48\textwidth}
      \begin{itemize}
        \item Conversiones de unidades
        \item Ejemplos en Almería
        \item Problema paso a paso
        \item Ejercicios Bloom
        \item Evidencia del día
      \end{itemize}
    \end{column}
  \end{columns}
  \vspace{0.4cm}
  \begin{notabox}[Niveles de Bloom]
    \bloomtag{Recordar}\quad
    \bloomtag{Comprender}\quad
    \bloomtag{Aplicar}
  \end{notabox}
\end{frame}

% ---------------------------------------------------------------
\seccionframe{¿Qué es el área?}
% ---------------------------------------------------------------

% ---------------------------------------------------------------
%  FRAME 3 — Definición
% ---------------------------------------------------------------
\begin{frame}{Área: la superficie interior}
  \begin{columns}[c]
    \begin{column}{0.54\textwidth}
      \begin{definicionbox}{Área}
        El \textbf{área} de una figura es la medida de la
        \textbf{superficie encerrada} en su interior.
        Indica cuánto espacio ocupa la figura en el plano.
      \end{definicionbox}

      \vspace{0.35cm}
      \begin{notabox}[Analogía]
        \AlmFontSmall El área es como el \textbf{césped} a segar dentro
        del parque: mides cuánto verde hay,
        no la longitud de la valla que lo rodea.
      \end{notabox}

      \vspace{0.3cm}
      \textbf{Unidades:} m$^2$, km$^2$, cm$^2$, mm$^2$, ha, a \\
      \AlmFontSmall (dos dimensiones $\Rightarrow$ el exponente es 2)
    \end{column}
    \begin{column}{0.42\textwidth}
      \begin{center}
      \begin{tikzpicture}[scale=0.80]
        % Parque con cuadrícula de área
        \draw[AlmAzulC!30, very thin, step=0.5]
          (0,0) grid (3.5,2.5);
        \draw[AlmTerra, line width=2.5pt] (0,0) rectangle (3.5,2.5);
        \fill[AlmVerde!22] (0,0) rectangle (3.5,2.5);
        \draw[AlmAzulC!30, very thin, step=0.5]
          (0,0) grid (3.5,2.5);
        \node[font=\AlmFontSmall\bfseries, color=AlmVerde!70!black]
          at (1.75,1.25) {ÁREA};
        \node[font=\AlmFontTiny\bfseries, color=AlmTerra]
          at (1.75,-0.3) {Perímetro (borde)};
        % Cota
        \draw[acota] (0,-0.65) -- (3.5,-0.65)
          node[midway,below,font=\AlmFontTiny]{base $b$};
        \draw[acota] (3.9,0) -- (3.9,2.5)
          node[midway,right,font=\AlmFontTiny]{altura $h$};
      \end{tikzpicture}
      \end{center}
    \end{column}
  \end{columns}
\end{frame}

% ---------------------------------------------------------------
%  FRAME 4 — Unidades de superficie
% ---------------------------------------------------------------
\begin{frame}{Unidades de superficie y conversiones}
  \begin{columns}[c]
    \begin{column}{0.54\textwidth}
      \begin{center}
      \AlmFontSmall
      \begin{tabular}{>{\bfseries\color{AlmAzulM}}l l l}
        \toprule
        Unidad & Símbolo & Equivalencia \\
        \midrule
        milímetro cuadrado & mm$^2$ & \\
        centímetro cuadrado & cm$^2$ & 1 cm$^2$ = 100 mm$^2$ \\
        metro cuadrado & m$^2$ & 1 m$^2$ = 10{,}000 cm$^2$ \\
        área & a & 1 a = 100 m$^2$ \\
        hectárea & ha & 1 ha = 10{,}000 m$^2$ \\
        kilómetro cuadrado & km$^2$ & 1 km$^2$ = 1{,}000{,}000 m$^2$ \\
        \bottomrule
      \end{tabular}
      \end{center}
    \end{column}
    \begin{column}{0.42\textwidth}
      \begin{formulabox}[AlmAmbar!60!black]
        \AlmFontSmall
        1 ha = 10{,}000 m$^2$\\[4pt]
        1 km$^2$ = 100 ha
      \end{formulabox}
      \vspace{0.3cm}
      \begin{notabox}[Uso práctico]
        \AlmFontSmall
        \begin{itemize}
          \item Parcelas/campos: \textbf{ha}
          \item Pisos/plazas: \textbf{m$^2$}
          \item Comunidades/países: \textbf{km$^2$}
        \end{itemize}
      \end{notabox}
    \end{column}
  \end{columns}
\end{frame}

% ---------------------------------------------------------------
%  FRAME 5 — Fórmulas: cuadrado y rectángulo
% ---------------------------------------------------------------
\begin{frame}{Fórmulas de área (I): cuadrado y rectángulo}
  \begin{columns}[c]
    \begin{column}{0.50\textwidth}
      \begin{center}
      \begin{tikzpicture}[scale=0.80]
        \fill[AlmVerde!15] (0,0) rectangle (2.0,2.0);
        \draw[rectov, line width=1.5pt] (0,0) rectangle (2.0,2.0);
        \draw[acota] (0,-0.5) -- (2.0,-0.5)
          node[midway,below,font=\AlmFontScript]{$l$};
        \draw[acota] (2.4,0) -- (2.4,2.0)
          node[midway,right,font=\AlmFontScript]{$l$};
        \node[font=\AlmFontSmall\bfseries, color=AlmVerde!70!black]
          at (1.0,1.0) {$A = l^2$};
        \node[below, font=\AlmFontScript\bfseries, color=AlmVerde]
          at (1.0,-0.9) {Cuadrado};
      \end{tikzpicture}
      \end{center}
      \formuladest[AlmVerde]{A_{\text{cuad}} = l^2}
    \end{column}
    \begin{column}{0.46\textwidth}
      \begin{center}
      \begin{tikzpicture}[scale=0.80]
        \fill[AlmAzulC!18] (0,0) rectangle (2.8,1.6);
        \draw[zona, line width=1.5pt] (0,0) rectangle (2.8,1.6);
        \draw[acota] (0,-0.5) -- (2.8,-0.5)
          node[midway,below,font=\AlmFontScript]{$b$};
        \draw[acota] (3.2,0) -- (3.2,1.6)
          node[midway,right,font=\AlmFontScript]{$h$};
        \node[font=\AlmFontSmall\bfseries, color=AlmAzulM]
          at (1.4,0.8) {$A = b \times h$};
        \node[below, font=\AlmFontScript\bfseries, color=AlmAzulM]
          at (1.4,-0.9) {Rectángulo};
      \end{tikzpicture}
      \end{center}
      \formuladest[AlmAzulM]{A_{\text{rect}} = b \times h}
    \end{column}
  \end{columns}
\end{frame}

% ---------------------------------------------------------------
%  FRAME 6 — Fórmulas: triángulo y paralelogramo
% ---------------------------------------------------------------
\begin{frame}{Fórmulas de área (II): triángulo y paralelogramo}
  \begin{columns}[c]
    \begin{column}{0.50\textwidth}
      \begin{center}
      \begin{tikzpicture}[scale=0.80]
        \fill[AlmTerra!12] (0,0) -- (3.0,0) -- (2.0,1.8) -- cycle;
        \draw[rectot, line width=1.5pt] (0,0) -- (3.0,0) -- (2.0,1.8) -- cycle;
        % Altura
        \draw[AlmGris, dashed, thin] (2.0,0) -- (2.0,1.8);
        \draw[acota] (-0.4,0) -- (-0.4,1.8)
          node[midway,left,font=\AlmFontScript]{$h$};
        \draw[acota] (0,-0.5) -- (3.0,-0.5)
          node[midway,below,font=\AlmFontScript]{$b$};
        \node[font=\AlmFontSmall\bfseries, color=AlmTerra]
          at (1.5,0.6) {$\dfrac{b \times h}{2}$};
        \node[below, font=\AlmFontScript\bfseries, color=AlmTerra]
          at (1.5,-0.9) {Triángulo};
      \end{tikzpicture}
      \end{center}
      \formuladest[AlmTerra]{A_{\text{trián}} = \dfrac{b \times h}{2}}
    \end{column}
    \begin{column}{0.46\textwidth}
      \begin{center}
      \begin{tikzpicture}[scale=0.80]
        \fill[BAnalizar!10]
          (0.5,0) -- (3.1,0) -- (2.6,1.6) -- (0.0,1.6) -- cycle;
        \draw[rectov, line width=1.5pt]
          (0.5,0) -- (3.1,0) -- (2.6,1.6) -- (0.0,1.6) -- cycle;
        \draw[AlmGris, dashed, thin] (0.5,0) -- (0.5,1.6);
        \draw[acota] (-0.4,0) -- (-0.4,1.6)
          node[midway,left,font=\AlmFontScript]{$h$};
        \draw[acota] (0.5,-0.5) -- (3.1,-0.5)
          node[midway,below,font=\AlmFontScript]{$b$};
        \node[font=\AlmFontSmall\bfseries, color=AlmVerde]
          at (1.6,0.8) {$A = b \times h$};
        \node[below, font=\AlmFontScript\bfseries, color=AlmVerde]
          at (1.6,-0.9) {Paralelogramo};
      \end{tikzpicture}
      \end{center}
      \formuladest[BAnalizar]{A_{\text{paral}} = b \times h}
    \end{column}
  \end{columns}
\end{frame}

% ---------------------------------------------------------------
%  FRAME 7 — Fórmulas: trapecio y círculo
% ---------------------------------------------------------------
\begin{frame}{Fórmulas de área (III): trapecio y círculo}
  \begin{columns}[c]
    \begin{column}{0.50\textwidth}
      \begin{center}
      \begin{tikzpicture}[scale=0.80]
        \fill[AlmAmbar!15]
          (0.4,0) -- (3.2,0) -- (2.8,1.6) -- (0.8,1.6) -- cycle;
        \draw[AlmAmbar!70!black, line width=1.5pt]
          (0.4,0) -- (3.2,0) -- (2.8,1.6) -- (0.8,1.6) -- cycle;
        \draw[acota] (0.4,-0.5) -- (3.2,-0.5)
          node[midway,below,font=\AlmFontTiny]{$B$ (base mayor)};
        \draw[acota] (0.8,1.9) -- (2.8,1.9)
          node[midway,above,font=\AlmFontTiny]{$b$ (base menor)};
        \draw[AlmGris, dashed, thin] (0.8,0) -- (0.8,1.6);
        \draw[acota] (-0.35,0) -- (-0.35,1.6)
          node[midway,left,font=\AlmFontTiny]{$h$};
        \node[below, font=\AlmFontScript\bfseries, color=AlmAmbar!80!black]
          at (1.8,-0.9) {Trapecio};
      \end{tikzpicture}
      \end{center}
      \formuladest[AlmAmbar!70!black]{A_{\text{trap}} =
        \dfrac{(B+b)}{2} \times h}
    \end{column}
    \begin{column}{0.46\textwidth}
      \begin{center}
      \begin{tikzpicture}[scale=0.80]
        \fill[AlmAzulC!20] (1.5,1.5) circle (1.4);
        \draw[AlmAzulM, line width=1.5pt] (1.5,1.5) circle (1.4);
        \draw[AlmTerra, line width=1pt] (1.5,1.5) -- (2.9,1.5)
          node[midway,above,font=\AlmFontScript]{$r$};
        \node[font=\AlmFontSmall\bfseries, color=AlmAzulM]
          at (1.5,1.5) {$\pi r^2$};
        \node[below, font=\AlmFontScript\bfseries, color=AlmAzulM]
          at (1.5,0.0) {Círculo};
      \end{tikzpicture}
      \end{center}
      \formuladest[AlmAzulM]{A_{\text{círc}} = \pi \times r^2}
    \end{column}
  \end{columns}
\end{frame}

% ---------------------------------------------------------------
\seccionframe[BAplicar]{Áreas en Almería}
% ---------------------------------------------------------------

% ---------------------------------------------------------------
%  FRAME 8 — Ejemplos en Almería
% ---------------------------------------------------------------
\begin{frame}{Áreas de zonas verdes de Almería}
  \begin{columns}[t]
    \begin{column}{0.54\textwidth}
      \begin{ejemplobox}
        \AlmFontSmall
        \textbf{Parque Nicolás Salmerón}
        (rectángulo 400 m $\times$ 80 m):
        \[A = 400 \times 80 = \mathbf{32{,}000\ m^2} = \mathbf{3{,}2\ ha}\]

        \vspace{0.1cm}
        \textbf{Plaza triangular} (base 120 m, altura 95 m):
        \[A = \frac{120 \times 95}{2} = \frac{11{,}400}{2}
          = \mathbf{5{,}700\ m^2}\]

        \vspace{0.1cm}
        \textbf{Alcazaba de Almería}
        (área aproximada):
        \[A \approx \mathbf{43{,}000\ m^2} = \mathbf{4{,}3\ ha}\]
      \end{ejemplobox}
    \end{column}
    \begin{column}{0.42\textwidth}
      \begin{notabox}[¿Qué es una hectárea?]
        \AlmFontSmall
        1 ha = 10{,}000 m$^2$.
        \\[4pt]
        Equivale a un cuadrado de
        \textbf{100 m} $\times$ \textbf{100 m}.
        \\[4pt]
        Se usa en agricultura, urbanismo
        y cartografía.
        \\[4pt]
        La Alcazaba ocupa más de 4 campos
        de fútbol (un campo $\approx$ 0{,}7 ha).
      \end{notabox}
    \end{column}
  \end{columns}
\end{frame}

% ---------------------------------------------------------------
%  FRAME 9 — Problema paso a paso
% ---------------------------------------------------------------
\begin{frame}{Problema resuelto: manzana trapezoidal}
  \begin{definicionbox}{Enunciado}
    \AlmFontSmall Una manzana de Almería tiene forma de trapecio con
    bases paralelas de \textbf{200 m} y \textbf{140 m},
    y altura \textbf{90 m}.
    Calcula su área en m$^2$ y en hectáreas.
  \end{definicionbox}

  \vspace{0.15cm}
  \begin{columns}[t]
    \begin{column}{0.50\textwidth}
      \begin{pasobox}[1]
        \AlmFontSmall Identificar la fórmula del trapecio:
        \[A = \frac{(B + b)}{2} \times h\]
      \end{pasobox}
      \vspace{0.2cm}
      \begin{pasobox}[2]
        \AlmFontSmall Sustituir valores:
        \[A = \frac{(200 + 140)}{2} \times 90
          = \frac{340}{2} \times 90\]
      \end{pasobox}
    \end{column}
    \begin{column}{0.46\textwidth}
      \begin{pasobox}[3]
        \AlmFontSmall Calcular:
        \[A = 170 \times 90 = \mathbf{15{,}300\ m^2}\]
      \end{pasobox}
      \vspace{0.2cm}
      \begin{pasobox}[4]
        \AlmFontSmall Convertir a hectáreas:
        \[A = \frac{15{,}300}{10{,}000} = \mathbf{1{,}53\ ha}\]
      \end{pasobox}
    \end{column}
  \end{columns}
\end{frame}

% ---------------------------------------------------------------
\seccionframe[BRecordar]{Ejercicios — RECORDAR}
% ---------------------------------------------------------------

% ---------------------------------------------------------------
%  FRAME 10 — Ejercicio RECORDAR
% ---------------------------------------------------------------
\begin{frame}{Ejercicio 1 — RECORDAR \dificultad{1}}
  \begin{recordarbox}
    \begin{enumerate}
      \item Escribe la fórmula del área de:
            rectángulo, triángulo y círculo.
      \item Relaciona cada figura con su fórmula:

      \vspace{0.2cm}
      \begin{center}
      \AlmFontSmall
      \begin{tabular}{|l|l|}
        \hline
        \rowcolor{AlmAzul!15}
        \textbf{Figura} & \textbf{Fórmula} \\
        \hline
        Cuadrado & $A = \pi r^2$ \\
        \hline
        Triángulo & $A = l^2$ \\
        \hline
        Trapecio & $A = b \times h / 2$ \\
        \hline
        Círculo & $A = \tfrac{(B+b)}{2} \times h$ \\
        \hline
      \end{tabular}
      \end{center}
    \end{enumerate}
  \end{recordarbox}
  \vspace{0.15cm}
  \begin{notabox}[Pista]
    \AlmFontSmall Para el triángulo, la altura debe ser perpendicular a la base.
  \end{notabox}
\end{frame}

% ---------------------------------------------------------------
\seccionframe[BComprender]{Ejercicios — COMPRENDER}
% ---------------------------------------------------------------

% ---------------------------------------------------------------
%  FRAME 11 — Ejercicio COMPRENDER
% ---------------------------------------------------------------
\begin{frame}{Ejercicio 2 — COMPRENDER \dificultad{2}}
  \begin{comprenderbox}
    Un \textbf{cuadrado} de lado 10 m y un \textbf{rectángulo}
    de 15 m $\times$ 5 m tienen el mismo perímetro (40 m).
    \begin{enumerate}
      \item Calcula el área de cada uno.
      \item ¿Cuál tiene mayor área?
      \item ¿Qué conclusión extraes sobre la relación
            entre perímetro y área?
    \end{enumerate}
  \end{comprenderbox}

  \vspace{0.2cm}
  \begin{columns}[c]
    \begin{column}{0.40\textwidth}
      \begin{center}
      \begin{tikzpicture}[scale=0.60]
        \fill[AlmVerde!15] (0,0) rectangle (2.0,2.0);
        \draw[rectov] (0,0) rectangle (2.0,2.0);
        \draw[acota] (0,-0.5) -- (2.0,-0.5)
          node[midway,below,font=\AlmFontTiny]{10 m};
        \node[font=\AlmFontScript\bfseries, color=AlmVerde] at (1.0,1.0) {?};
      \end{tikzpicture}
      \end{center}
    \end{column}
    \begin{column}{0.55\textwidth}
      \begin{center}
      \begin{tikzpicture}[scale=0.60]
        \fill[AlmAzulC!18] (0,0) rectangle (3.0,1.0);
        \draw[zona] (0,0) rectangle (3.0,1.0);
        \draw[acota] (0,-0.5) -- (3.0,-0.5)
          node[midway,below,font=\AlmFontTiny]{15 m};
        \draw[acota] (3.4,0) -- (3.4,1.0)
          node[midway,right,font=\AlmFontTiny]{5 m};
        \node[font=\AlmFontScript\bfseries, color=AlmAzulM] at (1.5,0.5) {?};
      \end{tikzpicture}
      \end{center}
    \end{column}
  \end{columns}
\end{frame}

% ---------------------------------------------------------------
\seccionframe[BAplicar]{Ejercicios — APLICAR}
% ---------------------------------------------------------------

% ---------------------------------------------------------------
%  FRAME 12 — Ejercicio APLICAR
% ---------------------------------------------------------------
\begin{frame}{Ejercicio 3 — APLICAR \dificultad{3}}
  \begin{aplicarbox}
    \AlmFontSmall Se quiere crear una \textbf{zona verde triangular}
    en la playa de Almería con base \textbf{300 m}
    y altura \textbf{150 m}.
    \begin{enumerate}
      \item Calcula el área en m$^2$.
      \item Exprésala en hectáreas (ha).
      \item Calcula el coste de sembrar césped
            si el precio es \textbf{8 €/m$^2$}.
    \end{enumerate}
  \end{aplicarbox} 
  \begin{columns}[c]
    \begin{column}{0.42\textwidth}
      \begin{center}
      \begin{tikzpicture}[scale=0.55]
        \fill[AlmVerde!18] (0,0) -- (4.5,0) -- (2.25,2.25) -- cycle;
        \draw[rectov, line width=1.5pt]
          (0,0) -- (4.5,0) -- (2.25,2.25) -- cycle;
        \draw[AlmGris, dashed, thin] (2.25,0) -- (2.25,2.25);
        \draw[acota] (0,-0.6) -- (4.5,-0.6)
          node[midway,below,font=\AlmFontTiny]{300 m};
        \draw[acota] (2.6,0) -- (2.6,2.25)
          node[midway,right,font=\AlmFontTiny]{150 m};
        \node[font=\AlmFontScript\bfseries, color=AlmVerde!70!black]
          at (2.25,0.8) {Zona verde};
      \end{tikzpicture}
      \end{center}
      
            \begin{pasobox}[1]
              \AlmFontScript $A = \dfrac{b \times h}{2}
                = \dfrac{300 \times 150}{2} = ?\ \text{m}^2$
            \end{pasobox}
    \end{column}
    \begin{column}{0.54\textwidth}
      \begin{pasobox}[2]
        \AlmFontScript $A_{\text{ha}} = \dfrac{A_{\text{m}^2}}{10{,}000}$
      \end{pasobox}
      \vspace{0.15cm}
      \begin{pasobox}[3]
        \AlmFontScript $\text{Coste} = A_{\text{m}^2} \times 8$~€/m$^2$
      \end{pasobox}
    \end{column}
  \end{columns}
\end{frame}

% ---------------------------------------------------------------
%  FRAME 13 — Error frecuente
% ---------------------------------------------------------------
\begin{frame}{¡Cuidado con este error!}
  \begin{errorbox}
    \begin{itemize}
      \item \incorrecto\
            \textbf{«El parque tiene 32{,}000 metros»}
            \\[6pt]
            \correcto\
            El área se mide en metros \textbf{cuadrados}:
            «El parque tiene 32{,}000 \textbf{m$^2$}» o «3{,}2 ha».
            \\
            La palabra «metros» sola indica longitud (perímetro, distancia),
            no superficie.

      \vspace{0.4cm}
      \item \incorrecto\
            Usar la fórmula del rectángulo en un triángulo:
            \\
            $A = 120 \times 95 = 11{,}400$~m$^2$ (sin dividir entre 2).
            \\[6pt]
            \correcto\
            $A = \dfrac{120 \times 95}{2} = 5{,}700$~m$^2$
            \\
            El triángulo es \textbf{la mitad} del rectángulo que lo contiene.
    \end{itemize}
  \end{errorbox}
\end{frame}

% ---------------------------------------------------------------
%  FRAME 14 — Resumen
% ---------------------------------------------------------------
\begin{frame}{Resumen de la sesión}
  \begin{resumenbox}
    \begin{itemize}
      \item El \textbf{área} mide la superficie interior:
            como el césped a segar.
      \item Unidades: m$^2$, cm$^2$, km$^2$, \textbf{ha}
            (1 ha = 10{,}000 m$^2$).
      \item Fórmulas clave:
        $A_{\text{cuad}} = l^2$, \quad
        $A_{\text{rect}} = b \times h$, \quad
        $A_{\text{trián}} = \tfrac{bh}{2}$, \\
        $A_{\text{paral}} = bh$, \quad
        $A_{\text{trap}} = \tfrac{B+b}{2} \times h$, \quad
        $A_{\text{círc}} = \pi r^2$.
      \item La altura debe ser \textbf{perpendicular} a la base.
      \item Parque Salmerón: $A = 32{,}000$~m$^2 = 3{,}2$~ha.
    \end{itemize}
  \end{resumenbox}
\end{frame}

% ---------------------------------------------------------------
%  FRAME 15 — Evidencia del día
% ---------------------------------------------------------------
\begin{frame}{Evidencia del día}
  \begin{evidenciabox}
    \textbf{Mapa de zonas verdes con áreas calculadas}

    \vspace{0.2cm}
    Investiga \textbf{3 parques o plazas} de Almería.
    Para cada uno:
    \begin{enumerate}
      \item Dibuja su forma aproximada con medidas.
      \item Identifica qué figura geométrica modela mejor su planta.
      \item Aplica la fórmula correcta y calcula el área en m$^2$.
      \item Convierte a hectáreas.
      \item Calcula el coste hipotético de regar el parque
            a razón de 2~€/m$^2$.
    \end{enumerate}
    \vspace{0.2cm}
    \textbf{Criterio de éxito:} los 3 cálculos son correctos,
    las fórmulas están indicadas y las unidades son adecuadas.
  \end{evidenciabox}
\end{frame}

\end{document}




% % ================================================================
%  SESIÓN 09 — Perímetro y área: no son lo mismo
%  Almería en Cifras y Formas · Matemáticas 1.º ESO
%  IES Al-Ándalus · 2025-2026
% ================================================================
\documentclass[aspectratio=169, 10pt, spanish]{beamer}
\usepackage{almeria-beamer}

\renewcommand{\SessionNumber}{09}
\renewcommand{\SessionTitle}{Perímetro y área: no son lo mismo}
\renewcommand{\SessionName}{Sesión 09}
\renewcommand{\SessionObjective}{%
  Distinguir conceptual y numéricamente el perímetro del área,
  comprendiendo que son medidas independientes y que una puede
  cambiar sin afectar a la otra.%
}
\renewcommand{\BlockName}{Bloque 2: Geometría urbana}

\begin{document}

% ---------------------------------------------------------------
%  PORTADA
% ---------------------------------------------------------------
\begin{frame}[plain]
  \titlepage
\end{frame}

% ---------------------------------------------------------------
%  ÍNDICE
% ---------------------------------------------------------------
\begin{frame}{Contenidos de hoy}
  \begin{columns}[t]
    \begin{column}{0.48\textwidth}
      \begin{itemize}
        \item Tabla comparativa P vs.\ A
        \item Mismos P, distintos A
        \item Mismos A, distintos P
        \item Problema isoperimérico
        \item Implicaciones en urbanismo
      \end{itemize}
    \end{column}
    \begin{column}{0.48\textwidth}
      \begin{itemize}
        \item Almería: Rambla completa
        \item Cálculo paralelo P y A
        \item Ejercicios Bloom
        \item Error frecuente
        \item Evidencia del día
      \end{itemize}
    \end{column}
  \end{columns}
  \vspace{0.4cm}
  \begin{notabox}[Niveles de Bloom]
    \bloomtag{Recordar}\quad
    \bloomtag{Comprender}\quad
    \bloomtag{Aplicar}
  \end{notabox}
\end{frame}

% ---------------------------------------------------------------
\seccionframe{Comparando perímetro y área}
% ---------------------------------------------------------------

% ---------------------------------------------------------------
%  FRAME 3 — Tabla comparativa
% ---------------------------------------------------------------
\begin{frame}{Perímetro vs. Área: tabla comparativa}
  \begin{center}
  \AlmFontSmall
  \begin{tabular}{|>{\bfseries\color{AlmAzulM}}l|>{\centering\arraybackslash}p{3.8cm}|>{\centering\arraybackslash}p{3.8cm}|}
    \hline
    \rowcolor{AlmAzul!20}
    & \textbf{\color{AlmTerra}Perímetro} & \textbf{\color{AlmVerde}Área} \\
    \hline
    ¿Qué mide? &
      El \textbf{contorno} &
      La \textbf{superficie}\\ &
      (la valla) &(el césped) \\
    \hline
    Dimensión &
      \textbf{1D} (lineal) &
      \textbf{2D} (superficial) \\
    \hline
    Unidades &
      m, km, cm\ldots &
      m$^2$, km$^2$, ha\ldots \\
    \hline
    Fórmula rect. &
      $P = 2(l + w)$ &
      $A = l \times w$ \\
    \hline
    Uso &
      Vallas, bordes, caminos &
      Suelos, pinturas, césped \\
    \hline
  \end{tabular}
  \end{center}

  \vspace{0.2cm}
  \begin{notabox}[Idea clave]
    \AlmFontSmall Perímetro y área miden cosas \textbf{completamente distintas}.
    Una puede ser grande mientras la otra es pequeña.
  \end{notabox}
\end{frame}

% ---------------------------------------------------------------
%  FRAME 4 — Mismo P, distinto A
% ---------------------------------------------------------------
\begin{frame}{Mismo perímetro, distinto área}
  \begin{columns}[c]
    \begin{column}{0.62\textwidth}
      \begin{center}
      \begin{tikzpicture}[scale=0.78]
        % Rectángulo 6x2
        \fill[AlmTerra!12] (0,0) rectangle (4.8,1.6);
        \draw[rectot, line width=1.8pt] (0,0) rectangle (4.8,1.6);
        \draw[acota] (0,-0.55) -- (4.8,-0.55)
          node[midway,below,font=\AlmFontScript]{6};
        \draw[acota] (5.15,0) -- (5.15,1.6)
          node[midway,right,font=\AlmFontScript]{2};
        \node[font=\AlmFontScript\bfseries, color=AlmTerra] at (2.4,0.8)
          {$P = 16$ \quad $A = 12$};

        % Rectángulo 4x4 (más abajo)
        \fill[AlmVerde!15] (0,-3.0) rectangle (3.2,-3.0+3.2);
        \draw[rectov, line width=1.8pt] (0,-3.0) rectangle (3.2,-3.0+3.2);
        \draw[acota] (0,-3.6) -- (3.2,-3.6)
          node[midway,below,font=\AlmFontScript]{4};
        \draw[acota] (3.55,-3.0) -- (3.55,-3.0+3.2)
          node[midway,right,font=\AlmFontScript]{4};
        \node[font=\AlmFontScript\bfseries, color=AlmVerde] at (1.6,-1.4)
          {$P = 16$ \quad $A = 16$};

        % Etiqueta
        \node[font=\AlmFontScript\bfseries, color=AlmAmbar!80!black]
          at (2.5,-0.1) {};
      \end{tikzpicture}
      \end{center}
    \end{column}
    \begin{column}{0.34\textwidth}
      \begin{teoriabox}
        \AlmFontSmall
        Ambos rectángulos tienen
        $P = 16$ unidades.
        \\[6pt]
        Pero sus áreas son distintas:
        \begin{itemize}
          \item 6$\times$2: $A = 12$
          \item 4$\times$4: $A = 16$
        \end{itemize}
        \vspace{4pt}
        El cuadrado es la figura
        de \textbf{mayor área} para
        un perímetro fijo.
        \\[4pt]
        \AlmFontScript(Problema isoperimérico)
      \end{teoriabox}
    \end{column}
  \end{columns}
\end{frame}

% ---------------------------------------------------------------
%  FRAME 5 — Mismo A, distinto P
% ---------------------------------------------------------------
\begin{frame}{Mismo área, distinto perímetro}
  \begin{columns}[c]
    \begin{column}{0.62\textwidth}
      \begin{center}
      \begin{tikzpicture}[scale=0.78]
        % Rectángulo 3x4 → A=12, P=14
        \fill[AlmAzulC!18] (0,0) rectangle (2.4,3.2);
        \draw[zona, line width=1.8pt] (0,0) rectangle (2.4,3.2);
        \draw[acota] (0,-0.55) -- (2.4,-0.55)
          node[midway,below,font=\AlmFontScript]{3};
        \draw[acota] (2.75,0) -- (2.75,3.2)
          node[midway,right,font=\AlmFontScript]{4};
        \node[font=\AlmFontScript\bfseries, color=AlmAzulM] at (1.2,1.6)
          {$P = 14$\quad$A = 12$};

        % Rectángulo 6x2 → A=12, P=16
        \fill[AlmAmbar!15] (4.2,0.8) rectangle (4.2+4.8,0.8+1.6);
        \draw[AlmAmbar!70!black, line width=1.8pt]
          (4.2,0.8) rectangle (4.2+4.8,0.8+1.6);
        \draw[acota] (4.2,0.25) -- (9.0,0.25)
          node[midway,below,font=\AlmFontScript]{6};
        \draw[acota] (9.35,0.8) -- (9.35,2.4)
          node[midway,right,font=\AlmFontScript]{2};
        \node[font=\AlmFontScript\bfseries, color=AlmAmbar!80!black] at (6.6,1.6)
          {$P = 16$\quad$A = 12$};
      \end{tikzpicture}
      \end{center}
    \end{column}
    \begin{column}{0.34\textwidth}
      \begin{teoriabox}
        \AlmFontSmall
        Ambos rectángulos tienen
        $A = 12$ unidades cuadradas.
        \\[6pt]
        Pero sus perímetros son distintos:
        \begin{itemize}
          \item 3$\times$4: $P = 14$
          \item 6$\times$2: $P = 16$
        \end{itemize}
        \vspace{4pt}
        La figura más «cuadrada»
        tiene menor perímetro para
        una misma área.
      \end{teoriabox}
    \end{column}
  \end{columns}
\end{frame}

% ---------------------------------------------------------------
%  FRAME 6 — Implicaciones en urbanismo
% ---------------------------------------------------------------
\begin{frame}{¿Cuándo usar cada medida?}
  \begin{columns}[t]
    \begin{column}{0.48\textwidth}
      \begin{definicionbox}{Necesitas el PERÍMETRO}
        \AlmFontSmall
        \begin{itemize}
          \item Colocar una \textbf{valla} alrededor de un parque.
          \item Poner \textbf{bordillo} en el borde de una acera.
          \item Calcular la distancia de una \textbf{ruta a pie}.
          \item Instalar \textbf{luminarias} en el contorno de una plaza.
        \end{itemize}
      \end{definicionbox}
    \end{column}
    \begin{column}{0.48\textwidth}
      \begin{definicionbox}{Necesitas el ÁREA}
        \AlmFontSmall
        \begin{itemize}
          \item Comprar \textbf{baldosas} para pavimentar.
          \item Calcular el \textbf{césped} a sembrar.
          \item Valorar el \textbf{precio} de un terreno.
          \item Pintar \textbf{la fachada} de un edificio.
        \end{itemize}
      \end{definicionbox}
    \end{column}
  \end{columns}
  \vspace{0.3cm}
  \begin{notabox}[Resumen]
    \AlmFontSmall \textbf{Valla, bordillo, camino} $\rightarrow$ Perímetro (metros).
    \quad
    \textbf{Suelo, césped, pintura} $\rightarrow$ Área (m$^2$).
  \end{notabox}
\end{frame}

% ---------------------------------------------------------------
\seccionframe[BAplicar]{Almería: la Rambla}
% ---------------------------------------------------------------

% ---------------------------------------------------------------
%  FRAME 7 — Ejemplo completo: Rambla de Almería
% ---------------------------------------------------------------
\begin{frame}{Ejemplo completo: Rambla de Almería}
  \begin{columns}[c]
    \begin{column}{0.52\textwidth}
      \begin{ejemplobox}
        \AlmFontSmall La \textbf{Rambla de Almería} es aproximadamente
        rectangular con \textbf{370 m} de largo y \textbf{50 m}
        de ancho.
        \begin{align*}
          P &= 2(370 + 50) = 2 \times 420 = \mathbf{840\ m} \\
          A &= 370 \times 50 = \mathbf{18{,}500\ m^2}
            = \mathbf{1{,}85\ ha}
        \end{align*}
        \textbf{P} sirve para: vallar todo el contorno. \\
        \textbf{A} sirve para: pavimentar la superficie.
      \end{ejemplobox}
    \end{column}
    \begin{column}{0.44\textwidth}
      \begin{center}
      \begin{tikzpicture}[scale=0.70]
        % Rambla (rectángulo alargado)
        \fill[AlmAmbar!15] (0,0) rectangle (5.5,0.75);
        \draw[AlmAmbar!70!black, line width=2pt] (0,0) rectangle (5.5,0.75);
        % Valla (borde rojo)
        \draw[AlmTerra, line width=3pt, dashed] (0,0) rectangle (5.5,0.75);
        % Árboles (zona verde interior simplificada)
        \fill[AlmVerde!20] (0.2,0.1) rectangle (5.3,0.65);
        \draw[acota] (0,-0.5) -- (5.5,-0.5)
          node[midway,below,font=\AlmFontTiny]{370 m};
        \draw[acota] (5.9,0) -- (5.9,0.75)
          node[midway,right,font=\AlmFontTiny]{50 m};
        \node[font=\AlmFontTiny\bfseries, color=AlmVerde!70!black] at (2.75,0.375)
          {Rambla de Almería};
        \node[font=\AlmFontTiny\bfseries, color=AlmTerra] at (2.75,-0.95)
          {$P = 840$ m};
        \node[font=\AlmFontTiny\bfseries, color=AlmVerde!70!black] at (2.75,-1.35)
          {$A = 18{.}500$ m$^2$};
      \end{tikzpicture}
      \end{center}
    \end{column}
  \end{columns}
\end{frame}

% ---------------------------------------------------------------
\seccionframe[BRecordar]{Ejercicios — RECORDAR}
% ---------------------------------------------------------------

% ---------------------------------------------------------------
%  FRAME 8 — Ejercicio RECORDAR
% ---------------------------------------------------------------
\begin{frame}{Ejercicio 1 — RECORDAR \dificultad{1}}
  \begin{recordarbox}
    Para un rectángulo de \textbf{12 m} de largo y \textbf{5 m}
    de ancho, calcula:
    \begin{enumerate}
      \item El \textbf{perímetro}. Indica su unidad.
      \item El \textbf{área}. Indica su unidad.
    \end{enumerate}
    ¿Cuál es la diferencia entre las unidades de perímetro
    y área?
  \end{recordarbox}
  \vspace{0.25cm}
  \begin{center}
  \begin{tikzpicture}[scale=0.80]
    \fill[AlmAzulC!12] (0,0) rectangle (4.8,2.0);
    \draw[zona, line width=1.8pt] (0,0) rectangle (4.8,2.0);
    \draw[acota] (0,-0.55) -- (4.8,-0.55)
      node[midway,below,font=\AlmFontScript]{12 m};
    \draw[acota] (5.15,0) -- (5.15,2.0)
      node[midway,right,font=\AlmFontScript]{5 m};
    \node[font=\AlmFontSmall, color=AlmTerra] at (1.8,1.0) {$P = ?$\ m};
    \node[font=\AlmFontSmall, color=AlmVerde] at (3.4,1.0) {$A = ?$\ m$^2$};
  \end{tikzpicture}
  \end{center}
\end{frame}

% ---------------------------------------------------------------
\seccionframe[BComprender]{Ejercicios — COMPRENDER}
% ---------------------------------------------------------------

% ---------------------------------------------------------------
%  FRAME 9 — Ejercicio COMPRENDER (problema isoperimérico)
% ---------------------------------------------------------------
\begin{frame}{Ejercicio 2 — COMPRENDER \dificultad{3}}
  \begin{comprenderbox}
    Tienes \textbf{24 m} de valla para cercar un huerto rectangular.
    Explora todas las posibilidades (lados enteros) y encuentra
    la que da \textbf{mayor área}:

    \vspace{0.15cm}
    \AlmFontSmall
    \begin{center}
    \begin{tabular}{|c|c|c|c|}
      \hline
      \rowcolor{AlmAzul!15}
      \textbf{Largo} & \textbf{Ancho} & \textbf{Perímetro} & \textbf{Área} \\
      \hline
      1 m & 11 m & 24 m & \phantom{000} m$^2$ \\
      2 m & 10 m & 24 m & \phantom{000} m$^2$ \\
      3 m &  9 m & 24 m & \phantom{000} m$^2$ \\
      4 m &  8 m & 24 m & \phantom{000} m$^2$ \\
      5 m &  7 m & 24 m & \phantom{000} m$^2$ \\
      6 m &  6 m & 24 m & \phantom{000} m$^2$ \\
      \hline
    \end{tabular}
    \end{center}

    ¿Qué forma geométrica da el área máxima?
  \end{comprenderbox}
\end{frame}

% ---------------------------------------------------------------
\seccionframe[BAplicar]{Ejercicios — APLICAR}
% ---------------------------------------------------------------

% ---------------------------------------------------------------
%  FRAME 10 — Ejercicio APLICAR
% ---------------------------------------------------------------
\begin{frame}{Ejercicio 3 — APLICAR \dificultad{3}}
  \begin{aplicarbox}
    \AlmFontSmall Un urbanista de Almería necesita:
    \begin{enumerate}
      \item \textbf{Vallar} un parque rectangular de 80 m $\times$ 60 m
            (precio: 45 €/m de valla).
      \item \textbf{Cubrir} su superficie con gravilla
            (precio: 12 €/m$^2$).
    \end{enumerate}
    Calcula:
    \begin{enumerate}[a)]
      \item El coste total de la valla.
      \item El coste total de la gravilla.
      \item ¿Cuál es mayor? ¿Qué dato usa cada cálculo (P o A)?
    \end{enumerate}
  \end{aplicarbox}
  \vspace{0.15cm}
  \begin{columns}[c]
    \begin{column}{0.50\textwidth}
      \begin{pasobox}[1]
        \AlmFontScript $P = 2(80 + 60) = ?$ m
        \\
        Coste valla $= P \times 45$~€
      \end{pasobox}
    \end{column}
    \begin{column}{0.46\textwidth}
      \begin{pasobox}[2]
        \AlmFontScript $A = 80 \times 60 = ?$~m$^2$
        \\
        Coste gravilla $= A \times 12$~€
      \end{pasobox}
    \end{column}
  \end{columns}
\end{frame}

% ---------------------------------------------------------------
%  FRAME 11 — Los tres casos TikZ (comparativa completa)
% ---------------------------------------------------------------
\begin{frame}{Tres rectángulos: misma P o misma A}
  \begin{center}
  \begin{tikzpicture}[scale=0.72]
    % --- Caso 1: 6x2, P=16, A=12 ---
    \fill[AlmTerra!12] (0,0) rectangle (4.8,1.6);
    \draw[rectot, line width=1.5pt] (0,0) rectangle (4.8,1.6);
    \node[above, font=\AlmFontTiny\bfseries, color=AlmTerra] at (2.4,1.6)
      {$6 \times 2$};
    \node[font=\AlmFontTiny, color=AlmTerra] at (2.4,0.8) {$P=16$};
    \node[font=\AlmFontTiny, color=AlmTerra] at (2.4,0.3) {$A=12$};

    % --- Caso 2: 4x4, P=16, A=16 ---
    \fill[AlmVerde!15] (5.8,0) rectangle (5.8+3.2,3.2);
    \draw[rectov, line width=1.5pt] (5.8,0) rectangle (5.8+3.2,3.2);
    \node[above, font=\AlmFontTiny\bfseries, color=AlmVerde] at (7.4,3.2)
      {$4 \times 4$};
    \node[font=\AlmFontTiny, color=AlmVerde] at (7.4,1.8) {$P=16$};
    \node[font=\AlmFontTiny, color=AlmVerde] at (7.4,1.2) {$A=16$};

    % Flecha entre 1 y 2
    \draw[->, AlmAmbar!80!black, line width=1.2pt]
      (5.0,0.8) -- (5.6,0.8);
    \node[above, font=\AlmFontTiny\bfseries, color=AlmAmbar!80!black]
      at (5.3,0.8) {P igual};

    % --- Caso 3: 3x4, P=14, A=12 ---
    \fill[AlmAzulC!18] (0,-3.5) rectangle (2.4,-3.5+3.2);
    \draw[zona, line width=1.5pt] (0,-3.5) rectangle (2.4,-3.5+3.2);
    \node[below, font=\AlmFontTiny\bfseries, color=AlmAzulM] at (1.2,-3.5)
      {$3 \times 4$};
    \node[font=\AlmFontTiny, color=AlmAzulM] at (1.2,-2.0) {$P=14$};
    \node[font=\AlmFontTiny, color=AlmAzulM] at (1.2,-2.5) {$A=12$};

    % Flecha entre 3 y 1 (misma A)
    \draw[->, AlmGris, dashed, line width=1.0pt]
      (2.6,-2.1) -- (4.4,-0.5);
    \node[right, font=\AlmFontTiny\bfseries, color=AlmGris]
      at (3.5,-1.4) {A igual};
  \end{tikzpicture}
  \end{center}
  \begin{notabox}[Conclusión]
    \AlmFontSmall Si cambias la forma de un rectángulo manteniendo el mismo
    perímetro, el área cambia. Y viceversa.
  \end{notabox}
\end{frame}

% ---------------------------------------------------------------
%  FRAME 12 — Error frecuente
% ---------------------------------------------------------------
\begin{frame}{¡Cuidado con este error!}
  \begin{errorbox}
    \begin{itemize}
      \item \incorrecto\
            \textbf{«El área del parque es 280 metros»}
            \\[6pt]
            \correcto\
            El área se mide en \textbf{m$^2$}:
            «El área es 280 \textbf{m$^2$}».
            Si dices solo «metros», estás indicando una longitud.

      \vspace{0.3cm}
      \item \incorrecto\
            \textbf{«El perímetro es 4{,}800 metros cuadrados»}
            \\[6pt]
            \correcto\
            El perímetro es una longitud (1D):
            «El perímetro es 4{,}800 \textbf{m}».
            Nunca lleva el exponente 2.

      \vspace{0.3cm}
      \item \incorrecto\
            Pensar que «mayor área $\Rightarrow$ mayor perímetro».
            \\[6pt]
            \correcto\
            Son medidas \textbf{independientes}.
            Un rectángulo puede tener mayor área
            y menor perímetro que otro.
    \end{itemize}
  \end{errorbox}
\end{frame}

% ---------------------------------------------------------------
%  FRAME 13 — Resumen
% ---------------------------------------------------------------
\begin{frame}{Resumen de la sesión}
  \begin{resumenbox}
    \begin{itemize}
      \item \textbf{Perímetro} $\neq$ \textbf{Área}: miden cosas distintas.
      \item P mide el \emph{contorno} (1D, en metros).
            A mide la \emph{superficie} (2D, en m$^2$).
      \item Dos figuras con el mismo P pueden tener distintas A,
            y viceversa.
      \item Con P fijo, el cuadrado maximiza el área
            (problema isoperimérico).
      \item Rambla de Almería: $P = 840$ m, $A = 18{,}500$ m$^2$.
      \item Usar P para vallas/bordillos;
            usar A para suelos/césped.
    \end{itemize}
  \end{resumenbox}
\end{frame}

% ---------------------------------------------------------------
%  FRAME 14 — Evidencia del día
% ---------------------------------------------------------------
\begin{frame}{Evidencia del día}
  \begin{evidenciabox}
    \textbf{Ficha doble: Perímetro y Área en paralelo}

    \vspace{0.2cm}
    Elige \textbf{3 figuras} de diferente forma (rectángulo, triángulo,
    trapecio). Para cada una:
    \begin{enumerate}
      \item Dibuja la figura con sus medidas.
      \item Calcula el perímetro — muestra todos los pasos.
      \item Calcula el área — muestra todos los pasos.
      \item Escribe las unidades correctas para cada resultado.
      \item Indica para qué sirve cada medida en un contexto real.
    \end{enumerate}
    \vspace{0.2cm}
    \textbf{Criterio de éxito:} los 6 cálculos (3 P + 3 A) son correctos
    y las unidades son las adecuadas para cada tipo de medida.
  \end{evidenciabox}
\end{frame}

\end{document}




% % ================================================================
%  sesion10.tex  —  Sesión 10 de 20
%  Rutas y distancias por Almería
%  Bloque 2: Geometría urbana · Matemáticas 1.º ESO
% ================================================================
\documentclass[aspectratio=169, 10pt, spanish]{beamer}
\usepackage{almeria-beamer}

\renewcommand{\SessionNumber}{10}
\renewcommand{\SessionTitle}{Rutas y distancias por Almería}
\renewcommand{\SessionName}{Sesión 10}
\renewcommand{\SessionObjective}{%
  Calcular distancias reales entre puntos turísticos de Almería
  usando el Teorema de Pitágoras y la escala del mapa,
  y diseñar rutas turísticas optimizadas.}
\renewcommand{\BlockName}{Bloque 2: Geometría urbana}

\begin{document}

% ---------------------------------------------------------------
\begin{frame}[plain]\titlepage\end{frame}

% ---------------------------------------------------------------
\begin{frame}{Índice}
  \begin{columns}[t]
    \begin{column}{0.5\textwidth}
      \begin{enumerate}
        \item Dos tipos de distancia
        \item El Teorema de Pitágoras
        \item Fórmula de la distancia entre dos puntos
        \item Escala del mapa
        \item Ejemplo paso a paso
      \end{enumerate}
    \end{column}
    \begin{column}{0.5\textwidth}
      \begin{enumerate}\setcounter{enumi}{5}
        \item Ruta turística por Almería
        \item Ejercicios multinivel (Bloom)
        \item Evidencia del día
        \item Error frecuente
        \item Resumen
      \end{enumerate}
    \end{column}
  \end{columns}
\end{frame}

% ---------------------------------------------------------------
\seccionframe{1. Dos tipos de distancia}
% ---------------------------------------------------------------

\begin{frame}{Distancia en línea recta vs.\ distancia real}
  \begin{columns}[c]
    \begin{column}{0.55\textwidth}
      \begin{definicionbox}{Distancia euclídea (línea recta)}
        La longitud del segmento más corto que une dos puntos.
        Es la distancia \textit{en el aire}, sin obstáculos.
        Se calcula con el Teorema de Pitágoras.
      \end{definicionbox}
      \vspace{6pt}
      \begin{definicionbox}{Distancia real (por la calle)}
        El camino que recorremos siguiendo las calles o rutas
        existentes. Siempre es \textbf{mayor o igual} que
        la distancia en línea recta.
      \end{definicionbox}
    \end{column}
    \begin{column}{0.42\textwidth}
      \begin{tikzpicture}[scale=0.92]
        \draw[AlmFondo!60, step=0.6cm] (0,0) grid (4.8,3.6);
        % Puntos
        \node[punto] (A) at (0.6,0.6) {};
        \node[punto] (B) at (4.2,3.0) {};
        \node[etiqueta, above left] at (A) {A};
        \node[etiqueta, above right] at (B) {B};
        % Línea recta (euclídea)
        \draw[recto, AlmAzulM] (A) -- (B)
          node[midway, below right, font=\AlmFontTiny\bfseries, color=AlmAzulM]
          {distancia euclídea};
        % Ruta Manhattan
        \draw[rectot, dashed, line width=1.3pt]
          (A) -- (4.2,0.6) -- (B)
          node[midway, right, font=\AlmFontTiny, color=AlmTerra]
          {ruta calles};
        % Acotaciones
        \draw[acota] (0.6,0.2) -- (4.2,0.2)
          node[midway, below, font=\AlmFontTiny] {$\Delta x$};
        \draw[acota] (4.5,0.6) -- (4.5,3.0)
          node[midway, right, font=\AlmFontTiny] {$\Delta y$};
      \end{tikzpicture}
    \end{column}
  \end{columns}
\end{frame}

% ---------------------------------------------------------------
\seccionframe{2. El Teorema de Pitágoras}
% ---------------------------------------------------------------

\begin{frame}{Teorema de Pitágoras}
  \begin{columns}[c]
    \begin{column}{0.52\textwidth}
      \begin{teoriabox}
        En todo triángulo \textbf{rectángulo}, el cuadrado
        de la hipotenusa es igual a la suma de los cuadrados
        de los catetos:
        \[
          c^2 = a^2 + b^2
          \qquad \Longrightarrow \qquad
          c = \sqrt{a^2 + b^2}
        \]
        \vspace{2pt}
        \begin{itemize}
          \item $a, b$ = catetos (lados que forman el ángulo recto)
          \item $c$ = hipotenusa (lado opuesto al ángulo recto,
                el más largo)
        \end{itemize}
      \end{teoriabox}
    \end{column}
    \begin{column}{0.44\textwidth}
      \begin{center}
      \begin{tikzpicture}[scale=1.1]
        % Triángulo rectángulo
        \coordinate (O) at (0,0);
        \coordinate (A) at (3,0);
        \coordinate (B) at (0,2.2);
        \fill[zona] (O) -- (A) -- (B) -- cycle;
        \draw[recto] (O) -- (A) node[midway, below, font=\AlmFontSmall] {$a$};
        \draw[recto] (O) -- (B) node[midway, left, font=\AlmFontSmall] {$b$};
        \draw[rectot, line width=2pt] (A) -- (B)
          node[midway, right, font=\AlmFontSmall\bfseries, color=AlmTerra] {$c$};
        % Ángulo recto
        \draw[rectoang] (0.28,0) -- (0.28,0.28) -- (0,0.28);
        % Cuadrados visuales
        \fill[AlmAzulC, opacity=0.25]
          (0,-0.55) rectangle (3,-0.25) node[midway, font=\AlmFontTiny, color=AlmAzul] {$a^2$};
        \fill[BRecordar, opacity=0.20]
          (-0.6,0) rectangle (-0.3,2.2) node[midway, font=\AlmFontTiny, color=BRecordar, rotate=90] {$b^2$};
      \end{tikzpicture}
      \end{center}
      \vspace{4pt}
      \begin{formulabox}[AlmTerra]
        $c = \sqrt{a^2 + b^2}$
      \end{formulabox}
    \end{column}
  \end{columns}
\end{frame}

% ---------------------------------------------------------------
\seccionframe{3. Distancia entre dos puntos}
% ---------------------------------------------------------------

\begin{frame}{Fórmula de la distancia entre dos puntos}
  \begin{columns}[c]
    \begin{column}{0.52\textwidth}
      \begin{teoriabox}
        Dados dos puntos $A(x_1,\,y_1)$ y $B(x_2,\,y_2)$:
        \[
          d(A,B) = \sqrt{(x_2-x_1)^2 + (y_2-y_1)^2}
        \]
        Los desplazamientos $\Delta x = x_2 - x_1$ y
        $\Delta y = y_2 - y_1$ son los catetos;
        la distancia es la hipotenusa.
      \end{teoriabox}
      \vspace{6pt}
      \begin{notabox}[Tripletas pitagóricas]
        \AlmFontFootnote
        $3$-$4$-$5$;\quad $5$-$12$-$13$;\quad $8$-$15$-$17$\\
        ¡Útiles para comprobar resultados mentalmente!
      \end{notabox}
    \end{column}
    \begin{column}{0.44\textwidth}
      \begin{tikzpicture}[scale=0.85]
        \draw[AlmFondo!80, step=1cm] (-0.5,-0.5) grid (5.5,5.5);
        \draw[->, AlmAzulM, line width=1pt] (-0.3,0) -- (5.3,0)
          node[right, font=\AlmFontSmall] {$x$};
        \draw[->, AlmAzulM, line width=1pt] (0,-0.3) -- (0,5.3)
          node[above, font=\AlmFontSmall] {$y$};
        \coordinate (A) at (1,1);
        \coordinate (B) at (4,5);
        \node[punto] at (A) {};
        \node[punto] at (B) {};
        \node[etiqueta, below left, font=\AlmFontSmall] at (A) {$A(1,1)$};
        \node[etiqueta, above left, font=\AlmFontSmall] at (B) {$B(4,5)$};
        % Triángulo auxiliar
        \coordinate (C) at (4,1);
        \draw[AlmGris, dashed] (A) -- (C) -- (B);
        \draw[rectoang] (3.72,1) -- (3.72,1.28) -- (4,1.28);
        \draw[rectot, line width=2pt] (A) -- (B);
        % Catetos
        \node[below, font=\AlmFontTiny, color=AlmGris] at (2.5,1) {$\Delta x=3$};
        \node[right, font=\AlmFontTiny, color=AlmGris] at (4,3) {$\Delta y=4$};
        \node[above left, font=\AlmFontSmall\bfseries, color=AlmTerra] at (2.5,3)
          {$c=5$};
      \end{tikzpicture}
    \end{column}
  \end{columns}
\end{frame}

% ---------------------------------------------------------------
\seccionframe{4. Escala del mapa}
% ---------------------------------------------------------------

\begin{frame}{Escala del mapa}
  \begin{columns}[T]
    \begin{column}{0.52\textwidth}
      \begin{definicionbox}{Escala}
        Relación entre la medida en el mapa y la medida real.
        Si la escala es $1{:}E$, entonces:
        \[
          d_{\text{real}} = d_{\text{mapa}} \times E
        \]
        \textbf{Ejemplo:} Escala 1:5\,000\\
        $\Rightarrow$ 1 cm en el mapa = 5\,000 cm = \textbf{50 m} en realidad.
      \end{definicionbox}
      \vspace{8pt}
      En nuestro mapa de Almería usamos:
      \begin{center}
        \tcbox[colback=AlmAmbar!15, colframe=AlmAmbar,
               fontupper=\bfseries\AlmFontSmall, arc=4pt]{%
          \textbf{1 unidad} del plano $=$ \textbf{200 m} reales}
      \end{center}
    \end{column}
    \begin{column}{0.44\textwidth}
      \begin{tabular}{lll}
        \toprule
        \textbf{Mapa} & \textbf{Escala} & \textbf{Real}\\
        \midrule
        1 cm & 1:1\,000  & 10 m\\
        1 cm & 1:5\,000  & 50 m\\
        1 cm & 1:10\,000 & 100 m\\
        1 cm & 1:50\,000 & 500 m\\
        1 cm & 1:100\,000 & 1 km\\
        \bottomrule
      \end{tabular}
      \vspace{8pt}
      \begin{ejemplobox}
        \AlmFontFootnote
        Dos puntos en el mapa distan 3 unidades.
        Con escala 1 u = 200 m:\\
        $d_{\text{real}} = 3 \times 200 = \mathbf{600\ \text{m}}$
      \end{ejemplobox}
    \end{column}
  \end{columns}
\end{frame}

% ---------------------------------------------------------------
\seccionframe{5. Ejemplo paso a paso}
% ---------------------------------------------------------------

\begin{frame}{Ejemplo: Alcazaba $\rightarrow$ Catedral}
  \framesubtitle{Escala: 1 unidad = 200 m}
  \begin{columns}[c]
    \begin{column}{0.44\textwidth}
      \begin{tikzpicture}[scale=0.82]
        \draw[AlmFondo!80, step=1cm] (0,0) grid (5,5);
        \draw[->, AlmAzulM, line width=0.8pt] (0,0)--(5.3,0)
          node[right,font=\AlmFontTiny]{$x$};
        \draw[->, AlmAzulM, line width=0.8pt] (0,0)--(0,5.3)
          node[above,font=\AlmFontTiny]{$y$};
        % Puntos
        \coordinate (Alc) at (1,4);
        \coordinate (Cat) at (2,2);
        \coordinate (Mer) at (3,2);
        \coordinate (Pto) at (4,1);
        \node[punto] at (Alc) {};
        \node[puntot] at (Cat) {};
        \node[puntov] at (Mer) {};
        \node[punto] at (Pto) {};
        \node[font=\AlmFontTiny\bfseries, above right, color=AlmAzulM] at (Alc)
          {Alcazaba (1,4)};
        \node[font=\AlmFontTiny\bfseries, above right, color=AlmTerra] at (Cat)
          {Catedral (2,2)};
        \node[font=\AlmFontTiny\bfseries, above right, color=AlmVerde] at (Mer)
          {Mercado (3,2)};
        \node[font=\AlmFontTiny\bfseries, below right, color=AlmAzulM] at (Pto)
          {Puerto (4,1)};
        % Ruta
        \draw[rectot, line width=2pt, ->] (Alc) -- (Cat);
        \draw[rectov, line width=2pt, ->] (Cat) -- (Mer);
        \draw[recto, line width=2pt, ->]  (Mer) -- (Pto);
      \end{tikzpicture}
    \end{column}
    \begin{column}{0.52\textwidth}
      \begin{pasobox}[1]
        \AlmFontFootnote
        $A(1,4)\to B(2,2)$:\\
        $\Delta x=1,\;\Delta y=-2$\\
        $d=\sqrt{1^2+2^2}=\sqrt{5}\approx 2{,}24\ \text{u}$\\
        $\rightarrow 2{,}24\times 200 \approx \mathbf{448\ \text{m}}$
      \end{pasobox}
      \vspace{4pt}
      \begin{pasobox}[2]
        \AlmFontFootnote
        $B(2,2)\to C(3,2)$:\\
        $\Delta x=1,\;\Delta y=0$\\
        $d=\sqrt{1^2+0^2}=1\ \text{u} = \mathbf{200\ \text{m}}$
      \end{pasobox}
      \vspace{4pt}
      \begin{pasobox}[3]
        \AlmFontFootnote
        $C(3,2)\to D(4,1)$:\\
        $\Delta x=1,\;\Delta y=-1$\\
        $d=\sqrt{1^2+1^2}=\sqrt{2}\approx 1{,}41\ \text{u}$\\
        $\rightarrow 1{,}41\times 200\approx \mathbf{283\ \text{m}}$
      \end{pasobox}
      \vspace{4pt}
      \begin{pasobox}[4]
        \AlmFontFootnote
        \textbf{Distancia total:}
        $448+200+283 = \mathbf{931\ \text{m}}$\\[2pt]
        Tiempo a 5 km/h $= \frac{0{,}931}{5}\times 60
        \approx \mathbf{11\ \text{min}}$
      \end{pasobox}
    \end{column}
  \end{columns}
\end{frame}

% ---------------------------------------------------------------
\seccionframe{6. Ruta turística completa}
% ---------------------------------------------------------------

\begin{frame}{Planifica tu ruta en Almería}
  \begin{columns}[T]
    \begin{column}{0.56\textwidth}
      \begin{ejemplobox}
        \AlmFontSmall
        \textbf{Puntos turísticos} (escala 1 u = 200 m):
        \begin{center}
        \begin{tabular}{lcc}
          \toprule
          Lugar & $x$ & $y$\\
          \midrule
          Alcazaba         & 1 & 4\\
          Catedral         & 2 & 2\\
          Mercado Central  & 3 & 2\\
          Puerto           & 4 & 1\\
          Playa de Almería & 5 & 2\\
          \bottomrule
        \end{tabular}
        \end{center}
      \end{ejemplobox}
      \vspace{4pt}
      \begin{notabox}[Tiempo de paseo]
        \AlmFontFootnote
        Velocidad media turista: $\approx 4{,}5\ \text{km/h}$\\
        Fórmula: $t = \dfrac{d}{v}$\\
        (resultado en horas; $\times 60$ para minutos)
      \end{notabox}
    \end{column}
    \begin{column}{0.41\textwidth}
      \begin{tabular}{lc}
        \toprule
        \textbf{Tramo} & $d$ (m)\\
        \midrule
        Alcazaba → Catedral  & 448\\
        Catedral → Mercado   & 200\\
        Mercado → Puerto     & 283\\
        Puerto → Playa       & 224\\
        \midrule
        \textbf{Total ruta}  & \textbf{1\,155}\\
        \bottomrule
      \end{tabular}
      \vspace{8pt}
      \begin{formulabox}[AlmVerde]
        $t = \dfrac{1{,}155\ \text{km}}{4{,}5\ \text{km/h}}
           \approx 15{,}4\ \text{min}$
      \end{formulabox}
    \end{column}
  \end{columns}
\end{frame}

% ---------------------------------------------------------------
\seccionframe{7. Ejercicios multinivel}
% ---------------------------------------------------------------

\begin{frame}{Ejercicios multinivel}
  \begin{columns}[T]
    \begin{column}{0.31\textwidth}
      \begin{aplicarbox}
        \AlmFontFootnote\dificultad{2}\\[4pt]
        Calcula la distancia entre los pares de puntos:\\[3pt]
        $P(0,0)$ y $Q(3,4)$\\
        $R(1,2)$ y $S(4,6)$\\
        $M(-2,1)$ y $N(2,4)$\\[4pt]
        ¿Cuál es la distancia más corta?
        Convierte a metros (escala 1 u = 200 m).
      \end{aplicarbox}
    \end{column}
    \begin{column}{0.31\textwidth}
      \begin{analizarbox}
        \AlmFontFootnote\dificultad{3}\\[4pt]
        La Alcazaba está en $(1,4)$ y el
        faro del Cabo de Gata en $(8,0)$
        (escala 1 u = 1 km).\\[4pt]
        La distancia en línea recta es
        $\approx 8{,}06$ km. Pero la carretera
        mide 12 km.\\[4pt]
        \textbf{Analiza:} ¿Por qué difieren?
        ¿Cuándo usarías cada medida?
        Argumenta en 5 líneas.
      \end{analizarbox}
    \end{column}
    \begin{column}{0.31\textwidth}
      \begin{evaluarbox}
        \AlmFontFootnote\dificultad{4}\\[4pt]
        Tienes \textbf{2 horas} y caminas a
        4 km/h. Diseña una ruta turística
        por Almería visitando \textbf{al menos
        5 lugares} de la lista anterior.\\[4pt]
        Calcula la distancia total y el
        tiempo. ¿Es factible?\\[4pt]
        Justifica tu elección de ruta:
        ¿optimizas distancia o número de
        lugares?
      \end{evaluarbox}
    \end{column}
  \end{columns}
\end{frame}

% ---------------------------------------------------------------
\begin{frame}{Evidencia del día}
  \begin{evidenciabox}
    \begin{itemize}
      \item Entrega tu \textbf{Mapa de ruta turística} con:
        \begin{enumerate}
          \item Los puntos turísticos marcados con sus coordenadas
          \item Las distancias calculadas en cada tramo (con operaciones)
          \item La distancia total de la ruta en metros y km
          \item El tiempo estimado de recorrido a pie
        \end{enumerate}
      \item Las operaciones deben aparecer escritas, no solo el resultado.
      \item Indica la escala usada.
    \end{itemize}
  \end{evidenciabox}
\end{frame}

% ---------------------------------------------------------------
\begin{frame}{¡Cuidado con este error!}
  \begin{errorbox}
    \begin{columns}[c]
      \begin{column}{0.48\textwidth}
        \incorrecto\ \textbf{Error grave}\\[4pt]
        Olvidar la raíz cuadrada:\\[4pt]
        $d = (x_2-x_1)^2 + (y_2-y_1)^2$\\
        $d = 3^2 + 4^2 = 9+16 = 25$\\[4pt]
        \textcolor{AlmTerra}{Eso es $d^2$, ¡no $d$!}
      \end{column}
      \begin{column}{0.48\textwidth}
        \correcto\ \textbf{Correcto}\\[4pt]
        Siempre aplica la raíz:\\[4pt]
        $d = \sqrt{(x_2-x_1)^2+(y_2-y_1)^2}$\\
        $d = \sqrt{3^2+4^2} = \sqrt{25} = \mathbf{5}$
      \end{column}
    \end{columns}
  \end{errorbox}
  \vspace{6pt}
  \begin{notabox}[Truco de comprobación]
    \AlmFontSmall
    La distancia siempre debe ser \textbf{mayor que cada cateto
    por separado} y \textbf{menor que su suma}.
    Si $a=3$ y $b=4$, entonces $3 < d < 7$. ¿Cumple $d=5$? ¡Sí!
  \end{notabox}
\end{frame}

% ---------------------------------------------------------------
\begin{frame}{Resumen de la sesión 10}
  \begin{resumenbox}
    \begin{itemize}
      \item La \textbf{distancia euclídea} (línea recta) entre
            $A(x_1,y_1)$ y $B(x_2,y_2)$ es:
            $d = \sqrt{(x_2-x_1)^2+(y_2-y_1)^2}$
      \item Es la aplicación directa del
            \textbf{Teorema de Pitágoras} al plano cartesiano.
      \item La \textbf{escala} convierte unidades del mapa
            a metros reales: $d_{\text{real}} = d_{\text{mapa}}\times E$
      \item La distancia por calles siempre es
            $\geq$ la distancia en línea recta.
      \item El \textbf{tiempo} de recorrido: $t = d / v$
    \end{itemize}
  \end{resumenbox}
  \vspace{6pt}
  \begin{center}
    {\AlmFontSmall\color{AlmAzulM}\faArrowCircleRight\enskip
    \textbf{Próximo bloque:}
    Datos y gráficas (Sesiones 11--16)}
  \end{center}
  \fuente{Datos de distancias: estimaciones didácticas sobre plano de Almería}
\end{frame}

\end{document}




% % ================================================================
%  SESIÓN 11 — Variables urbanas
%  Almería en Cifras y Formas · Matemáticas 1.º ESO
%  IES Al-Ándalus · 2025-2026
% ================================================================
\documentclass[aspectratio=169, 12pt, spanish]{beamer}
\usepackage{almeria-beamer}

\renewcommand{\SessionNumber}{11}
\renewcommand{\SessionTitle}{Variables urbanas}
\renewcommand{\SessionName}{Sesión 11}
\renewcommand{\SessionObjective}{Identificar variables cuantitativas y cualitativas en el contexto urbano de Almería, distinguiendo variables independientes de dependientes}
\renewcommand{\BlockName}{Bloque 3: Datos y gráficas}

\begin{document}

% ---------------------------------------------------------------
%  PORTADA
% ---------------------------------------------------------------
\begin{frame}[plain]
  \titlepage
\end{frame}

% ---------------------------------------------------------------
%  FRAME 1 — ¿Qué es una variable?
% ---------------------------------------------------------------
\begin{frame}{¿Qué es una variable?}
  \begin{definicionbox}{Variable}
    Una \textbf{variable} es una característica o cualidad que puede tomar
    \textbf{diferentes valores} en los distintos elementos de un conjunto.
    Si todos los elementos tuvieran el mismo valor, no sería variable.
  \end{definicionbox}

  \vspace{0.4cm}
  \begin{columns}[T]
    \begin{column}{0.48\textwidth}
      \begin{ejemplobox}
        \small
        \begin{itemize}
          \item \textbf{Temperatura} de cada día de julio en Almería
          \item \textbf{Número de visitantes} al Alcázar cada mes
          \item \textbf{Tipo de monumento} (árabe, renacentista, moderno\ldots)
          \item \textbf{Barrio} de residencia de cada alumno
        \end{itemize}
      \end{ejemplobox}
    \end{column}
    \begin{column}{0.48\textwidth}
      \begin{notabox}[Dato curioso]
        \small
        El término \textit{variable} viene del latín \textit{variabilis}:
        \textbf{que puede cambiar}.
        En estadística, cada columna de una tabla de datos
        suele representar una variable distinta.
      \end{notabox}
    \end{column}
  \end{columns}
\end{frame}

% ---------------------------------------------------------------
%  FRAME 2 — Clasificación de variables (árbol TikZ)
% ---------------------------------------------------------------
\begin{frame}{Tipos de variables: árbol de clasificación}
  \begin{center}
  \begin{tikzpicture}[
      every node/.style={font=\small},
      nodo/.style={rounded corners=4pt, draw, align=center,
                   inner sep=4pt, text width=2.8cm},
      nodov/.style={nodo, fill=AlmAzul!15, draw=AlmAzul, text=AlmAzul},
      nodoc/.style={nodo, fill=AlmAmbar!20, draw=AlmAmbar!70!black},
      nodod/.style={nodo, fill=AlmTerra!12, draw=AlmTerra},
      flecha/.style={-Stealth, AlmAzulM, line width=1pt},
      scale=0.92
    ]
    % Raíz
    \node[nodov] (var) at (0,0) {\textbf{VARIABLE}};
    % Hijos nivel 1
    \node[nodoc] (cual) at (-3.5,-1.6) {Cualitativa\\\small(describe una\\cualidad)};
    \node[nodoc] (cuant) at (3.5,-1.6) {Cuantitativa\\\small(valor numérico)};
    % Hijos nivel 2
    \node[nodod] (disc) at (2.0,-3.4) {Discreta\\\small números enteros,\\\small contables};
    \node[nodod] (cont) at (5.0,-3.4) {Continua\\\small cualquier valor\\\small en un intervalo};
    % Flechas
    \draw[flecha] (var) -- (cual);
    \draw[flecha] (var) -- (cuant);
    \draw[flecha] (cuant) -- (disc);
    \draw[flecha] (cuant) -- (cont);
    % Ejemplos
    \node[font=\tiny, text=AlmGris, align=center] at (-3.5,-2.9)
      {barrio, tipo de\\monumento, medio\\de transporte};
    \node[font=\tiny, text=AlmGris, align=center] at (2.0,-4.5)
      {turistas/día,\\nº de autobuses};
    \node[font=\tiny, text=AlmGris, align=center] at (5.0,-4.5)
      {temperatura,\\altura, distancia};
  \end{tikzpicture}
  \end{center}
\end{frame}

% ---------------------------------------------------------------
%  FRAME 3 — Cualitativa vs Cuantitativa (detalles)
% ---------------------------------------------------------------
\begin{frame}{Cualitativa y cuantitativa: diferencias clave}
  \begin{columns}[T]
    \begin{column}{0.48\textwidth}
      \begin{definicionbox}{Variable cualitativa}
        \small
        Describe una \textbf{cualidad o categoría}.
        \textbf{No es un número} (aunque a veces se codifique con números).\\[4pt]
        \textit{Ejemplos en Almería:}
        \begin{itemize}\setlength\itemsep{2pt}
          \item Barrio del municipio
          \item Estilo del monumento
          \item Medio de transporte usado
          \item Temporada turística
        \end{itemize}
      \end{definicionbox}
    \end{column}
    \begin{column}{0.48\textwidth}
      \begin{definicionbox}{Variable cuantitativa}
        \small
        Toma \textbf{valores numéricos} con los que tiene sentido operar.\\[4pt]
        \begin{itemize}\setlength\itemsep{2pt}
          \item \textbf{Discreta}: valores enteros y contables\\
            \textit{nº de autobuses de línea, visitantes por día}
          \item \textbf{Continua}: cualquier valor real en un rango\\
            \textit{temperatura, velocidad del viento, consumo de agua}
        \end{itemize}
      \end{definicionbox}
    \end{column}
  \end{columns}

  \vspace{0.3cm}
  \begin{notabox}[Truco para distinguirlas]
    \small
    Pregúntate: \textbf{«¿Tiene sentido calcular la media?»}
    Si sí (media de temperatura = 24{,}3\,°C \correcto),
    es cuantitativa.
    Si no («media de barrios» no tiene sentido \incorrecto),
    es cualitativa.
  \end{notabox}
\end{frame}

% ---------------------------------------------------------------
%  FRAME 4 — Variable independiente y dependiente
% ---------------------------------------------------------------
\begin{frame}{Variable independiente y dependiente}
  \begin{teoriabox}
    \begin{itemize}
      \item \textbf{Variable independiente} (\(x\)): la que controlamos o elegimos.
            Causa el cambio. Va en el \textbf{eje horizontal}.
      \item \textbf{Variable dependiente} (\(y\)): la que cambia \textit{como consecuencia} de \(x\).
            Va en el \textbf{eje vertical}.
    \end{itemize}
  \end{teoriabox}

  \vspace{0.3cm}
  \begin{center}
  \begin{tikzpicture}
    \node[rounded corners=5pt, fill=AlmAzulM!15, draw=AlmAzulM,
          minimum width=3.2cm, minimum height=1cm, font=\bfseries\small]
      (xi) at (0,0) {Mes del año\\{\tiny(variable independiente \(x\))}};
    \node[rounded corners=5pt, fill=AlmAmbar!25, draw=AlmAmbar!70!black,
          minimum width=3.6cm, minimum height=1cm, font=\bfseries\small]
      (yi) at (7,0) {Nº de turistas\\{\tiny(variable dependiente \(y\))}};
    \draw[-Stealth, AlmTerra, line width=2pt] (xi.east) -- node[above,font=\small,color=AlmTerra]
      {\textit{determina}} (yi.west);
  \end{tikzpicture}
  \end{center}

  \vspace{0.2cm}
  \begin{columns}[T]
    \begin{column}{0.55\textwidth}
      \begin{ejemplobox}
        \small
        \textbf{Pregunta clave}: ¿Qué \textit{causa} a qué?\\
        \begin{itemize}
          \item «A más temperatura \textrightarrow\ más consumo de agua»\\
            $x$ = temperatura; $y$ = consumo
          \item «A más lluvia \textrightarrow\ menos turistas»\\
            $x$ = precipitación; $y$ = turistas
        \end{itemize}
      \end{ejemplobox}
    \end{column}
    \begin{column}{0.42\textwidth}
      \begin{notabox}[Fórmula mental]
        \small
        \[\text{Cuanto más } x \;\Rightarrow\; \text{más/menos } y\]
        Si tiene sentido, \(x\) es la independiente.
      \end{notabox}
    \end{column}
  \end{columns}
\end{frame}

% ---------------------------------------------------------------
%  FRAME 5 — Tabla de variables urbanas de Almería
% ---------------------------------------------------------------
\begin{frame}{Variables urbanas de Almería: tabla de ejemplos}
  \begin{center}
  \small
  \renewcommand{\arraystretch}{1.25}
  \begin{tabular}{>{\bfseries\color{AlmAzulM}}l l l l}
    \toprule
    \rowcolor{AlmAzul!12}
    \multicolumn{1}{c}{\textbf{Variable}} & \textbf{Tipo} & \textbf{Rol} & \textbf{Se relaciona con\ldots} \\
    \midrule
    Mes del año & Cuant. discreta & $x$ (independiente) & Nº de visitantes ($y$) \\
    \rowcolor{AlmFondo}
    Temperatura media & Cuant. continua & $x$ & Consumo de agua ($y$) \\
    Barrio & Cualitativa & $x$ & Precio medio de vivienda ($y$) \\
    \rowcolor{AlmFondo}
    Hora del día & Cuant. discreta & $x$ & Intensidad de tráfico ($y$) \\
    Estación del año & Cualitativa & $x$ & Tipo de turista ($y$) \\
    \rowcolor{AlmFondo}
    Nº de bares en calle & Cuant. discreta & $y$ & Mes del año ($x$) \\
    \bottomrule
  \end{tabular}
  \end{center}

  \vspace{0.2cm}
  \begin{notabox}[Observa]
    \small
    Una variable puede ser \textbf{independiente en un par}
    y \textbf{dependiente en otro}.
    El rol depende de la \textit{relación} que estudiamos.
  \end{notabox}
\end{frame}

% ---------------------------------------------------------------
%  FRAME 6 — Gráfico de dispersión: turistas por mes
% ---------------------------------------------------------------
\begin{frame}{Turistas en Almería por mes: diagrama de puntos}
  \begin{center}
  \begin{tikzpicture}
    \begin{axis}[
      almeria,
      width=0.82\linewidth, height=0.55\linewidth,
      xlabel={Mes (1 = enero \ldots 12 = diciembre)},
      ylabel={Turistas (miles)},
      title={Turistas mensuales en Almería (datos aproximados)},
      xtick={1,2,3,4,5,6,7,8,9,10,11,12},
      xticklabels={E,F,M,A,M,J,J,A,S,O,N,D},
      ymin=0, ymax=700,
      ytick={0,100,200,300,400,500,600,700},
      yticklabels={0,100,200,300,400,500,600,700},
    ]
      \addplot[
        only marks,
        mark=*, mark size=3.5pt,
        color=AlmAzulM,
      ] coordinates {
        (1,60) (2,65) (3,80) (4,120) (5,150)
        (6,250) (7,500) (8,600) (9,300) (10,200)
        (11,100) (12,80)
      };
      \addplot[
        AlmTerra!70, dashed, line width=1pt, smooth
      ] coordinates {
        (1,60) (2,65) (3,80) (4,120) (5,150)
        (6,250) (7,500) (8,600) (9,300) (10,200)
        (11,100) (12,80)
      };
      \node[font=\tiny\bfseries, color=AlmTerra, above] at (axis cs:8,600)
        {Máximo: agosto};
      \node[font=\tiny\bfseries, color=AlmAzulM, below right] at (axis cs:1,60)
        {Mínimo: enero};
    \end{axis}
  \end{tikzpicture}
  \end{center}
  \vspace{-0.1cm}
  {\small $x$ = mes (variable independiente) \quad $y$ = turistas en miles (variable dependiente)}
\end{frame}

% ---------------------------------------------------------------
%  FRAME 7 — Ejercicio RECORDAR
% ---------------------------------------------------------------
\begin{frame}{Ejercicio: clasificar variables}
  \begin{recordarbox}
    Clasifica cada variable como \textbf{cualitativa},
    \textbf{cuantitativa discreta} o \textbf{cuantitativa continua}:

    \vspace{0.25cm}
    \renewcommand{\arraystretch}{1.3}
    \begin{tabular}{clc}
      \toprule
      \textbf{N.º} & \textbf{Variable} & \textbf{Tipo} \\
      \midrule
      1 & Altura de la Torre de la Catedral de Almería & \textcolor{AlmGris}{\textit{(completa)}} \\
      2 & Tipo de transporte utilizado & \\
      3 & Visitantes diarios al Alcázar & \\
      4 & Velocidad del viento en el Puerto & \\
      5 & Estilo artístico del monumento & \\
      6 & Consumo de agua por hogar (litros) & \\
      \bottomrule
    \end{tabular}
  \end{recordarbox}

  \vspace{0.2cm}
  \begin{notabox}[Pista]
    \small
    Pregunta 1: la altura puede ser 98{,}36\,m → admite decimales → \textbf{continua}.
    Para las demás, aplica el mismo razonamiento.
  \end{notabox}
\end{frame}

% ---------------------------------------------------------------
%  FRAME 8 — Ejercicio COMPRENDER
% ---------------------------------------------------------------
\begin{frame}{Ejercicio: temperatura y ventas de helados}
  \begin{comprenderbox}
    El ayuntamiento recoge datos sobre la \textbf{temperatura diaria}
    y las \textbf{ventas de helados} en el Paseo de Almería.
    Responde:

    \vspace{0.2cm}
    \begin{enumerate}
      \item ¿Cuál es la variable independiente (\(x\))?
            ¿Por qué?
      \item ¿Cuál es la variable dependiente (\(y\))?
      \item Explica con tus palabras la relación entre ambas.
      \item ¿Qué ocurrirá con las ventas de helados si la temperatura
            sube de 25\,°C a 38\,°C?
            ¿Y si baja a 10\,°C?
    \end{enumerate}
  \end{comprenderbox}

  \vspace{0.3cm}
  \begin{columns}[T]
    \begin{column}{0.48\textwidth}
      \begin{notabox}[Para reflexionar]
        \small
        ¿Podría ser que las ventas de helados
        \textit{causaran} el aumento de temperatura?
        ¿Tiene sentido invertir el par?
      \end{notabox}
    \end{column}
    \begin{column}{0.48\textwidth}
      \begin{notabox}[Vocabulario]
        \small
        «A mayor temperatura, mayores ventas» →
        relación \textbf{directa} o \textbf{positiva}.
      \end{notabox}
    \end{column}
  \end{columns}
\end{frame}

% ---------------------------------------------------------------
%  FRAME 9 — Ejercicio APLICAR
% ---------------------------------------------------------------
\begin{frame}{Ejercicio: inventa pares de variables urbanas}
  \begin{aplicarbox}
    Inventa \textbf{dos pares de variables} propias de Almería:

    \vspace{0.2cm}
    \textbf{Par A}: la variable independiente \(x\) es \textit{el tiempo}
    (hora, día, mes o año).\\
    \textbf{Par B}: la variable independiente \(x\) es
    \textit{una magnitud física} (distancia, temperatura, altura\ldots).

    \vspace{0.2cm}
    Para cada par indica:
    \begin{enumerate}
      \item Nombre y clasificación de \(x\)
      \item Nombre y clasificación de \(y\)
      \item Predicción: «Cuanto más \(x\)\ldots»
    \end{enumerate}
  \end{aplicarbox}

  \vspace{0.2cm}
  \begin{ejemplobox}
    \small
    \textbf{Ejemplo de Par A}: $x$ = mes del año (cuantitativa discreta);
    $y$ = litros de agua consumidos (cuantitativa continua).\\
    Predicción: «En verano, a mayor calor, mayor consumo».
  \end{ejemplobox}
\end{frame}

% ---------------------------------------------------------------
%  FRAME 10 — Ejercicio ANALIZAR
% ---------------------------------------------------------------
\begin{frame}{Ejercicio: una variable oculta}
  \begin{analizarbox}
    El siguiente conjunto de datos muestra turistas en Almería
    en el mes de \textbf{agosto} de varios años:

    \vspace{0.2cm}
    \begin{center}
    \small
    \begin{tabular}{cc}
      \toprule
      \textbf{Año (\(x\))} & \textbf{Turistas agosto (\(y\))} \\
      \midrule
      2019 & 620\,000 \\
      2020 & \textcolor{AlmTerra}{\textbf{180\,000}} \\
      2021 & 390\,000 \\
      2022 & 590\,000 \\
      \bottomrule
    \end{tabular}
    \end{center}

    \vspace{0.2cm}
    \begin{enumerate}
      \item Identifica \(x\) e \(y\) y clasifica cada una.
      \item En 2020 hay una \textbf{caída enorme}.
            El año por sí solo no la explica.
            ¿Qué \textit{variable oculta} podría haber causado ese descenso?
      \item ¿Podría haber más de una variable oculta?
            Razona la respuesta.
    \end{enumerate}
  \end{analizarbox}
\end{frame}

% ---------------------------------------------------------------
%  FRAME 11 — Error frecuente
% ---------------------------------------------------------------
\begin{frame}{¡Cuidado con este error!}
  \begin{errorbox}
    \textbf{Confundir el sentido de la causalidad}:

    \vspace{0.3cm}
    \begin{center}
    \begin{tikzpicture}
      \node[rounded corners=4pt, fill=AlmTerra!15, draw=AlmTerra,
            text width=5cm, align=center, font=\small] (mal)
        at (0,0) {«El \textbf{número de turistas}\\
                  hace que \textbf{cambie el mes}»};
      \node[right=0.5cm of mal, font=\Large, color=AlmTerra] {\incorrecto};
    \end{tikzpicture}

    \vspace{0.4cm}
    \begin{tikzpicture}
      \node[rounded corners=4pt, fill=AlmVerde!12, draw=AlmVerde,
            text width=5cm, align=center, font=\small] (bien)
        at (0,0) {«El \textbf{mes del año}\\
                  influye en el \textbf{número de turistas}»};
      \node[right=0.5cm of bien, font=\Large, color=AlmVerde] {\correcto};
    \end{tikzpicture}
    \end{center}

    \vspace{0.3cm}
    \small
    El \textbf{tiempo} (mes, año, hora) siempre es la variable
    \textit{independiente}: el tiempo avanza solo, no lo controla
    ningún dato de la tabla.
  \end{errorbox}
\end{frame}

% ---------------------------------------------------------------
%  FRAME 12 — Evidencia del día
% ---------------------------------------------------------------
\begin{frame}{Evidencia del día}
  \begin{evidenciabox}
    Completa en tu cuaderno la siguiente \textbf{tabla de variables urbanas}
    con al menos \textbf{5 pares}. Para cada uno indica:

    \vspace{0.2cm}
    \renewcommand{\arraystretch}{1.3}
    \begin{tabular}{>{\bfseries}c l l l l}
      \toprule
      \textbf{Par} & \textbf{Variable \(x\)} & \textbf{Tipo de \(x\)} &
      \textbf{Variable \(y\)} & \textbf{Tipo de \(y\)} \\
      \midrule
      1 & Mes del año & C. discreta & Nº visitantes Alcazaba & C. discreta \\
      2 & & & & \\
      3 & & & & \\
      4 & & & & \\
      5 & & & & \\
      \bottomrule
    \end{tabular}

    \vspace{0.3cm}
    \small
    Indica con \textbf{una flecha} la dirección de la dependencia:
    $x \longrightarrow y$.
    Escribe una frase prediciendo la relación («a mayor $x$\ldots»).
  \end{evidenciabox}
\end{frame}

% ---------------------------------------------------------------
%  FRAME 13 — Resumen de sesión
% ---------------------------------------------------------------
\begin{frame}{Resumen: lo que hemos aprendido hoy}
  \begin{resumenbox}
    \begin{columns}[T]
      \begin{column}{0.48\textwidth}
        \textbf{Tipos de variable:}
        \begin{itemize}
          \item \textbf{Cualitativa}: describe categorías
          \item \textbf{Cuantitativa discreta}: valores enteros
          \item \textbf{Cuantitativa continua}: cualquier valor real
        \end{itemize}
      \end{column}
      \begin{column}{0.48\textwidth}
        \textbf{Independiente vs. dependiente:}
        \begin{itemize}
          \item $x$ \textbf{independiente}: la que controla o causa
          \item $y$ \textbf{dependiente}: la que cambia por $x$
          \item Pregunta: ¿qué \textit{causa} a qué?
        \end{itemize}
      \end{column}
    \end{columns}

    \vspace{0.3cm}
    \separador
    \vspace{0.2cm}
    \small
    \textbf{Contexto de Almería:} mes del año ($x$) \textrightarrow\ turistas ($y$);
    temperatura ($x$) \textrightarrow\ consumo de agua ($y$);
    barrio ($x$) \textrightarrow\ precio vivienda ($y$).
  \end{resumenbox}

  \vspace{0.2cm}
  \begin{center}
    \bloomtag{Recordar}\quad
    \bloomtag{Comprender}\quad
    \bloomtag{Aplicar}\quad
    \bloomtag{Analizar}
  \end{center}
\end{frame}

\end{document}

% % ================================================================
%  SESIÓN 12 — De enunciado a tabla
%  Almería en Cifras y Formas · Matemáticas 1.º ESO
%  IES Al-Ándalus · 2025-2026
% ================================================================
\documentclass[aspectratio=169, 10pt, spanish]{beamer}
\usepackage{almeria-beamer}

\renewcommand{\SessionNumber}{12}
\renewcommand{\SessionTitle}{De enunciado a tabla}
\renewcommand{\SessionName}{Sesión 12}
\renewcommand{\SessionObjective}{Organizar datos de un enunciado textual en una tabla ordenada y funcional, identificando las variables y sus valores}
\renewcommand{\BlockName}{Bloque 3: Datos y gráficas}

\begin{document}

% ---------------------------------------------------------------
%  PORTADA
% ---------------------------------------------------------------
\begin{frame}[plain]
  \titlepage
\end{frame}

% ---------------------------------------------------------------
%  FRAME 1 — ¿Por qué necesitamos tablas?
% ---------------------------------------------------------------
\begin{frame}{¿Por qué organizar datos en una tabla?}
  \begin{columns}[T]
    \begin{column}{0.52\textwidth}
      \begin{teoriabox}
        \AlmFontSmall
        Un \textbf{texto con datos} es difícil de analizar:
        \begin{itemize}
          \item Los números están dispersos entre palabras.
          \item No se aprecia el orden ni la magnitud.
          \item Es imposible hacer una gráfica directamente.
        \end{itemize}
        \vspace{4pt}
        Una \textbf{tabla} organiza la información de forma
        \textit{sistemática, ordenada y visual}.
      \end{teoriabox}
    \end{column}
    \begin{column}{0.44\textwidth}
      \begin{notabox}[Dato de Almería]
        \AlmFontSmall
        El Patronato de Turismo de Almería publica sus estadísticas
        en tablas mensuales.
        Sin esa organización, los datos serían ilegibles.
      \end{notabox}

      \vspace{0.3cm}
      \begin{center}
      \begin{tikzpicture}
        \node[font=\AlmFontLargeUp, color=AlmTerra] (texto) at (0,0)
          {\faFileAlt};
        \node[font=\AlmFontSmall, color=AlmTerra, below=2pt of texto]
          {Texto disperso};
        \draw[-Stealth, AlmAzulM, line width=2pt] (0.6,0) -- (1.6,0);
        \node[font=\AlmFontLargeUp, color=AlmVerde] (tabla) at (2.2,0)
          {\faTable};
        \node[font=\AlmFontSmall, color=AlmVerde, below=2pt of tabla]
          {Tabla ordenada};
      \end{tikzpicture}
      \end{center}
    \end{column}
  \end{columns}
\end{frame}

% ---------------------------------------------------------------
%  FRAME 2 — Estructura de una tabla
% ---------------------------------------------------------------
\begin{frame}{Estructura de una tabla de datos}
  \begin{center}
  \begin{tikzpicture}[font=\AlmFontSmall, every node/.style={align=center}]
    % Tabla visual con anotaciones
    \draw[AlmAzulM, line width=1.5pt] (0,0) rectangle (8,3.4);
    % Cabecera
    \fill[AlmAzul!85] (0,2.8) rectangle (8,3.4);
    \node[text=white, font=\bfseries\AlmFontSmall] at (2.4,3.1)
      {Encabezado col.~1};
    \draw[AlmAzulM, line width=1pt] (4.8,0) -- (4.8,3.4);
    \node[text=white, font=\bfseries\AlmFontSmall] at (6.4,3.1)
      {Encabezado col.~2};
    % Filas
    \foreach \y/\col in {2.2/AlmFondo, 1.6/white, 1.0/AlmFondo, 0.4/white}{
      \fill[\col, opacity=0.6] (0.1,\y-0.4) rectangle (7.9,\y+0.2);
    }
    \node at (2.4,2.1) {Fila 1 · celda A}; \node at (6.4,2.1) {Fila 1 · celda B};
    \node at (2.4,1.5) {Fila 2 · celda A}; \node at (6.4,1.5) {Fila 2 · celda B};
    \node at (2.4,0.9) {Fila 3 · celda A}; \node at (6.4,0.9) {Fila 3 · celda B};
    \node at (2.4,0.3) {\ldots};           \node at (6.4,0.3) {\ldots};
    % Anotaciones externas
    \draw[acota] (-0.7,2.8) -- (-0.7,3.4)
      node[midway, left=2pt, font=\AlmFontTiny\bfseries, color=AlmAzulM]
      {\rotatebox{90}{Encabezados}};
    \draw[acota] (-0.7,0) -- (-0.7,2.8)
      node[midway, left=2pt, font=\AlmFontTiny\bfseries, color=AlmGris]
      {\rotatebox{90}{Filas de datos}};
    \draw[<->, >=Stealth, AlmTerra, line width=0.8pt]
      (0,3.6) -- (4.7,3.6)
      node[midway, above, font=\AlmFontTiny\bfseries, color=AlmTerra]
      {Columna \(x\) independiente};
    \draw[<->, >=Stealth, AlmVerde, line width=0.8pt]
      (4.9,3.6) -- (8,3.6)
      node[midway, above, font=\AlmFontTiny\bfseries, color=AlmVerde]
      {Columna \(y\) dependiente};
  \end{tikzpicture}
  \end{center}

  \vspace{0.1cm}
  \begin{center}
  \AlmFontSmall
  \textbf{Regla:} columna izquierda = variable independiente \(x\);
  columna derecha = variable dependiente \(y\).
  \end{center}
\end{frame}

% ---------------------------------------------------------------
%  FRAME 3 — Pasos para construir una tabla
% ---------------------------------------------------------------
\begin{frame}{Pasos para construir una tabla desde un texto}
  \begin{columns}[T]
    \begin{column}{0.48\textwidth}
      \begin{pasobox}[1]
        \AlmFontSmall Lee con atención y \textbf{subraya todos los números}.
      \end{pasobox}

      \begin{pasobox}[2]
        \AlmFontSmall Identifica \textbf{qué representa cada número}
        (¿es una fecha? ¿una cantidad? ¿una medida?).
      \end{pasobox}

      \begin{pasobox}[3]
        \AlmFontSmall Decide cuál es la \textbf{variable independiente \(x\)}
        y cuál la \textbf{dependiente \(y\)}.
      \end{pasobox}
    \end{column}
    \begin{column}{0.48\textwidth}
      \begin{pasobox}[4]
        \AlmFontSmall Crea los \textbf{encabezados} con nombre claro y unidades
        entre paréntesis: \textit{Mes}, \textit{Visitantes (miles)}.
      \end{pasobox}

      \begin{pasobox}[5]
        \AlmFontSmall Rellena los valores en \textbf{orden creciente de \(x\)}
        (de menor a mayor).
      \end{pasobox}

      \begin{pasobox}[6]
        \AlmFontSmall \textbf{Verifica}: ¿están todos los datos?
        ¿son consistentes las unidades?
      \end{pasobox}
    \end{column}
  \end{columns}

  \vspace{0.2cm}
  \begin{notabox}[Clave]
    \AlmFontSmall
    Ordenar por \(x\) no es un capricho: facilita la lectura
    y es imprescindible para trazar una gráfica correcta después.
  \end{notabox}
\end{frame}

% ---------------------------------------------------------------
%  FRAME 4 — Ejemplo: visitantes a la Alcazaba
% ---------------------------------------------------------------
\begin{frame}{Ejemplo 1: visitantes a la Alcazaba}
  \begin{columns}[T]
    \begin{column}{0.50\textwidth}
      \begin{ejemplobox}
        \AlmFontSmall\itshape
        «En enero la Alcazaba recibió 8\,500 visitantes.
        En febrero, 9\,200. Marzo registró 12\,000 visitantes.
        En abril hubo 18\,500. En mayo la cifra llegó a 22\,000 personas.»
      \end{ejemplobox}

      \vspace{0.3cm}
      \begin{pasobox}[Análisis]
        \AlmFontSmall
        \begin{itemize}
          \item \(x\) = mes (cuantitativa discreta)
          \item \(y\) = visitantes (cuantitativa discreta)
          \item Ordenar: enero → mayo
        \end{itemize}
      \end{pasobox}
    \end{column}
    \begin{column}{0.46\textwidth}
      \begin{center}
      \renewcommand{\arraystretch}{1.35}
      \begin{tabular}{>{\bfseries\color{AlmAzulM}}c
                      >{\centering\arraybackslash}c
                      r}
        \toprule
        \rowcolor{AlmAzul!12}
        \textbf{Mes} & \textbf{N.º mes} & \textbf{Visitantes} \\
        \midrule
        Enero    & 1 & 8\,500  \\
        \rowcolor{AlmFondo}
        Febrero  & 2 & 9\,200  \\
        Marzo    & 3 & 12\,000 \\
        \rowcolor{AlmFondo}
        Abril    & 4 & 18\,500 \\
        Mayo     & 5 & 22\,000 \\
        \bottomrule
      \end{tabular}
      \end{center}
      \vspace{0.2cm}
      \begin{notabox}[Observa]
        \AlmFontSmall
        ¡Los datos estaban desordenados en el texto!
        La tabla los ordena por el mes.
      \end{notabox}
    \end{column}
  \end{columns}
\end{frame}

% ---------------------------------------------------------------
%  FRAME 5 — Ejemplo 2: tarifa de taxi en Almería
% ---------------------------------------------------------------
\begin{frame}{Ejemplo 2: tarifa de taxi en Almería}
  \begin{columns}[T]
    \begin{column}{0.50\textwidth}
      \begin{ejemplobox}
        \AlmFontSmall\itshape
        «La bajada de bandera es 2{,}40\,€.
        Cada kilómetro recorrido cuesta 1{,}20\,€ adicionales.
        Calcula el precio para 1, 2, 3, 4 y 5\,km.»
      \end{ejemplobox}

      \vspace{0.25cm}
      \begin{formulabox}
        $y = 2{,}40 + 1{,}20 \cdot x$
      \end{formulabox}

      \vspace{0.2cm}
      {\AlmFontSmall $x$ = km recorridos \quad $y$ = precio total (€)}
    \end{column}
    \begin{column}{0.46\textwidth}
      \begin{center}
      \renewcommand{\arraystretch}{1.35}
      \begin{tabular}{>{\bfseries\color{AlmAzulM}\centering\arraybackslash}p{1.8cm}
                      >{\centering\arraybackslash}p{2.4cm}}
        \toprule
        \rowcolor{AlmAzul!12}
        \textbf{Km (\(x\))} & \textbf{Precio €\,(\(y\))} \\
        \midrule
        1 & 2{,}40 + 1{,}20 = 3{,}60 \\
        \rowcolor{AlmFondo}
        2 & 2{,}40 + 2{,}40 = 4{,}80 \\
        3 & 2{,}40 + 3{,}60 = 6{,}00 \\
        \rowcolor{AlmFondo}
        4 & 2{,}40 + 4{,}80 = 7{,}20 \\
        5 & 2{,}40 + 6{,}00 = 8{,}40 \\
        \bottomrule
      \end{tabular}
      \end{center}
    \end{column}
  \end{columns}
\end{frame}

% ---------------------------------------------------------------
%  FRAME 6 — Tabla de frecuencias
% ---------------------------------------------------------------
\begin{frame}{Tabla de frecuencias: cuando muchos datos se repiten}
  \begin{columns}[T]
    \begin{column}{0.50\textwidth}
      \begin{teoriabox}
        \AlmFontSmall
        Cuando hay \textbf{muchas repeticiones} del mismo valor,
        usamos una \textbf{tabla de frecuencias}:
        \begin{itemize}
          \item \textbf{Valor}: cada valor posible de la variable.
          \item \textbf{Marcas de conteo}: palotes para contar
                (5 palotes = \texttt{IIII\hspace{-1pt}/}).
          \item \textbf{Frecuencia}: número total de repeticiones.
        \end{itemize}
      \end{teoriabox}
    \end{column}
    \begin{column}{0.46\textwidth}
      \begin{center}
      \AlmFontSmall
      \textit{Ejemplo: tipo de transporte en Almería (encuesta a 30 alumnos)}\\[4pt]
      \renewcommand{\arraystretch}{1.25}
      \begin{tabular}{lcc}
        \toprule
        \rowcolor{AlmAzul!12}
        \textbf{Transporte} & \textbf{Marcas} & \textbf{Frec.} \\
        \midrule
        A pie      & \texttt{IIIII IIIII II}  & 12 \\
        \rowcolor{AlmFondo}
        Autobús    & \texttt{IIIII IIII}       &  9 \\
        Coche      & \texttt{IIIII I}          &  6 \\
        \rowcolor{AlmFondo}
        Bicicleta  & \texttt{III}              &  3 \\
        \midrule
        \textbf{Total} & & \textbf{30} \\
        \bottomrule
      \end{tabular}
      \end{center}
    \end{column}
  \end{columns}
\end{frame}

% ---------------------------------------------------------------
%  FRAME 7 — Gráfico de barras a partir de la tabla (pgfplots)
% ---------------------------------------------------------------
\begin{frame}{De la tabla al gráfico: Alcazaba enero–mayo}
  \begin{center}
  \begin{tikzpicture}
    \begin{axis}[
      almeria bar,
      width=0.82\linewidth, height=0.56\linewidth,
      title={Visitantes mensuales a la Alcazaba},
      xlabel={Mes},
      ylabel={Visitantes},
      symbolic x coords={Enero,Febrero,Marzo,Abril,Mayo},
      xtick=data,
      ymin=0, ymax=26000,
      nodes near coords style={font=\AlmFontTiny\bfseries, color=AlmAzul},
      bar width=0.6cm,
    ]
      \addplot[fill=AlmAzulM, draw=AlmAzul] coordinates {
        (Enero,8500)
        (Febrero,9200)
        (Marzo,12000)
        (Abril,18500)
        (Mayo,22000)
      };
    \end{axis}
  \end{tikzpicture}
  \end{center}
  \vspace{-0.2cm}
  {\AlmFontSmall\color{AlmGris} La tabla se convierte directamente en este gráfico
  gracias al orden correcto de los datos.}
\end{frame}

% ---------------------------------------------------------------
%  FRAME 8 — Ejercicio RECORDAR
% ---------------------------------------------------------------
\begin{frame}{Ejercicio: elementos de una tabla}
  \begin{recordarbox}
    Responde sin mirar los apuntes:
    \begin{enumerate}
      \item ¿Cuáles son los \textbf{4 elementos} que debe tener
            toda tabla bien estructurada?
      \item ¿Qué información debe aparecer en los
            \textbf{encabezados de columna}?
      \item ¿Qué variable va en la columna \textbf{izquierda}?
            ¿Y en la derecha?
      \item ¿Por qué hay que ordenar los datos
            de menor a mayor valor de \(x\)?
    \end{enumerate}
  \end{recordarbox}

  \vspace{0.3cm}
  \begin{center}
  \AlmFontSmall
  \begin{tabular}{>{\bfseries\color{AlmAzulM}}l l}
    \toprule
    Elemento & Descripción \\
    \midrule
    Encabezados & Nombre de variable + unidad \\
    Filas & Un dato por fila \\
    Columna \(x\) & Variable independiente (izquierda) \\
    Columna \(y\) & Variable dependiente (derecha) \\
    \bottomrule
  \end{tabular}
  \end{center}
\end{frame}

% ---------------------------------------------------------------
%  FRAME 9 — Ejercicio COMPRENDER
% ---------------------------------------------------------------
\begin{frame}{Ejercicio: detectar errores en una tabla}
  \begin{comprenderbox}
    La siguiente tabla tiene \textbf{4 errores}. Identifícalos
    y reescribe la tabla de forma correcta.

    \vspace{0.25cm}
    \begin{center}
    \AlmFontSmall
    \begin{tabular}{cc}
      \toprule
      \rowcolor{AlmTerra!10}
      \textbf{Datos} & \textbf{Cosas} \\
      \midrule
      Julio  & 280000 \\
      Enero  & 60 \\
      Marzo  & 80000 \\
      Agosto & mucho \\
      Febrero& 65000 \\
      \bottomrule
    \end{tabular}
    \end{center}

    \vspace{0.15cm}
    \textit{Pista}: busca problemas en encabezados, unidades,
    orden y consistencia de datos.
  \end{comprenderbox}
\end{frame}

% ---------------------------------------------------------------
%  FRAME 10 — Ejercicio APLICAR
% ---------------------------------------------------------------
\begin{frame}{Ejercicio: el autobús de Almería}
  \begin{aplicarbox}
    Lee el enunciado y construye la tabla correspondiente:

    \vspace{0.2cm}
    \begin{center}
    \begin{tcolorbox}[colback=AlmAmbar!12, colframe=AlmAmbar!60!black,
                      width=0.88\textwidth, arc=4pt, boxrule=0.8pt]
      \AlmFontSmall\itshape
      «El autobús urbano de Almería tarda 5 minutos en hacer 2 paradas.
      Para 4 paradas necesita 10 minutos.
      Con 6 paradas son 15 minutos.
      Las 8 paradas se completan en 20 minutos,
      y las 12 paradas en 30 minutos.»
    \end{tcolorbox}
    \end{center}

    \vspace{0.2cm}
    \begin{enumerate}
      \item Construye la tabla completa con encabezados y unidades.
      \item ¿Cuál es la variable independiente? ¿Y la dependiente?
      \item ¿Ves un patrón? ¿Cuánto tarda para \textbf{10 paradas}?
    \end{enumerate}
  \end{aplicarbox}
\end{frame}

% ---------------------------------------------------------------
%  FRAME 11 — Error frecuente
% ---------------------------------------------------------------
\begin{frame}{¡Cuidado con este error!}
  \begin{errorbox}
    \textbf{Meter los datos en orden aleatorio}:

    \vspace{0.3cm}
    \begin{columns}[T]
      \begin{column}{0.44\textwidth}
        \incorrecto\; \textbf{Tabla desordenada}
        \vspace{4pt}
        \AlmFontSmall
        \begin{tabular}{cc}
          \toprule
          \rowcolor{AlmTerra!10}
          Mes & Visitantes \\
          \midrule
          Julio & 500\,000 \\
          Enero & 60\,000 \\
          Abril & 120\,000 \\
          Marzo & 80\,000 \\
          \bottomrule
        \end{tabular}
      \end{column}
      \begin{column}{0.44\textwidth}
        \correcto\; \textbf{Tabla ordenada}
        \vspace{4pt}
        \AlmFontSmall
        \begin{tabular}{cc}
          \toprule
          \rowcolor{AlmVerde!10}
          Mes & Visitantes \\
          \midrule
          Enero & 60\,000 \\
          Marzo & 80\,000 \\
          Abril & 120\,000 \\
          Julio & 500\,000 \\
          \bottomrule
        \end{tabular}
      \end{column}
    \end{columns}

    \vspace{0.25cm}
    \AlmFontSmall
    Una tabla desordenada \textbf{impide ver la tendencia} y genera gráficas
    con puntos conectados de forma caótica.
    Ordena siempre de menor a mayor valor de \(x\).
  \end{errorbox}
\end{frame}

% ---------------------------------------------------------------
%  FRAME 12 — Evidencia del día
% ---------------------------------------------------------------
\begin{frame}{Evidencia del día}
  \begin{evidenciabox}
    Construye en tu cuaderno \textbf{dos tablas completas}
    a partir de estos dos párrafos:

    \vspace{0.2cm}
    \begin{tcolorbox}[colback=AlmAmbar!10, colframe=AlmAmbar!60!black,
                      arc=3pt, boxrule=0.8pt, fonttitle=\bfseries\AlmFontSmall,
                      title={Párrafo A — Temperaturas en Almería}]
      \AlmFontSmall\itshape
      «En enero la temperatura media es 12\,°C.
      En abril sube a 19\,°C. En julio alcanza 30\,°C.
      Octubre registra 21\,°C y diciembre baja a 13\,°C.»
    \end{tcolorbox}

    \vspace{0.2cm}
    \begin{tcolorbox}[colback=AlmAzulC!10, colframe=AlmAzulC!70!black,
                      arc=3pt, boxrule=0.8pt, fonttitle=\bfseries\AlmFontSmall,
                      title={Párrafo B — Precio del taxi}]
      \AlmFontSmall\itshape
      «Un taxi en Almería cobra 1{,}80\,€/km con bajada de bandera de 3\,€.
      Calcula el precio total para 1, 3, 5, 7 y 10\,km.»
    \end{tcolorbox}

    \vspace{0.2cm}
    \AlmFontSmall
    Cada tabla debe tener: \textbf{encabezados con unidades},
    \textbf{orden correcto de \(x\)} y \textbf{5 filas de datos}.
  \end{evidenciabox}
\end{frame}

% ---------------------------------------------------------------
%  FRAME 13 — Resumen de sesión
% ---------------------------------------------------------------
\begin{frame}{Resumen: lo que hemos aprendido hoy}
  \begin{resumenbox}
    \begin{columns}[T]
      \begin{column}{0.50\textwidth}
      \vspace{0.15cm}
        \textbf{Estructura de una tabla:}
        \begin{itemize}
          \item Encabezados con nombre y unidades
          \item Columna izquierda = \(x\) independiente
          \item Columna derecha = \(y\) dependiente
          \item Filas ordenadas de menor a mayor \(x\)
        \end{itemize}
      \end{column}
      \begin{column}{0.46\textwidth}
        \textbf{6 pasos para construir la tabla:}
        \begin{enumerate}
          \AlmFontSmall
          \item Subraya todos los números
          \item Identifica qué representa cada uno
          \item Decide \(x\) e \(y\)
          \item Crea encabezados claros
          \item Rellena en orden
          \item Verifica consistencia
        \end{enumerate}
      \end{column}
    \end{columns}

    \vspace{0.2cm}
    \separador
    \vspace{0.15cm}
    \AlmFontSmall
    \textbf{Error a evitar:}
    datos en orden aleatorio dificultan la lectura y hacen
    imposible trazar una gráfica correcta.
  \end{resumenbox}

  \vspace{0.2cm}
  \begin{center}
    \bloomtag{Recordar}\quad
    \bloomtag{Comprender}\quad
    \bloomtag{Aplicar}
  \end{center}
\end{frame}

\end{document}



% % ================================================================
%  SESIÓN 13 — De tabla a gráfica
%  Almería en Cifras y Formas · Matemáticas 1.º ESO
%  IES Al-Ándalus · 2025-2026
% ================================================================
\documentclass[aspectratio=169, 12pt, spanish]{beamer}
\usepackage{almeria-beamer}

\renewcommand{\SessionNumber}{13}
\renewcommand{\SessionTitle}{De tabla a gráfica}
\renewcommand{\SessionName}{Sesión 13}
\renewcommand{\SessionObjective}{Representar datos de una tabla en una gráfica (diagrama de barras y gráfica lineal) eligiendo el tipo adecuado según el tipo de variable}
\renewcommand{\BlockName}{Bloque 3: Datos y gráficas}

\begin{document}

% ---------------------------------------------------------------
%  PORTADA
% ---------------------------------------------------------------
\begin{frame}[plain]
  \titlepage
\end{frame}

% ---------------------------------------------------------------
%  FRAME 1 — ¿Por qué representar gráficamente?
% ---------------------------------------------------------------
\begin{frame}{¿Por qué usar gráficas?}
  \begin{columns}[T]
    \begin{column}{0.52\textwidth}
      \begin{teoriabox}
        \small
        Una tabla muestra los datos con \textbf{precisión},
        pero una gráfica los hace \textbf{visibles}:
        \begin{itemize}
          \item Revela \textbf{tendencias} (subida, bajada, estabilidad).
          \item Permite detectar \textbf{valores atípicos} de un vistazo.
          \item Facilita la \textbf{comparación} entre series.
          \item Comunica la información más rápido que un texto.
        \end{itemize}
      \end{teoriabox}
    \end{column}
    \begin{column}{0.44\textwidth}
      \begin{notabox}[Datos → Gráfica → Conclusión]
        \small
        \begin{center}
        \begin{tikzpicture}[node distance=0.7cm,
            caja/.style={rounded corners=4pt, draw=AlmAzulM,
                         fill=AlmFondo, text width=2.2cm,
                         align=center, font=\tiny\bfseries}]
          \node[caja] (d) {Tabla\\de datos};
          \node[caja, below=of d] (g) {Gráfica\\adecuada};
          \node[caja, below=of g] (c) {Conclusión\\estadística};
          \draw[-Stealth, AlmAzulM] (d)--(g);
          \draw[-Stealth, AlmAzulM] (g)--(c);
        \end{tikzpicture}
        \end{center}
      \end{notabox}
    \end{column}
  \end{columns}
\end{frame}

% ---------------------------------------------------------------
%  FRAME 2 — Tipos de gráficas y cuándo usarlas
% ---------------------------------------------------------------
\begin{frame}{Tipos de gráficas y cuándo usarlas}
  \small
  \renewcommand{\arraystretch}{1.4}
  \begin{tabular}{>{\bfseries\color{AlmAzulM}}l l l}
    \toprule
    \rowcolor{AlmAzul!12}
    \textbf{Tipo de gráfica} & \textbf{Cuándo usarla} & \textbf{Ejemplo en Almería} \\
    \midrule
    Diagrama de barras
      & Variable $x$ discreta o categórica
      & Visitantes por barrio \\
    \rowcolor{AlmFondo}
    Gráfica lineal
      & $x$ continua o temporal; muestra tendencia
      & Temperatura media mensual \\
    Diagrama de puntos
      & Dos variables cuantitativas; correlación
      & Temperatura vs.\ turistas \\
    \rowcolor{AlmFondo}
    Diagrama de sectores
      & Partes de un todo (porcentajes)
      & Tipo de transporte usado \\
    \bottomrule
  \end{tabular}

  \vspace{0.3cm}
  \begin{notabox}[Regla práctica]
    \small
    Si $x$ son \textbf{categorías o meses} que comparo → barras.\\
    Si $x$ es \textbf{tiempo continuo} y quiero ver la evolución → línea.
  \end{notabox}
\end{frame}

% ---------------------------------------------------------------
%  FRAME 3 — Elementos de una gráfica (TikZ anotado)
% ---------------------------------------------------------------
\begin{frame}{Elementos de una gráfica bien construida}
  \begin{columns}[T]
    \begin{column}{0.54\textwidth}
      \begin{center}
      \begin{tikzpicture}[font=\scriptsize, scale=0.88]
        % Ejes
        \draw[AlmAzulM, line width=1.2pt, -Stealth] (0,0)--(6.2,0)
          node[right, color=AlmAzulM, font=\tiny\bfseries] {Mes ($x$)};
        \draw[AlmAzulM, line width=1.2pt, -Stealth] (0,0)--(0,4.2)
          node[above, color=AlmAzulM, font=\tiny\bfseries,
               xshift=-0.3cm, rotate=90] {Visitantes ($y$)};
        % Barras de muestra
        \foreach \x/\h/\c in {
          0.6/0.9/AlmAzulC, 1.2/1.0/AlmAzulC, 1.8/1.3/AlmAzulC,
          2.4/1.9/AlmAzulM, 3.0/2.7/AlmAzulM}{
          \fill[\c, opacity=0.7] (\x,0) rectangle (\x+0.5,\h);
        }
        % Título
        \node[color=AlmAzul, font=\tiny\bfseries] at (3.0,4.0)
          {Visitantes mensuales a la Alcazaba};
        % Escala y
        \foreach \h/\v in {1.0/10000, 2.0/20000, 3.0/30000}{
          \draw[AlmGris!40, line width=0.4pt] (0,\h)--(5.8,\h);
          \node[left, color=AlmGris, font=\tiny] at (0,\h) {\v};
        }
        % Etiquetas x
        \foreach \x/\l in {0.85/E, 1.45/F, 2.05/M, 2.65/A, 3.25/M}{
          \node[below=2pt, color=AlmGris] at (\x,0) {\l};
        }
        % Anotaciones externas
        \draw[-Stealth, AlmTerra, thin] (5.4,3.8) -- (4.6,2.8)
          node[above right, color=AlmTerra, font=\tiny\bfseries,
               xshift=8pt, yshift=4pt] {Título};
        \draw[-Stealth, AlmVerde, thin] (5.4,0.5) -- (0.7,0.5)
          node[above right, color=AlmVerde, font=\tiny\bfseries,
               xshift=-24pt, yshift=6pt] {Barra};
        \node[AlmAzulM, font=\tiny\bfseries] at (-1.1,2.0)
          {\rotatebox{90}{Escala $y$}};
        \node[AlmAzulM, font=\tiny\bfseries] at (3.2,-0.5)
          {Etiquetas $x$};
      \end{tikzpicture}
      \end{center}
    \end{column}
    \begin{column}{0.42\textwidth}
      \begin{teoriabox}
        \small
        Todo gráfico debe tener:
        \begin{itemize}
          \item \textbf{Eje $x$}: variable independiente + unidad
          \item \textbf{Eje $y$}: variable dependiente + unidad
          \item \textbf{Escala}: intervalos regulares, comenzando en~0
          \item \textbf{Título}: qué se representa
          \item \textbf{Leyenda}: si hay más de una serie
        \end{itemize}
      \end{teoriabox}
    \end{column}
  \end{columns}
\end{frame}

% ---------------------------------------------------------------
%  FRAME 4 — Pasos para dibujar un gráfico
% ---------------------------------------------------------------
\begin{frame}{Pasos para trazar una gráfica}
  \begin{columns}[T]
    \begin{column}{0.48\textwidth}
      \begin{pasobox}[1]
        \small Decide el \textbf{tipo} de gráfica
        (¿barras o línea?).
      \end{pasobox}\vspace{3pt}
      \begin{pasobox}[2]
        \small Dibuja los \textbf{ejes} $x$ (horizontal) e $y$ (vertical).
      \end{pasobox}\vspace{3pt}
      \begin{pasobox}[3]
        \small \textbf{Etiqueta} cada eje: variable + unidad.
      \end{pasobox}\vspace{3pt}
      \begin{pasobox}[4]
        \small Elige la \textbf{escala}: mayor valor $\times\;1{,}1$
        para el máximo del eje.
      \end{pasobox}
    \end{column}
    \begin{column}{0.48\textwidth}
      \begin{pasobox}[5]
        \small \textbf{Representa} cada par $(x,y)$
        de la tabla como punto.
      \end{pasobox}\vspace{3pt}
      \begin{pasobox}[6]
        \small \textbf{Une} los puntos con línea (gráfica lineal)
        o rellena barras (diagrama de barras).
      \end{pasobox}\vspace{3pt}
      \begin{pasobox}[7]
        \small Añade el \textbf{título} descriptivo.
      \end{pasobox}

      \vspace{0.25cm}
      \begin{notabox}[Clave del paso 4]
        \small
        El eje $y$ \textbf{siempre empieza en 0}
        salvo razón justificada.
      \end{notabox}
    \end{column}
  \end{columns}
\end{frame}

% ---------------------------------------------------------------
%  FRAME 5 — Diagrama de barras: visitantes Almería ene–jun
% ---------------------------------------------------------------
\begin{frame}{Diagrama de barras: turistas en Almería (ene–jun)}
  \begin{center}
  \begin{tikzpicture}
    \begin{axis}[
      almeria bar,
      width=0.82\linewidth, height=0.56\linewidth,
      title={Turistas mensuales en Almería (enero–junio)},
      xlabel={Mes},
      ylabel={Turistas},
      symbolic x coords={Enero,Febrero,Marzo,Abril,Mayo,Junio},
      xtick=data,
      xticklabel style={font=\tiny, rotate=15, anchor=north east},
      ymin=0, ymax=330000,
      ytick={0,50000,100000,150000,200000,250000,300000},
      yticklabels={0,50\,k,100\,k,150\,k,200\,k,250\,k,300\,k},
      nodes near coords style={font=\tiny\bfseries, color=AlmAzul,
                                rotate=90, anchor=west},
    ]
      \addplot[fill=AlmAzulM!75, draw=AlmAzul, fill opacity=0.85] coordinates {
        (Enero,60000)
        (Febrero,65000)
        (Marzo,85000)
        (Abril,130000)
        (Mayo,180000)
        (Junio,280000)
      };
    \end{axis}
  \end{tikzpicture}
  \end{center}
  {\small\color{AlmGris}
   Usamos barras porque los meses son \textbf{categorías discretas}
   que queremos \textit{comparar}.}
\end{frame}

% ---------------------------------------------------------------
%  FRAME 6 — Gráfica lineal: temperatura mensual en Almería
% ---------------------------------------------------------------
\begin{frame}{Gráfica lineal: temperatura media mensual en Almería}
  \begin{center}
  \begin{tikzpicture}
    \begin{axis}[
      almeria line,
      width=0.82\linewidth, height=0.54\linewidth,
      title={Temperatura media mensual en Almería (°C)},
      xlabel={Mes},
      ylabel={Temperatura (°C)},
      xtick={1,2,3,4,5,6,7,8,9,10,11,12},
      xticklabels={E,F,M,A,M,J,J,A,S,O,N,D},
      ymin=8, ymax=34,
      ytick={10,15,20,25,30},
    ]
      \addplot[
        color=AlmTerra, line width=2pt, mark=*, mark size=2.5pt,
        mark options={fill=AlmTerra!70, draw=AlmTerra}
      ] coordinates {
        (1,12)(2,13)(3,16)(4,19)(5,23)(6,27)
        (7,30)(8,30)(9,26)(10,21)(11,16)(12,13)
      };
      \node[font=\tiny\bfseries, color=AlmTerra, above=4pt] at (axis cs:7,30)
        {30\,°C};
      \node[font=\tiny\bfseries, color=AlmAzulM, below=4pt] at (axis cs:1,12)
        {12\,°C};
    \end{axis}
  \end{tikzpicture}
  \end{center}
  {\small\color{AlmGris}
   Usamos línea porque la temperatura varía de forma \textbf{continua}
   a lo largo del tiempo y queremos ver la \textit{tendencia}.}
\end{frame}

% ---------------------------------------------------------------
%  FRAME 7 — Barras vs. línea: ¿cuándo usar cada una?
% ---------------------------------------------------------------
\begin{frame}{Barras vs. línea: la elección correcta}
  \begin{columns}[T]
    \begin{column}{0.48\textwidth}
      \begin{definicionbox}{Diagrama de barras — úsalo cuando:}
        \small
        \begin{itemize}
          \item $x$ son \textbf{categorías separadas}\\
                (tipos, barrios, meses como etiquetas)
          \item Quieres \textbf{comparar magnitudes}
          \item No hay continuidad entre valores de $x$
        \end{itemize}
        \vspace{4pt}
        \textit{Ejemplo}: nº de plazas hoteleras por barrio de Almería
      \end{definicionbox}
    \end{column}
    \begin{column}{0.48\textwidth}
      \begin{definicionbox}{Gráfica lineal — úsala cuando:}
        \small
        \begin{itemize}
          \item $x$ es \textbf{continua o temporal}
                (tiempo, distancia, temperatura)
          \item Quieres mostrar una \textbf{tendencia} o evolución
          \item Tiene sentido hablar de valores intermedios
        \end{itemize}
        \vspace{4pt}
        \textit{Ejemplo}: temperatura media horaria en un día de verano
      \end{definicionbox}
    \end{column}
  \end{columns}

  \vspace{0.3cm}
  \begin{notabox}[¿Y si tengo dudas?]
    \small
    Pregúntate: «¿Tendría sentido un valor \textit{entre} dos puntos del eje $x$?»
    Si sí → línea. Si no → barras.
  \end{notabox}
\end{frame}

% ---------------------------------------------------------------
%  FRAME 8 — Ejercicio RECORDAR
% ---------------------------------------------------------------
\begin{frame}{Ejercicio: tipos de gráficas}
  \begin{recordarbox}
    \begin{enumerate}
      \item Nombra \textbf{4 tipos de gráficas} que conoces.
      \item Para cada una, escribe \textbf{un ejemplo concreto}
            relacionado con Almería en el que sea apropiada.
      \item Completa la tabla:
    \end{enumerate}

    \vspace{0.2cm}
    \small
    \renewcommand{\arraystretch}{1.3}
    \begin{tabular}{lcc}
      \toprule
      \rowcolor{AlmAzul!12}
      \textbf{Situación} & \textbf{¿Barras o línea?} & \textbf{¿Por qué?} \\
      \midrule
      Visitantes por mes (etiquetas) & & \\
      \rowcolor{AlmFondo}
      Temperatura a lo largo del día & & \\
      Cuota de turistas por país & & \\
      \rowcolor{AlmFondo}
      Evolución anual de pernoctaciones & & \\
      \bottomrule
    \end{tabular}
  \end{recordarbox}
\end{frame}

% ---------------------------------------------------------------
%  FRAME 9 — Ejercicio COMPRENDER
% ---------------------------------------------------------------
\begin{frame}{Ejercicio: ¿barras o línea?}
  \begin{comprenderbox}
    Se recogen estos datos sobre alquileres en Almería:

    \vspace{0.2cm}
    \begin{center}
    \small
    \begin{tabular}{cc}
      \toprule
      \rowcolor{AlmAzul!12}
      \textbf{Distancia al centro (km)} & \textbf{Precio alquiler (€/m²)} \\
      \midrule
      0{,}5 & 12 \\
      \rowcolor{AlmFondo}
      1{,}0 & 10 \\
      1{,}5 & 8{,}5 \\
      \rowcolor{AlmFondo}
      2{,}0 & 7 \\
      2{,}5 & 6 \\
      \bottomrule
    \end{tabular}
    \end{center}

    \vspace{0.2cm}
    \begin{enumerate}
      \item ¿Deberías usar un diagrama de barras o una gráfica lineal?
            \textbf{Justifica tu respuesta}.
      \item ¿Cuál es la variable independiente? ¿Y la dependiente?
      \item Describe en una frase la tendencia que observas.
    \end{enumerate}
  \end{comprenderbox}
\end{frame}

% ---------------------------------------------------------------
%  FRAME 10 — Ejercicio APLICAR
% ---------------------------------------------------------------
\begin{frame}{Ejercicio: graficar los datos del autobús}
  \begin{aplicarbox}
    Retoma la tabla del autobús de Almería (sesión 12):

    \vspace{0.2cm}
    \begin{center}
    \small
    \begin{tabular}{cc}
      \toprule
      \rowcolor{AlmAzul!12}
      \textbf{Paradas (\(x\))} & \textbf{Tiempo min (\(y\))} \\
      \midrule
      2 & 5 \\ \rowcolor{AlmFondo}
      4 & 10 \\
      6 & 15 \\ \rowcolor{AlmFondo}
      8 & 20 \\
      12 & 30 \\
      \bottomrule
    \end{tabular}
    \end{center}

    \vspace{0.2cm}
    \begin{enumerate}
      \item Dibuja la gráfica lineal correspondiente en papel cuadriculado.
            Etiqueta ambos ejes con nombre y unidad.
            Elige una escala adecuada para el eje $y$ (máx. 30~min).
      \item ¿Qué \textbf{tendencia} observas?
      \item ¿Qué representa la \textbf{inclinación} de la línea?
    \end{enumerate}
  \end{aplicarbox}
\end{frame}

% ---------------------------------------------------------------
%  FRAME 11 — Error frecuente
% ---------------------------------------------------------------
\begin{frame}{¡Cuidado con este error!}
  \begin{errorbox}
    \textbf{Empezar el eje \(y\) en un valor distinto de cero}
    para exagerar las diferencias:

    \vspace{0.3cm}
    \begin{columns}[T]
      \begin{column}{0.46\textwidth}
        \incorrecto\; \textbf{Eje \(y\) desde 95\,000}\\[4pt]
        \begin{center}
        \begin{tikzpicture}[scale=0.7]
          \begin{axis}[
            width=0.9\linewidth, height=0.7\linewidth,
            axis lines=left, ymin=95000, ymax=105000,
            xtick={1,2,3}, xticklabels={E,F,M},
            ytick={95000,100000,105000},
            yticklabels={95k,100k,105k},
            tick label style={font=\tiny},
            label style={font=\tiny},
            ylabel style={font=\tiny},
            ylabel={Visitantes},
            xlabel={Mes},
            ybar, bar width=0.4cm,
            nodes near coords, nodes near coords style={font=\tiny},
          ]
            \addplot[fill=AlmTerra!60] coordinates
              {(1,97000)(2,100000)(3,104000)};
          \end{axis}
        \end{tikzpicture}
        \end{center}
        \small\color{AlmTerra}
        ¡Parece que la diferencia es enorme!
      \end{column}
      \begin{column}{0.46\textwidth}
        \correcto\; \textbf{Eje \(y\) desde 0}\\[4pt]
        \begin{center}
        \begin{tikzpicture}[scale=0.7]
          \begin{axis}[
            width=0.9\linewidth, height=0.7\linewidth,
            axis lines=left, ymin=0, ymax=115000,
            xtick={1,2,3}, xticklabels={E,F,M},
            ytick={0,50000,100000},
            yticklabels={0,50k,100k},
            tick label style={font=\tiny},
            label style={font=\tiny},
            ylabel style={font=\tiny},
            ylabel={Visitantes},
            xlabel={Mes},
            ybar, bar width=0.4cm,
            nodes near coords, nodes near coords style={font=\tiny},
          ]
            \addplot[fill=AlmVerde!60] coordinates
              {(1,97000)(2,100000)(3,104000)};
          \end{axis}
        \end{tikzpicture}
        \end{center}
        \small\color{AlmVerde}
        La diferencia es real pero pequeña.
      \end{column}
    \end{columns}
  \end{errorbox}
\end{frame}

% ---------------------------------------------------------------
%  FRAME 12 — Evidencia del día
% ---------------------------------------------------------------
\begin{frame}{Evidencia del día}
  \begin{evidenciabox}
    Dibuja en tu cuaderno \textbf{dos gráficas} a partir de las tablas
    que construiste en la sesión anterior:

    \vspace{0.15cm}
    \begin{enumerate}
      \item \textbf{Gráfica A} — Temperaturas medias de Almería:
            \begin{itemize}
              \small
              \item Elige el tipo correcto (barras o línea).
              \item Etiqueta los ejes: «Mes» y «Temperatura (°C)».
              \item Escala: 0–35\,°C.
              \item Título: «Temperatura media mensual en Almería».
            \end{itemize}
      \item \textbf{Gráfica B} — Precio del taxi en Almería:
            \begin{itemize}
              \small
              \item Elige el tipo correcto.
              \item Etiqueta: «Km recorridos» y «Precio (€)».
              \item Escala acorde al máximo de la tabla.
            \end{itemize}
    \end{enumerate}

    \vspace{0.15cm}
    \small
    Para cada gráfica: justifica en una frase por qué elegiste
    ese tipo de representación.
  \end{evidenciabox}
\end{frame}

% ---------------------------------------------------------------
%  FRAME 13 — Resumen de sesión
% ---------------------------------------------------------------
\begin{frame}{Resumen: lo que hemos aprendido hoy}
  \begin{resumenbox}
    \begin{columns}[T]
      \begin{column}{0.48\textwidth}
        \textbf{Elegir el tipo de gráfica:}
        \begin{itemize}
          \item \textbf{Barras}: categorías, comparar
          \item \textbf{Línea}: variable continua, tendencia
          \item \textbf{Puntos}: correlación entre dos variables
          \item \textbf{Sectores}: partes de un todo
        \end{itemize}
      \end{column}
      \begin{column}{0.48\textwidth}
        \textbf{Elementos imprescindibles:}
        \begin{itemize}
          \item Ejes etiquetados con unidades
          \item Escala regular comenzando en 0
          \item Título claro y descriptivo
          \item Leyenda si hay varias series
        \end{itemize}
      \end{column}
    \end{columns}

    \vspace{0.25cm}
    \separador
    \vspace{0.15cm}
    \small
    \textbf{Error crítico:} empezar el eje $y$ en un valor
    distinto de cero distorsiona visualmente las diferencias y
    puede ser \textit{engañoso}.
  \end{resumenbox}

  \vspace{0.2cm}
  \begin{center}
    \bloomtag{Recordar}\quad
    \bloomtag{Comprender}\quad
    \bloomtag{Aplicar}
  \end{center}
\end{frame}

\end{document}

% % ================================================================
%  sesion14.tex  —  Sesión 14 de 20
%  Leer e interpretar gráficas
%  Bloque 3: Datos y gráficas · Matemáticas 1.º ESO
% ================================================================
\documentclass[aspectratio=169, 12pt, spanish]{beamer}
\usepackage{almeria-beamer}

\renewcommand{\SessionNumber}{14}
\renewcommand{\SessionTitle}{Leer e interpretar gráficas}
\renewcommand{\SessionName}{Sesión 14}
\renewcommand{\SessionObjective}{%
  Leer coordenadas de una gráfica, identificar tendencias,
  máximos, mínimos y valores clave, y extraer conclusiones
  sobre la realidad de Almería.}
\renewcommand{\BlockName}{Bloque 3: Datos y gráficas}

\begin{document}

\begin{frame}[plain]\titlepage\end{frame}

\begin{frame}{Índice}
  \begin{columns}[t]
    \begin{column}{0.5\textwidth}
      \begin{enumerate}
        \item Cómo leer un valor de una gráfica
        \item Elementos clave: máximo, mínimo, tendencia
        \item Interpretación de la pendiente
        \item Gráfica de temperatura en Almería
        \item Gráfica de visitantes 2017--2022
      \end{enumerate}
    \end{column}
    \begin{column}{0.5\textwidth}
      \begin{enumerate}\setcounter{enumi}{5}
        \item Preguntas tipo sobre gráficas
        \item Ejercicios multinivel (Bloom)
        \item Evidencia del día
        \item Error frecuente
        \item Resumen
      \end{enumerate}
    \end{column}
  \end{columns}
\end{frame}

% ---------------------------------------------------------------
\seccionframe{1. Cómo leer un valor de una gráfica}
% ---------------------------------------------------------------

\begin{frame}{Protocolo de lectura de una gráfica}
  \begin{columns}[c]
    \begin{column}{0.5\textwidth}
      \begin{pasobox}[1]
        \footnotesize
        Localiza el valor de \textbf{$x$} que te interesa
        en el eje horizontal.
      \end{pasobox}
      \vspace{4pt}
      \begin{pasobox}[2]
        \footnotesize
        Sube en vertical hasta tocar la
        \textbf{curva o barra}.
      \end{pasobox}
      \vspace{4pt}
      \begin{pasobox}[3]
        \footnotesize
        Desplázate en horizontal hacia el
        \textbf{eje vertical} y lee el valor $y$.
      \end{pasobox}
      \vspace{4pt}
      \begin{pasobox}[4]
        \footnotesize
        Escribe el resultado con sus
        \textbf{unidades} (no olvides el eje $y$).
      \end{pasobox}
    \end{column}
    \begin{column}{0.46\textwidth}
      \begin{tikzpicture}[scale=0.82]
        \begin{axis}[
          almeria line,
          width=6.8cm, height=5.2cm,
          xmin=0, xmax=7, ymin=0, ymax=30,
          xtick={1,2,3,4,5,6},
          xticklabels={Ene,Feb,Mar,Abr,May,Jun},
          ytick={0,10,20,30},
          ylabel={Temperatura (°C)},
          xlabel={Mes},
        ]
          \addplot[AlmTerra, line width=2pt, mark=*]
            coordinates {(1,12)(2,13)(3,16)(4,19)(5,23)(6,27)};
          % Líneas de lectura para Abril (x=4, y=19)
          \draw[AlmVerde, dashed, line width=1pt]
            (axis cs:4,0) -- (axis cs:4,19) -- (axis cs:0,19);
          \node[fill=AlmVerde!20, draw=AlmVerde, rounded corners,
                font=\tiny\bfseries, inner sep=2pt]
            at (axis cs:2.5,19) {$y=19\,°C$};
          \node[fill=AlmVerde!20, draw=AlmVerde, rounded corners,
                font=\tiny\bfseries, inner sep=2pt]
            at (axis cs:4,5) {$x=\text{Abr}$};
        \end{axis}
      \end{tikzpicture}
    \end{column}
  \end{columns}
\end{frame}

% ---------------------------------------------------------------
\seccionframe{2. Elementos clave de una gráfica}
% ---------------------------------------------------------------

\begin{frame}{Máximo, mínimo, tendencia y punto de inflexión}
  \begin{columns}[T]
    \begin{column}{0.42\textwidth}
      \begin{definicionbox}{Valor máximo}
        Punto más alto de la gráfica.
        Mayor valor de $y$.
      \end{definicionbox}
      \vspace{5pt}
      \begin{definicionbox}{Valor mínimo}
        Punto más bajo de la gráfica.
        Menor valor de $y$.
      \end{definicionbox}
      \vspace{5pt}
      \begin{definicionbox}{Tendencia}
        Dirección general de la curva:\\
        $\nearrow$ creciente \quad
        $\searrow$ decreciente \quad
        $\rightarrow$ constante
      \end{definicionbox}
      \vspace{5pt}
      \begin{definicionbox}{Punto de inflexión}
        Donde la tendencia cambia de sentido.
      \end{definicionbox}
    \end{column}
    \begin{column}{0.54\textwidth}
      \begin{tikzpicture}
        \begin{axis}[
          almeria line,
          width=8.5cm, height=5.5cm,
          xmin=1, xmax=12, ymin=10, ymax=33,
          xtick={1,2,3,4,5,6,7,8,9,10,11,12},
          xticklabels={E,F,M,A,M,J,J,A,S,O,N,D},
          ylabel={°C}, xlabel={Mes},
          title={Temperatura media mensual en Almería},
        ]
          \addplot[AlmTerra, line width=2pt, mark=*,
                   mark options={fill=AlmTerra}]
            coordinates {
              (1,12)(2,13)(3,16)(4,19)(5,23)(6,27)
              (7,30)(8,30)(9,26)(10,21)(11,16)(12,13)
            };
          % Máximo
          \node[above, font=\tiny\bfseries, color=AlmTerra]
            at (axis cs:7.5,30) {\faArrowUp\ Máx: 30°C};
          % Mínimo
          \node[below right, font=\tiny\bfseries, color=AlmAzulM]
            at (axis cs:1,12) {\faArrowDown\ Mín: 12°C};
          % Tendencia creciente
          \draw[->, AlmVerde, line width=1.2pt]
            (axis cs:2,14) -- (axis cs:6,26)
            node[midway, above left, font=\tiny, color=AlmVerde]
            {creciente};
          % Tendencia decreciente
          \draw[->, BEvaluar, line width=1.2pt]
            (axis cs:8,29) -- (axis cs:11,17)
            node[midway, above right, font=\tiny, color=BEvaluar]
            {decreciente};
        \end{axis}
      \end{tikzpicture}
    \end{column}
  \end{columns}
\end{frame}

% ---------------------------------------------------------------
\seccionframe{3. Interpretación de la pendiente}
% ---------------------------------------------------------------

\begin{frame}{¿Qué nos dice la pendiente de una curva?}
  \begin{columns}[c]
    \begin{column}{0.5\textwidth}
      \begin{teoriabox}
        La \textbf{pendiente} (inclinación) indica la
        \textit{velocidad de cambio}:
        \begin{itemize}
          \item Pendiente \textbf{pronunciada}: el valor
                cambia \textbf{rápido}.
          \item Pendiente \textbf{suave}: el valor cambia
                \textbf{lentamente}.
          \item Línea \textbf{horizontal}: el valor
                \textbf{no cambia}.
        \end{itemize}
      \end{teoriabox}
    \end{column}
    \begin{column}{0.46\textwidth}
      \begin{tikzpicture}[scale=0.85]
        \begin{axis}[
          almeria line,
          width=6.8cm, height=4.8cm,
          xmin=0, xmax=6, ymin=0, ymax=12,
          axis lines=left,
          xtick={1,2,3,4,5},
          ytick={0,4,8,12},
          ylabel={$y$}, xlabel={$x$},
        ]
          % Pendiente pronunciada
          \addplot[AlmTerra, line width=2pt]
            coordinates {(0,0)(1,4)(2,8)(3,12)};
          \node[right, font=\tiny\bfseries, color=AlmTerra]
            at (axis cs:3,12) {pronunciada};
          % Pendiente suave
          \addplot[AlmVerde, line width=2pt]
            coordinates {(0,0)(1,1.5)(2,3)(3,4.5)(5,7.5)};
          \node[right, font=\tiny\bfseries, color=AlmVerde]
            at (axis cs:5,7.5) {suave};
          % Horizontal
          \addplot[AlmAzulM, line width=2pt, dashed]
            coordinates {(0,2)(5,2)};
          \node[right, font=\tiny\bfseries, color=AlmAzulM]
            at (axis cs:4,2.5) {constante};
        \end{axis}
      \end{tikzpicture}
    \end{column}
  \end{columns}
\end{frame}

% ---------------------------------------------------------------
\seccionframe{4. Temperatura en Almería}
% ---------------------------------------------------------------

\begin{frame}{Análisis de la gráfica de temperatura}
  \begin{columns}[c]
    \begin{column}{0.55\textwidth}
      \begin{tikzpicture}
        \begin{axis}[
          almeria line,
          width=9cm, height=5.5cm,
          xmin=1, xmax=12, ymin=8, ymax=34,
          xtick={1,...,12},
          xticklabels={E,F,M,A,M,J,J,A,S,O,N,D},
          ylabel={Temperatura (°C)},
          xlabel={Mes},
          title={Temperatura media · Almería},
          ytick={10,15,20,25,30},
        ]
          \addplot[AlmTerra, line width=2.5pt, mark=*,
                   mark options={fill=AlmTerra, scale=1.2}]
            coordinates {
              (1,12)(2,13)(3,16)(4,19)(5,23)(6,27)
              (7,30)(8,30)(9,26)(10,21)(11,16)(12,13)
            };
          % Línea de referencia 20°C
          \addplot[AlmAmbar!70, dashed, line width=1pt]
            coordinates {(1,20)(12,20)}
            node[right, font=\tiny, color=AlmAmbar!80!black]
            {20°C};
        \end{axis}
      \end{tikzpicture}
    \end{column}
    \begin{column}{0.42\textwidth}
      Respondemos las preguntas clave:
      \begin{itemize}
        \item \textbf{Máximo:} julio--agosto, $30\,°C$
        \item \textbf{Mínimo:} enero, $12\,°C$
        \item \textbf{Rango:} $30-12 = 18\,°C$
        \item \textbf{Meses $>20\,°C$:} junio, julio, agosto, septiembre
        \item \textbf{Tendencia} ene→ago: \textit{creciente}
        \item \textbf{Tendencia} ago→dic: \textit{decreciente}
        \item \textbf{Simetría:} casi simétrica respecto a agosto
      \end{itemize}
      \vspace{6pt}
      \begin{notabox}
        \footnotesize
        Almería tiene uno de los climas
        más cálidos y secos de Europa.
        ¡Más de 3\,000 horas de sol al año!
      \end{notabox}
    \end{column}
  \end{columns}
\end{frame}

% ---------------------------------------------------------------
\seccionframe{5. Visitantes 2017--2022}
% ---------------------------------------------------------------

\begin{frame}{Turismo en Almería: gráfica de barras 2017--2022}
  \begin{columns}[c]
    \begin{column}{0.58\textwidth}
      \begin{tikzpicture}
        \begin{axis}[
          almeria bar,
          width=9.5cm, height=5.5cm,
          symbolic x coords={2017,2018,2019,2020,2021,2022},
          xtick=data,
          ylabel={Visitantes (millones)},
          xlabel={Año},
          title={Visitantes a Almería (millones)},
          ymin=0, ymax=3.5,
          nodes near coords style={font=\tiny\bfseries},
          bar width=0.55cm,
        ]
          \addplot[fill=AlmAzulM, fill opacity=0.85]
            coordinates {
              (2017,2.5)(2018,2.7)(2019,2.9)
              (2020,1.2)(2021,1.8)(2022,2.8)
            };
          % Línea de tendencia general
          \addplot[AlmTerra, line width=1.5pt, dashed, no marks]
            coordinates {(2017,2.5)(2018,2.7)(2019,2.9)
                          (2020,1.2)(2021,1.8)(2022,2.8)};
        \end{axis}
      \end{tikzpicture}
    \end{column}
    \begin{column}{0.38\textwidth}
      Datos relevantes:
      \begin{itemize}
        \item Máximo: 2019 → 2,9 M
        \item Mínimo: 2020 → 1,2 M
        \item Caída 2019→2020:\\
          $\dfrac{2{,}9-1{,}2}{2{,}9}\times 100
          \approx \mathbf{59\,\%}$
        \item Recuperación 2021→2022:\\
          $\dfrac{2{,}8-1{,}8}{1{,}8}\times 100
          \approx \mathbf{56\,\%}$
      \end{itemize}
      \vspace{6pt}
      \begin{analizarbox}
        \footnotesize
        ¿Qué evento de 2020 explica
        la caída drástica de visitantes?
        ¿Qué nos indica la recuperación
        de 2022?
      \end{analizarbox}
    \end{column}
  \end{columns}
\end{frame}

% ---------------------------------------------------------------
\seccionframe{6. Preguntas tipo sobre gráficas}
% ---------------------------------------------------------------

\begin{frame}{Preguntas que siempre debes poder responder}
  \begin{columns}[T]
    \begin{column}{0.48\textwidth}
      \textbf{Lectura de valores:}
      \begin{itemize}
        \item «¿Cuál es el valor de $y$ cuando $x = \ldots$?»
        \item «¿En qué mes/año fue el valor más alto?»
        \item «¿Entre qué valores oscila $y$?»
      \end{itemize}
      \vspace{8pt}
      \textbf{Tendencias:}
      \begin{itemize}
        \item «¿La gráfica es creciente, decreciente
              o constante entre $x_1$ y $x_2$?»
        \item «¿Dónde está el punto de inflexión?»
        \item «¿Cuál es la tendencia general?»
      \end{itemize}
    \end{column}
    \begin{column}{0.48\textwidth}
      \textbf{Comparación:}
      \begin{itemize}
        \item «¿Entre qué dos puntos fue mayor
              el crecimiento?»
        \item «¿Cuál es la variación total de $y$?»
      \end{itemize}
      \vspace{8pt}
      \textbf{Interpretación:}
      \begin{itemize}
        \item «¿Qué puede explicar la anomalía en $x = \ldots$?»
        \item «¿Qué predices para el siguiente valor?»
        \item «¿Qué conclusión puedes extraer del gráfico?»
      \end{itemize}
      \vspace{8pt}
      \begin{notabox}[Consejo]
        \footnotesize
        Siempre lee el \textbf{título}, los
        \textbf{ejes} y la \textbf{unidad} antes
        de responder cualquier pregunta.
      \end{notabox}
    \end{column}
  \end{columns}
\end{frame}

% ---------------------------------------------------------------
\seccionframe{7. Ejercicios multinivel}
% ---------------------------------------------------------------

\begin{frame}{Ejercicios multinivel}
  \begin{columns}[T]
    \begin{column}{0.31\textwidth}
      \begin{comprenderbox}
        \footnotesize\dificultad{1}\\[4pt]
        Usando la gráfica de temperatura de Almería:\\[4pt]
        (a) ¿Temperatura en abril?\\
        (b) ¿Mes más caluroso?\\
        (c) ¿Cuántos meses supera 20°C?\\
        (d) Describe la tendencia de enero a agosto
        en una frase.
      \end{comprenderbox}
    \end{column}
    \begin{column}{0.31\textwidth}
      \begin{aplicarbox}
        \footnotesize\dificultad{2}\\[4pt]
        Con la gráfica de visitantes:\\[4pt]
        (a) Mejor año y número de visitantes.\\
        (b) Peor año y número de visitantes.\\
        (c) Calcula el \% de cambio de 2021 a 2022.\\
        (d) Si la tendencia continúa, ¿cuántos
        visitantes habría en 2023?
      \end{aplicarbox}
    \end{column}
    \begin{column}{0.31\textwidth}
      \begin{analizarbox}
        \footnotesize\dificultad{3}\\[4pt]
        La gráfica muestra el número de terrazas
        abiertas en Almería por mes: máximo en
        agosto, mínimo en enero.\\[4pt]
        \textbf{Analiza:} ¿Qué relación tiene con
        la temperatura? ¿Con el turismo?
        ¿Qué variable \textit{causa} la apertura
        de terrazas? Argumenta en 6 líneas.
      \end{analizarbox}
    \end{column}
  \end{columns}
\end{frame}

% ---------------------------------------------------------------
\begin{frame}{Evidencia del día}
  \begin{evidenciabox}
    Escribe un \textbf{análisis de dos gráficas} (temperatura y
    visitantes) que incluya:
    \begin{enumerate}
      \item Los valores leídos en al menos 3 puntos de cada gráfica
      \item Identificación del máximo, mínimo y tendencia general
      \item Una anomalía o punto de inflexión con su posible explicación
      \item Una conclusión de 3--4 líneas sobre lo que la gráfica revela
            sobre Almería
    \end{enumerate}
    \vspace{4pt}
    \textit{No basta con leer los números: debes interpretarlos.}
  \end{evidenciabox}
\end{frame}

% ---------------------------------------------------------------
\begin{frame}{¡Cuidado con este error!}
  \begin{errorbox}
    \begin{columns}[c]
      \begin{column}{0.48\textwidth}
        \incorrecto\ Leer los ejes al revés:\\[4pt]
        «En el mes de julio, la temperatura
        es 7» ← Leyó el eje $x$ donde debía
        leer el eje $y$.
      \end{column}
      \begin{column}{0.48\textwidth}
        \correcto\ Protocolo correcto:\\[4pt]
        1. Localiza julio en el eje \textbf{horizontal}.\\
        2. Sube hasta la curva.\\
        3. Lee el valor en el eje \textbf{vertical}: $30\,°C$.
      \end{column}
    \end{columns}
  \end{errorbox}
  \vspace{6pt}
  \begin{errorbox}
    \incorrecto\ «En 2020 bajaron los visitantes porque
    no había buenos hoteles.»\\[4pt]
    \correcto\ Una gráfica muestra \textbf{qué} pasó,
    pero no siempre explica \textbf{por qué}.
    Hay que buscar la explicación con otras fuentes o razonar
    con más datos (en 2020 hubo pandemia mundial de COVID-19).
  \end{errorbox}
\end{frame}

% ---------------------------------------------------------------
\begin{frame}{Resumen de la sesión 14}
  \begin{resumenbox}
    \begin{itemize}
      \item Para leer un valor: localiza $x$ → sube a la curva
            → lee $y$ en el eje vertical.
      \item Identificamos \textbf{máximo}, \textbf{mínimo},
            \textbf{tendencia} (creciente/decreciente/constante)
            y \textbf{punto de inflexión}.
      \item La pendiente indica la \textbf{velocidad de cambio}.
      \item La gráfica de temperatura de Almería es \textit{creciente}
            de enero a julio--agosto, y \textit{decreciente} de agosto
            a diciembre.
      \item Anomalías (como 2020) requieren \textbf{interpretación
            contextual}, no solo matemática.
    \end{itemize}
  \end{resumenbox}
  \vspace{6pt}
  \begin{center}
    {\small\color{AlmAzulM}\faArrowCircleRight\enskip
    \textbf{Sesión 15:} Comparar gráficas y detectar correlaciones}
  \end{center}
\end{frame}

\end{document}

% % ================================================================
%  sesion15.tex  —  Sesión 15 de 20
%  Comparar gráficas
%  Bloque 3: Datos y gráficas · Matemáticas 1.º ESO
% ================================================================
\documentclass[aspectratio=169, 12pt, spanish]{beamer}
\usepackage{almeria-beamer}

\renewcommand{\SessionNumber}{15}
\renewcommand{\SessionTitle}{Comparar gráficas}
\renewcommand{\SessionName}{Sesión 15}
\renewcommand{\SessionObjective}{%
  Superponer y comparar dos gráficas sobre el mismo sistema de ejes,
  identificar correlaciones positivas y negativas, y extraer
  conclusiones comparativas sobre datos urbanos de Almería.}
\renewcommand{\BlockName}{Bloque 3: Datos y gráficas}

\begin{document}

\begin{frame}[plain]\titlepage\end{frame}

\begin{frame}{Índice}
  \begin{columns}[t]
    \begin{column}{0.5\textwidth}
      \begin{enumerate}
        \item ¿Por qué comparar gráficas?
        \item Tipos de correlación
        \item Temperatura vs.\ turismo en Almería
        \item Lluvia vs.\ aforo en playas
        \item Correlación $\neq$ causalidad
      \end{enumerate}
    \end{column}
    \begin{column}{0.5\textwidth}
      \begin{enumerate}\setcounter{enumi}{5}
        \item Cómo «miente» una gráfica
        \item Ejercicios multinivel (Bloom)
        \item Evidencia del día
        \item Error frecuente
        \item Resumen
      \end{enumerate}
    \end{column}
  \end{columns}
\end{frame}

% ---------------------------------------------------------------
\seccionframe{1. ¿Por qué comparar gráficas?}
% ---------------------------------------------------------------

\begin{frame}{La comparación revela relaciones}
  \begin{columns}[c]
    \begin{column}{0.52\textwidth}
      \begin{teoriabox}
        Superponer dos gráficas en los mismos ejes permite:
        \begin{itemize}
          \item Detectar si dos variables
                \textbf{suben y bajan juntas}
                (correlación positiva).
          \item Detectar si una \textbf{sube cuando
                la otra baja} (correlación negativa).
          \item Identificar \textbf{retrasos} (lag):
                una variable responde a la otra con demora.
          \item Descubrir períodos de
                \textbf{comportamiento anómalo}.
        \end{itemize}
      \end{teoriabox}
    \end{column}
    \begin{column}{0.44\textwidth}
      \begin{tikzpicture}[every node/.style={font=\scriptsize}]
        \begin{axis}[
          almeria line,
          width=6.5cm, height=4.5cm,
          xmin=1, xmax=6, ymin=0, ymax=11,
          xtick={1,...,6},
          xticklabels={E,F,M,A,M,J},
          ylabel={Valor (escalado)},
          xlabel={Mes},
          legend pos=north west,
          legend style={font=\tiny},
        ]
          \addplot[AlmTerra, line width=2pt, mark=*]
            coordinates {(1,2)(2,3)(3,5)(4,7)(5,9)(6,10)};
          \addlegendentry{Variable A}
          \addplot[AlmAzulM, line width=2pt, mark=square*]
            coordinates {(1,8)(2,7)(3,6)(4,4)(5,3)(6,2)};
          \addlegendentry{Variable B}
        \end{axis}
      \end{tikzpicture}
      \begin{center}
        \small\color{AlmGris}
        A sube, B baja → correlación \textbf{negativa}
      \end{center}
    \end{column}
  \end{columns}
\end{frame}

% ---------------------------------------------------------------
\seccionframe{2. Tipos de correlación}
% ---------------------------------------------------------------

\begin{frame}{Correlación: positiva, negativa o nula}
  \begin{columns}[T]
    \begin{column}{0.31\textwidth}
    \centering
    {\color{AlmVerde}\bfseries Correlación positiva}\\[4pt]
    \begin{tikzpicture}
      \begin{axis}[
        width=4.5cm, height=3.8cm,
        axis lines=left, xmin=0, xmax=6, ymin=0, ymax=10,
        xtick=\empty, ytick=\empty,
        tick style={draw=none},
      ]
        \addplot[AlmVerde, line width=2pt, no marks]
          coordinates {(0,1)(5,9)};
        \addplot[only marks, mark=*, AlmVerde!80, mark size=2pt]
          coordinates {(1,1.8)(2,3.5)(3,5.1)(4,6.8)(5,8.5)};
      \end{axis}
    \end{tikzpicture}\\[4pt]
    \footnotesize Cuando $x$ $\uparrow$, también $y$ $\uparrow$\\
    Ejemplo: temperatura $\uparrow$ → turistas $\uparrow$
    \end{column}
    \begin{column}{0.31\textwidth}
    \centering
    {\color{AlmTerra}\bfseries Correlación negativa}\\[4pt]
    \begin{tikzpicture}
      \begin{axis}[
        width=4.5cm, height=3.8cm,
        axis lines=left, xmin=0, xmax=6, ymin=0, ymax=10,
        xtick=\empty, ytick=\empty,
        tick style={draw=none},
      ]
        \addplot[AlmTerra, line width=2pt, no marks]
          coordinates {(0,9)(5,1)};
        \addplot[only marks, mark=*, AlmTerra!80, mark size=2pt]
          coordinates {(1,8.2)(2,6.5)(3,4.9)(4,3.2)(5,1.5)};
      \end{axis}
    \end{tikzpicture}\\[4pt]
    \footnotesize Cuando $x$ $\uparrow$, $y$ $\downarrow$\\
    Ejemplo: lluvia $\uparrow$ → aforo playa $\downarrow$
    \end{column}
    \begin{column}{0.31\textwidth}
    \centering
    {\color{AlmGris}\bfseries Sin correlación}\\[4pt]
    \begin{tikzpicture}
      \begin{axis}[
        width=4.5cm, height=3.8cm,
        axis lines=left, xmin=0, xmax=6, ymin=0, ymax=10,
        xtick=\empty, ytick=\empty,
        tick style={draw=none},
      ]
        \addplot[only marks, mark=*, AlmGris, mark size=2pt]
          coordinates {
            (0.5,2)(1,8)(1.5,4)(2,6)(2.5,1)(3,9)
            (3.5,3)(4,7)(4.5,2)(5,5)(5.5,8)
          };
      \end{axis}
    \end{tikzpicture}\\[4pt]
    \footnotesize Los puntos no siguen ningún patrón\\
    Ejemplo: n.º de farolas $\leftrightarrow$ turistas
    \end{column}
  \end{columns}
\end{frame}

% ---------------------------------------------------------------
\seccionframe{3. Temperatura vs.\ turismo}
% ---------------------------------------------------------------

\begin{frame}{Temperatura y turismo en Almería — dos ejes}
  \begin{columns}[c]
    \begin{column}{0.62\textwidth}
      \begin{tikzpicture}
        \begin{axis}[
          width=10cm, height=5.8cm,
          axis y line*=left,
          axis x line=bottom,
          xmin=1, xmax=12, ymin=8, ymax=33,
          xtick={1,...,12},
          xticklabels={E,F,M,A,M,J,J,A,S,O,N,D},
          ylabel={\color{AlmTerra}Temperatura (°C)},
          xlabel={Mes},
          tick label style={font=\tiny},
          label style={font=\small},
          title={Temperatura (°C) y turistas ($\times 10^4$)},
          title style={font=\small\bfseries, color=AlmAzul},
          grid=both,
          grid style={line width=0.3pt, AlmAzulC!25},
        ]
          \addplot[AlmTerra, line width=2.5pt, mark=*,
                   mark options={fill=AlmTerra}]
            coordinates {
              (1,12)(2,13)(3,16)(4,19)(5,23)(6,27)
              (7,30)(8,30)(9,26)(10,21)(11,16)(12,13)
            };
          \addlegendentry{\footnotesize Temperatura}
        \end{axis}
        \begin{axis}[
          width=10cm, height=5.8cm,
          axis y line*=right,
          axis x line=none,
          xmin=1, xmax=12, ymin=0, ymax=80,
          ylabel={\color{AlmAzulM}Turistas ($\times 10^4$/mes)},
          tick label style={font=\tiny},
          label style={font=\small},
          ytick={0,20,40,60,80},
        ]
          \addplot[AlmAzulM, line width=2.5pt, mark=square*,
                   mark options={fill=AlmAzulM}, dashed]
            coordinates {
              (1,6)(2,6.5)(3,8)(4,12)(5,15)(6,25)
              (7,50)(8,60)(9,30)(10,20)(11,10)(12,8)
            };
          \addlegendentry{\footnotesize Turistas}
        \end{axis}
      \end{tikzpicture}
    \end{column}
    \begin{column}{0.34\textwidth}
      \textbf{Observaciones:}
      \begin{itemize}
        \item Ambas curvas tienen el \textbf{pico en julio--agosto}.
        \item Ambas tienen el \textbf{mínimo en enero}.
        \item Cuando la temperatura $\uparrow$,
              los turistas también $\uparrow$.
        \item \textbf{Correlación positiva} clara.
      \end{itemize}
      \vspace{6pt}
      \begin{notabox}[Doble eje $y$]
        \footnotesize
        Se usan dos ejes $y$ cuando las
        variables tienen \textbf{escalas muy
        diferentes} (°C vs.\ decenas de miles).
      \end{notabox}
    \end{column}
  \end{columns}
\end{frame}

% ---------------------------------------------------------------
\seccionframe{4. Lluvia vs.\ aforo en playas}
% ---------------------------------------------------------------

\begin{frame}{Lluvia y aforo de playas — correlación negativa}
  \begin{columns}[c]
    \begin{column}{0.58\textwidth}
      \begin{tikzpicture}
        \begin{axis}[
          almeria bar,
          width=9.5cm, height=5.2cm,
          symbolic x coords={E,F,M,A,M,J,J,A,S,O,N,D},
          xtick=data,
          ylabel={Lluvia (mm) / Aforo (\%)},
          xlabel={Mes},
          title={Lluvia mensual (barras) y aforo de playa (línea)},
          ymin=0, ymax=100,
          bar width=0.38cm,
          legend pos=north west,
          legend style={font=\tiny},
          nodes near coords=\empty,
        ]
          \addplot[fill=AlmAzulC!60, fill opacity=0.8]
            coordinates {
              (E,25)(F,22)(M,18)(A,12)(M,8)(J,3)
              (J,1)(A,1)(S,8)(O,20)(N,28)(D,30)
            };
          \addlegendentry{Lluvia (mm)}
          \addplot[AlmTerra, line width=2.5pt, mark=*,
                   mark options={fill=AlmTerra}]
            coordinates {
              (E,15)(F,18)(M,22)(A,35)(M,50)(J,70)
              (J,95)(A,98)(S,65)(O,35)(N,18)(D,12)
            };
          \addlegendentry{Aforo playa (\%)}
        \end{axis}
      \end{tikzpicture}
    \end{column}
    \begin{column}{0.38\textwidth}
      \textbf{Análisis:}
      \begin{itemize}
        \item Lluvia alta (otoño--invierno) →
              aforo de playa \textbf{bajo}.
        \item Lluvia baja (verano) →
              aforo de playa \textbf{alto}.
        \item \textbf{Correlación negativa}: cuando
              una variable sube, la otra baja.
      \end{itemize}
      \vspace{6pt}
      \begin{analizarbox}
        \footnotesize
        ¿La lluvia es la \textit{causa} de que
        las playas estén vacías, o existen
        otras variables explicativas?
        (temperatura, vacaciones, etc.)
      \end{analizarbox}
    \end{column}
  \end{columns}
\end{frame}

% ---------------------------------------------------------------
\seccionframe{5. Correlación $\neq$ Causalidad}
% ---------------------------------------------------------------

\begin{frame}{¡Atención! Correlación no implica causalidad}
  \begin{columns}[T]
    \begin{column}{0.55\textwidth}
      \begin{errorbox}
        \textbf{Error de razonamiento frecuente:}\\[4pt]
        «Como la temperatura y el turismo suben juntos,
        \textit{la temperatura es la causa del turismo.}»\\[6pt]
        Pero también hay \textbf{otra variable oculta}:
        las \textit{vacaciones escolares} coinciden con
        el verano. Son las vacaciones (no solo el calor)
        las que permiten a las familias viajar.
      \end{errorbox}
      \vspace{6pt}
      \begin{notabox}[Variable oculta o confusora]
        \footnotesize
        Una \textbf{variable confusora} es aquella que
        influye en $x$ e $y$ simultáneamente, haciendo
        que parezca que $x$ causa $y$ cuando en realidad
        ambas están causadas por una tercera variable.
      \end{notabox}
    \end{column}
    \begin{column}{0.42\textwidth}
      \textbf{Ejemplo clásico:}\\[6pt]
      \begin{center}
      \begin{tikzpicture}[
        nodo/.style={
          rounded corners=4pt, draw=AlmAzulM,
          fill=AlmFondo, font=\small, inner sep=5pt,
          text width=2.5cm, align=center
        },
        causa/.style={->, AlmTerra, line width=1.5pt}
      ]
        \node[nodo] (V) at (0,2.5) {Vacaciones\\de verano};
        \node[nodo] (T) at (-2,0) {Temperatura\\alta (°C)};
        \node[nodo] (Tu) at (2,0) {Turistas\\$\uparrow$};
        \draw[causa] (V) -- (T);
        \draw[causa] (V) -- (Tu);
        \draw[AlmGris, dashed, ->, line width=1pt] (T) -- (Tu)
          node[midway, below, font=\tiny, color=AlmGris]
          {correlación\,$\neq$\,causa};
      \end{tikzpicture}
      \end{center}
      \vspace{4pt}
      \footnotesize\color{AlmGris}
      Las vacaciones causan TANTO que haya mucha gente
      en la playa COMO que observemos días calurosos
      en los registros turísticos.
    \end{column}
  \end{columns}
\end{frame}

% ---------------------------------------------------------------
\seccionframe{6. Cómo «miente» una gráfica}
% ---------------------------------------------------------------

\begin{frame}{La escala puede engañar}
  \begin{columns}[T]
    \begin{column}{0.48\textwidth}
      \centering
      {\color{AlmTerra}\bfseries\small Gráfica «alarmista»}\\
      (eje $y$ truncado: 0 omitido)\\[4pt]
      \begin{tikzpicture}
        \begin{axis}[
          width=6cm, height=4.2cm,
          axis lines=left,
          xmin=2019.5, xmax=2022.5,
          ymin=2.7, ymax=3.1,
          xtick={2020,2021,2022},
          ytick={2.7,2.8,2.9,3.0,3.1},
          xlabel={Año}, ylabel={Mill.},
          tick label style={font=\tiny},
          label style={font=\tiny},
          title style={font=\tiny},
          title={¡Caída catastrófica!},
        ]
          \addplot[AlmTerra, line width=2.5pt, mark=*]
            coordinates {(2020,2.75)(2021,2.85)(2022,3.0)};
        \end{axis}
      \end{tikzpicture}
    \end{column}
    \begin{column}{0.48\textwidth}
      \centering
      {\color{AlmVerde}\bfseries\small Gráfica honesta}\\
      (eje $y$ desde 0)\\[4pt]
      \begin{tikzpicture}
        \begin{axis}[
          width=6cm, height=4.2cm,
          axis lines=left,
          xmin=2019.5, xmax=2022.5,
          ymin=0, ymax=3.5,
          xtick={2020,2021,2022},
          ytick={0,1,2,3},
          xlabel={Año}, ylabel={Mill.},
          tick label style={font=\tiny},
          label style={font=\tiny},
          title style={font=\tiny},
          title={Crecimiento moderado},
        ]
          \addplot[AlmVerde, line width=2.5pt, mark=*]
            coordinates {(2020,2.75)(2021,2.85)(2022,3.0)};
        \end{axis}
      \end{tikzpicture}
    \end{column}
  \end{columns}
  \vspace{6pt}
  \begin{center}
    \small ¡Los mismos datos, escalas diferentes,
    impresiones \textbf{completamente distintas}!
    Siempre comprueba si el eje $y$ empieza en $0$.
  \end{center}
\end{frame}

% ---------------------------------------------------------------
\seccionframe{7. Ejercicios multinivel}
% ---------------------------------------------------------------

\begin{frame}{Ejercicios multinivel}
  \begin{columns}[T]
    \begin{column}{0.31\textwidth}
      \begin{analizarbox}
        \footnotesize\dificultad{2}\\[4pt]
        Tienes un gráfico dual: temperatura (°C)
        y ventas de helados en Almería por mes.\\[4pt]
        \textbf{Analiza:}\\
        (a) Tipo de correlación.\\
        (b) ¿Cuándo coinciden los picos?\\
        (c) ¿Hay un desfase (lag) entre ambas?\\
        Argumenta en 5 líneas.
      \end{analizarbox}
    \end{column}
    \begin{column}{0.31\textwidth}
      \begin{evaluarbox}
        \footnotesize\dificultad{3}\\[4pt]
        Tienes tres comparaciones:\\[3pt]
        A: temperatura vs.\ turistas\\
        B: vacaciones escolares vs.\ turistas\\
        C: precio vuelos vs.\ turistas\\[4pt]
        \textbf{Evalúa} cuál de los tres factores
        explica mejor el turismo en Almería.
        ¿Son independientes o están relacionados
        entre sí? Escribe un argumento de 8 líneas.
      \end{evaluarbox}
    \end{column}
    \begin{column}{0.31\textwidth}
      \begin{crearbox}
        \footnotesize\dificultad{4}\\[4pt]
        Inventa \textbf{dos variables urbanas}
        de Almería que tengan:
        \begin{itemize}
          \item Una correlación \textbf{positiva}
          \item Una correlación \textbf{negativa}
        \end{itemize}
        Dibuja a mano un gráfico esquemático
        de cada par. Explica la posible variable
        confusora en cada caso.
      \end{crearbox}
    \end{column}
  \end{columns}
\end{frame}

% ---------------------------------------------------------------
\begin{frame}{Evidencia del día}
  \begin{evidenciabox}
    Escribe un \textbf{análisis comparativo} de dos gráficas de
    datos de Almería. Debe incluir:
    \begin{enumerate}
      \item Identificación del tipo de correlación (positiva/negativa/nula)
            con justificación basada en los valores leídos
      \item Al menos un punto donde ambas variables cambian juntas
            y uno donde se separan — explica por qué
      \item Discusión de si la correlación implica causalidad:
            ¿existe una variable oculta?
      \item Conclusión de 4--5 líneas sobre la relación entre
            las dos variables en el contexto de Almería
    \end{enumerate}
  \end{evidenciabox}
\end{frame}

% ---------------------------------------------------------------
\begin{frame}{¡Cuidado con este error!}
  \begin{errorbox}
    \textbf{Correlación $\neq$ Causalidad}\\[6pt]
    \incorrecto\ «El turismo aumenta porque hace más calor.
    Por tanto, si calentamos la ciudad artificialmente,
    vendrán más turistas.»\\[6pt]
    \correcto\ La correlación entre temperatura y turismo se debe en
    parte a las vacaciones de verano, al sol en la playa, y a
    patrones culturales. El calor \textit{solo} no causaría más turismo.
    Identificar la variable confusora es fundamental para
    razonar correctamente.
  \end{errorbox}
\end{frame}

% ---------------------------------------------------------------
\begin{frame}{Resumen de la sesión 15}
  \begin{resumenbox}
    \begin{itemize}
      \item Comparar gráficas permite detectar
            \textbf{correlaciones} entre variables.
      \item \textbf{Correlación positiva}: ambas variables
            suben/bajan juntas.
      \item \textbf{Correlación negativa}: una sube cuando
            la otra baja.
      \item Se usan \textbf{dos ejes $y$} cuando las escalas
            son muy distintas.
      \item \textbf{Correlación $\neq$ causalidad}:
            siempre busca posibles variables confusoras.
      \item Una gráfica puede \textbf{engañar} si el eje $y$
            no empieza en $0$.
    \end{itemize}
  \end{resumenbox}
  \vspace{6pt}
  \begin{center}
    {\small\color{AlmAzulM}\faArrowCircleRight\enskip
    \textbf{Sesión 16:} Relaciones funcionales y proporcionalidad directa}
  \end{center}
\end{frame}

\end{document}

% % ================================================================
%  SESIÓN 16 — Relaciones funcionales y proporcionalidad
%  Almería en Cifras y Formas · Matemáticas 1.º ESO · IES Al-Ándalus
%  Compilar: pdflatex sesion16.tex
% ================================================================
\documentclass[aspectratio=169, 12pt, spanish]{beamer}
\usepackage{almeria-beamer}

\renewcommand{\SessionNumber}{16}
\renewcommand{\SessionTitle}{Relaciones funcionales\\y proporcionalidad}
\renewcommand{\SessionName}{Sesión 16}
\renewcommand{\SessionObjective}{Identificar y representar relaciones de proporcionalidad directa mediante tablas y gráficas, reconociendo la constante de proporcionalidad y aplicándola a situaciones reales de Almería}
\renewcommand{\BlockName}{Bloque 3: Datos y gráficas}

\begin{document}

% ── PORTADA ────────────────────────────────────────────────────────
\begin{frame}[plain]
  \titlepage
\end{frame}

% ── ÍNDICE ─────────────────────────────────────────────────────────
\begin{frame}{Contenidos de hoy}
  \begin{columns}[t]
    \begin{column}{0.48\textwidth}
      \begin{enumerate}
        \item Relación vs.\ función
        \item Proporcionalidad directa
        \item Cómo encontrar la constante $k$
        \item Gráfica de $y = kx$
        \item Ejemplos en Almería
      \end{enumerate}
    \end{column}
    \begin{column}{0.48\textwidth}
      \begin{enumerate}\setcounter{enumi}{5}
        \item Tabla de la proporcionalidad
        \item Lo que \textbf{no} es proporcional
        \item Propiedades clave
        \item Ejercicios por niveles Bloom
      \end{enumerate}
    \end{column}
  \end{columns}
  \vspace{0.4cm}
  \begin{center}
    \bloomtag{Comprender}\bloomtag{Aplicar}\bloomtag{Analizar}\bloomtag{Evaluar}
  \end{center}
\end{frame}

% ── SECCIÓN 1 ──────────────────────────────────────────────────────
\seccionframe[AlmAzulM]{Relación vs.\ Función}

% ── FRAME 3: Relación y función ────────────────────────────────────
\begin{frame}{Relación y función: ¿cuál es la diferencia?}
  \begin{columns}[t]
    \begin{column}{0.48\textwidth}
      \begin{definicionbox}{Relación}
        Cualquier \textbf{conexión} entre dos conjuntos de valores. A cada $x$ pueden corresponder \emph{varios} valores de $y$.
        \[
          x = 4 \;\longrightarrow\; y \in \{2,\,-2\}
        \]
      \end{definicionbox}
    \end{column}
    \begin{column}{0.48\textwidth}
      \begin{definicionbox}{Función}
        Una regla que asigna a cada valor de entrada \textbf{exactamente un} valor de salida.
        \[
          \text{Para cada } x,\; \exists!\; y = f(x)
        \]
        Se lee: «$y$ es función de $x$».
      \end{definicionbox}
    \end{column}
  \end{columns}
  \vspace{0.4cm}
  \begin{center}
  \begin{tikzpicture}[scale=0.85]
    % RELACIÓN (izq)
    \node[draw=AlmAzulM, ellipse, minimum width=1.2cm, minimum height=2.2cm,
          fill=AlmAzulC!15] at (0,0) {};
    \foreach \v/\y in {2/0.6, 4/0, 6/-0.6}
      \node[font=\small] at (0,\y) {\v};
    \node[draw=AlmAzulM, ellipse, minimum width=1.2cm, minimum height=2.8cm,
          fill=AlmAzulC!15] at (2.8,0) {};
    \foreach \v/\y in {-2/0.9, 1/0.3, 2/-0.3, 3/-0.9}
      \node[font=\small] at (2.8,\y) {\v};
    \draw[->, AlmAzulM, thick] (0.6, 0.6) -- (2.2, 0.9);
    \draw[->, AlmAzulM, thick] (0.6, 0.6) -- (2.2,-0.3);
    \draw[->, AlmAzulM, thick] (0.6, 0) -- (2.2, 0.3);
    \draw[->, AlmTerra, thick] (0.6,-0.6) -- (2.2, -0.9);
    \node[above, font=\tiny\bfseries, color=AlmGris] at (1.4,1.4) {RELACIÓN};
    % FUNCIÓN (der)
    \node[draw=AlmVerde, ellipse, minimum width=1.2cm, minimum height=2.2cm,
          fill=AlmVerde!10] at (5.4,0) {};
    \foreach \v/\y in {1/0.6, 2/0, 3/-0.6}
      \node[font=\small] at (5.4,\y) {\v};
    \node[draw=AlmVerde, ellipse, minimum width=1.2cm, minimum height=2.2cm,
          fill=AlmVerde!10] at (8.2,0) {};
    \foreach \v/\y in {3/0.6, 6/0, 9/-0.6}
      \node[font=\small] at (8.2,\y) {\v};
    \draw[->, AlmVerde, thick] (6.0, 0.6) -- (7.6, 0.6);
    \draw[->, AlmVerde, thick] (6.0, 0.0) -- (7.6, 0.0);
    \draw[->, AlmVerde, thick] (6.0,-0.6) -- (7.6,-0.6);
    \node[above, font=\tiny\bfseries, color=AlmGris] at (6.8,1.4) {FUNCIÓN};
  \end{tikzpicture}
  \end{center}
\end{frame}

% ── SECCIÓN 2 ──────────────────────────────────────────────────────
\seccionframe[BComprender]{Proporcionalidad directa}

% ── FRAME 4: Definición proporcionalidad ───────────────────────────
\begin{frame}{Proporcionalidad directa: definición}
  \begin{teoriabox}
    Dos magnitudes $x$ e $y$ son \textbf{directamente proporcionales} si existe una constante $k \neq 0$ tal que:
    \[
      \boxed{y = k \cdot x}
    \]
    $k$ se llama \textbf{constante de proporcionalidad}.
  \end{teoriabox}
  \vspace{0.3cm}
  \begin{columns}[t]
    \begin{column}{0.32\textwidth}
      \begin{notabox}[Duplicar]
        Si $x$ se duplica,\\$y$ también se duplica.
      \end{notabox}
    \end{column}
    \begin{column}{0.32\textwidth}
      \begin{notabox}[Triplicar]
        Si $x$ se triplica,\\$y$ también se triplica.
      \end{notabox}
    \end{column}
    \begin{column}{0.32\textwidth}
      \begin{notabox}[Cociente]
        El cociente\\$\dfrac{y}{x} = k$ es \textbf{siempre constante}.
      \end{notabox}
    \end{column}
  \end{columns}
  \vspace{0.35cm}
  \begin{center}
    \begin{tikzpicture}
      \node[draw=AlmAzulM, fill=AlmAzulC!12, rounded corners=5pt,
            font=\small\bfseries, inner sep=6pt] {
        La gráfica es una \textcolor{AlmTerra}{recta que pasa por el origen} $(0,0)$};
    \end{tikzpicture}
  \end{center}
\end{frame}

% ── FRAME 5: Encontrar k ───────────────────────────────────────────
\begin{frame}{¿Cómo encontrar la constante $k$?}
  \begin{columns}[c]
    \begin{column}{0.55\textwidth}
      \begin{pasobox}[1]
        Escribe la fórmula $y = k \cdot x$.
      \end{pasobox}
      \vspace{0.2cm}
      \begin{pasobox}[2]
        Despeja $k$: $\quad k = \dfrac{y}{x}$
      \end{pasobox}
      \vspace{0.2cm}
      \begin{pasobox}[3]
        Sustituye \textbf{cualquier par} $(x,\,y)$ conocido.
      \end{pasobox}
      \vspace{0.2cm}
      \begin{pasobox}[4]
        Verifica con \textbf{otro par}: si $k$ es siempre igual, hay proporcionalidad directa.
      \end{pasobox}
    \end{column}
    \begin{column}{0.42\textwidth}
      \begin{ejemplobox}
        \textbf{Taxi en Almería:}\\
        \SI{2}{\km} cuestan \SI{5.00}{\euro}\\[4pt]
        $k = \dfrac{5{,}00}{2} = \SI{2{,}50}{\euro/\km}$\\[4pt]
        Verificamos con $x=4$:\\
        $y = 2{,}50 \times 4 = \SI{10{,}00}{\euro}$ \correcto
      \end{ejemplobox}
    \end{column}
  \end{columns}
\end{frame}

% ── FRAME 6: Gráfica de y = 2.5x ──────────────────────────────────
\begin{frame}{Gráfica de proporcionalidad directa: tarifa de taxi en Almería}
  \begin{columns}[c]
    \begin{column}{0.60\textwidth}
      \begin{tikzpicture}
        \begin{axis}[
          almeria line,
          width=\linewidth,
          height=0.75\linewidth,
          xmin=0, xmax=5.5,
          ymin=0, ymax=15,
          xlabel={Distancia (km)},
          ylabel={Coste (\euro)},
          title={Tarifa taxi Almería: $y = 2{,}5x$},
          xtick={0,1,2,3,4,5},
          ytick={0,2.5,5,7.5,10,12.5},
          yticklabels={0, 2{,}5, 5, 7{,}5, 10, 12{,}5},
        ]
          % línea de proporcionalidad
          \addplot[AlmTerra, line width=2pt, domain=0:5.5] {2.5*x};
          % puntos marcados
          \addplot[only marks, mark=*, mark size=3pt, AlmAzulM]
            coordinates {(0,0)(1,2.5)(2,5)(3,7.5)(4,10)(5,12.5)};
          % etiqueta de la recta
          \node[font=\tiny\bfseries, color=AlmTerra, anchor=west]
            at (axis cs: 3.5, 10.5) {$y = 2{,}5\,x$};
          % pendiente k
          \draw[<->, AlmVerde, thick]
            (axis cs: 2, 5) -- (axis cs: 3, 5)
            node[midway, below, font=\tiny\bfseries, color=AlmVerde] {$\Delta x = 1$};
          \draw[<->, AlmVerde, thick]
            (axis cs: 3, 5) -- (axis cs: 3, 7.5)
            node[midway, right, font=\tiny\bfseries, color=AlmVerde] {$k = 2{,}5$};
        \end{axis}
      \end{tikzpicture}
    \end{column}
    \begin{column}{0.38\textwidth}
      \small
      \begin{tabular}{@{}cc@{}}
        \toprule
        \textbf{km} & \textbf{\euro} \\
        \midrule
        0 & 0{,}00 \\
        1 & 2{,}50 \\
        2 & 5{,}00 \\
        3 & 7{,}50 \\
        4 & 10{,}00 \\
        5 & 12{,}50 \\
        \bottomrule
      \end{tabular}
      \vspace{0.4cm}

      \begin{formulabox}
        $k = \dfrac{y}{x} = 2{,}5$
      \end{formulabox}
      \vspace{0.2cm}
      {\tiny La recta \textbf{pasa por el origen} y tiene pendiente $k = 2{,}5$.}
    \end{column}
  \end{columns}
\end{frame}

% ── FRAME 7: Ejemplos almería ──────────────────────────────────────
\begin{frame}{Proporcionalidad directa en Almería}
  \begin{columns}[t]
    \begin{column}{0.32\textwidth}
      \begin{ejemplobox}
        \textbf{\faMapMarkerAlt\ Bus urbano}\\[4pt]
        3 viajes $= 3 \times 1{,}40 = \SI{4{,}20}{\euro}$\\[4pt]
        \[y = 1{,}40\,x\]
        $k = \SI{1{,}40}{\euro/\text{viaje}}$
      \end{ejemplobox}
    \end{column}
    \begin{column}{0.32\textwidth}
      \begin{ejemplobox}
        \textbf{\faWalking\ Caminando}\\[4pt]
        Velocidad $= 5$ km/h\\[4pt]
        \[d = 5\,t\]
        $k = 5$ km/h
      \end{ejemplobox}
    \end{column}
    \begin{column}{0.32\textwidth}
      \begin{ejemplobox}
        \textbf{\faPaintRoller\ Baldosas plaza}\\[4pt]
        1 m² $\to$ 9 baldosas\\[4pt]
        \[y = 9\,x\]
        20 m² $\to$ 180 baldosas
      \end{ejemplobox}
    \end{column}
  \end{columns}
  \vspace{0.4cm}
  \begin{center}
    \begin{tikzpicture}
      \node[draw=AlmAmbar, fill=AlmAmbar!10, rounded corners=4pt,
            font=\small, inner sep=8pt, text width=10cm, align=center] {
        En todos los casos: si $x$ se multiplica por $n$, $y$ también se multiplica por $n$,\\
        y el cociente $y/x$ es siempre la misma constante $k$.
      };
    \end{tikzpicture}
  \end{center}
\end{frame}

% ── FRAME 8: Tabla y = 1.40x ──────────────────────────────────────
\begin{frame}{Tabla de proporcionalidad: viajes en bus}
  \begin{columns}[c]
    \begin{column}{0.55\textwidth}
      \renewcommand{\arraystretch}{1.3}
      \begin{tabular}{|>{\bfseries\color{AlmAzulM}}c|c|c|}
        \hline
        \rowcolor{AlmAzul!15}
        \textbf{Viajes ($x$)} & \textbf{Coste ($y$, \euro)} & \textbf{$y/x$} \\
        \hline
        1  & 1{,}40  & 1{,}40 \\
        2  & 2{,}80  & 1{,}40 \\
        3  & 4{,}20  & 1{,}40 \\
        4  & 5{,}60  & 1{,}40 \\
        5  & 7{,}00  & 1{,}40 \\
        10 & 14{,}00 & 1{,}40 \\
        \hline
      \end{tabular}
    \end{column}
    \begin{column}{0.42\textwidth}
      \begin{teoriabox}
        \textbf{Propiedades que verificar:}
        \begin{itemize}
          \item[\correcto] $y/x = k$ siempre igual
          \item[\correcto] Gráfica: recta por $(0,0)$
          \item[\correcto] Fórmula: $y = 1{,}40\,x$
        \end{itemize}
      \end{teoriabox}
      \vspace{0.3cm}
      \begin{formulabox}[AlmAmbar]
        $y = 1{,}40\,x$
      \end{formulabox}
    \end{column}
  \end{columns}
\end{frame}

% ── FRAME 9: No proporcional vs proporcional ───────────────────────
\begin{frame}{Proporcional vs.\ lineal: ¡no confundir!}
  \begin{columns}[c]
    \begin{column}{0.48\textwidth}
      \begin{block}{\correcto\ Proporcional: $y = 2{,}5\,x$}
        Taxi \textbf{sin tarifa de bajada:}\\
        $y = 2{,}5\,x$
        \begin{itemize}
          \item Pasa por el origen $(0,\,0)$
          \item $y/x = 2{,}5$ (constante)
        \end{itemize}
      \end{block}
    \end{column}
    \begin{column}{0.48\textwidth}
      \begin{alertblock}{\incorrecto\ Lineal (no proporcional): $y = 2{,}40 + 1{,}20\,x$}
        Taxi \textbf{con tarifa de bajada} €2{,}40:
        \begin{itemize}
          \item \textbf{No} pasa por el origen
          \item $y/x \neq$ constante
          \item Sí es lineal (recta)
        \end{itemize}
      \end{alertblock}
    \end{column}
  \end{columns}
  \vspace{0.3cm}
  \begin{center}
  \begin{tikzpicture}[scale=0.85]
    \begin{axis}[
      almeria,
      width=0.65\linewidth,
      height=0.38\linewidth,
      xmin=0, xmax=5.5,
      ymin=0, ymax=12,
      xlabel={km}, ylabel={\euro},
      title={Comparación de tarifas},
      legend pos=north west,
      xtick={0,1,2,3,4,5},
    ]
      \addplot[AlmVerde, line width=2pt, domain=0:5.5] {2.5*x};
      \addlegendentry{$y=2{,}5x$ (proporcional)}
      \addplot[AlmTerra, line width=2pt, dashed, domain=0:5.5] {2.4 + 1.2*x};
      \addlegendentry{$y=2{,}4+1{,}2x$ (lineal)}
    \end{axis}
  \end{tikzpicture}
  \end{center}
\end{frame}

% ── FRAME 10: Propiedades clave ────────────────────────────────────
\begin{frame}{Tres formas de reconocer la proporcionalidad directa}
  \begin{columns}[t]
    \begin{column}{0.32\textwidth}
      \begin{block}{\faTable\ Por la tabla}
        \[
          \frac{y}{x} = k \;\text{(siempre igual)}
        \]
      \end{block}
    \end{column}
    \begin{column}{0.32\textwidth}
      \begin{block}{\faChartLine\ Por la gráfica}
        Recta que pasa por el \textbf{origen} $(0,0)$
      \end{block}
    \end{column}
    \begin{column}{0.32\textwidth}
      \begin{block}{\faCalculator\ Por la fórmula}
        \[y = k \cdot x\]
        sin término independiente
      \end{block}
    \end{column}
  \end{columns}
  \vspace{0.5cm}
  \begin{errorbox}
    $y = k\,x$ \textbf{sí} es proporcionalidad directa.\\
    $y = mx + b$ con $b \neq 0$ es \textbf{lineal pero NO proporcional}.\\
    El término $b$ «levanta» la recta y la separa del origen.
  \end{errorbox}
\end{frame}

% ── FRAME 11: Ejercicio COMPRENDER ────────────────────────────────
\begin{frame}{Ejercicio — \bloomtag{Comprender}}
  \begin{comprenderbox}
    \textbf{Fuente de la Alameda (Almería):} llenado a \textbf{40 litros/minuto}.

    \vspace{0.3cm}
    \begin{columns}[t]
      \begin{column}{0.55\textwidth}
        \textbf{a)} Completa la tabla:
        \vspace{0.2cm}

        \renewcommand{\arraystretch}{1.2}
        \begin{tabular}{|c|c|c|}
          \hline
          \rowcolor{AlmAzulC!20}
          $t$ (min) & $V$ (litros) & $V/t$ \\
          \hline
          1 & \phantom{40}  & \phantom{40} \\
          2 & \phantom{80}  & \phantom{40} \\
          3 & \phantom{120} & \phantom{40} \\
          4 & \phantom{160} & \phantom{40} \\
          5 & \phantom{200} & \phantom{40} \\
          6 & \phantom{240} & \phantom{40} \\
          \hline
        \end{tabular}
      \end{column}
      \begin{column}{0.42\textwidth}
        \textbf{b)} ¿Es proporcionalidad directa?

        \vspace{0.3cm}
        Justifica usando el cociente $V/t$.

        \vspace{0.3cm}
        \textbf{c)} ¿Cuántos litros habrá tras 15 minutos?
      \end{column}
    \end{columns}
  \end{comprenderbox}
\end{frame}

% ── FRAME 12: Ejercicio APLICAR ───────────────────────────────────
\begin{frame}{Ejercicio — \bloomtag{Aplicar}}
  \begin{aplicarbox}
    \textbf{Entrada a la Alcazaba de Almería:} €3 por adulto.

    \vspace{0.25cm}
    \begin{columns}[t]
      \begin{column}{0.60\textwidth}
        \begin{enumerate}
          \item[\textbf{a)}] Escribe la fórmula de proporcionalidad.
          \item[\textbf{b)}] Construye una tabla para 1 a 10 personas.
          \item[\textbf{c)}] Dibuja la gráfica (ejes correctamente etiquetados).
          \item[\textbf{d)}] ¿Cuántas personas necesitan €21?
        \end{enumerate}
      \end{column}
      \begin{column}{0.37\textwidth}
        \begin{notabox}[Pista]
          $\text{Coste} = k \times \text{personas}$\\[4pt]
          Despeja «personas» para el apartado d).
        \end{notabox}
        \vspace{0.2cm}
        \dificultad{2}
      \end{column}
    \end{columns}
  \end{aplicarbox}
\end{frame}

% ── FRAME 13: Ejercicio ANALIZAR ──────────────────────────────────
\begin{frame}{Ejercicio — \bloomtag{Analizar}}
  \begin{analizarbox}
    \textbf{Dos planes de dispensador de agua:}

    \vspace{0.2cm}
    \begin{columns}[t]
      \begin{column}{0.47\textwidth}
        \begin{block}{Plan A}
          €0{,}80 por litro\\
          \[y_A = 0{,}80\,x\]
        \end{block}
      \end{column}
      \begin{column}{0.47\textwidth}
        \begin{alertblock}{Plan B}
          €5 fijos + €0{,}40 por litro\\
          \[y_B = 5 + 0{,}40\,x\]
        \end{alertblock}
      \end{column}
    \end{columns}
    \vspace{0.3cm}
    \begin{enumerate}
      \item ¿Cuál de los dos es proporcionalidad directa? ¿Por qué?
      \item ¿Para qué cantidad son igual de caros?\\
            Resuelve: $0{,}80\,x = 5 + 0{,}40\,x$
      \item ¿Cuándo conviene el Plan A? ¿Cuándo el Plan B?
    \end{enumerate}
    \dificultad{3}
  \end{analizarbox}
\end{frame}

% ── FRAME 14: Ejercicio EVALUAR ──────────────────────────────────
\begin{frame}{Ejercicio — \bloomtag{Evaluar}}
  \begin{evaluarbox}
    \textbf{Dos ciclistas salen de Almería al mismo tiempo:}
    \begin{itemize}
      \item Ciclista A: 15 km/h \quad $\Rightarrow$ \quad $d_A = 15\,t$
      \item Ciclista B: 20 km/h \quad $\Rightarrow$ \quad $d_B = 20\,t$
    \end{itemize}
    \vspace{0.3cm}
    \begin{columns}[t]
      \begin{column}{0.55\textwidth}
        \begin{enumerate}
          \item ¿Son ambas distancias proporcionales al tiempo? Razona.
          \item Tras 2 horas, ¿cuántos km les separan?
          \item La diferencia $d_B - d_A = 5t$, ¿es también proporcional al tiempo?
          \item ¿Cuál sería la constante de esa diferencia?
        \end{enumerate}
      \end{column}
      \begin{column}{0.42\textwidth}
        \begin{notabox}[Reflexión]
          «La suma o diferencia de dos funciones proporcionales, ¿sigue siendo proporcional?»
        \end{notabox}
        \dificultad{4}
      \end{column}
    \end{columns}
  \end{evaluarbox}
\end{frame}

% ── FRAME 15: Error frecuente ────────────────────────────────────
\begin{frame}{¡Cuidado! El error más habitual}
  \begin{errorbox}
    \textbf{Confundir proporcionalidad directa con función lineal:}
    \vspace{0.2cm}

    \begin{columns}[t]
      \begin{column}{0.47\textwidth}
        \begin{center}
          \correcto\; \textbf{Proporcionalidad directa}\\[4pt]
          $y = 3x$\\[4pt]
          {\small Pasa por el origen $(0,0)$}\\
          {\small $y/x = 3$ siempre}
        \end{center}
      \end{column}
      \begin{column}{0.47\textwidth}
        \begin{center}
          \incorrecto\; \textbf{Solo lineal (no proporcional)}\\[4pt]
          $y = 3x + 5$\\[4pt]
          {\small No pasa por el origen}\\
          {\small $y/x \neq$ constante}
        \end{center}
      \end{column}
    \end{columns}
    \vspace{0.3cm}
    El término independiente $b$ en $y = mx + b$ «rompe» la proporcionalidad. Si $b = 0$, entonces sí es proporcional.
  \end{errorbox}
\end{frame}

% ── FRAME 16: Evidencia del día ──────────────────────────────────
\begin{frame}{Evidencia del día}
  \begin{evidenciabox}
    \textbf{Ficha de proporcionalidad} (entregar al final de la sesión):

    \vspace{0.3cm}
    \begin{columns}[t]
      \begin{column}{0.55\textwidth}
        La ficha debe incluir \textbf{los cuatro elementos}:
        \begin{enumerate}
          \item \textbf{Tabla} de valores con columna $y/x$
          \item \textbf{Fórmula} $y = k \cdot x$ con valor de $k$
          \item \textbf{Gráfica} con ejes etiquetados y línea que pase por el origen
          \item \textbf{Verificación} del cociente constante $y/x = k$
        \end{enumerate}
      \end{column}
      \begin{column}{0.42\textwidth}
        \begin{notabox}[Situación a elegir]
          \begin{itemize}
            \item Bus urbano ($k = 1{,}40$)
            \item Alcazaba ($k = 3{,}00$)
            \item Baldosas ($k = 9$)
            \item Tu propio ejemplo
          \end{itemize}
        \end{notabox}
      \end{column}
    \end{columns}
  \end{evidenciabox}
\end{frame}

% ── FRAME 17: Resumen ─────────────────────────────────────────────
\begin{frame}{Resumen de la sesión 16}
  \begin{resumenbox}
    \begin{columns}[t]
      \begin{column}{0.48\textwidth}
        \begin{itemize}
          \item Una \textbf{función} asigna exactamente un $y$ a cada $x$.
          \item La \textbf{proporcionalidad directa}: $y = k\,x$.
          \item $k$ es la \textbf{constante de proporcionalidad}: $k = y/x$.
          \item La gráfica es una \textbf{recta que pasa por el origen}.
        \end{itemize}
      \end{column}
      \begin{column}{0.48\textwidth}
        \begin{itemize}
          \item Si $b \neq 0$, $y = kx + b$ es \textbf{lineal} pero \textbf{no proporcional}.
          \item En Almería: taxi ($k=2{,}5$), bus ($k=1{,}40$), baldosas ($k=9$).
          \item \textbf{Verificar siempre}: tabla, gráfica y fórmula.
        \end{itemize}
      \end{column}
    \end{columns}
  \end{resumenbox}

  \vspace{0.3cm}
  \begin{center}
    \bloomtag{Comprender}\ identificar relaciones de proporcionalidad\\[3pt]
    \bloomtag{Aplicar}\ construir tablas y gráficas\\[3pt]
    \bloomtag{Analizar}\ comparar planes y situaciones\\[3pt]
    \bloomtag{Evaluar}\ razonar sobre propiedades de funciones\\
  \end{center}
\end{frame}

\end{document}

% % ================================================================
%  SESIÓN 17 — Seleccionar la información de la guía
%  Almería en Cifras y Formas · Matemáticas 1.º ESO · IES Al-Ándalus
%  Compilar: pdflatex sesion17.tex
% ================================================================
\documentclass[aspectratio=169, 12pt, spanish]{beamer}
\usepackage{almeria-beamer}

\renewcommand{\SessionNumber}{17}
\renewcommand{\SessionTitle}{Seleccionar la información\\de la guía}
\renewcommand{\SessionName}{Sesión 17}
\renewcommand{\SessionObjective}{Aplicar criterios matemáticos y comunicativos para seleccionar, verificar y organizar la información que formará parte de la guía turística matemática de Almería}
\renewcommand{\BlockName}{Bloque 4: Producto final}

\begin{document}

% ── PORTADA ────────────────────────────────────────────────────────
\begin{frame}[plain]
  \titlepage
\end{frame}

% ── ÍNDICE ─────────────────────────────────────────────────────────
\begin{frame}{Contenidos de hoy}
  \begin{columns}[t]
    \begin{column}{0.48\textwidth}
      \begin{enumerate}
        \item Contexto: ¿qué hemos recogido?
        \item Criterios de selección
        \item Lista de verificación de datos
        \item Problemas de calidad habituales
        \item Árbol de decisión
      \end{enumerate}
    \end{column}
    \begin{column}{0.48\textwidth}
      \begin{enumerate}\setcounter{enumi}{5}
        \item Categorías de la guía
        \item Evaluar datos concretos
        \item Ejercicios ANALIZAR / EVALUAR / CREAR
        \item Evidencia del día
      \end{enumerate}
    \end{column}
  \end{columns}
  \vspace{0.4cm}
  \begin{center}
    \bloomtag{Analizar}\bloomtag{Evaluar}\bloomtag{Crear}
  \end{center}
\end{frame}

% ── SECCIÓN 1 ──────────────────────────────────────────────────────
\seccionframe[BAnalizar]{De la exploración a la selección}

% ── FRAME 3: Contexto ──────────────────────────────────────────────
\begin{frame}{Ya tenemos datos: ¿cuáles elegimos?}
  \begin{columns}[c]
    \begin{column}{0.52\textwidth}
      \begin{teoriabox}
        En las \textbf{16 sesiones anteriores} hemos recogido datos matemáticos sobre Almería:
        \begin{itemize}
          \item Coordenadas de monumentos
          \item Áreas, perímetros y alturas
          \item Tablas estadísticas y gráficas
          \item Rutas y distancias
          \item Relaciones de proporcionalidad
        \end{itemize}
        \vspace{0.2cm}
        Ahora toca \textbf{seleccionar lo mejor} para la guía.
      \end{teoriabox}
    \end{column}
    \begin{column}{0.45\textwidth}
      \begin{center}
      \begin{tikzpicture}[scale=0.9]
        % Triángulo de selección
        \node[draw=AlmAzul, fill=AlmAzulC!15, rounded corners=6pt,
              minimum width=3.5cm, minimum height=0.65cm,
              font=\scriptsize\bfseries, color=AlmAzul] (a) at (0, 2.2)
          {Muchos datos (sesiones 1-16)};
        \node[draw=AlmAmbar!80!black, fill=AlmAmbar!15, rounded corners=6pt,
              minimum width=2.6cm, minimum height=0.65cm,
              font=\scriptsize\bfseries, color=AlmAzul] (b) at (0, 1.1)
          {Datos candidatos};
        \node[draw=AlmVerde, fill=AlmVerde!15, rounded corners=6pt,
              minimum width=2.0cm, minimum height=0.65cm,
              font=\scriptsize\bfseries, color=AlmVerde] (c) at (0, 0)
          {Datos seleccionados};
        \draw[->, thick, AlmAzulM] (a) -- (b);
        \draw[->, thick, AlmAmbar!70!black] (b) -- (c);
        \node[right=0.1cm of a, font=\tiny, color=AlmGris] {100\%};
        \node[right=0.1cm of b, font=\tiny, color=AlmGris] {50\%};
        \node[right=0.1cm of c, font=\tiny, color=AlmGris] {25\%};
      \end{tikzpicture}
      \end{center}
      \begin{notabox}[Principio]
        \textbf{Calidad} sobre cantidad.
      \end{notabox}
    \end{column}
  \end{columns}
\end{frame}

% ── FRAME 4: Criterios de selección ───────────────────────────────
\begin{frame}{Los cinco criterios de selección}
  \begin{columns}[t]
    \begin{column}{0.48\textwidth}
      \begin{block}{\faCheckDouble\ R — Relevancia}
        ¿Es este dato \textbf{matemáticamente interesante}? ¿Aporta algo al turista?
      \end{block}
      \vspace{0.2cm}
      \begin{block}{\faSearch\ E — Exactitud}
        ¿Es el dato \textbf{preciso y verificable}? ¿Tiene fuente?
      \end{block>}
      \vspace{0.2cm}
      \begin{block}{\faComments\ C — Claridad}
        ¿Se puede explicar a un turista \textbf{sin formación matemática}?
      \end{block}
    \end{column}
    \begin{column}{0.48\textwidth}
      \begin{block}{\faMapMarkerAlt\ R — Representatividad}
        ¿Representa bien a Almería? ¿Es \textbf{característico} de la ciudad?
      \end{block}
      \vspace{0.2cm}
      \begin{block}{\faStar\ O — Originalidad}
        ¿Es \textbf{sorprendente o único}? ¿Algo que el turista no esperaría?
      \end{block}
      \vspace{0.3cm}
      \begin{formulabox}[BAnalizar]
        \textbf{R\,·\,E\,·\,C\,·\,R\,·\,O}
      \end{formulabox}
    \end{column}
  \end{columns}
\end{frame}

% ── FRAME 5: Lista de verificación ────────────────────────────────
\begin{frame}{Lista de verificación para cada dato}
  \begin{columns}[c]
    \begin{column}{0.55\textwidth}
      \renewcommand{\arraystretch}{1.4}
      \begin{tabular}{|c|p{6.5cm}|}
        \hline
        \rowcolor{AlmAzul!15}
        \textbf{Check} & \textbf{Pregunta} \\
        \hline
        $\square$ & Fuente verificada (no inventada) \\
        $\square$ & Unidades correctas (m, no m² para distancias) \\
        $\square$ & Cálculos comprobados dos veces \\
        $\square$ & Visualizable (mapa, tabla o gráfica) \\
        $\square$ & Interesante para un turista \\
        \hline
      \end{tabular}
    \end{column}
    \begin{column}{0.42\textwidth}
      \begin{notabox}[Regla de los 5/5]
        Un dato de \textbf{calidad} supera \textbf{todos} los checks.\\[6pt]
        Un dato con 3 o menos checks debe \textbf{revisarse o descartarse}.
      \end{notabox}
      \vspace{0.2cm}
      \begin{notabox}[Unidades]
        Distancia $\to$ m, km\\
        Área $\to$ m², ha, km²\\
        Tiempo $\to$ min, h\\
        Temperatura $\to$ °C
      \end{notabox}
    \end{column}
  \end{columns}
\end{frame}

% ── FRAME 6: Problemas de calidad ─────────────────────────────────
\begin{frame}{Errores habituales en la calidad de los datos}
  \begin{columns}[t]
    \begin{column}{0.48\textwidth}
      \begin{errorbox}
        \textbf{Tipos de datos problemáticos:}
        \begin{itemize}
          \item[\incorrecto] \textbf{Obsoletos:} «La población de Almería en 2000 era...»
          \item[\incorrecto] \textbf{Sin fuente:} «Dicen que la Alcazaba tiene 43.000 m²»
          \item[\incorrecto] \textbf{Sin unidades:} «La torre mide 55»
          \item[\incorrecto] \textbf{Cálculo incorrecto:} perímetro calculado como área
        \end{itemize}
      \end{errorbox}
    \end{column}
    \begin{column}{0.48\textwidth}
      \begin{exampleblock}{\correcto\ Datos de calidad:}
        \begin{itemize}
          \item «La Alcazaba ocupa \textbf{4,3 ha} = 43.000 m²\\ (Fuente: Junta de Andalucía, 2023)»
          \item «El Cabo de Gata tiene \textbf{32 m} de altura (IGN, 2022)»
          \item «Almería recibió \textbf{1,2 millones} de turistas en 2023 (INE)»
        \end{itemize}
      \end{exampleblock}
    \end{column}
  \end{columns}
\end{frame}

% ── FRAME 7: Árbol de decisión TikZ ───────────────────────────────
\begin{frame}{Árbol de decisión: ¿incluyo este dato?}
  \begin{center}
  \begin{tikzpicture}[
    decision/.style={diamond, draw=AlmAzulM, fill=AlmAzulC!20,
                     text width=2.6cm, align=center, font=\tiny\bfseries,
                     inner sep=1pt, aspect=1.8},
    result/.style={rounded corners=4pt, draw=AlmVerde, fill=AlmVerde!15,
                   font=\tiny\bfseries, text width=2cm, align=center, inner sep=3pt},
    reject/.style={rounded corners=4pt, draw=AlmTerra, fill=AlmTerra!10,
                   font=\tiny\bfseries, text width=1.8cm, align=center, inner sep=3pt},
    arr/.style={->, >=Stealth, thick, AlmAzulM},
    scale=0.9, every node/.style={font=\tiny}
  ]
    \node[decision] (Q1) at (0,0) {¿Tiene fuente verificable?};
    \node[reject] (R1) at (-3.2,-1.5) {Descartar o buscar fuente};
    \node[decision] (Q2) at (2.0,-1.5) {¿Tiene unidades correctas?};
    \node[reject] (R2) at (4.8,-1.5) {Corregir unidades};
    \node[decision] (Q3) at (2.0,-3.0) {¿Es relevante para el turista?};
    \node[reject] (R3) at (5.0,-3.0) {Descartar};
    \node[result] (OK) at (2.0,-4.5) {\correcto\ Incluir en la guía};

    \draw[arr] (Q1) -- node[above,sloped]{\tiny No} (R1);
    \draw[arr] (Q1) -- node[above]{\tiny Sí} (Q2);
    \draw[arr] (Q2) -- node[above]{\tiny No} (R2);
    \draw[arr] (Q2) -- node[right]{\tiny Sí} (Q3);
    \draw[arr] (Q3) -- node[above]{\tiny No} (R3);
    \draw[arr] (Q3) -- node[right]{\tiny Sí} (OK);
  \end{tikzpicture}
  \end{center}
\end{frame}

% ── FRAME 8: Categorías de la guía ────────────────────────────────
\begin{frame}{Categorías de la guía turística matemática}
  \begin{columns}[t]
    \begin{column}{0.48\textwidth}
      \begin{definicionbox}{Estructura de la guía}
        \begin{enumerate}
          \item \textbf{Mapa de la ciudad} — coordenadas, escala, puntos clave
          \item \textbf{Monumentos principales} — áreas, alturas, fechas, visitantes
          \item \textbf{Rutas turísticas} — distancias, tiempos, dificultad
        \end{enumerate}
      \end{definicionbox}
    \end{column}
    \begin{column}{0.48\textwidth}
      \begin{definicionbox}{Estructura (cont.)}
        \begin{enumerate}\setcounter{enumi}{3}
          \item \textbf{Datos curiosos} — estadísticas sorprendentes de Almería
          \item \textbf{Gráficas estadísticas} — turismo, temperatura, población
          \item \textbf{Glosario matemático} — términos usados en la guía
        \end{enumerate}
      \end{definicionbox}
    \end{column}
  \end{columns}
  \vspace{0.35cm}
  \begin{center}
  \begin{tikzpicture}
    \foreach \i/\cat/\col in {
      1/Mapa/AlmAzulM,
      2/Monumentos/AlmTerra,
      3/Rutas/AlmVerde,
      4/Curiosidades/AlmAmbar!80!black,
      5/Gráficas/BAnalizar,
      6/Glosario/BCrear}{
      \node[draw=\col, fill=\col!12, rounded corners=3pt,
            font=\tiny\bfseries\color{\col}, inner sep=4pt,
            minimum width=1.8cm, minimum height=0.55cm]
        at ({(\i-1)*2.2},0) {\cat};
    }
  \end{tikzpicture}
  \end{center}
\end{frame>}

% ── FRAME 9: Evaluar datos concretos ──────────────────────────────
\begin{frame}{Evaluando datos concretos}
  \begin{columns}[t]
    \begin{column}{0.48\textwidth}
      \begin{block}{Dato A}
        «La Alcazaba ocupa aproximadamente 43.000 m²»
        \vspace{0.2cm}

        \begin{tabular}{lc}
          Fuente verificable & $\square$ \\
          Unidades correctas & \correcto \\
          Relevante & \correcto \\
          Verificable & \correcto \\
        \end{tabular}
        \vspace{0.2cm}
        \textcolor{AlmAmbar!80!black}{\textbf{Candidato: buscar fuente}}
      \end{block}
    \end{column}
    \begin{column}{0.48\textwidth}
      \begin{alertblock}{Dato B}
        «La Alcazaba tiene 100 torres»
        \vspace{0.2cm}

        \begin{tabular}{lc}
          Fuente verificable & \incorrecto \\
          Unidades correctas & \correcto \\
          Relevante & \correcto \\
          Plausible & \incorrecto \\
        \end{tabular}
        \vspace{0.2cm}
        \textcolor{AlmTerra}{\textbf{Rechazado: valor sospechoso}}
      \end{alertblock}
    \end{column}
  \end{columns}
\end{frame}

% ── FRAME 10: Ejercicio ANALIZAR ──────────────────────────────────
\begin{frame}{Ejercicio — \bloomtag{Analizar}}
  \begin{analizarbox}
    \textbf{8 candidatos para la guía.} Aplica la lista de verificación a cada uno y decide cuáles de los 5 incluirías:

    \vspace{0.2cm}
    \begin{columns}[t]
      \begin{column}{0.48\textwidth}
        \begin{enumerate}
          \item «El Paseo de Almería mide 1,2 km de largo» (sin fuente)
          \item «El faro del Cabo de Gata: 32 m» (IGN 2022)
          \item «Almería tiene muchos días de sol» (sin número)
          \item «La temperatura media en julio es 27°C» (AEMET 2023)
        \end{enumerate}
      \end{column>}
      \begin{column}{0.48\textwidth}
        \begin{enumerate}\setcounter{enumi}{4}
          \item «El Cable Inglés mide 1.100 m» (Ayto. 2021)
          \item «La Alcazaba es muy grande» (sin dato)
          \item «En 2023 llegaron 1.200.000 turistas» (INE)
          \item «El estadio Juegos Mediterráneos: 15.000 localidades» (sin año)
        \end{enumerate}
      \end{column}
    \end{columns}

    \vspace{0.2cm}
    \textbf{¿Cuáles 5 incluyes? ¿Por qué?} Usa la rúbrica R-E-C-R-O.
  \end{analizarbox}
\end{frame>}

% ── FRAME 11: Ejercicio EVALUAR ───────────────────────────────────
\begin{frame}{Ejercicio — \bloomtag{Evaluar}}
  \begin{evaluarbox}
    \textbf{Tu equipo tiene los siguientes materiales:}
    \begin{columns}[t]
      \begin{column}{0.48\textwidth}
        \begin{itemize}
          \item 3 mapas con coordenadas
          \item 2 gráficas estadísticas
          \item 5 listas de coordenadas
          \item 1 tabla de cálculos verificados
        \end{itemize}
      \end{column}
      \begin{column}{0.48\textwidth}
        Decide la \textbf{estructura óptima} de tu capítulo de guía:
        \begin{enumerate}
          \item ¿Qué materiales incluyes?
          \item ¿Cuáles descartas? ¿Por qué?
          \item ¿Qué sección del capítulo encabeza cada material?
          \item ¿Falta algo? ¿Qué añadirías?
        \end{enumerate}
      \end{column}
    \end{columns}
    \vspace{0.2cm}
    Justifica \textbf{cada} decisión con criterios matemáticos.
    \dificultad{3}
  \end{evaluarbox}
\end{frame>}

% ── FRAME 12: Ejercicio CREAR ─────────────────────────────────────
\begin{frame}{Ejercicio — \bloomtag{Crear}}
  \begin{crearbox}
    \textbf{Diseña la portada de tu guía turística matemática:}

    \vspace{0.3cm}
    \begin{columns}[t]
      \begin{column}{0.55\textwidth}
        La portada debe incluir:
        \begin{enumerate}
          \item \textbf{Título} de la guía
          \item \textbf{Subtítulo} que describa el contenido matemático
          \item \textbf{Una estadística clave} sobre Almería (bien elegida y verificada)
          \item \textbf{Miniatura de mapa} con coordenadas de 3 puntos turísticos
          \item Nombre del equipo y fecha
        \end{enumerate}
      \end{column}
      \begin{column}{0.42\textwidth}
        \begin{notabox}[Diseño]
          Recuerda:
          \begin{itemize}
            \item El número debe ser \textbf{preciso} y tener \textbf{unidades}
            \item El mapa debe tener \textbf{ejes} y \textbf{escala}
            \item El título debe mencionar \textbf{matemáticas}
          \end{itemize}
        \end{notabox}
        \dificultad{4}
      \end{column}
    \end{columns}
  \end{crearbox}
\end{frame}

% ── FRAME 13: Error frecuente ─────────────────────────────────────
\begin{frame}{¡Cuidado! «Más datos = mejor guía»}
  \begin{errorbox}
    \textbf{Error habitual:} incluir \textbf{todo} lo recopilado porque «más información es mejor».

    \vspace{0.3cm}
    \begin{columns}[t]
      \begin{column}{0.47\textwidth}
        \incorrecto\ \textbf{Guía con exceso:}
        \begin{itemize}
          \item 47 coordenadas sin clasificar
          \item Datos sin fuente mezclados con datos verificados
          \item Gráficas redundantes
          \item Números sin contexto
        \end{itemize}
      \end{column}
      \begin{column}{0.47\textwidth}
        \correcto\ \textbf{Guía bien seleccionada:}
        \begin{itemize}
          \item 10--15 datos \textbf{seleccionados y verificados}
          \item Cada dato tiene fuente y unidades
          \item Gráficas complementarias, no repetidas
          \item Números con contexto y comparación
        \end{itemize}
      \end{column}
    \end{columns}
    \vspace{0.2cm}
    \begin{center}
      \textit{Una buena guía tiene datos \textbf{seleccionados}, \textbf{verificados} y \textbf{relevantes}.}
    \end{center}
  \end{errorbox}
\end{frame}

% ── FRAME 14: Evidencia del día ───────────────────────────────────
\begin{frame}{Evidencia del día}
  \begin{evidenciabox}
    \textbf{Lista de datos seleccionados} (entregar al final):

    \vspace{0.3cm}
    La lista debe contener \textbf{exactamente 10 datos}, organizados por categoría, con:

    \begin{columns}[t]
      \begin{column}{0.55\textwidth}
        Para cada dato:
        \begin{itemize}
          \item El dato exacto (con número, unidades y fuente)
          \item La categoría de la guía a la que pertenece
          \item Puntuación en los 5 criterios R-E-C-R-O
          \item Justificación en una frase de por qué se incluye
        \end{itemize}
      \end{column}
      \begin{column}{0.42\textwidth}
        \begin{notabox}[Distribución mínima]
          \begin{itemize}
            \item 2 datos de Mapa
            \item 2 datos de Monumentos
            \item 2 datos de Rutas
            \item 2 datos de Gráficas
            \item 2 datos de Curiosidades
          \end{itemize}
        \end{notabox}
      \end{column}
    \end{columns}
  \end{evidenciabox}
\end{frame}

% ── FRAME 15: Resumen ─────────────────────────────────────────────
\begin{frame}{Resumen de la sesión 17}
  \begin{resumenbox}
    \begin{columns}[t]
      \begin{column}{0.48\textwidth}
        \begin{itemize}
          \item Hemos aplicado los \textbf{5 criterios} R-E-C-R-O para evaluar datos.
          \item Hemos aprendido a detectar \textbf{datos problemáticos}: sin fuente, sin unidades, obsoletos o implausibles.
          \item El \textbf{árbol de decisión} ayuda a tomar decisiones sistemáticas.
        \end{itemize}
      \end{column}
      \begin{column}{0.48\textwidth}
        \begin{itemize}
          \item La guía tiene \textbf{6 categorías} bien definidas.
          \item La \textbf{selección rigurosa} es parte del proceso matemático.
          \item En la próxima sesión: \textbf{construimos la guía completa}.
        \end{itemize}
        \vspace{0.3cm}
        \begin{formulabox}[BCrear]
          Calidad $>$ Cantidad
        \end{formulabox}
      \end{column}
    \end{columns}
  \end{resumenbox}
\end{frame}

\end{document}

% % ================================================================
%  sesion18.tex  —  Sesión 18 de 20
%  Construimos la guía — Almería en Cifras y Formas
%  Bloque 4: Producto final · Matemáticas 1.º ESO
% ================================================================
\documentclass[aspectratio=169, 12pt, spanish]{beamer}
\usepackage{almeria-beamer}

\renewcommand{\SessionNumber}{18}
\renewcommand{\SessionTitle}{Construimos la guía turística matemática}
\renewcommand{\SessionName}{Sesión 18}
\renewcommand{\SessionObjective}{%
  Elaborar el borrador completo de la guía turística matemática
  integrando mapa con coordenadas, tablas, gráficas, medidas
  y rutas calculadas de Almería.}
\renewcommand{\BlockName}{Bloque 4: Producto final}

\begin{document}

\begin{frame}[plain]\titlepage\end{frame}

\begin{frame}{Índice}
  \begin{columns}[t]
    \begin{column}{0.5\textwidth}
      \begin{enumerate}
        \item Estructura de la guía (7 secciones)
        \item Principios de diseño profesional
        \item Requisitos del mapa de coordenadas
        \item Requisitos de tablas y gráficas
        \item Lista de verificación matemática
      \end{enumerate}
    \end{column}
    \begin{column}{0.5\textwidth}
      \begin{enumerate}\setcounter{enumi}{5}
        \item Datos curiosos de Almería
        \item Ejercicios: construir la guía
        \item Lista de cotejo entre equipos
        \item Error frecuente
        \item Resumen
      \end{enumerate}
    \end{column}
  \end{columns}
\end{frame}

% ---------------------------------------------------------------
\seccionframe{1. Estructura de la guía}
% ---------------------------------------------------------------

\begin{frame}{Las 7 secciones de vuestra guía}
  \begin{columns}[T]
    \begin{column}{0.5\textwidth}
      \begin{enumerate}
        \item[\faBookOpen] \textbf{Portada}\\
          {\footnotesize Título, autores, estadística llamativa, miniatura del mapa}
        \vspace{3pt}
        \item[\faMapMarkerAlt] \textbf{Mapa con coordenadas}\\
          {\footnotesize $\geq 10$ puntos turísticos con $(x,y)$, escala, rosa de los vientos}
        \vspace{3pt}
        \item[\faTable] \textbf{Monumentos en cifras}\\
          {\footnotesize Tabla: nombre, coordenadas, área/altura, visitantes/año}
        \vspace{3pt}
        \item[\faRoute] \textbf{Rutas matemáticas}\\
          {\footnotesize 2--3 rutas: distancias calculadas, tiempo estimado}
      \end{enumerate}
    \end{column}
    \begin{column}{0.5\textwidth}
      \begin{enumerate}
        \item[\faChartLine] \textbf{Gráficas de Almería}\\
          {\footnotesize $\geq 2$ gráficas (temperatura, turismo\ldots) con análisis}
        \vspace{3pt}
        \item[\faLightbulb] \textbf{Datos curiosos}\\
          {\footnotesize 5--8 hechos matemáticos sorprendentes de Almería}
        \vspace{3pt}
        \item[\faBook] \textbf{Glosario matemático}\\
          {\footnotesize $\geq 10$ términos usados en la guía con definición}
      \end{enumerate}
      \vspace{8pt}
      \begin{notabox}[Extensión orientativa]
        \footnotesize
        Mínimo 8 páginas (papel o digital).\\
        Cada sección debe llevar título claro
        y numeración de página.
      \end{notabox}
    \end{column}
  \end{columns}
\end{frame}

% ---------------------------------------------------------------
\seccionframe{2. Principios de diseño profesional}
% ---------------------------------------------------------------

\begin{frame}{Diseño que comunica matemáticas}
  \begin{columns}[T]
    \begin{column}{0.48\textwidth}
      \begin{teoriabox}
        \textbf{4 principios de diseño}\\[4pt]
        \begin{enumerate}
          \item \textbf{Coherencia visual}\\
            Mismos colores, misma tipografía\\
            en todos los apartados.
          \vspace{4pt}
          \item \textbf{Equilibrio texto--imagen}\\
            Aprox.\ 50\% texto, 50\% figuras
            (mapas, tablas, gráficas).
          \vspace{4pt}
          \item \textbf{Jerarquía de información}\\
            Título $>$ subtítulo $>$ cuerpo de texto\\
            (tamaños distintos, negrita).
          \vspace{4pt}
          \item \textbf{Accesibilidad}\\
            Ejes etiquetados, unidades presentes,
            fórmulas explicadas.
        \end{enumerate}
      \end{teoriabox}
    \end{column}
    \begin{column}{0.48\textwidth}
      % Maqueta esquemática de doble página
      \begin{tikzpicture}[scale=0.62]
        % Página izquierda
        \fill[AlmAzul] (0,8) rectangle (6,8.6);
        \node[text=AlmBlanco, font=\tiny\bfseries] at (3,8.3)
          {ALMERÍA EN CIFRAS Y FORMAS};
        \fill[AlmAmbar!30, rounded corners=2pt]
          (0.3,5.8) rectangle (5.7,7.8);
        \node[font=\tiny, text=AlmNegro, align=center] at (3,6.8)
          {MAPA DE COORDENADAS\\(10+ puntos etiquetados)};
        \foreach \y in {5.1,4.7,4.3,3.9}{
          \fill[AlmGris!20] (0.3,\y) rectangle (5.7,\y+0.25);
        }
        \node[font=\tiny, color=AlmGris] at (3,3.6) {texto · texto · texto};
        % Página derecha
        \fill[AlmAzul!12] (6.5,3.5) rectangle (12.5,8.6);
        \fill[AlmAzulC!25, rounded corners=2pt]
          (6.8,5.5) rectangle (12.2,8.3);
        \node[font=\tiny, text=AlmAzul, align=center] at (9.5,6.9)
          {GRÁFICA: turistas/mes\\(pgfplots / mano)};
        \foreach \y in {4.8,4.4,4.0,3.7}{
          \fill[AlmGris!20] (6.8,\y) rectangle (12.2,\y+0.22);
        }
        % Títulos de sección
        \node[font=\tiny\bfseries, color=AlmAzul] at (3,8.0)
          {Sección 2: Mapa};
        \node[font=\tiny\bfseries, color=AlmAzul] at (9.5,8.0)
          {Sección 5: Gráficas};
      \end{tikzpicture}
    \end{column}
  \end{columns}
\end{frame}

% ---------------------------------------------------------------
\seccionframe{3. Requisitos del mapa}
% ---------------------------------------------------------------

\begin{frame}{El mapa de coordenadas: qué debe incluir}
  \begin{columns}[c]
    \begin{column}{0.52\textwidth}
      \begin{itemize}
        \item[\faCheck] \textbf{Escala indicada}
              (ej.: 1 unidad = 200 m)
        \item[\faCheck] \textbf{Rosa de los vientos} (Norte)
        \item[\faCheck] \textbf{Ejes etiquetados} ($x$ e $y$)
        \item[\faCheck] \textbf{$\geq 10$ puntos} con nombre
              y coordenadas $(x,y)$
        \item[\faCheck] \textbf{Código de color} por tipo:\\
          {\footnotesize
          \textcolor{AlmAzulM}{$\bullet$} Monumentos \quad
          \textcolor{AlmAzulC}{$\bullet$} Playas \\
          \textcolor{AlmVerde}{$\bullet$} Parques \quad
          \textcolor{AlmTerra}{$\bullet$} Gastronomía}
        \item[\faCheck] \textbf{Leyenda} con los colores
        \item[\faCheck] \textbf{Al menos una ruta} dibujada
              con flechas
      \end{itemize}
    \end{column}
    \begin{column}{0.44\textwidth}
      \begin{tikzpicture}[scale=0.75]
        \draw[AlmFondo!70, step=1cm] (0,0) grid (6,5.5);
        \draw[->, AlmAzulM] (0,0)--(6.4,0) node[right,font=\tiny]{$x$};
        \draw[->, AlmAzulM] (0,0)--(0,5.8) node[above,font=\tiny]{$y$};
        % Puntos
        \foreach \x/\y/\c/\n in {
          1/4/AlmAzulM/Alcazaba,
          2/2/AlmAzulM/Catedral,
          3/2/AlmTerra/Mercado,
          4/1/AlmAzulC/Puerto,
          5/2/AlmAzulC/Playa,
          1/2/AlmVerde/Parque}{
          \node[circle,fill=\c,inner sep=2pt] at (\x,\y) {};
          \node[font=\tiny, above, color=\c] at (\x,\y) {\n};
        }
        % Ruta
        \draw[AlmAmbar, line width=1.2pt, ->]
          (1,4)--(2,2)--(3,2)--(4,1)--(5,2);
        % Escala
        \draw[acota] (0,-0.5) -- (1,-0.5)
          node[midway,below,font=\tiny]{200 m};
        % Norte
        \node[font=\large, color=AlmAzulM] at (5.5,5.2) {$\uparrow$N};
      \end{tikzpicture}
    \end{column}
  \end{columns}
\end{frame}

% ---------------------------------------------------------------
\seccionframe{4. Tablas y gráficas}
% ---------------------------------------------------------------

\begin{frame}{Requisitos de tablas y gráficas}
  \begin{columns}[T]
    \begin{column}{0.48\textwidth}
      \textbf{Tablas} — lista de verificación:
      \begin{itemize}
        \item[\faCheck] Encabezados claros \textbf{con unidades}
        \item[\faCheck] Datos \textbf{ordenados} por la variable $x$
        \item[\faCheck] Fuente de los datos indicada
        \item[\faCheck] Mínimo \textbf{8 filas} de datos
        \item[\faCheck] Sin celdas vacías sin justificar
      \end{itemize}
      \vspace{6pt}
      \begin{tabular}{lrr}
        \toprule
        \textbf{Monumento} & \textbf{Área (m²)} & \textbf{Visit./año}\\
        \midrule
        Alcazaba & 43\,000 & 300\,000\\
        Catedral & 5\,800  & 180\,000\\
        Cable Inglés & 2\,200 & 95\,000\\
        Puerto & 350\,000 & --\\
        \bottomrule
      \end{tabular}
    \end{column}
    \begin{column}{0.48\textwidth}
      \textbf{Gráficas} — lista de verificación:
      \begin{itemize}
        \item[\faCheck] \textbf{Título} informativo
        \item[\faCheck] Eje $x$ etiquetado con unidad
        \item[\faCheck] Eje $y$ etiquetado con unidad
        \item[\faCheck] Eje $y$ comienza en \textbf{0} (salvo excepción justificada)
        \item[\faCheck] Escala \textbf{regular} en ambos ejes
        \item[\faCheck] Leyenda si hay más de una serie
        \item[\faCheck] Tipo de gráfica \textbf{adecuado} al tipo de variable
      \end{itemize}
      \vspace{6pt}
      \begin{notabox}[Recuerda]
        \footnotesize
        Variables discretas → diagrama de \textbf{barras}\\
        Variables continuas / tiempo → gráfica \textbf{lineal}\\
        Partes de un todo → diagrama de \textbf{sectores}
      \end{notabox}
    \end{column}
  \end{columns}
\end{frame}

% ---------------------------------------------------------------
\seccionframe{5. Lista de verificación matemática}
% ---------------------------------------------------------------

\begin{frame}{Lista de verificación antes de entregar}
  \begin{columns}[T]
    \begin{column}{0.48\textwidth}
      \begin{tcolorbox}[
        enhanced, colback=AlmVerde!8, colframe=AlmVerde,
        title={\faCheckDouble\enskip Verificación matemática},
        fonttitle=\bfseries\small\color{AlmBlanco},
        attach boxed title to top left={yshift=-2.2mm, xshift=4mm},
        boxed title style={colback=AlmVerde, arc=3pt, boxrule=0pt},
        arc=4pt]
        \footnotesize
        \begin{itemize}
          \item[$\square$] Todas las \textbf{distancias} calculadas con
                $d=\sqrt{\Delta x^2+\Delta y^2}$
          \item[$\square$] Todas las \textbf{áreas} con la fórmula correcta
                y en m²
          \item[$\square$] Los \textbf{perímetros} en m (no m²)
          \item[$\square$] Las \textbf{proporcionalidades} verificadas
                ($y/x = k$ constante)
          \item[$\square$] Las \textbf{coordenadas} escritas como $(x,y)$,
                no $(y,x)$
          \item[$\square$] Los \textbf{gráficos} con ejes etiquetados y
                escala que empieza en 0
          \item[$\square$] Todos los \textbf{tiempos} calculados
                con $t = d/v$
        \end{itemize}
      \end{tcolorbox}
    \end{column}
    \begin{column}{0.48\textwidth}
      \begin{tcolorbox}[
        enhanced, colback=AlmAmbar!10, colframe=AlmAmbar!70!black,
        title={\faStickyNote\enskip Verificación de presentación},
        fonttitle=\bfseries\small\color{AlmAzul},
        attach boxed title to top left={yshift=-2.2mm, xshift=4mm},
        boxed title style={colback=AlmAmbar, coltext=AlmAzul, arc=3pt},
        arc=4pt]
        \footnotesize
        \begin{itemize}
          \item[$\square$] Portada con título, autores y estadística
          \item[$\square$] Todas las páginas numeradas
          \item[$\square$] Índice o tabla de contenidos
          \item[$\square$] Fuentes de datos citadas
          \item[$\square$] Glosario con $\geq 10$ términos
          \item[$\square$] Sin faltas de ortografía
          \item[$\square$] Diseño coherente (mismos colores y fuentes)
        \end{itemize}
      \end{tcolorbox}
    \end{column}
  \end{columns}
\end{frame}

% ---------------------------------------------------------------
\seccionframe{6. Datos curiosos de Almería}
% ---------------------------------------------------------------

\begin{frame}{Datos curiosos para vuestra guía}
  \begin{columns}[T]
    \begin{column}{0.5\textwidth}
      \begin{notabox}[¿Sabías que\ldots?]
        \begin{itemize}
          \item La \textbf{Alcazaba} ocupa \SI{43000}{m^2}
                $\approx$ 6 campos de fútbol.
          \item El \textbf{Cable Inglés} mide \SI{1.1}{km}
                y pesa \num{1400} toneladas.
          \item Almería tiene \textbf{3\,000 horas de sol}
                al año — es la ciudad más soleada de España.
          \item El \textbf{faro del Cabo de Gata} tiene
                \SI{32}{m} de altura sobre el mar.
          \item La \textbf{Rambla de Almería} mide \SI{370}{m}
                de longitud y hasta \SI{50}{m} de anchura.
        \end{itemize}
      \end{notabox}
    \end{column}
    \begin{column}{0.46\textwidth}
      \begin{notabox}[Más datos\ldots]
        \begin{itemize}
          \item El \textbf{Puerto de Almería} ocupa
                \SI{35}{ha} = \SI{350000}{m^2}.
          \item La \textbf{catedral} (1524) tiene
                \SI{55}{m} de altura en su torre.
          \item El Parque Natural del Cabo de Gata
                tiene \SI{38000}{ha}.
          \item Almería recibe $\approx$\,\SI{3}{millones}
                de turistas al año.
          \item La temperatura media anual es \SI{18.5}{\celsius}.
        \end{itemize}
      \end{notabox}
    \end{column}
  \end{columns}
\end{frame}

% ---------------------------------------------------------------
\seccionframe{7. Ejercicios: construir la guía}
% ---------------------------------------------------------------

\begin{frame}{Construye la guía — Taller de aula}
  \begin{columns}[T]
    \begin{column}{0.48\textwidth}
      \begin{aplicarbox}
        \footnotesize
        \textbf{Sección «Monumentos en cifras»}\\[4pt]
        Completa la tabla con \textbf{5 monumentos} de Almería.
        Para cada uno indica:
        \begin{enumerate}
          \item Nombre y breve descripción
          \item Coordenadas $(x,y)$ en el mapa del proyecto
          \item Área o altura (con unidades)
          \item Dato histórico (año de construcción, etc.)
          \item Un hecho matemático curioso
        \end{enumerate}
      \end{aplicarbox}
    \end{column}
    \begin{column}{0.48\textwidth}
      \begin{crearbox}
        \footnotesize
        \textbf{Sección «Ruta de 3 horas»}\\[4pt]
        Diseña una ruta turística visitando
        \textbf{6 lugares} del mapa. Para cada tramo:
        \begin{enumerate}
          \item Anota los puntos inicial y final con coordenadas
          \item Calcula la distancia con Pitágoras
          \item Multiplica por la escala (1 u = 200 m)
          \item Calcula el tiempo a 4{,}5 km/h
        \end{enumerate}
        Suma todas las distancias y tiempos.
        ¿Cabe en 3 horas? Ajusta si es necesario.
      \end{crearbox}
    \end{column}
  \end{columns}
\end{frame}

% ---------------------------------------------------------------
\begin{frame}{Evidencia del día}
  \begin{evidenciabox}
    Entrega el \textbf{borrador de la guía} con al menos
    \textbf{4 secciones completas}:
    \begin{enumerate}
      \item Portada con título, autores y una estadística de Almería
      \item Mapa con $\geq 8$ puntos etiquetados, escala y norte
      \item Tabla de monumentos con $\geq 5$ filas y datos verificados
      \item Ruta turística con distancias calculadas y tiempo total
    \end{enumerate}
    \vspace{4pt}
    \textit{Recuerda: toda información numérica debe
    tener su unidad. No existe el «200 metros cuadrados
    de perímetro».}
  \end{evidenciabox}
\end{frame}

% ---------------------------------------------------------------
\begin{frame}{¡Cuidado con este error!}
  \begin{errorbox}
    \textbf{Una guía matemática no es solo una guía con números.}\\[6pt]
    \incorrecto\ «La Alcazaba es un monumento histórico muy bonito
    situado en Almería. Tiene muchas salas y vistas preciosas.»\\
    {\footnotesize ← Esto es una guía turística clásica, no matemática.}\\[8pt]
    \correcto\ «La Alcazaba (coordenadas $(1,4)$ en nuestro mapa,
    escala 1\,u = 200\,m) ocupa \SI{43000}{m^2} $\approx$ \SI{4.3}{ha}
    y registra \num{300000} visitantes anuales. Su muro de mayor altura
    mide aproximadamente \SI{15}{m}.»\\
    {\footnotesize ← Mismo lugar, pero con contenido matemático real.}
  \end{errorbox}
\end{frame}

% ---------------------------------------------------------------
\begin{frame}{Resumen de la sesión 18}
  \begin{resumenbox}
    \begin{itemize}
      \item La guía tiene \textbf{7 secciones}: portada, mapa,
            monumentos en cifras, rutas, gráficas, datos curiosos
            y glosario.
      \item Todo dato matemático necesita su \textbf{unidad}.
      \item Las gráficas necesitan \textbf{título, ejes etiquetados
            y escala desde 0}.
      \item Las tablas deben estar \textbf{ordenadas por $x$}.
      \item El mapa debe tener \textbf{escala, norte y $\geq 10$ puntos}.
      \item Calidad $>$ cantidad: datos \textbf{verificados y relevantes}.
    \end{itemize}
  \end{resumenbox}
  \vspace{6pt}
  \begin{center}
    {\small\color{AlmAzulM}\faArrowCircleRight\enskip
    \textbf{Sesión 19:} Revisión entre pares y mejora del borrador}
  \end{center}
\end{frame}

\end{document}

% % ================================================================
%  sesion19.tex  —  Sesión 19 de 20
%  Ensayo y mejora
%  Bloque 4: Producto final · Matemáticas 1.º ESO
% ================================================================
\documentclass[aspectratio=169, 12pt, spanish]{beamer}
\usepackage{almeria-beamer}

\renewcommand{\SessionNumber}{19}
\renewcommand{\SessionTitle}{Ensayo y mejora}
\renewcommand{\SessionName}{Sesión 19}
\renewcommand{\SessionObjective}{%
  Revisar y mejorar la guía turística matemática mediante
  autoevaluación y coevaluación estructurada, identificando
  errores conceptuales y de presentación.}
\renewcommand{\BlockName}{Bloque 4: Producto final}

\begin{document}

\begin{frame}[plain]\titlepage\end{frame}

\begin{frame}{Índice}
  \begin{columns}[t]
    \begin{column}{0.5\textwidth}
      \begin{enumerate}
        \item ¿Por qué revisar?
        \item Tipos de errores a buscar
        \item Rúbrica de autoevaluación
        \item Protocolo de coevaluación
        \item Errores más comunes del proyecto
      \end{enumerate}
    \end{column}
    \begin{column}{0.5\textwidth}
      \begin{enumerate}\setcounter{enumi}{5}
        \item El ciclo de mejora
        \item Ejercicios: evaluar y corregir
        \item Evidencia del día
        \item Error de actitud
        \item Resumen
      \end{enumerate}
    \end{column}
  \end{columns}
\end{frame}

% ---------------------------------------------------------------
\seccionframe{1. ¿Por qué revisar?}
% ---------------------------------------------------------------

\begin{frame}{La revisión es parte del proceso matemático}
  \begin{columns}[c]
    \begin{column}{0.52\textwidth}
      \begin{teoriabox}
        Ningún trabajo matemático se entrega sin revisión.
        Los matemáticos, ingenieros y científicos revisan
        sus trabajos \textbf{varias veces} antes de publicarlos.\\[6pt]
        El primer borrador \textbf{siempre} contiene:
        \begin{itemize}
          \item Errores de cálculo
          \item Unidades incorrectas
          \item Gráficas incompletas
          \item Datos sin verificar
          \item Explicaciones poco claras
        \end{itemize}
      \end{teoriabox}
    \end{column}
    \begin{column}{0.44\textwidth}
      \begin{tikzpicture}[
        nodo/.style={
          rounded corners=5pt,
          draw=AlmAzulM, fill=AlmFondo,
          font=\small\bfseries,
          text width=2.2cm, align=center,
          inner sep=5pt
        },
        flecha/.style={->, AlmAmbar, line width=1.5pt, shorten >=2pt}
      ]
        \node[nodo, fill=BCrear!15] (c) at (0,3.5) {Crear};
        \node[nodo] (r) at (2.5,2)   {Revisar};
        \node[nodo] (i) at (2.5,0)   {Identificar\\error};
        \node[nodo] (co)at (0,-1.2)  {Corregir};
        \node[nodo] (v) at (-2.5,0)  {Verificar};
        \node[nodo, fill=AlmVerde!15] (p) at (-2.5,2) {Presentar};
        \draw[flecha] (c)  -- (r);
        \draw[flecha] (r)  -- (i);
        \draw[flecha] (i)  -- (co);
        \draw[flecha] (co) -- (v);
        \draw[flecha] (v)  -- (p);
        \draw[flecha, dashed] (r) to[bend left=20] (c)
          node[midway, above right, font=\tiny, color=AlmGris]
          {¿Todo bien?};
      \end{tikzpicture}
    \end{column}
  \end{columns}
\end{frame}

% ---------------------------------------------------------------
\seccionframe{2. Tipos de errores}
% ---------------------------------------------------------------

\begin{frame}{Cuatro tipos de errores a buscar}
  \begin{columns}[T]
    \begin{column}{0.48\textwidth}
      \begin{tcolorbox}[enhanced, colback=BRecordar!8,
        colframe=BRecordar, fonttitle=\bfseries\small\color{AlmBlanco},
        title={\faCalculator\enskip Errores matemáticos},
        attach boxed title to top left={yshift=-2.2mm, xshift=4mm},
        boxed title style={colback=BRecordar, arc=3pt},
        arc=4pt, leftrule=4pt]
        \footnotesize
        \begin{itemize}
          \item Fórmula incorrecta ($A = b + h$ en vez de $b \times h$)
          \item Operación aritmética errónea
          \item Unidad equivocada (m² en un perímetro)
          \item Raíz cuadrada olvidada en Pitágoras
          \item Coordenadas invertidas $(y,x)$ en vez de $(x,y)$
        \end{itemize}
      \end{tcolorbox}
      \vspace{6pt}
      \begin{tcolorbox}[enhanced, colback=BAnalizar!8,
        colframe=BAnalizar, fonttitle=\bfseries\small\color{AlmBlanco},
        title={\faDatabase\enskip Errores de datos},
        attach boxed title to top left={yshift=-2.2mm, xshift=4mm},
        boxed title style={colback=BAnalizar, arc=3pt},
        arc=4pt, leftrule=4pt]
        \footnotesize
        \begin{itemize}
          \item Dato inventado sin fuente
          \item Dato obsoleto (más de 5 años)
          \item Dato sin unidades
          \item Valor implausible (Alcazaba 430\,m²)
        \end{itemize}
      \end{tcolorbox}
    \end{column}
    \begin{column}{0.48\textwidth}
      \begin{tcolorbox}[enhanced, colback=BComprender!8,
        colframe=BComprender, fonttitle=\bfseries\small\color{AlmBlanco},
        title={\faChartBar\enskip Errores de representación},
        attach boxed title to top left={yshift=-2.2mm, xshift=4mm},
        boxed title style={colback=BComprender, arc=3pt},
        arc=4pt, leftrule=4pt]
        \footnotesize
        \begin{itemize}
          \item Ejes sin etiqueta o sin unidades
          \item Escala no regular (saltos irregulares)
          \item Eje $y$ no empieza en 0 sin justificación
          \item Tipo de gráfica inadecuado
          \item Tabla no ordenada por $x$
        \end{itemize}
      \end{tcolorbox}
      \vspace{6pt}
      \begin{tcolorbox}[enhanced, colback=BEvaluar!8,
        colframe=BEvaluar, fonttitle=\bfseries\small\color{AlmBlanco},
        title={\faCommentAlt\enskip Errores de comunicación},
        attach boxed title to top left={yshift=-2.2mm, xshift=4mm},
        boxed title style={colback=BEvaluar, arc=3pt},
        arc=4pt, leftrule=4pt]
        \footnotesize
        \begin{itemize}
          \item Descripción sin datos numéricos
          \item Cálculo sin explicación de los pasos
          \item Conclusión vaga («hay muchos turistas»)
          \item Vocabulario matemático incorrecto
        \end{itemize}
      \end{tcolorbox}
    \end{column}
  \end{columns}
\end{frame}

% ---------------------------------------------------------------
\seccionframe{3. Rúbrica de autoevaluación}
% ---------------------------------------------------------------

\begin{frame}{Rúbrica de autoevaluación}
  \begin{center}
  \footnotesize
  \begin{tabular}{p{2.8cm} p{2.6cm} p{2.6cm} p{2.6cm} p{2.6cm}}
    \toprule
    \textbf{Criterio}
      & \textbf{4 — Excelente}
      & \textbf{3 — Bien}
      & \textbf{2 — Regular}
      & \textbf{1 — Mejora}\\
    \midrule
    Exactitud\\matemática
      & Todos los cálculos correctos
      & 1--2 errores menores
      & 3--4 errores
      & Errores graves\\[4pt]
    Unidades
      & Siempre correctas
      & Casi siempre
      & Algunas mal
      & Frecuentemente mal\\[4pt]
    Gráficas
      & Completas y claras
      & Falta 1 elemento
      & Faltan 2 elementos
      & Incompletas\\[4pt]
    Selección\\de datos
      & Verificados y relevantes
      & Mayoría correctos
      & Algunos irrelevantes
      & Sin criterio\\[4pt]
    Comunicación
      & Clara y matemática
      & Casi siempre clara
      & A veces confusa
      & Vaga o incorrecta\\
    \bottomrule
  \end{tabular}
  \end{center}
  \vspace{4pt}
  \begin{center}
    \small\textbf{Puntuación máxima: 20 puntos}\\
    {\footnotesize 18--20: Sobresaliente · 14--17: Notable ·
    10--13: Bien · $<$10: Necesita revisión}
  \end{center}
\end{frame}

% ---------------------------------------------------------------
\seccionframe{4. Protocolo de coevaluación}
% ---------------------------------------------------------------

\begin{frame}{Coevaluación: revisión entre equipos}
  \begin{columns}[c]
    \begin{column}{0.55\textwidth}
      \begin{teoriabox}
        \textbf{Protocolo «Estrella + Deseo»}\\[6pt]
        \begin{enumerate}
          \item \textbf{Intercambia} la guía con otro equipo.
          \item Lee con atención \textbf{3 secciones concretas}.
          \item Para cada sección anota:\\[4pt]
            \fastar\ \textbf{Estrella} (1): algo que funciona muy bien\\
            \fastar\ \textbf{Estrella} (2): otro punto positivo\\
            \faCommentAlt\ \textbf{Deseo} (1): algo a mejorar\\
            (sé específico: «la tabla de la pág.~3
            no tiene unidades en la cabecera»)
          \item Devuelve la guía con tus notas escritas.
          \item Lee el feedback recibido y \textbf{decide}
                qué cambios aplicar.
        \end{enumerate}
      \end{teoriabox}
    \end{column}
    \begin{column}{0.41\textwidth}
      \textbf{Normas de la coevaluación:}
      \begin{itemize}
        \item Sé \textbf{específico}: no «está bien» sino
              «la gráfica de barras tiene los ejes
              correctamente etiquetados con unidades».
        \item Sé \textbf{constructivo}: «En la pág.~4 falta
              la unidad del área. Añade m² al resultado
              para que sea correcto.»
        \item No es una crítica personal: estás
              evaluando el \textbf{trabajo matemático},
              no a la persona.
        \item El receptor del feedback puede
              \textbf{aceptar o rechazar} cada sugerencia
              con justificación matemática.
      \end{itemize}
    \end{column}
  \end{columns}
\end{frame}

% ---------------------------------------------------------------
\seccionframe{5. Errores más comunes}
% ---------------------------------------------------------------

\begin{frame}{Los 8 errores más frecuentes en este proyecto}
  \begin{columns}[T]
    \begin{column}{0.48\textwidth}
      \begin{enumerate}
        \item \incorrecto\ Perímetro en m²\\
              \correcto\ Perímetro en m (lineal)
        \vspace{4pt}
        \item \incorrecto\ Coordenadas como $(y,x)$\\
              \correcto\ Siempre $(x,y)$: primero horizontal
        \vspace{4pt}
        \item \incorrecto\ Diagrama de barras para temperatura\\
              \correcto\ Gráfica lineal para variables continuas
        \vspace{4pt}
        \item \incorrecto\ Tabla desordenada\\
              \correcto\ Ordenada por la variable $x$
      \end{enumerate}
    \end{column}
    \begin{column}{0.48\textwidth}
      \begin{enumerate}\setcounter{enumi}{4}
        \item \incorrecto\ $d = 3^2 + 4^2 = 25$ (falta $\sqrt{}$)\\
              \correcto\ $d = \sqrt{3^2+4^2} = 5$
        \vspace{4pt}
        \item \incorrecto\ «La Alcazaba tiene 43.000»\\
              \correcto\ «\SI{43000}{m^2}» (siempre con unidad)
        \vspace{4pt}
        \item \incorrecto\ Eje $y$ empieza en 2.000.000\\
              \correcto\ Empieza en 0 (o justifica por qué no)
        \vspace{4pt}
        \item \incorrecto\ $A = b + h$\\
              \correcto\ $A = b \times h$ para el rectángulo
      \end{enumerate}
    \end{column}
  \end{columns}
\end{frame}

% ---------------------------------------------------------------
\seccionframe{7. Ejercicios: evaluar y corregir}
% ---------------------------------------------------------------

\begin{frame}{Ejercicios multinivel}
  \begin{columns}[T]
    \begin{column}{0.48\textwidth}
      \begin{evaluarbox}
        \footnotesize\dificultad{3}\\[4pt]
        Aplica la rúbrica de autoevaluación
        a \textbf{tu propia guía}.\\[4pt]
        \begin{enumerate}
          \item Puntúa cada criterio (1--4).
          \item Calcula la puntuación total.
          \item Identifica tus \textbf{3 puntos más débiles}.
          \item Escribe un plan concreto de mejora
                para cada punto: qué cambiarás y cómo.
        \end{enumerate}
      \end{evaluarbox}
    \end{column}
    \begin{column}{0.48\textwidth}
      \begin{crearbox}
        \footnotesize\dificultad{4}\\[4pt]
        Basándote en el \textbf{feedback del equipo revisor},
        aplica \textbf{3 mejoras concretas} a tu guía.\\[4pt]
        Para cada mejora documenta:
        \begin{enumerate}
          \item \textbf{Versión original} (copia textual)
          \item \textbf{Problema identificado}
                (tipo de error y criterio de la rúbrica)
          \item \textbf{Versión corregida}
                (con el cambio aplicado)
        \end{enumerate}
        Adjunta la documentación de mejoras
        a tu guía final.
      \end{crearbox}
    \end{column}
  \end{columns}
\end{frame}

% ---------------------------------------------------------------
\begin{frame}{Evidencia del día}
  \begin{evidenciabox}
    Entrega:
    \begin{enumerate}
      \item \textbf{Rúbrica de autoevaluación} completada
            con puntuación numérica en cada criterio y
            reflexión escrita sobre los 3 puntos más débiles.
      \item \textbf{Guía revisada} con las correcciones
            aplicadas marcadas claramente (usa un color diferente
            para las correcciones).
      \item \textbf{Hoja de mejoras} documentando las 3 correcciones
            (versión original → problema → versión corregida).
    \end{enumerate}
    \vspace{4pt}
    \textit{Una guía bien revisada vale mucho más que una
    guía apresurada con muchos errores. La calidad matemática
    es lo que hace profesional vuestro trabajo.}
  \end{evidenciabox}
\end{frame}

% ---------------------------------------------------------------
\begin{frame}{¡Cuidado con este error de actitud!}
  \begin{errorbox}
    \textbf{«Mi cálculo está bien porque yo lo hice»}\\[6pt]
    \incorrecto\ Resistirse al feedback y defender los errores
    por orgullo. La revisión crítica es debilidad.\\[8pt]
    \correcto\ Los mejores matemáticos y científicos del mundo
    \textbf{dan la bienvenida a las correcciones} porque les
    permiten mejorar. En matemáticas hay una respuesta correcta
    y es \textbf{verificable}. Si alguien señala un error, es
    una oportunidad, no un ataque.\\[6pt]
    «\textit{Un error descubierto hoy es mejor que
    un error entregado mañana.}»
  \end{errorbox}
\end{frame}

% ---------------------------------------------------------------
\begin{frame}{Resumen de la sesión 19}
  \begin{resumenbox}
    \begin{itemize}
      \item La revisión es una \textbf{parte fundamental} del proceso
            matemático, no un extra opcional.
      \item Hay 4 tipos de errores: \textbf{matemáticos},
            de datos, de representación y de comunicación.
      \item La \textbf{rúbrica} permite autoevaluar
            objetivamente con criterios claros.
      \item La \textbf{coevaluación} con el protocolo
            «Estrella + Deseo» da feedback constructivo.
      \item Cada corrección debe documentarse:
            original → problema → versión corregida.
    \end{itemize}
  \end{resumenbox}
  \vspace{6pt}
  \begin{center}
    {\small\color{AlmAzulM}\faArrowCircleRight\enskip
    \textbf{Sesión 20 (última):}
    Exposición final y cierre del proyecto}
  \end{center}
\end{frame}

\end{document}

% % ================================================================
%  sesion20.tex  —  Sesión 20 de 20
%  Exposición y cierre del proyecto
%  Bloque 4: Producto final · Matemáticas 1.º ESO
% ================================================================
\documentclass[aspectratio=169, 12pt, spanish]{beamer}
\usepackage{almeria-beamer}

\renewcommand{\SessionNumber}{20}
\renewcommand{\SessionTitle}{Exposición y cierre del proyecto}
\renewcommand{\SessionName}{Sesión 20}
\renewcommand{\SessionObjective}{%
  Presentar la guía turística matemática ante la clase,
  defender las decisiones matemáticas tomadas, y reflexionar
  sobre los aprendizajes adquiridos a lo largo del proyecto.}
\renewcommand{\BlockName}{Bloque 4: Producto final}

\begin{document}

\begin{frame}[plain]\titlepage\end{frame}

\begin{frame}{Índice}
  \begin{columns}[t]
    \begin{column}{0.5\textwidth}
      \begin{enumerate}
        \item Estructura de la exposición oral
        \item Consejos para presentar matemáticas
        \item Vocabulario matemático requerido
        \item Rúbrica de evaluación oral
        \item Síntesis de los 4 bloques
      \end{enumerate}
    \end{column}
    \begin{column}{0.5\textwidth}
      \begin{enumerate}\setcounter{enumi}{5}
        \item Lo que hemos construido
        \item Preguntas de reflexión final
        \item Ejercicios (evaluar y crear)
        \item Error final: el mito
        \item Cierre del proyecto
      \end{enumerate}
    \end{column}
  \end{columns}
\end{frame}

% ---------------------------------------------------------------
\seccionframe{1. Estructura de la exposición}
% ---------------------------------------------------------------

\begin{frame}{Exposición oral — 5 minutos por equipo}
  \begin{columns}[c]
    \begin{column}{0.55\textwidth}
      \begin{tikzpicture}[
        bloque/.style={
          rounded corners=4pt,
          draw=AlmAzulM, fill=AlmFondo,
          text width=6.5cm, align=left,
          inner sep=7pt, font=\small
        },
        tiempo/.style={
          circle, fill=AlmAmbar, text=AlmAzul,
          font=\bfseries\small, inner sep=3pt,
          minimum size=0.8cm
        }
      ]
        \node[tiempo] (t1) at (0,4.5) {1'};
        \node[bloque, right=0.3cm of t1] {
          \textbf{Introducción}\\
          {\footnotesize Equipo, objetivo del proyecto,
          monumento o zona de Almería elegida como tema central.}
        };
        \node[tiempo] (t2) at (0,3.1) {2'};
        \node[bloque, right=0.3cm of t2] {
          \textbf{2 resultados matemáticos clave}\\
          {\footnotesize Un cálculo completo (área, distancia, etc.)
          y una gráfica explicada. Mostrad vuestros pasos.}
        };
        \node[tiempo] (t3) at (0,1.7) {1'};
        \node[bloque, right=0.3cm of t3] {
          \textbf{Descubrimiento más sorprendente}\\
          {\footnotesize El dato de Almería que más os ha sorprendido
          y por qué tiene importancia matemática.}
        };
        \node[tiempo] (t4) at (0,0.3) {1'};
        \node[bloque, right=0.3cm of t4] {
          \textbf{Reflexión}\\
          {\footnotesize Una cosa que habéis aprendido y una
          que haríais diferente en un segundo intento.}
        };
        % Flecha temporal
        \draw[AlmAmbar, line width=1.5pt, ->]
          (t1.south) -- (t2.north);
        \draw[AlmAmbar, line width=1.5pt, ->]
          (t2.south) -- (t3.north);
        \draw[AlmAmbar, line width=1.5pt, ->]
          (t3.south) -- (t4.north);
      \end{tikzpicture}
    \end{column}
    \begin{column}{0.41\textwidth}
      \begin{notabox}[Materiales permitidos]
        \begin{itemize}
          \item La guía impresa o en pantalla
          \item Fichas de notas (no leer todo)
          \item El mapa de coordenadas para señalar
          \item Las gráficas de la guía en grande
        \end{itemize}
      \end{notabox}
      \vspace{6pt}
      \begin{notabox}[Preguntas del profesor]
        \footnotesize
        Tras la exposición habrá 1--2 preguntas
        que \textbf{cualquier miembro del equipo}
        debe poder responder. No son trampa:
        solo verifican que entendéis lo que
        habéis hecho.
      \end{notabox}
    \end{column}
  \end{columns}
\end{frame}

% ---------------------------------------------------------------
\seccionframe{2. Consejos para presentar}
% ---------------------------------------------------------------

\begin{frame}{Consejos para presentar matemáticas en público}
  \begin{columns}[T]
    \begin{column}{0.48\textwidth}
      \textbf{\color{AlmVerde}\faCheck\ Sí hacer:}
      \begin{itemize}
        \item Mirar a la audiencia, no al papel.
        \item Señalar los números en el mapa o la gráfica
              mientras los explicas.
        \item Usar vocabulario matemático correcto:
              «\textit{calculamos el área con $b \times h$}»
              en vez de «\textit{lo multiplicamos}».
        \item Justificar cada decisión:
              «Elegimos la gráfica lineal porque
              la temperatura es variable continua.»
        \item Hablar despacio en los cálculos:
              la audiencia necesita seguir los pasos.
      \end{itemize}
    \end{column}
    \begin{column}{0.48\textwidth}
      \textbf{\color{AlmTerra}\faTimes\ No hacer:}
      \begin{itemize}
        \item Leer el trabajo en voz alta sin mirarnos.
        \item Decir «este número» sin señalar cuál.
        \item Usar «esto», «aquello», «la cosa»
              en lugar del nombre matemático.
        \item Ignorar las unidades en la presentación oral.
        \item Excusarse continuamente: «Sé que no está bien pero\ldots»\\
              {\footnotesize Si sabes que hay un error,
              corrígelo antes de presentar.}
      \end{itemize}
      \vspace{6pt}
      \begin{notabox}[Ante preguntas difíciles]
        \footnotesize
        «Buena pregunta. Según nuestros datos\ldots»\\
        «No lo hemos calculado, pero podríamos\ldots»\\
        «No estoy seguro/a, lo podría verificar con\ldots»
      \end{notabox}
    \end{column}
  \end{columns}
\end{frame}

% ---------------------------------------------------------------
\seccionframe{3. Vocabulario matemático}
% ---------------------------------------------------------------

\begin{frame}{Vocabulario matemático: ¡úsalo!}
  \begin{columns}[T]
    \begin{column}{0.48\textwidth}
      Debéis usar \textbf{al menos 10} de estos términos
      durante la exposición:
      \vspace{4pt}
      \begin{center}
      \begin{tabular}{ll}
        \toprule
        \textbf{Término} & \textbf{Contexto de uso}\\
        \midrule
        Coordenadas & Localización en el mapa\\
        Perímetro   & Borde de una zona\\
        Área        & Superficie de un espacio\\
        Función     & Relación x→y\\
        Proporcionalidad & k = y/x constante\\
        Variable    & Lo que medimos\\
        Tendencia   & Dirección de una gráfica\\
        Correlación & Relación entre dos variables\\
        Escala      & Conversión mapa--realidad\\
        Magnitud    & Lo que se mide (longitud, etc.)\\
        Perpendicular & Ángulo recto entre calles\\
        Polígono    & Figura cerrada de lados rectos\\
        \bottomrule
      \end{tabular}
      \end{center}
    \end{column}
    \begin{column}{0.48\textwidth}
      \textbf{Ejemplos de frases bien construidas:}\\[6pt]
      \begin{tcolorbox}[colback=AlmVerde!8, colframe=AlmVerde,
                        arc=4pt, boxrule=0.5pt, fontupper=\small]
        «La \textbf{distancia} entre la Alcazaba y el Puerto
        es $d = \sqrt{3^2+3^2} = \sqrt{18} \approx 4{,}2$
        unidades, que con nuestra \textbf{escala}
        de 1\,u = 200\,m equivale a \SI{840}{m}.»
      \end{tcolorbox}
      \vspace{6pt}
      \begin{tcolorbox}[colback=AlmVerde!8, colframe=AlmVerde,
                        arc=4pt, boxrule=0.5pt, fontupper=\small]
        «La gráfica de turistas muestra una
        \textbf{tendencia creciente} de enero a agosto,
        lo que indica una \textbf{correlación positiva}
        con la temperatura.»
      \end{tcolorbox}
    \end{column}
  \end{columns}
\end{frame}

% ---------------------------------------------------------------
\seccionframe{4. Rúbrica de evaluación oral}
% ---------------------------------------------------------------

\begin{frame}{Rúbrica de evaluación de la exposición oral}
  \begin{center}
  \footnotesize
  \begin{tabular}{p{3.5cm} p{2.2cm} p{2.2cm} p{2.2cm} c}
    \toprule
    \textbf{Criterio} & \textbf{Excelente (3)} & \textbf{Bien (2)}
      & \textbf{Mejora (1)} & \textbf{Pts.}\\
    \midrule
    Exactitud matemática\\de los datos
      & Sin errores
      & 1 error menor
      & 2+ errores
      & /3\\[4pt]
    Claridad en la\\explicación de cálculos
      & Paso a paso,\\muy claro
      & Mayormente claro
      & Confuso o\\incompleto
      & /3\\[4pt]
    Vocabulario\\matemático
      & $\geq 10$ términos\\correctos
      & 6--9 términos
      & $<6$ términos
      & /2\\[4pt]
    Calidad visual\\(gráficas, mapas, tablas)
      & Todos completos\\y correctos
      & Algún elemento\\incompleto
      & Varios\\incompletos
      & /2\\
    \midrule
    \multicolumn{4}{r}{\textbf{TOTAL}} & \textbf{/10}\\
    \bottomrule
  \end{tabular}
  \end{center}
  \vspace{4pt}
  \begin{center}
    {\small 9--10: Sobresaliente \quad 7--8: Notable \quad
    5--6: Bien \quad $<5$: Necesita revisión}
  \end{center}
\end{frame}

% ---------------------------------------------------------------
\seccionframe{5. Síntesis: 4 bloques, 20 sesiones}
% ---------------------------------------------------------------

\begin{frame}{Lo que hemos aprendido — Los 4 bloques}
  \begin{columns}[T]
    \begin{column}{0.48\textwidth}
      \begin{tcolorbox}[enhanced, colback=AlmAzulM!12,
        colframe=AlmAzulM, arc=5pt,
        title={\faMapMarkerAlt\enskip Bloque 1: Mapa y coordenadas},
        fonttitle=\bfseries\small\color{AlmBlanco},
        attach boxed title to top left={yshift=-2.2mm, xshift=4mm},
        boxed title style={colback=AlmAzulM, arc=3pt}]
        \footnotesize
        Sesiones 1--4\\
        \textbf{Clave:} localizar cualquier punto con $(x,y)$,
        escala, tipos de información (Ubicación/Medida/Dato)
      \end{tcolorbox}
      \vspace{6pt}
      \begin{tcolorbox}[enhanced, colback=AlmTerra!10,
        colframe=AlmTerra, arc=5pt,
        title={\faCity\enskip Bloque 2: Geometría urbana},
        fonttitle=\bfseries\small\color{AlmBlanco},
        attach boxed title to top left={yshift=-2.2mm, xshift=4mm},
        boxed title style={colback=AlmTerra, arc=3pt}]
        \footnotesize
        Sesiones 5--10\\
        \textbf{Clave:} perímetro, área, polígonos,
        Teorema de Pitágoras para distancias
      \end{tcolorbox}
    \end{column}
    \begin{column}{0.48\textwidth}
      \begin{tcolorbox}[enhanced, colback=AlmVerde!10,
        colframe=AlmVerde, arc=5pt,
        title={\faChartLine\enskip Bloque 3: Datos y gráficas},
        fonttitle=\bfseries\small\color{AlmBlanco},
        attach boxed title to top left={yshift=-2.2mm, xshift=4mm},
        boxed title style={colback=AlmVerde, arc=3pt}]
        \footnotesize
        Sesiones 11--16\\
        \textbf{Clave:} variables, tablas, gráficas,
        proporcionalidad directa $y=kx$
      \end{tcolorbox}
      \vspace{6pt}
      \begin{tcolorbox}[enhanced, colback=BCrear!10,
        colframe=BCrear, arc=5pt,
        title={\faStar\enskip Bloque 4: Producto final},
        fonttitle=\bfseries\small\color{AlmBlanco},
        attach boxed title to top left={yshift=-2.2mm, xshift=4mm},
        boxed title style={colback=BCrear, arc=3pt}]
        \footnotesize
        Sesiones 17--20\\
        \textbf{Clave:} selección, construcción, revisión
        y presentación de la guía turística matemática
      \end{tcolorbox}
    \end{column}
  \end{columns}
\end{frame}

% ---------------------------------------------------------------
\seccionframe{6. Lo que hemos construido}
% ---------------------------------------------------------------

\begin{frame}{¿Qué hemos construido?}
  \begin{columns}[c]
    \begin{column}{0.55\textwidth}
      \begin{resumenbox}
        \textbf{Una guía turística matemática real de Almería:}
        \begin{itemize}
          \item \faMapMarkerAlt\ Mapa con coordenadas cartesianas
                y $\geq 10$ puntos de interés
          \item \faRuler\ Áreas y perímetros de zonas urbanas
          \item \faRoute\ Rutas calculadas con el Teorema de Pitágoras
          \item \faChartBar\ Gráficas de datos reales de Almería
          \item \faTable\ Tablas organizadas con datos verificados
          \item \faEquals\ Proporcionalidades y funciones lineales
          \item \faBookOpen\ Glosario de términos matemáticos
          \item \faLightbulb\ Datos curiosos con significado matemático
        \end{itemize}
      \end{resumenbox}
    \end{column}
    \begin{column}{0.41\textwidth}
      \begin{tikzpicture}[scale=0.7]
        % Guía representada como libro
        \fill[AlmAzul, rounded corners=5pt]
          (0,0) rectangle (5,7);
        \fill[AlmAmbar]
          (0,6.2) rectangle (5,7);
        \node[text=AlmAzul, font=\tiny\bfseries] at (2.5,6.6)
          {ALMERÍA EN CIFRAS Y FORMAS};
        \fill[AlmAzulC!30] (0.3,4.6) rectangle (4.7,6.0);
        \node[text=AlmAzul, font=\tiny, align=center] at (2.5,5.3)
          {mapa + coordenadas};
        \fill[AlmAmbar!20] (0.3,3.5) rectangle (2.2,4.4);
        \node[text=AlmNegro, font=\tiny, align=center] at (1.25,3.95)
          {tabla\\monum.};
        \fill[AlmVerde!20] (2.5,3.5) rectangle (4.7,4.4);
        \node[text=AlmNegro, font=\tiny, align=center] at (3.6,3.95)
          {gráfica\\turistas};
        \foreach \y in {2.7,2.3,1.9,1.5,1.1,0.7}{
          \fill[AlmBlanco!20] (0.3,\y) rectangle (4.7,\y+0.22);
        }
        \node[text=AlmAmbar, font=\tiny] at (2.5,0.3)
          {Matemáticas 1.º ESO · IES Al-Ándalus};
      \end{tikzpicture}
    \end{column}
  \end{columns}
\end{frame}

% ---------------------------------------------------------------
\seccionframe{7. Reflexión final}
% ---------------------------------------------------------------

\begin{frame}{Preguntas para la reflexión final}
  \begin{columns}[T]
    \begin{column}{0.48\textwidth}
      \begin{tcolorbox}[
        enhanced, colback=AlmAmbar!10, colframe=AlmAmbar!70!black,
        title={\faBrain\enskip Reflexiones matemáticas},
        fonttitle=\bfseries\small\color{AlmAzul},
        attach boxed title to top left={yshift=-2.2mm, xshift=4mm},
        boxed title style={colback=AlmAmbar, coltext=AlmAzul, arc=3pt},
        arc=4pt]
        \small
        \begin{enumerate}
          \item ¿Qué herramienta matemática
                (coordenadas, geometría, gráficas)
                ha resultado más \textbf{útil} para
                describir Almería?
          \vspace{4pt}
          \item ¿Qué dato de Almería os ha
                \textbf{sorprendido más} y por qué
                tiene importancia matemática?
          \vspace{4pt}
          \item Si tuvieras que guiar a un turista
                por Almería usando \textbf{solo matemáticas},
                ¿qué tres cosas le mostrarías?
        \end{enumerate}
      \end{tcolorbox}
    \end{column}
    \begin{column}{0.48\textwidth}
      \begin{tcolorbox}[
        enhanced, colback=AlmAzulM!10, colframe=AlmAzulM,
        title={\faStar\enskip Tarjeta de cierre},
        fonttitle=\bfseries\small\color{AlmBlanco},
        attach boxed title to top left={yshift=-2.2mm, xshift=4mm},
        boxed title style={colback=AlmAzulM, arc=3pt},
        arc=4pt]
        \small
        Completa estas frases en tu cuaderno:
        \begin{enumerate}
          \item «La herramienta matemática más útil
                para una guía turística es
                \underline{\hspace{2cm}}
                porque \underline{\hspace{2cm}}.»
          \item «Me sorprendió descubrir que
                \underline{\hspace{3cm}}.»
          \item «La próxima vez mejoraría
                \underline{\hspace{3cm}}.»
          \item «Las matemáticas en una ciudad
                son \underline{\hspace{3cm}}.»
        \end{enumerate}
      \end{tcolorbox}
    \end{column}
  \end{columns}
\end{frame}

% ---------------------------------------------------------------
\seccionframe[BCrear]{8. Ejercicios}
% ---------------------------------------------------------------

\begin{frame}{Ejercicios finales}
  \begin{columns}[T]
    \begin{column}{0.48\textwidth}
      \begin{evaluarbox}
        \footnotesize\dificultad{3}\\[4pt]
        Después de escuchar \textbf{tres exposiciones}
        de otros equipos, evalúa cada una con la rúbrica
        (puntuación numérica + 2 líneas de justificación
        por criterio).\\[6pt]
        ¿Qué aspectos destacan de los mejores trabajos?
        ¿Qué diferencia a una guía «notable» de una
        «sobresaliente»?
      \end{evaluarbox}
    \end{column}
    \begin{column}{0.48\textwidth}
      \begin{crearbox}
        \footnotesize\dificultad{4}\\[4pt]
        \textbf{Proyecto ampliado (opcional):}\\[6pt]
        ¿Y si aplicáramos el mismo método a
        \textbf{otra ciudad}? Elige una ciudad
        (tu ciudad natal, Madrid, Sevilla\ldots)
        y diseña las \textbf{primeras 2 páginas}
        de una guía turística matemática para ella,
        usando exactamente las mismas herramientas:
        mapa con coordenadas, tabla de monumentos
        y una gráfica de datos locales.
      \end{crearbox}
    \end{column}
  \end{columns}
\end{frame}

% ---------------------------------------------------------------
\begin{frame}{Evidencia del día}
  \begin{evidenciabox}
    Entrega final del proyecto — 3 elementos:
    \begin{enumerate}
      \item \textbf{Guía turística matemática} completa
            (versión final revisada con todas las correcciones)
      \item \textbf{Tarjeta de cierre} con las 4 frases completadas
            (reflexión personal)
      \item \textbf{Evaluación de 3 exposiciones} de otros equipos
            con rúbrica completada y justificación escrita
    \end{enumerate}
    \vspace{6pt}
    \begin{center}
      \large\color{AlmAmbar}
      {\Large\faAward}\quad
      \textbf{¡Enhorabuena por completar el proyecto!}
    \end{center}
  \end{evidenciabox}
\end{frame}

% ---------------------------------------------------------------
\begin{frame}{¡Cuidado con este error final!}
  \begin{errorbox}
    \textbf{«Las matemáticas son abstractas y no tienen
    conexión con la vida real»}\\[8pt]
    \incorrecto\ Esta afirmación es el error más grande
    de todos los que hemos visto en 20 sesiones.\\[8pt]
    \correcto\ Hemos demostrado exactamente lo contrario:\\
    Las matemáticas \textbf{sí están} en las calles de Almería,
    en las distancias entre monumentos, en las gráficas
    de turistas, en las áreas de los parques, en la
    proporcionalidad del precio de los transportes y en
    las coordenadas que nos llevan de un lugar a otro.\\[6pt]
    \textit{La ciudad que acabáis de describir matemáticamente
    es la misma ciudad en la que vivís cada día.}
  \end{errorbox}
\end{frame}

% ---------------------------------------------------------------
\begin{frame}{Cierre del proyecto — Almería en Cifras y Formas}
  \begin{center}
    \begin{tikzpicture}
      \fill[AlmAzul, rounded corners=8pt]
        (-5.5,-2.5) rectangle (5.5,2.5);
      \fill[AlmAmbar]
        (-5.5,1.8) rectangle (5.5,2.5);
      \node[text=AlmBlanco, font=\Large\bfseries] at (0,0.6)
        {Almería en Cifras y Formas};
      \node[text=AlmAmbar, font=\normalsize] at (0,-0.1)
        {20 sesiones · 4 bloques · 1 guía turística matemática};
      \node[text=AlmBlanco!70, font=\small] at (0,-0.8)
        {Matemáticas 1.º ESO · \InstituteName};
      \node[text=AlmBlanco!55, font=\footnotesize] at (0,-1.4)
        {\AuthorName · \SchoolYear};
      % Estrellas decorativas
      \foreach \a in {0,60,...,300}{
        \node[text=AlmAmbar, opacity=0.35, font=\scriptsize]
          at (\a:4.2) {$\bigstar$};
      }
    \end{tikzpicture}
  \end{center}
  \vspace{4pt}
  \begin{center}
    {\color{AlmAzulM}\large
    ¡Gracias por vuestro trabajo durante todo el proyecto!}\\[4pt]
    {\color{AlmGris}\footnotesize
    La matemática no es solo un libro de ejercicios.
    La matemática eres tú describiendo el mundo que te rodea.}
  \end{center}
\end{frame}

\end{document}


\end{document}
